% 纲要标题：# 第五卷　代数几何
% 纲要位置：几何学计划纲要.md，第 4728 行。
% 正文状态：未撰写；本文件目前是框架。

% 纲要内容开始
% 
% > 以方程、局部结构、相对几何和模问题为主线。读者具备基础环、域及线性代数；必要交换代数和同调方法在几何问题中系统建立。深度通过拆开长证明、重建依赖和完成算例实现，不通过提前堆入所有现代术语实现。
% 
% 纲要内容结束

% !TeX TXS-program:compile = cmd /c build-Book5.bat -NoPause
% !TeX TXS-program:view = txs:///view-pdf-internal --embedded "Book5.pdf"

\documentclass[UTF8,zihao=-4,openany,fontset=fandol]{ctexbook}
\newcommand{\BookTitle}{代数几何}
\newcommand{\BookVolumeName}{第五卷}
\newcommand{\BookNumber}{5}
\newcommand{\BookDraftDate}{2026年10月}
% 《几何学》五卷共用的基础排版与语义宏。
% 各卷主文件须先定义 \BookTitle、\BookVolumeName 和 \BookDraftDate。
\newcommand{\BookAuthor}{R. EverStarine}
% 稿件状态由各卷元数据指定；尚未开写的卷沿用框架状态。
\providecommand{\BookStatus}{框架预览·正文未撰写}

\usepackage[a4paper,inner=28mm,outer=24mm,top=27mm,bottom=28mm,
  headheight=16pt,headsep=8mm,footskip=12mm]{geometry}

\usepackage{amsmath}
\usepackage{mathtools}
\usepackage{amsthm}
\usepackage{aliascnt}
\usepackage{unicode-math}
% 数学字母使用默认数学字体内的对应字形；禁止 mathrm 等默认继承正文族。
% 运算符（sin、log、operatorname 等）随 mathrm 使用数学正体。
\unimathsetup{mathrm=sym,mathit=sym,mathbf=sym,mathsf=sym,mathtt=sym}
\setmainfont{texgyrepagella-regular.otf}[
  BoldFont=texgyrepagella-bold.otf,
  ItalicFont=texgyrepagella-italic.otf,
  BoldItalicFont=texgyrepagella-bolditalic.otf
]
\setsansfont{texgyreheros-regular.otf}[
  BoldFont=texgyreheros-bold.otf,
  ItalicFont=texgyreheros-italic.otf,
  BoldItalicFont=texgyreheros-bolditalic.otf
]
\setmonofont{lmmonolt10-regular.otf}[
  BoldFont=lmmonolt10-bold.otf,
  ItalicFont=lmmonolt10-oblique.otf,
  BoldItalicFont=lmmonolt10-boldoblique.otf
]
% 人物注释中的古希腊文、斯拉夫文字、日文异体字与希伯来文原名使用有相应字形的字体。
\newfontfamily\GeometryOriginalNameFont{Times New Roman}
\NewDocumentCommand{\OriginalName}{m}{{\GeometryOriginalNameFont #1}}
\newCJKfontfamily\GeometryOriginalHanFont{SimSun}
\NewDocumentCommand{\OriginalHanName}{m}{{\GeometryOriginalHanFont #1}}
\newfontfamily\GeometryHebrewFont{Times New Roman}[Script=Hebrew]
\TeXXeTstate=1
\NewDocumentCommand{\OriginalHebrewName}{m}{{\GeometryHebrewFont\beginR#1\endR}}
% 行内、行间及标题内公式均使用 unicode-math 默认 Latin Modern Math。
% 公式之外的正文西文仍由 setmainfont 使用 Pagella；黑板粗体单独采用 AMS。
\let\mathbb\relax
\DeclareMathAlphabet{\mathbb}{U}{msb}{m}{n}
% unicode-math 的 Latin Modern Math 将 \varnothing 映为瘦长椭圆；空集单独
% 采用传统 AMS MSBM 的圆形斜杠字形，不改变其余数学符号的字体。
\DeclareSymbolFont{GeometryAMSb}{U}{msb}{m}{n}
\DeclareMathSymbol{\GeometryRoundEmptySet}{\mathord}{GeometryAMSb}{"3F}
% 当前 Latin Modern Math 缺少横截符号 U+22D4；仅此关系符使用 AMSa 的现有字形。
\DeclareSymbolFont{GeometryAMSa}{U}{msa}{m}{n}
\DeclareMathSymbol{\GeometryTransverse}{\mathrel}{GeometryAMSa}{"74}
\AtBeginDocument{%
  \let\varnothing\GeometryRoundEmptySet
  \let\emptyset\GeometryRoundEmptySet
  \let\pitchfork\GeometryTransverse
}

\usepackage{microtype}
\usepackage{xcolor}
\usepackage{graphicx}
\usepackage{booktabs}
\usepackage{longtable}
\usepackage{array}
\usepackage{tabularx}
\usepackage{algorithm}
\usepackage{algpseudocode}
\usepackage{enumitem}
\usepackage{fancyhdr}
% 双面排版补齐奇数页时，空白页不残留上一章的页眉或页码。
\usepackage{emptypage}
\usepackage[most]{tcolorbox}
\usepackage{tikz-cd}
\usepackage{csquotes}
\usepackage[
  backend=biber,
  style=numeric-comp,
  sorting=nyt,
  giveninits=true,
  maxbibnames=99,
  doi=true,
  url=true,
  isbn=true
]{biblatex}
\addbibresource{../Shared/References.bib}
% 骨架没有引文时保留书目入口；写入真实条目后由完整构建自动排入书目。
\newif\ifGeometryHasBibliography
\AtDataInput{\global\GeometryHasBibliographytrue}
\NewDocumentCommand{\PrintBookBibliography}{}{%
  \ifGeometryHasBibliography
    \printbibliography[heading=bibintoc,title={参考文献}]%
  \else
    \chapter*{参考文献}%
    \addcontentsline{toc}{chapter}{参考文献}%
    \noindent 待撰写\par
  \fi
}
% 访问日期保留在书目库供内部核验，成品参考文献不显示 visited on。
\renewbibmacro*{url+urldate}{\usebibmacro{url}}
% biblatex 在书目内重绑定格式宏；在其作用域中将强调统一为正体，
% 同时避免书名中的中文及标点仍继承旧的斜体样式。
\AtBeginBibliography{\let\emph\textup\let\textit\textup}
% 中文条目的责任者连接符：不可套用西文的 “and”，按 GB/T 7714—2015 的著录
% 惯例改用逗号分隔。中文条目在 Shared/References.bib 中以 langid = {chinese}
% 标注（含译本，因其著录信息为中文形式），西文条目保持 “and” 不变。
% 按条目切换须在每条书目项处理时判定；\DeclareDelimFormat 的可选参数不是
% 语言选择器（实测无效），故用 \AtEveryBibitem 配合 \iffieldequalstr。
\AtEveryBibitem{%
  \iffieldequalstr{langid}{chinese}%
    {\DeclareDelimFormat{finalnamedelim}{\addcomma\space}%
     \DeclareDelimFormat{multinamedelim}{\addcomma\space}}%
    {\DeclareDelimFormat{finalnamedelim}{\addspace\bibstring{and}\space}%
     \DeclareDelimFormat{multinamedelim}{\addcomma\space}}%
}
\usepackage{etoolbox}
\usepackage{xparse}
\usepackage[noautomatic]{imakeidx}

\definecolor{GeometryBlue}{HTML}{244A64}
\definecolor{GeometryBlueLight}{HTML}{EEF4F7}
\definecolor{GeometryGray}{HTML}{4F5962}
\definecolor{GeometryGrayLight}{HTML}{F4F5F6}

\ctexset{
  part = {
    % 编题页分两行：第一行“第X编”，与封面“第一卷”同款（大号无衬线灰色）；
    % 第二行编名，与封面“几何学”同款（大号无衬线蓝色粗体），字号保持原样。
    % 第一行与行距由 number／aftername 给出，format 只负责编名的字体，
    % 以免字号被重置或被重复输出。
    format = {\centering\fontsize{30}{38}\selectfont\sffamily\bfseries\color{GeometryBlue}},
    name = {,},
    number = {\centering\normalfont\normalsize{\fontsize{24}{30}\selectfont\sffamily\color{GeometryGray}第\chinese{part}编}},
    aftername = {\par\vspace{0.7cm}},
    beforeskip = 0pt,
    afterskip = 2.5\baselineskip,
    pagestyle = plain
  },
  chapter = {
    format = \raggedright\Huge\bfseries\color{GeometryBlue},
    name = {第,章},
    number = \chinese{chapter},
    aftername = \quad,
    beforeskip = 0pt,
    afterskip = 30pt,
    pagestyle = plain
  },
  section = {
    format = \Large\bfseries\color{GeometryBlue},
    beforeskip = 2.4ex plus .6ex minus .2ex,
    afterskip = 1.1ex plus .2ex
  },
  subsection = {
    format = \large\bfseries\color{GeometryGray},
    beforeskip = 2ex plus .4ex minus .2ex,
    afterskip = .8ex plus .2ex
  }
}

\setlength{\parindent}{2em}
\setlength{\parskip}{0pt}
\linespread{1.18}
% 仅用于两个连续的普通正文段落；列表、公式、定理、标题等结构边界不用。
% 段间额外留白统一由这一长度控制，避免改动全局 \parskip 后扰乱各类环境。
\newlength{\GeometryProseParagraphGap}
\AtBeginDocument{%
  \setlength{\GeometryProseParagraphGap}{0.40\baselineskip}%
}
\NewDocumentCommand{\ProseParagraphBreak}{}{%
  \par\addvspace{\GeometryProseParagraphGap}%
}
% 与代数学保持一致：小节编号并进入目录，便于按具体主题定位。
\setcounter{secnumdepth}{3}
\setcounter{tocdepth}{2}
\allowdisplaybreaks[2]
\emergencystretch=2em
\clubpenalty=10000
\widowpenalty=10000
\displaywidowpenalty=10000

% 固定列表内部段距，避免 flushbottom 将默认的弹性 \parsep 拉大到超过正文段距。
\setlist[itemize]{leftmargin=2.2em,itemsep=.35ex,parsep=0pt,topsep=.6ex}
\setlist[enumerate]{leftmargin=2.5em,itemsep=.35ex,parsep=0pt,topsep=.6ex,
  font=\normalfont\upshape}
\setlist[enumerate,1]{label=\arabic*.,ref=\arabic*}
\setlist[enumerate,2]{label=\arabic*),ref=\arabic*}
\setlist[enumerate,3]{label=\alph*),ref=\alph*}
\setlist[description]{leftmargin=2.5em,labelindent=0pt,itemsep=.35ex,parsep=0pt,topsep=.6ex}

% “下列命题等价”专用列表；普通罗列仍使用 enumerate。
\newlist{equivalences}{enumerate}{2}
\setlist[equivalences,1]{leftmargin=2.8em,itemsep=.35ex,parsep=0pt,topsep=.6ex,
  label=\Roman*.,ref=\Roman*,font=\normalfont\upshape}
\setlist[equivalences,2]{leftmargin=2.8em,itemsep=.35ex,parsep=0pt,topsep=.6ex,
  label=\roman*),ref=\roman*,font=\normalfont\upshape}

\pagestyle{fancy}
\fancyhf{}
\fancyhead[LE]{\small\color{GeometryGray}\thepage\quad\leftmark}
% 奇数页采用页前最后有效的节标记，避免页中开始新节时提前更换页眉。
% 新节或习题的标记从后续页生效；章首页使用 plain，不显示上一章的标记。
\fancyhead[RO]{\small\color{GeometryGray}\TopMark{2e-right}\quad\thepage}
\renewcommand{\headrulewidth}{0.35pt}
% 保留数学标题的原始大小写，避免 book 类的 \MakeUppercase
% 与 unicode-math 命令在页眉标记中发生冲突。
% 偶数页页眉显示“章序＋章标题”（如“第一章 代数对象、映射与基本构造”）；
% 章序与章标题之间以 \quad 分隔，间距与奇数页节眉“节号＋节标题”之间一致；
% 附录章的章序由 \CTEXthechapter 按本卷附录字母（如 A）给出。
\renewcommand{\chaptermark}[1]{\markboth{\CTEXthechapter\quad #1}{}}
\renewcommand{\sectionmark}[1]{\markright{\thesection\quad #1}}
% 章末习题不编号，但进入习题页后应覆盖上一节留下的奇数页页眉；
% 同时写入目录项。PDF 书签由 hyperref 对 \section*{习题} 自动生成，
% 不再另行 \pdfbookmark，避免“习题”下再嵌套一个“习题”。
\newcommand{\BookExercises}{%
  \section*{习题}%
  \markright{习题}%
  \addcontentsline{toc}{section}{习题}%
}
\fancypagestyle{plain}{
  \fancyhf{}
  \fancyfoot[C]{\small\color{GeometryGray}\thepage}
  \renewcommand{\headrulewidth}{0pt}
}

\usepackage[unicode,colorlinks=true,linkcolor=GeometryBlue,
  citecolor=GeometryBlue,urlcolor=GeometryBlue,pdfborder={0 0 0}]{hyperref}
\usepackage{bookmark}
% PDF 书签沿用目录条目中的编号：正文章为“第几章”，附录章为“附录 A”，
% 节与小节分别保留 x.x、x.x.x；无编号的导读、习题和卷末单元仍不加号。
\bookmarksetup{numbered}

% unicode-math 会调整原始 Unicode 希腊字母的类别码。若把这些字母直接写进
% \texorpdfstring 的书签分支，辅助文件回读时类别码不同，LaTeX 会在每一轮都
% 误报标签发生变化。书签分支统一写成下列可展开宏，既保留希腊字母显示，
% 又使写入和回读的标签文本稳定。
\newcommand{\GeometryPdfAlpha}{α}
\newcommand{\GeometryPdfGamma}{Γ}
\newcommand{\GeometryPdfLambda}{λ}
\newcommand{\GeometryPdfMu}{μ}
\newcommand{\GeometryPdfPi}{π}
\newcommand{\GeometryPdfSigma}{σ}
\newcommand{\GeometryPdfTheta}{θ}
\newcommand{\GeometryPdfXi}{ξ}
\newcommand{\GeometryPdfZeta}{ζ}
\newcommand{\GeometryPdfCapitalLambda}{Λ}
\newcommand{\GeometryPdfOmega}{Ω}
\newcommand{\GeometryPdfCapitalSigma}{Σ}
\newcommand{\GeometryPdfDelta}{δ}
\newcommand{\GeometryPdfEpsilon}{ε}
\newcommand{\GeometryPdfIota}{ι}
\newcommand{\GeometryPdfSmallOmega}{ω}
\newcommand{\GeometryPdfEll}{ℓ}

\usepackage[nameinlink,noabbrev]{cleveref}
\newcommand{\crefpairconjunction}{和}
\newcommand{\crefmiddleconjunction}{、}
\newcommand{\creflastconjunction}{和}

% CTeX 2.5.10 的标题前缀宏带一个参数，而 TeX Live 2025 中 cleveref
% 仍按普通前缀宏展开它们，会把节号写成含 \endcsname 的损坏元数据。
% 改为恒等包装后，CTeX 传入的标题编号保持不变，\ref 与 \cref 均可正常工作。
\makeatletter
\def\p@part#1{#1}
\def\p@chapter#1{#1}
\def\p@section#1{#1}
\def\p@subsection#1{#1}
\def\p@subsubsection#1{#1}
\def\p@paragraph#1{#1}
\def\p@subparagraph#1{#1}
\makeatother
% 非数学西文字母统一正体；不改变括号类别、标点或数学字体。
% 2026-09-11：非数学模式的西文字母统一使用 Pagella 正体，保留粗细区别。
% 保留字母级边界，不重分类数字、括号、标点及中文。
% 此设置也覆盖列表序号字母和术语原文；数学模式保持原样。
\newfontfamily\GeometryLetterFont{texgyrepagella-regular.otf}[
  BoldFont=texgyrepagella-bold.otf,
  ItalicFont=texgyrepagella-italic.otf,
  BoldItalicFont=texgyrepagella-bolditalic.otf
]
\ExplSyntaxOn
\bool_new:N \l_geometry_text_letters_bool
\bool_set_true:N \l_geometry_text_letters_bool
\cs_new_protected:Npn \GeometryPreserveMathText
  { \bool_set_false:N \l_geometry_text_letters_bool }
\cs_new_protected:Npn \geometry_letter_group_begin:
  {
    \group_begin:
    \bool_if:NT \l_geometry_text_letters_bool
      { \GeometryLetterFont \upshape }
  }
% 继承 Default 的全部 xeCJK 转移，保留中西文间距及标点处理；
% 不能用只有开关分组的转移覆盖 Boundary 或 CJK 的原处理。
\xeCJK_new_class:n { GeometryLatinLetter }
\seq_map_inline:Nn \g__xeCJK_class_seq
  {
    \str_if_eq:nnF {#1} { GeometryLatinLetter }
      {
        \xeCJK_copy_inter_class_toks:nnnn
          {#1} {GeometryLatinLetter} {#1} {Default}
        \xeCJK_copy_inter_class_toks:nnnn
          {GeometryLatinLetter} {#1} {Default} {#1}
        \xeCJK_app_inter_class_toks:nnn {#1} {GeometryLatinLetter}
          { \geometry_letter_group_begin: }
        \xeCJK_pre_inter_class_toks:nnn {GeometryLatinLetter} {#1}
          { \group_end: }
      }
  }
% ASCII 字母及常用拉丁扩展字母；明确跳过乘号、除号，不重分类标点。
\xeCJKDeclareCharClass { GeometryLatinLetter }
  { "0041->"005A, "0061->"007A, "00C0->"00D6, "00D8->"00F6,
    "00F8->"024F, "1E00->"1EFF }
\ExplSyntaxOff
% 数学组内局部关闭字母切换，连公式中 \text{...} 的字体也保持原样。
% 只追加，不覆盖 amsmath、unicode-math、xeCJK 已有的数学初始化。
\AtBeginDocument{%
  \everymath\expandafter{\the\everymath\GeometryPreserveMathText}%
  \everydisplay\expandafter{\the\everydisplay\GeometryPreserveMathText}%
}


% 定理类环境（定理、定义、命题、引理、推论、断言、猜想、开放问题、注记、
% 习题）的正文：中文用楷体，西文保持正体，数学字体不受影响。
% 注意 \normalfont 只重置西文字族，中文字族仍随默认（宋体），故须显式 \kaishu；
% 不能退回 \itshape——那会把西文一并变斜。
\newtheoremstyle{geometryplain}
  {1.2ex}{1.2ex}{\normalfont\upshape\kaishu}{}
  {\color{GeometryBlue}\bfseries}{.}{0.6em}{}
\newtheoremstyle{geometrydefinition}
  {1.2ex}{1.2ex}{\normalfont\kaishu}{}
  {\color{GeometryBlue}\bfseries}{.}{0.6em}{}
\newtheoremstyle{geometryremark}
  {1.0ex}{1.0ex}{\normalfont\kaishu}{}
  {\color{GeometryGray}\bfseries}{.}{0.6em}{}

\theoremstyle{geometryplain}
\newtheorem{theorem}{定理}[section]
% 正文默认按“章.节.流水号”编号；未分节的附录省略中间的 0。
\renewcommand{\thetheorem}{%
  \ifnum\value{section}=0
    \thechapter.\arabic{theorem}%
  \else
    \thesection.\arabic{theorem}%
  \fi
}
\newaliascnt{lemma}{theorem}
\newtheorem{lemma}[lemma]{引理}
\aliascntresetthe{lemma}
\newaliascnt{proposition}{theorem}
\newtheorem{proposition}[proposition]{命题}
\aliascntresetthe{proposition}
\newaliascnt{corollary}{theorem}
\newtheorem{corollary}[corollary]{推论}
\aliascntresetthe{corollary}
\newaliascnt{claim}{theorem}
\newtheorem{claim}[claim]{断言}
\aliascntresetthe{claim}
\newaliascnt{conjecture}{theorem}
\newtheorem{conjecture}[conjecture]{猜想}
\aliascntresetthe{conjecture}
\newaliascnt{openproblem}{theorem}
\newtheorem{openproblem}[openproblem]{开放问题}
\aliascntresetthe{openproblem}

\theoremstyle{geometrydefinition}
\newaliascnt{definition}{theorem}
\newtheorem{definition}[definition]{定义}
\aliascntresetthe{definition}
\newaliascnt{example}{theorem}
\newtheorem{example}[example]{例}
\aliascntresetthe{example}
% 习题按章单独编号，不占用定理类环境的流水号。
\newtheorem{exercise}{习题}[chapter]

\theoremstyle{geometryremark}
\newaliascnt{remark}{theorem}
\newtheorem{remark}[remark]{注记}
\aliascntresetthe{remark}
\renewcommand{\proofname}{证明}
% 仅在证明紧接定理类陈述时收紧间距，并按陈述结尾是否为行间公式调整。
% 证明标题（证明／解／证明概要）的中文用楷体，
% 西文保持正体；正文不施加楷体，仍为默认宋体。不能退回 \itshape。
\newif\ifGeometryDirectProof
\newif\ifGeometryCoupledProof
\newif\ifGeometryStatementEndsInDisplay
\newskip\GeometryProofTrailingSkip
% 行间公式结束后，TeX 留下空的水平列表；在定理环境结束前辨认这一状态。
% 证明跟在公式之后时不再减去定理尾距，否则公式下方会比上方更紧。
\newcommand{\GeometryRecordStatementTail}{%
  \ifhmode
    \ifnum\lastnodetype=-1\relax
      \global\GeometryStatementEndsInDisplaytrue
    \else
      \global\GeometryStatementEndsInDisplayfalse
    \fi
  \else
    \global\GeometryStatementEndsInDisplayfalse
  \fi
}
\AddToHook{para/begin}{\global\GeometryDirectProoffalse\global\GeometryStatementEndsInDisplayfalse}
\AddToHook{env/theorem/end}{\GeometryRecordStatementTail}
\AddToHook{env/lemma/end}{\GeometryRecordStatementTail}
\AddToHook{env/proposition/end}{\GeometryRecordStatementTail}
\AddToHook{env/corollary/end}{\GeometryRecordStatementTail}
\AddToHook{env/claim/end}{\GeometryRecordStatementTail}
\AddToHook{env/exercise/end}{\GeometryRecordStatementTail}
\AddToHook{env/theorem/after}{\global\GeometryDirectProoftrue}
\AddToHook{env/lemma/after}{\global\GeometryDirectProoftrue}
\AddToHook{env/proposition/after}{\global\GeometryDirectProoftrue}
\AddToHook{env/corollary/after}{\global\GeometryDirectProoftrue}
\AddToHook{env/claim/after}{\global\GeometryDirectProoftrue}
\AddToHook{env/exercise/after}{\global\GeometryDirectProoftrue}
\AddToHook{env/figure/before}{\global\GeometryDirectProoffalse}
\AddToHook{env/table/before}{\global\GeometryDirectProoffalse}
\makeatletter
\renewenvironment{proof}[1][\proofname]{\par
  \pushQED{\qed}%
  \normalfont \topsep6\p@\@plus6\p@\relax
  \GeometryCoupledProoffalse
  \ifGeometryDirectProof
    \GeometryCoupledProoftrue
    \GeometryProofTrailingSkip=\topsep
    \advance\GeometryProofTrailingSkip by \partopsep
    \ifGeometryStatementEndsInDisplay\else
      \kern-5\p@% 只在陈述以普通文字结尾时抵消部分尾距。
    \fi
    \topsep2\p@\partopsep\z@\relax
  \fi
  \global\GeometryDirectProoffalse
  \global\GeometryStatementEndsInDisplayfalse
  \trivlist
  \item[\hskip\labelsep\upshape\kaishu #1\@addpunct{.}]\ignorespaces
}{%
  \popQED
  \ifGeometryCoupledProof\@topsepadd\GeometryProofTrailingSkip\fi
  \endtrivlist\@endpefalse
}
\makeatother
\newenvironment{solution}
  {\begin{proof}[解]}
  {\end{proof}}
\newenvironment{proofsketch}
  {\begin{proof}[证明概要]}
  {\end{proof}}

% 正文选读内容的星号属于编号而不属于标题；第二参数统一接收标签。
\newcommand{\GeometryOptionalMark}{\texorpdfstring{\textsuperscript{*}}{*}}
% 选做习题沿用同一编号星号。可选参数仍只表示习题自身的简短题名，
% 不得再写“选读”或“选做”；星号由本环境自动附在题号右侧。
\NewDocumentEnvironment{optionalexercise}{o}{%
  \begingroup
  \let\GeometrySavedTheExercise\theexercise
  \renewcommand{\theexercise}{\GeometrySavedTheExercise\GeometryOptionalMark}%
  \IfNoValueTF{#1}{\begin{exercise}}{\begin{exercise}[#1]}%
}{%
  \end{exercise}%
  \endgroup
}
% 可选参数只供目录与 PDF 书签使用；正文题名可借此按语义手动换行，
% 又不把换行命令写入目录或书签。
\NewDocumentCommand{\OptionalChapter}{O{}mm}{%
  \begingroup
  \let\GeometrySavedTheChapter\thechapter
  \renewcommand{\thechapter}{\GeometrySavedTheChapter\GeometryOptionalMark}%
  \ctexset{chapter={name={第,章\GeometryOptionalMark},
    number=\chinese{chapter},aftername=\quad}}%
  \ifstrempty{#1}{\chapter{#2}}{\chapter[#1]{#2}}\label{#3}%
  \endgroup
}
\NewDocumentCommand{\OptionalSection}{mm}{%
  \begingroup
  \let\GeometrySavedTheSection\thesection
  \renewcommand{\thesection}{\GeometrySavedTheSection\GeometryOptionalMark}%
  \ctexset{section={number=\thesection}}%
  \section{#1}\label{#2}%
  \endgroup
}
\NewDocumentCommand{\OptionalSubsection}{mm}{%
  \begingroup
  \let\GeometrySavedTheSubsection\thesubsection
  \renewcommand{\thesubsection}{\GeometrySavedTheSubsection\GeometryOptionalMark}%
  \ctexset{subsection={number=\thesubsection}}%
  \subsection{#1}\label{#2}%
  \endgroup
}

% 编题页必须从奇数页开始，题页仅在页脚显示页码，其后保留一张空白页。
\makeatletter
\NewDocumentCommand{\BookPart}{mm}{%
  \cleardoublepage
  \begingroup
  \def\GeometryPartLabel{#2}%
  % \part 会在返回前翻页，故须在 \@endpart 开始处记录题页标签。
  \pretocmd{\@endpart}{\label{\GeometryPartLabel}}{}%
    {\PackageError{geometry}{无法为编题页写入标签}{请检查当前文档类的 \string\@endpart 定义。}}%
  \part{#1}%
  \endgroup
  \thispagestyle{empty}%
  \null
  \clearpage
}
% 各编的“本编导读”是无编号章，在总目录和 PDF 书签中与所属编的各章并列。
% addcontentsline 同时建立目录项和书签，不再另加同名 pdfbookmark。
\NewDocumentCommand{\BookPartGuide}{O{本编导读}}{%
  \chapter*{#1}%
  \addcontentsline{toc}{chapter}{#1}%
}

% “附录”单独排无编号编题页，使用与正文编名相同的字号、字体和颜色。
% 题页从奇数页开始，其后保留一张空白页，再进入第一个附录章。
% startatroot 使根级书签与书后顶层节点分离；phantomsection 须先于
% addcontentsline，目录中的“附录”行才能获得独立锚点而非沿用前一节的
% section*.N。根级“附录”书签连同其子章由 addcontentsline 的 part 行自动
% 建立，不得再手写 \pdfbookmark（实测会产生重复的“附录”节点）。
\NewDocumentCommand{\BookAppendixPart}{m}{%
  \cleardoublepage
  \begingroup
  \def\GeometryPartLabel{#1}%
  % 在 part* 的题页尚未翻走时建立唯一的目录与书签目标。
  \pretocmd{\@endpart}{%
    \bookmarksetup{startatroot}%
    \phantomsection
    \addcontentsline{toc}{part}{附录}%
    \label{\GeometryPartLabel}%
  }{}{\PackageError{geometry}{无法为附录编题页写入目录与标签}{请检查当前文档类的 \string\@endpart 定义。}}%
  \part*{附录}%
  \endgroup
  \thispagestyle{empty}%
  \null
  \clearpage
}
\makeatother

% 列表项的括号与编号整体承担链接，避免只让中间的编号可点击。
\NewDocumentCommand{\itemref}{m}{\hyperref[#1]{(\ref*{#1})}}

% 章、节、小节的引用：整个“第 X 章／第 X 节／第 X 小节”都是跳转链接。
% 链接必须走 hyperref 的标签解析机制（目标名 chapter.N／section.N.M）；
% 早先用 \hyperlink{裸标签名} 指向不存在的目标，链接全部失效。
%
% 章号用汉字（“第五章”）且必须取“被引用那一章”的编号；\thechapter 是当前
% 计数器，在编导读、前言或跨章引用时会给出错误章号（表现为“第零章”）。
% hyperref 下 \ref 受保护、\getrefnumber 在 \edef 中又失真，故改由 zref 记录：
% 各章在 \chapter 之后调用 \BookChapterNumber，把汉字号存为 zref 属性；
% \zref@extract 可展开，能在 \edef 中安全取出。
\usepackage[user,hyperref]{zref}
\makeatletter
\zref@newprop{geometrychapternum}{}
% 在 \chapter 之后调用，参数即该章标签；把本章汉字号存为 zref 属性
\NewDocumentCommand{\BookChapterNumber}{m}{%
  \zref@setcurrent{geometrychapternum}{\chinese{chapter}}%
  \zref@labelbyprops{zch@#1}{geometrychapternum}%
}
% 取出某章标签对应的汉字号（\zref@extract 可展开，能在 \edef 中安全使用）
\NewDocumentCommand{\GeometryChineseChapterOf}{m}{%
  \zref@extract{zch@#1}{geometrychapternum}}
\makeatother
% \ref* 的分组会遮住编号末端的西文字符类别；编号前的间距可由 xeCJK
% 自动产生，编号后的同类间距须在完整链接内部显式补回。
\NewDocumentCommand{\secref}{m}{%
  \hyperref[#1]{第\ref*{#1}\CJKecglue 节}}
\NewDocumentCommand{\subsecref}{m}{%
  \hyperref[#1]{第\ref*{#1}\CJKecglue 小节}}
\NewDocumentCommand{\chapref}{m}{%
  \begingroup
  \xeCJKsetup{CJKecglue={}}%
  \edef\GeometryChapterDisp{第\GeometryChineseChapterOf{#1}章}%
  \hyperref[#1]{\GeometryChapterDisp}%
  \endgroup}

% 算法采用带编号的步骤，并明确标出输入与输出。
\floatname{algorithm}{算法}
\renewcommand{\listalgorithmname}{算法目录}
\algrenewcommand\algorithmicrequire{\textbf{输入：}}
\algrenewcommand\algorithmicensure{\textbf{输出：}}
\numberwithin{algorithm}{chapter}
% 完整中文编号进入题注、各自目录和引用；附录自动使用字母章号。
% 内部 \nobreakspace 与 cleveref 默认的定理引用一致；引用整串置于同一链接，
% 链接后的中文间距由下方 \cref 包装补齐。
\renewcommand{\thealgorithm}{算法\nobreakspace\thechapter-\arabic{algorithm}}
\renewcommand{\thefigure}{图\nobreakspace\thechapter-\arabic{figure}}
\renewcommand{\thetable}{表\nobreakspace\thechapter-\arabic{table}}
\renewcommand{\figurename}{}
\renewcommand{\tablename}{}
\makeatletter
\renewcommand{\fnum@figure}{\thefigure}
\renewcommand{\fnum@table}{\thetable}
\renewcommand{\fnum@algorithm}{\thealgorithm}
\makeatother
\NewDocumentEnvironment{algorithmsteps}{mm}
  {\begin{algorithmic}[1]\Require #1\Ensure #2}
  {\end{algorithmic}}

\newtcolorbox{learninggoals}{
  enhanced,breakable,colback=GeometryBlueLight,colframe=GeometryBlue,
  boxrule=.45pt,arc=1mm,left=1.2mm,right=1.2mm,top=1mm,bottom=1mm,
  title=本章目标,fonttitle=\bfseries
}
\newtcolorbox{partguide}{
  enhanced,breakable,colback=GeometryGrayLight,colframe=GeometryGray,
  boxrule=.35pt,arc=1mm,left=1.2mm,right=1.2mm,top=1mm,bottom=1mm,
  title=本编导读,fonttitle=\bfseries
}
\newtcolorbox{chaptersummary}{
  enhanced,breakable,colback=GeometryBlueLight,colframe=GeometryBlue,
  boxrule=.45pt,arc=1mm,left=1.2mm,right=1.2mm,top=1mm,bottom=1mm,
  title=本章小结,fonttitle=\bfseries,
  % 与阅读提示框一致：仅本框内容为多段散文，须恢复首行缩进。
  before upper={\parindent=2em\parskip=0pt}
}
% 仅用于尚待高阶数学复核的核心断言；成稿前必须处理或移除。
\newtcolorbox{coreverification}{
  enhanced,breakable,colback=yellow!10,colframe=orange!85!black,
  boxrule=.55pt,arc=1mm,left=1.2mm,right=1.2mm,top=1mm,bottom=1mm,
  title=待数学复核,fonttitle=\bfseries
}
% 兼容旧稿的占位命令：只输出普通文字，不另加方框。
\NewDocumentCommand{\GeometryPlaceholder}{m}{待撰写}
\NewDocumentEnvironment{chapterabstract}{}
  {\small\list{}{\rightmargin\leftmargin
    \listparindent=2\ccwd\itemindent=\listparindent}%
    \item\relax\color{GeometryGray}}
  {\endlist}
% 引诗与短引文统一居中，中文使用定理类环境的楷体；章摘要已在上方独立定义，
% 不继承这里的字形与对齐方式。
\RenewDocumentEnvironment{quote}{}
  {\list{}{\rightmargin\leftmargin}\item\relax\centering\kaishu}
  {\endlist}
% 引诗环境保留为语义清楚的兼容入口，版式与 quote 完全一致。
\NewDocumentEnvironment{poemquote}{}
  {\begin{quote}}
  {\end{quote}}
% 引诗末尾的作者或出处在引诗区域内右对齐。
\NewDocumentCommand{\quoteattribution}{m}{%
  \par\smallskip
  {\raggedleft #1\par}%
}

\makeindex[name=chinese,title=中文索引,columns=2,intoc,options=-q]
\makeindex[name=foreign,title=外文索引,columns=2,intoc,options=-q]
\makeindex[name=symbols,title=符号索引,columns=2,intoc,options=-q]
% 索引采用稍小字号与自然栏底，字母组间仅留少量空白；末页平衡不得越过正文高度。
\indexsetup{othercode={\small\raggedcolumns\maxbalancingoverflow=0pt
  \renewcommand{\indexspace}{\par\addvspace{0.2\baselineskip}}}}

% 有条目时读入 makeindex 生成的真实索引；尚无条目时仍保留书后结构与目录入口。
% 一键脚本会在 .idx 为空时删除旧 .ind，故这里不会误排上一次构建的陈旧索引。
\NewDocumentCommand{\PrintBookIndex}{m m}{%
  \IfFileExists{#1.ind}{%
    % 双栏输出可能丢失页眉标记；在索引范围内直接绑定标题并完成分页。
    \cleardoublepage
    \begingroup
      \fancyhead[LE]{\small\color{GeometryGray}\thepage\quad #2}%
      \fancyhead[RO]{\small\color{GeometryGray}#2\quad\thepage}%
      \printindex[#1]%
      \clearpage
    \endgroup
  }{%
    \chapter*{#2}%
    \addcontentsline{toc}{chapter}{#2}%
    \noindent 待撰写\par
  }%
}

% 所有术语在每卷首次正式定义或引入时都显示完整原文，并登记中外文索引。
% 星号形式与缺失原文均视为错误，防止旧的隐藏原文规则被重新带回源码。
% 定理类正文处于楷体；Fandol 楷体没有粗体字形，直接 \textbf 会退回普通楷体。
% 术语主称因此局部恢复正文的宋体字族，再使用与普通正文一致的粗体。
\NewDocumentCommand{\term}{s O{} m o}{%
  \IfBooleanT{#1}{%
    \PackageError{geometry}{\string\term* 已取消}{请改用 \string\term，并保留末尾的原文注解。}%
  }%
  {\songti\textbf{\textup{#3}}}%
  \IfNoValueTF{#4}{%
    \PackageError{geometry}{术语缺少原文注解}{每个术语首次正式定义或引入时都必须写成 \string\term[排序键]{中文主称}[原文]。}%
  }{%
    \textnormal{(\textup{#4})}%
  }%
  \ifstrempty{#2}{\index[chinese]{#3}}{\index[chinese]{#2@#3}}%
  \IfNoValueF{#4}{\index[foreign]{#4}}%
}
\NewDocumentCommand{\termidx}{O{}m}{%
  \ifstrempty{#1}{\index[chinese]{#2}}{\index[chinese]{#1@#2}}%
}
% 别称与主称一样用 \termidx 登记实际页码，不生成“见某某”的转引条目。
\NewDocumentCommand{\foreignidx}{O{}m}{%
  \ifstrempty{#1}{\index[foreign]{#2}}{\index[foreign]{#1@#2}}%
}
\newcounter{geometrySymbolIndexOrder}
\NewDocumentCommand{\symidx}{m}{%
  \stepcounter{geometrySymbolIndexOrder}%
  \edef\GeometrySymbolSortKey{%
    \number\numexpr10000+\value{geometrySymbolIndexOrder}\relax}%
  \index[symbols]{\GeometrySymbolSortKey@#1}%
}

% 一般强调与术语粗体分离。
\NewDocumentCommand{\keypoint}{m}{\CJKunderdot{#1}}
\NewDocumentCommand{\foreignkeypoint}{m}{\underline{\textup{#1}}}


\crefname{chapter}{章}{章}
\crefname{part}{编}{编}
\crefname{section}{节}{节}
\crefname{subsection}{小节}{小节}
\crefname{theorem}{定理}{定理}
\crefname{lemma}{引理}{引理}
\crefname{proposition}{命题}{命题}
\crefname{corollary}{推论}{推论}
\crefname{claim}{断言}{断言}
\crefname{conjecture}{猜想}{猜想}
\crefname{openproblem}{开放问题}{开放问题}
\crefname{definition}{定义}{定义}
\crefname{example}{例}{例}
\crefname{remark}{注记}{注记}
\crefname{exercise}{习题}{习题}
\crefname{equation}{式}{式}
\crefname{figure}{}{}
\crefname{table}{}{}
\crefname{algorithm}{}{}
\crefformat{figure}{#2#1#3}
\Crefformat{figure}{#2#1#3}
\crefformat{table}{#2#1#3}
\Crefformat{table}{#2#1#3}
\crefformat{algorithm}{#2#1#3}
\Crefformat{algorithm}{#2#1#3}

% cleveref 会在自己的分组中排出“定义 3.1.5”一类完整链接，类别名与编号之间
% 使用 \nobreakspace。分组结束后，xeCJK 无法再从末端识别出西文数字，
% 因而不会在紧随其后的汉字前插入 \CJKecglue。这里只在下一个源码记号确为
% CJK 正文字时补回同一种间距；空格、控制序列以及中英文标点均不处理，
% 从而不在“定义3.1.5，”或“定义3.1.5.”前制造额外空白。
\NewCommandCopy{\GeometryOriginalCref}{\cref}
\NewCommandCopy{\GeometryOriginalCrefCapital}{\Cref}
\ExplSyntaxOn
\cs_new_protected:Npn \geometry_cref_maybe_cjk_glue:
  { \peek_after:Nw \geometry_cref_maybe_cjk_glue_aux: }
\cs_new_protected:Npn \geometry_cref_maybe_cjk_glue_aux:
  {
    \token_if_space:NF \l_peek_token
      {
        \token_if_cs:NF \l_peek_token
          {
            \xeCJK_if_same_class:NNT \l_peek_token 中
              { \CJKecglue }
          }
      }
  }
\RenewDocumentCommand{\cref}{s m}
  {
    \IfBooleanTF{#1}
      { \GeometryOriginalCref*{#2} }
      { \GeometryOriginalCref{#2} }
    \geometry_cref_maybe_cjk_glue:
  }
\RenewDocumentCommand{\Cref}{s m}
  {
    \IfBooleanTF{#1}
      { \GeometryOriginalCrefCapital*{#2} }
      { \GeometryOriginalCrefCapital{#2} }
    \geometry_cref_maybe_cjk_glue:
  }
% 附录标签沿用字母编号；名称与编号整体进入同一链接，并与 \cref
% 采用相同的链接后中西文间距规则。
\NewDocumentCommand{\appref}{m}
  {
    \hyperref[#1]{附录\CJKecglue\ref*{#1}}
    \geometry_cref_maybe_cjk_glue:
  }
\ExplSyntaxOff

\newcommand{\MakeBookTitle}[1][封面]{%
  % 书名页在 PDF 侧栏中作为顶级入口；卷级文件可传入更具辨识性的书签名。
  \bookmarksetup{startatroot}%
  \pdfbookmark[0]{#1}{GeometryCover}%
  \begin{titlepage}
    \centering
    \vspace*{2.2cm}
    {\fontsize{36}{46}\selectfont\sffamily 几何学\par}
    \vspace{0.7cm}
    {\Large\sffamily\color{GeometryGray}\BookVolumeName\par}
    \vspace{1.0cm}
    {\fontsize{24}{32}\selectfont\bfseries\color{GeometryBlue}\BookTitle\par}
    \vfill
    {\large \BookAuthor\ 编写\par}
    \vspace{0.5cm}
    {\normalsize\BookDraftDate\par}
    \vspace{0.7cm}
    {\large\sffamily\bfseries\color{GeometryGray}\BookStatus\par}
    \vspace*{1.8cm}
  \end{titlepage}
}

% 目录也应有独立的顶级 PDF 书签。使用公共命令避免五卷分别遗漏。
\newcommand{\MakeBookContents}{%
  \cleardoublepage
  \bookmarksetup{startatroot}%
  \pdfbookmark[0]{目录}{GeometryContents}%
  \tableofcontents
}

% 图、表、算法目录不进入总目录，也不建立 PDF 书签；各自页面使用原有章级标题。
% 只另起一页而不补单双页。
\makeatletter
\newcommand{\GeometryListStart}{%
  \clearpage
}
% 各列表命令自带 \chapter*，局部消除其中的再次换页，保留原有标题与列表排版。
\newcommand{\GeometryListOnCurrentPage}[1]{%
  \begingroup
    \let\clearpage\relax
    \let\cleardoublepage\relax
    #1%
  \endgroup
}
\newcommand{\MakeBookListOfFigures}{%
  \GeometryListStart
  \GeometryListOnCurrentPage{\listoffigures}%
}
\newcommand{\MakeBookListOfTables}{%
  \begingroup
    \renewcommand{\listtablename}{表目录}%
    \GeometryListStart
    \GeometryListOnCurrentPage{\listoftables}%
  \endgroup
}
\newcommand{\MakeBookListOfAlgorithms}{%
  \begingroup
    \GeometryListStart
    % float.sty 在 listof 内重设编号栏；在标题回调中再加宽它。
    \let\GeometryOriginalFloatListHead\float@listhead
    \renewcommand{\float@listhead}[1]{%
      \renewcommand*\l@algorithm{\@dottedtocline{1}{1.5em}{5em}}%
      \GeometryOriginalFloatListHead{##1}%
    }%
    \GeometryListOnCurrentPage{\listofalgorithms}%
  \endgroup
}
\makeatother

% 无编号章（序言、目录、编导读、后记、参考文献、索引等）同样写入页眉标记，
% 使后续偶数页页眉显示该单元标题。置于 hyperref 之后，保证保留 hyperref 的章锚点。
\makeatletter
\let\Geometry@orig@schapter\@schapter
\def\@schapter#1{%
  \markboth{#1}{#1}%
  \Geometry@orig@schapter{#1}%
}
% 目录按层级着色：编为黑色、小节为深灰；章与节保持默认链接蓝。
\pretocmd{\l@part}{\hypersetup{linkcolor=black}}{}{%
  \PackageWarning{geometry}{无法为目录编条目着色}}
\pretocmd{\l@chapter}{\hypersetup{linkcolor=GeometryBlue}}{}{%
  \PackageWarning{geometry}{无法为目录章条目着色}}
\pretocmd{\l@section}{\hypersetup{linkcolor=GeometryBlue}}{}{%
  \PackageWarning{geometry}{无法为目录节条目着色}}
\pretocmd{\l@subsection}{\hypersetup{linkcolor=GeometryGray}}{}{%
  \PackageWarning{geometry}{无法为目录小节条目着色}}
% 编的目录条目与 PDF 书签须用纯文本“第X编”。ctex 的 \CTEX@part@tocline
% 取 \CTEXthepart，而该宏带上编题页的颜色与字体设置，会把格式代码写进目录，
% 并使书签出现 “=GeometryGray第一编” 之类乱码。此覆盖须置于 \ctexset 之后，
% 否则会被 ctex 重新定义；\CTEXifname 是 ctex 的内部开关，此处沿用。
\def\CTEX@part@tocline#1#2{%
  \CTEXifname
    {第\chinese{part}编\hspace{1em}}%
    {}%
  #2}
\makeatother

\hypersetup{pdftitle={几何学·\BookVolumeName：\BookTitle},pdfauthor={\BookAuthor},pdfsubject={几何学书系框架预览}}

\begin{document}
\frontmatter
% !TeX root = ../Book5.tex
\MakeBookTitle[\BookTitle]
\clearpage
\thispagestyle{empty}
\null
\clearpage
% !TeX root = ../Book5.tex
\chapter*{序言}
\addcontentsline{toc}{chapter}{序言}
本卷为《几何学》\BookVolumeName《\BookTitle》的框架预览。目录依据《几何学计划纲要》建立，选读标记与原纲要一致。

待撰写

\MakeBookContents
\MakeBookListOfFigures
\MakeBookListOfTables
\MakeBookListOfAlgorithms


\mainmatter
% !TeX root = ../Book5.tex
% 纲要标题：# 第一编　方程、局部化与局部环的结构
% 纲要位置：几何学计划纲要.md，第 5026 行。
% 正文状态：未撰写；本文件目前是框架。

% 纲要内容开始
% 
% 纲要内容结束

\BookPart{方程、局部化与局部环的结构}{g5:part:01}
\BookPartGuide
待撰写

% !TeX root = ../../Book5.tex
% 纲要标题：## 第1章　仿射、射影与有理几何的模型
% 纲要位置：几何学计划纲要.md，第 5028 行。
% 正文状态：初稿；尚未正式审核。

% 纲要内容开始
% 
% > 本章目标与证明范围：把经典几何入口合成有实质内容的一章：方程—函数—映射—消去—射影紧化。初始取代数闭域，进入概形后再系统改变基域；不把投影像等同于消去得到的闭包。
% 
% 纲要内容结束

\chapter{仿射、射影与有理几何的模型}
\label{g5:ch:01}
\BookChapterNumber{g5:ch:01}

\begin{chapterabstract}
	多项式方程既给出点集，也给出点上及点附近可以使用的函数。本章先证明多项式环的有限性与零点定理，建立仿射代数集和约化坐标环的反变对应，再以正则函数、函数域和稠密开集比较几何模型。射影坐标把乘积、子空间及曲线参数实现为齐次方程；多元除法、Gröbner 基和消去计算则从方程及参数中提取映射像的闭包。最后用结点、尖点、无穷远分支及重合纤维，说明点集、函数域与局部代数各自保留的信息。
\end{chapterabstract}

\begin{learninggoals}
	\begin{itemize}
		\item 证明多项式环的理想有限生成，使用弱、强零点定理把零点条件转换成理想条件，并找出不可约分量。
		\item 从局部分式判断函数与映射的正则性，以坐标环同态和函数域嵌入分别描述正则映射与支配有理映射。
		\item 写出射影仿射图、Segre 与 Veronese 嵌入及 Plücker 关系，并由像坐标恢复原来的几何对象。
		\item 完成多元除法和 Buchberger 算法，计算消去理想、饱和与有限商代数，区分实际投影像及其闭包。
		\item 在具体曲线与纤维中辨认分支、延拓障碍和幂零信息，说明局部环及保留方程的纤维为何必要。
	\end{itemize}
\end{learninggoals}

方程 \(xy=1\) 的每个非零横坐标都对应唯一纵坐标 \(y=1/x\)，但横坐标 \(0\) 没有原像。若只寻找所有横坐标共同满足的多项式方程，结果却是零理想：一个在所有非零元素处消失的一元多项式必为零。因此，方程可以准确描述这条曲线，却不能只靠横坐标上的方程描述其投影像；从方程消去变量时，必须知道所得对象与原来的点集有什么关系。

\ProseParagraphBreak

另一个变化发生在 \(t\mapsto(t^2,t^3)\)。这个参数把直线上的点一一映到 \(y^2=x^3\)，而逆参数 \(y/x\) 在原点附近不能用不消失的分母表示。两空间有相同的有理函数域，却有不同的局部函数。再看 \(x=0\) 和 \(x^2=0\)：二者零点相同，方程取商得到的环却不同。由方程到函数、再由函数理解映射，能够同时处理这些现象。

\ProseParagraphBreak

本章所需基础是域、理想、商环、域扩张以及有限维线性代数，并使用一元多项式带余除法、因式分解和行列式。开集、闭集、连续性、不可约性会在所研究的空间中说明。初始基域 \(k\) 取代数闭域；纯多项式算法及若干代数引理在陈述中给出较一般的基域条件，结点和纤维计算则逐处注明特征。仿射代数集允许可约，非空不可约者称为仿射簇。局部化、模和概形的一般理论随后建立，本章先通过可直接计算的模型引出它们。

\ProseParagraphBreak

仿射与射影簇、正则映射和有理映射可对照\cite[第I章第1--4节]{Hartshorne1977}；典型射影嵌入还可对照\cite[第1章第4节]{Shafarevich2013Basic1}。计算方法的阅读位置随\secref{g5:sec:0104}各项构造列出。本章完整证明所用的有限性、零点对应、映射对应及消去结论，曲线与有限商代数的算例将贯穿这些证明。

% !TeX root = ../../Book5.tex
% 纲要标题：### 1.1 方程与坐标环
% 纲要位置：几何学计划纲要.md，第 4738 行。
% 正文状态：未撰写；本文件目前是框架。

% 纲要内容开始
% 
% * Zariski 拓扑、理想和根式
% * Hilbert 基定理的有限性证明
% * 弱零点定理与 Zariski 引理
% * 强零点定理的 Rabinowitsch 论证
% * 不可约分支及反变对应
% 
% 纲要内容结束

\section{方程与坐标环}
\label{g5:sec:0101}

待撰写

\subsection{Zariski 拓扑、理想和根式}
\label{g5:subsec:0101:01}
待撰写

\subsection{Hilbert 基定理的有限性证明}
\label{g5:subsec:0101:02}
待撰写

\subsection{弱零点定理与 Zariski 引理}
\label{g5:subsec:0101:03}
待撰写

\subsection{强零点定理的 Rabinowitsch 论证}
\label{g5:subsec:0101:04}
待撰写

\subsection{不可约分支及反变对应}
\label{g5:subsec:0101:05}
待撰写
% 纲要条目排为编号小节标题，并进入目录与书签；说明段仍保留于纲要注释中。

% !TeX root = ../../Book5.tex
% 纲要标题：### 1.2 正则与有理映射
% 纲要位置：几何学计划纲要.md，第 5040 行。
% 正文状态：未撰写；本文件目前是框架。

% 纲要内容开始
% 
% * 正则函数和局部表达
% * 函数域及支配有理映射
% * 双有理等价和稠密开集
% * 图像、未定义点及延拓障碍
% 
% 纲要内容结束

\section{正则与有理映射}
\label{g5:sec:0102}

待撰写

\subsection{正则函数和局部表达}
\label{g5:subsec:0102:01}
待撰写

\subsection{函数域及支配有理映射}
\label{g5:subsec:0102:02}
待撰写

\subsection{双有理等价和稠密开集}
\label{g5:subsec:0102:03}
待撰写

\subsection{图像、未定义点及延拓障碍}
\label{g5:subsec:0102:04}
待撰写
% 纲要条目排为编号小节标题，并进入目录与书签；说明段仍保留于纲要注释中。

% !TeX root = ../../Book5.tex
% 纲要标题：### 1.3 射影空间与典型嵌入
% 纲要位置：几何学计划纲要.md，第 5047 行。
% 正文状态：初稿；尚未正式审核。

% 纲要内容开始
% 
% * 齐次方程和仿射图
% * Segre、Veronese 及其方程
% * Grassmann 与 Plücker 关系
% * 射影曲线及有理参数
% 
% 纲要内容结束

\section{射影空间与典型嵌入}
\label{g5:sec:0103}

在仿射平面中，双曲线 \(xy=1\) 的投影漏掉原点；一条直线的斜率也不能记录竖直方向。把坐标改成比例，可以把这些遗漏的方向作为点保存。新的问题随即出现：比例坐标满足什么方程，如何与仿射计算衔接，又怎样用一个点记录一对子空间或一个高次表达式？本节仍取代数闭域 \(k\)，从齐次方程建立射影模型，再把乘积、子空间和曲线参数写成可恢复的坐标。

\subsection{齐次方程和仿射图}
\label{g5:subsec:0103:01}

\begin{definition}
\label{g5:def:projective-space}
	设 \(k\) 为代数闭域，\(n\geq0\)。非零向量 \(k^{n+1}\setminus\{0\}\) 在关系 \(x\sim\lambda x\)（\(\lambda\in k^\times\)）下的商集称为 \(k\) 上的 \(n\) 维\term[sheyingkongjian]{射影空间}[projective space]，记为 \(\mathbb P^n_k\)。点写作 \([x_0:\cdots:x_n]\)，其中的 \(x_i\) 称为\term[qicizuobiao]{齐次坐标}[homogeneous coordinates]。
	\symidx{\(\mathbb P^n_k\)：域 \(k\) 上的射影空间}
\end{definition}

一个射影点等价地是一条通过原点的向量直线。这里共同倍数可以改变，单独一个坐标的值没有不依赖代表的意义。若 \(F\in k[x_0,\ldots,x_n]\) 是 \(d\) 次齐次多项式，则 \(F(\lambda x)=\lambda^dF(x)\)，所以条件 \(F(x)=0\) 却与代表无关。对一族齐次多项式 \(\mathcal F\)，定义
\[
	V_+(\mathcal F)=\{[x]\in\mathbb P^n_k:F(x)=0\text{ 对所有 }F\in\mathcal F\}.
\]
这些集合满足闭集公理：任意交对应合并方程，两个零点集的并对应两族方程的所有乘积；空集由 \(x_0=\cdots=x_n=0\) 给出，全集由零多项式给出。由此得到射影空间的 Zariski 拓扑。其闭集称为\term[sheyingdaishuji]{射影代数集}[projective algebraic set]；非空不可约的射影代数集称为\term[sheyingcu]{射影簇}[projective variety]。

把多项式按总次数分解成齐次部分，就有 \(k[x_0,\ldots,x_n]=\bigoplus_{d\geq0}S_d\)。由齐次多项式生成的理想称为\term[qicilixiang]{齐次理想}[homogeneous ideal]。每个元素的齐次部分仍在这个理想中，因为把生成式中的系数也逐次分解即可。射影方程因此可记为 \(V_+(J)\)，其中 \(J\) 是齐次理想。Hilbert 基定理保证方程可取有限个齐次生成元：先取有限生成元，再取它们的全部齐次部分。

\begin{proposition}
\label{g5:prop:projective-affine-charts}
	设 \(k\) 为代数闭域，\(n\geq1\)。开集 \(U_i=\{x_i\ne0\}\subseteq\mathbb P^n_k\) 覆盖 \(\mathbb P^n_k\)，而映射
	\[
		\phi_i:U_i\longrightarrow\mathbb A^n_k,\qquad
		[x_0:\cdots:x_n]\longmapsto(x_0/x_i,\ldots,\widehat{x_i/x_i},\ldots,x_n/x_i)
	\]
	是同胚。在两张图的交上，坐标变换及其逆都由分母不为零的多项式商给出。
\end{proposition}
\begin{proof}
	至少有一个坐标非零，所以这些开集覆盖全部点。\(\phi_i\) 的逆在第 \(i\) 个位置插入 \(1\)，故它是双射。齐次方程 \(F(x)=0\) 在 \(U_i\) 中成为令 \(x_i=1\) 得到的多项式方程。反过来，仿射坐标中的多项式 \(f\) 若次数为 \(d\)，则
	\[
		F(x)=x_i^d f(x_0/x_i,\ldots,\widehat{x_i/x_i},\ldots,x_n/x_i)
	\]
	是齐次多项式，且其在 \(U_i\) 上的零点恰对应 \(V(f)\)。这证明双向连续性。在 \(U_i\cap U_j\) 上写 \(u_\ell=x_\ell/x_i\)，补记 \(u_i=1\)，则第 \(j\) 张图中的坐标为 \(u_\ell/u_j\)，而 \(u_j\ne0\)。交换 \(i,j\) 即得逆变换。
\end{proof}

\(U_i\) 称为一张\term[fangshetu]{仿射图}[affine chart]。这使\secref{g5:sec:0102}的正则函数定义可以用于射影代数集：函数在每张仿射图中局部为多项式商，且分母在所讨论点不为零。图间变换正则，因而此定义不依赖图的选择。射影簇的开子集称为\term[zhunsheyingcu]{准射影簇}[quasi-projective variety]；允许可约时，同样讨论射影代数集的开子集及其正则函数。映射在源、靶的仿射图中为正则映射时，称为正则映射；有正则逆的双射称为同构。

例如，\(\mathbb P^1_k=\mathbb A^1_k\cup\{[0:1]\}\)：在 \(x_0\ne0\) 的图中用 \(t=x_1/x_0\)，另一图用 \(s=x_0/x_1\)，交上的关系是 \(s=1/t\)。每张 \(U_i\) 都在 \(\mathbb P^n_k\) 中稠密。事实上，若一个齐次多项式在 \(U_i\) 上为零，其去齐次化在 \(k^n\) 上恒为零；无限域上的多项式恒等原理表明该多项式为零，再齐次化就得到原式为零。仿射空间不可约，因此其稠密闭包 \(\mathbb P^n_k\) 也不可约。

若 \(F_0,\ldots,F_m\) 是同次齐次多项式，则
\[
	[x]\longmapsto[F_0(x):\cdots:F_m(x)]
\]
在这些多项式不同时为零的开集上定义正则映射：在 \(F_j\ne0\) 的靶图中，坐标恰是 \(F_i/F_j\)。共同零点处的公式不能直接赋值；能否换用别的局部表达来延拓，要另行检查。本节末的曲线会给出一个可以延拓的例子。

上述齐次方程与仿射图的构造可对照\cite[第I章，第2节，命题2.1--2.2]{Hartshorne1977}。

\subsection{Segre、Veronese 及其方程}
\label{g5:subsec:0103:02}

考虑一对射影点 \([x]\in\mathbb P^m_k\)、\([y]\in\mathbb P^n_k\)。把两组坐标直接并排，不能消去它们各自独立的缩放；取乘积 \(x_i y_j\) 则使两次缩放合成一个共同倍数。乘积 \(\mathbb P^m_k\times\mathbb P^n_k\) 用图 \(U_a\times U_b\cong\mathbb A^{m+n}_k\) 赋予代数几何结构，正则性按这些图检验。两组坐标分别换图，所以转换公式仍为分母非零的多项式商。这里的 Zariski 拓扑由仿射图中的多项式方程确定。

\begin{proposition}
\label{g5:prop:segre-embedding}
	设 \(k\) 为代数闭域，\(m,n\geq0\)。\term[Segreqianru]{Segre 嵌入}[Segre embedding]
	\[
		\sigma:\mathbb P^m_k\times\mathbb P^n_k\longrightarrow\mathbb P^{(m+1)(n+1)-1}_k,
		\qquad([x],[y])\longmapsto[x_i y_j]_{i,j}
	\]
	将其定义域同构到由全部方程
	\[
		z_{ij}z_{ab}-z_{ib}z_{aj}=0
		\qquad(0\leq i,a\leq m,\ 0\leq j,b\leq n)
	\]
	定义的射影代数集。
\end{proposition}

这里嵌入指与一个闭代数子集的同构，因而同时要求像的闭性与逆映射的正则性。Segre 的乘积坐标把这两项要求化为秩一矩阵的恢复。\footnote{Corrado Segre（1863--1924），意大利数学家，研究高维射影几何与代数曲面。}

\begin{proof}
	分别缩放 \(x,y\) 只会共同缩放全部乘积，且至少一个乘积非零，所以 \(\sigma\) 良定义。像矩阵 \((x_i y_j)\) 的所有二阶子式为零。反过来，设 \(Z=(z_{ij})\ne0\) 满足这些方程，选 \(z_{ab}\ne0\)，则
	\[
		\frac{z_{ij}}{z_{ab}}
		=\frac{z_{ib}}{z_{ab}}\frac{z_{aj}}{z_{ab}}.
	\]
	令 \(x_a=y_b=1\)、\(x_i=z_{ib}/z_{ab}\)、\(y_j=z_{aj}/z_{ab}\)，就恢复了原来的两个射影点。这些比值还说明恢复唯一。源图 \(x_a y_b\ne0\) 到靶图 \(z_{ab}\ne0\) 的坐标是多项式，逆坐标是上述比值，故双向正则。
\end{proof}

对 \(m=n=1\)，得到二次曲面
\[
	Q=V_+(z_{00}z_{11}-z_{01}z_{10})\subseteq\mathbb P^3_k.
\]
固定 \([x]\) 让 \([y]\) 变化，或固定 \([y]\) 让 \([x]\) 变化，分别得到一族直线。每族内不同直线不交，两族各取一条恰交于一个点，因为原像中这正是固定一个坐标因子的两个条件。\(Q\) 上的对角线像 \(\sigma(\{([x],[x])\})\) 又由 \(z_{01}=z_{10}\) 切出。乘积的代数几何 Zariski 拓扑因而不同于两个 Zariski 拓扑的普通积拓扑：\(\mathbb P^1_k\) 的每个非空开集都是有限点的补集，任意两个非空开集相交，所以对角线在普通积拓扑中稠密，却在这里是一个真闭集。

另一种构造只用一个点，却同时记录它的全部 \(d\) 次单项式。它称为\term[Veroneseqianru]{Veronese 嵌入}[Veronese embedding]。\footnote{Giuseppe Veronese（1854--1917），意大利数学家，研究高维射影几何与几何基础。}记 \(e_0,\ldots,e_n\) 为标准多重指标，\(\alpha=(\alpha_0,\ldots,\alpha_n)\) 的各分量非负，\(|\alpha|=\sum\alpha_i\)，\(x^\alpha=\prod x_i^{\alpha_i}\)。

\begin{proposition}
\label{g5:prop:veronese-embedding}
	设 \(k\) 为代数闭域，\(n\geq1\)、\(d\geq1\)，\(N=\binom{n+d}{d}-1\)。映射
	\[
		\nu_d:\mathbb P^n_k\longrightarrow\mathbb P^N_k,
		\qquad[x]\longmapsto[x^\alpha]_{|\alpha|=d}
	\]
	是到闭代数子集的同构。用 \(z_\alpha\) 表示靶坐标，其像可由以下有限族齐次方程刻画：
	\begin{align*}
		z_\alpha^d&=\prod_{j=0}^n z_{de_j}^{\alpha_j},\\
		z_\alpha z_{de_i}^{d-1}
		&=\prod_{j=0}^n z_{(d-1)e_i+e_j}^{\alpha_j}
		\qquad(0\leq i\leq n,\ |\alpha|=d).
	\end{align*}
\end{proposition}
\begin{proof}
	坐标中包括每个 \(x_i^d\)，所以它们不会同时为零；所有坐标同次，故映射良定义且正则。代入 \(z_\alpha=x^\alpha\) 直接验证所列方程。

	设一个射影点 \([z]\) 满足这些方程。选 \(z_\alpha\ne0\)，第一族方程保证至少一个纯幂坐标 \(z_{de_i}\ne0\)。将全部坐标除以它，令
	\[
		x_j=\frac{z_{(d-1)e_i+e_j}}{z_{de_i}},\qquad x_i=1.
	\]
	第二族方程给出 \(z_\alpha/z_{de_i}=\prod_j x_j^{\alpha_j}\)，因而 \([z]=\nu_d([x])\)。在这张图上，\(x_j=x_j/x_i\) 的恢复公式唯一且正则；各图覆盖整个像，证明了同构。
\end{proof}

这些方程给出的是像的零点集描述。就计算而言，二次情形还有更简洁的方程：把坐标 \(z_{ij}\)（\(i\leq j\)）对称地补成矩阵 \(Z\)，则 \(\nu_2\) 的像恰由 \(Z\) 的全部二阶子式为零定义。若 \(Z\ne0\) 且全部二阶子式为零，某个 \(z_{aa}\) 必非零，否则 \(z_{ij}^2=z_{ii}z_{jj}=0\) 强迫 \(Z=0\)。随后
\[
	\frac{z_{ij}}{z_{aa}}
	=\frac{z_{ia}}{z_{aa}}\frac{z_{ja}}{z_{aa}}
\]
恢复 \(x_i/x_a=z_{ia}/z_{aa}\)。这个证明在特征 \(2\) 也成立，不需要取平方根。

Veronese 嵌入把高次方程变为线性截面：若 \(F(x)=\sum_{|\alpha|=d}c_\alpha x^\alpha\)，则 \(F(x)=0\) 等价于像点满足超平面方程 \(\sum c_\alpha z_\alpha=0\)。Segre 嵌入则把双线性方程变为超平面截面。这些构造可对照\cite[第I章，第2节，习题2.12--2.15]{Hartshorne1977}；Veronese 的二次关系式还可参看\cite[第1章，第4节，例1.28]{Shafarevich2013Basic1}。

\subsection{Grassmann 与 Plücker 关系}
\label{g5:subsec:0103:03}

射影空间以一维向量子空间为点。若希望把一个 \(r\) 维子空间也当作一个几何点，直接记录一组基会多记一次任意换基。全部最大子式会在换基时乘以同一个行列式，恰好可以作为射影坐标。

\begin{definition}
\label{g5:def:grassmannian}
	设 \(k\) 为代数闭域，\(V=k^n\)，\(1\leq r\leq n\)。所有 \(r\) 维向量子空间组成的集合称为\term[Grassmannkongjian]{Grassmann 空间}[Grassmannian]，记为 \(\operatorname{Gr}(r,V)\) 或 \(\operatorname{Gr}(r,n)\)。
	\symidx{\(\operatorname{Gr}(r,V)\)：\(V\) 的 \(r\) 维子空间的参数空间}
\end{definition}

Grassmann 空间将“一个子空间怎样变化”变成了点的坐标问题。\footnote{Hermann Günter Grassmann（1809--1877），德国数学家，发展了向量运算与外代数。}选子空间的一组基作为 \(r\times n\) 矩阵 \(A\) 的行，令 \(p_{i_1\cdots i_r}\) 为第 \(i_1,\ldots,i_r\) 列组成的子式。先对递增指标记录坐标；一般指标按交替规则延伸，有重复指标时规定为零。换基将所有子式乘以同一个非零行列式，且矩阵满秩保证它们不全为零。这些\term[Pluckerzuobiao]{Plücker 坐标}[Plücker coordinates]给出映射
\[
	\iota:\operatorname{Gr}(r,n)\longrightarrow\mathbb P^{\binom nr-1}_k,
	\qquad W\longmapsto[p_{i_1\cdots i_r}].
\]
Plücker 坐标特别适于描述直线族。\footnote{Julius Plücker（1801--1868），德国数学家、物理学家，研究解析几何、直线几何与光谱。}

\begin{theorem}[Plücker 关系]
\label{g5:thm:plucker-relations}
	设 \(k\) 为代数闭域，\(1\leq r\leq n\)。由基矩阵的全部 \(r\) 阶子式定义的映射 \(\iota:\operatorname{Gr}(r,n)\rightarrow\mathbb P^{\binom nr-1}_k\) 为单射，其像恰由齐次二次方程
	\begin{equation}
	\label{g5:eq:plucker-relations}
		\sum_{t=1}^{r+1}(-1)^{t-1}
		p_{i_1\cdots i_{r-1}j_t}
		p_{j_1\cdots\widehat{j_t}\cdots j_{r+1}}=0
	\end{equation}
	给出，其中指标取自 \(\{1,\ldots,n\}\)，并采用交替约定。像中每个开集 \(p_{i_1\cdots i_r}\ne0\) 同构于 \(\mathbb A^{r(n-r)}_k\)，同构由将相应 \(r\) 列化为单位矩阵给出。
\end{theorem}

\begin{proof}
	必要性来自一个行列式恒等式。对 \(r\) 维空间中的任意 \(r+1\) 个列向量 \(v_{j_t}\)，按余子式展开得到
	\[
		\sum_{t=1}^{r+1}(-1)^{t-1}
		\det(v_{j_1},\ldots,\widehat{v_{j_t}},\ldots,v_{j_{r+1}})v_{j_t}=0.
	\]
	可逐个分量检验：每个分量都是一个含两行相同的 \((r+1)\) 阶行列式的展开。对这个向量等式作用线性函数 \(v\mapsto\det(v_{i_1},\ldots,v_{i_{r-1}},v)\)，便得到\eqref{g5:eq:plucker-relations}。

	为证明充分性，需要从满足方程的 \([p]\) 恢复矩阵，并证明它的所有子式都等于给定坐标。重排列后可设 \(p_{1\cdots r}\ne0\)，再共同缩放使 \(p_{1\cdots r}=1\)。定义矩阵 \(A\) 的前 \(r\) 列为单位矩阵，其余元素为
	\[
		A_{aj}=p_{1\cdots(a-1)j(a+1)\cdots r}
		\qquad(1\leq a\leq r,\ j>r).
	\]
	设 \(q_J\) 是这个矩阵的子式。对仅含基指标 \(1,\ldots,r\) 的 \(r\) 元组，交替性给出 \(q_J=p_J\)。一般情形按 \(J\) 中基指标的缺失数归纳。取缺失的 \(a\in\{1,\ldots,r\}\)，在\eqref{g5:eq:plucker-relations}中代入
	\[
		(i_1,\ldots,i_{r-1})=(1,\ldots,\widehat a,\ldots,r),
		\qquad(j_1',\ldots,j_{r+1}')=(a,j_1,\ldots,j_r).
	\]
	利用 \(p_{1\cdots\widehat a\cdots r a}=(-1)^{r-a}\)，得到
	\[
		p_{j_1\cdots j_r}
		=\sum_{t=1}^r(-1)^{t-1}A_{a j_t}\,
		p_{a j_1\cdots\widehat{j_t}\cdots j_r}.
	\]
	若 \(j_t\) 是基指标，则 \(j_t\ne a\)，从而 \(A_{a j_t}=0\)。其余非零项中的子式少缺一个基指标，故可用归纳假设把 \(p\) 换成 \(q\)。换后右边恰是 \(\det A_J\) 沿第 \(a\) 行的展开：\(q_{a,J\setminus j_t}\) 的首列是 \(e_a\)，其展开再贡献符号 \((-1)^{a+1}\)。所以右边等于 \(q_J\)，充分性得证。

	矩阵 \((I_r\ B)\) 的 \(r(n-r)\) 个自由元素可以任取，全部子式是它们的多项式；反向恢复则是 Plücker 坐标除以选定的非零子式。这既证明每张图与仿射空间的同构，也证明行空间由 \([p]\) 唯一确定。换用别的非零子式时，恢复的是同一行空间，图间变换因而相容。
\end{proof}

由这些方程及仿射图，\(\operatorname{Gr}(r,n)\) 获得射影代数集的结构，\(\iota\) 称为\term[Pluckerqianru]{Plücker 嵌入}[Plücker embedding]。它的不可约性也可直接检查：满秩矩阵构成 \(\mathbb A^{rn}_k\) 的非空开集，故不可约；取全部子式给出到 Grassmann 空间的连续满射，而不可约空间的连续像不可约。此处满射来自每个子空间都能选基，连续性则由齐次方程的多项式代入给出。因此这个参数空间是一射影簇。上述构造与关系式可对照\cite[第1章，第4节，例1.24，公式(1.24)]{Shafarevich2013Basic1}。

\begin{example}
\label{g5:exm:grassmann-24}
	设 \(k\) 为代数闭域。\(\operatorname{Gr}(2,4)\) 的六个坐标满足
	\[
		p_{12}p_{34}-p_{13}p_{24}+p_{14}p_{23}=0.
	\]
	这一个方程已经足够。在 \(p_{12}\ne0\) 的图中归一化 \(p_{12}=1\)，恢复矩阵
	\[
		A=\begin{pmatrix}
			1&0&-p_{23}&-p_{24}\\
			0&1&p_{13}&p_{14}
		\end{pmatrix}.
	\]
	前五个子式就是给定坐标，最后一个是 \(p_{13}p_{24}-p_{14}p_{23}=p_{34}\)。任意非零子式都可通过重排列移到 \(p_{12}\)，所以这些图覆盖全部方程解。这个二次超曲面参数化 \(\mathbb P^3_k\) 中的射影直线；图中的四个自由参数就是上述矩阵右侧四个元素。整个计算在任意特征下成立。
\end{example}

\subsection{射影曲线及有理参数}
\label{g5:subsec:0103:04}

仿射曲线参数常在某个参数值出现分母为零。射影参数把这个值放到另一张仿射图中，因而可以检查那里究竟是遗漏了一个点，还是映射存在真实的延拓障碍。不可约准射影簇上的有理映射仍按稠密开集上正则映射的等价类定义；函数域可在任一非空仿射开集中计算，因为交上的正则函数给出相同的分式域。若其函数域的超越次数为 \(1\)，我们称它为曲线；这一描述把本节的参数计算与后续维数理论联系起来。

\begin{example}
\label{g5:exm:projective-conic-parameter}
	设 \(k\) 为代数闭域。二次曲线
	\[
		C=V_+(XZ-Y^2)\subseteq\mathbb P^2_k
	\]
	具有处处定义的参数化
	\[
		\nu:\mathbb P^1_k\longrightarrow C,\qquad[s:t]\longmapsto[s^2:st:t^2].
	\]
	它是二次 Veronese 嵌入，其逆还可直接写成
	\[
		\rho([X:Y:Z])=
		\begin{cases}
			[X:Y],&X\ne0,\\
			[Y:Z],&Z\ne0.
		\end{cases}
	\]
	这两个开集覆盖 \(C\)，因为 \(X=Z=0\) 会由曲线方程推出 \(Y=0\)。在交上，\([X:Y]=[Y:Z]\) 正是 \(XZ=Y^2\)，所以局部公式相容。代入可知 \(\rho\nu\) 和 \(\nu\rho\) 均为恒等映射，且各局部表达正则。这证明 \(C\cong\mathbb P^1_k\)，并给出 \(k(C)=k(t/s)\)。
\end{example}

从点 \(Q=[0:0:1]\) 作投影，可把这个参数化看成线与二次曲线的相交。经过 \(Q\) 的直线由比例 \([X:Y]=[s:t]\) 标记；在 \(s\ne0\) 时，除去固定交点 \(Q\)，另一个交点为 \([s^2:st:t^2]\)。当 \(s=0\) 时，这个交点也移到 \(Q\)。下面的图记录投影公式与其延拓的关系：
\[
	\begin{tikzcd}[column sep=large,row sep=large]
		C\setminus\{Q\}\arrow[r,"{[X:Y]}"]\arrow[d,hook]
		&\mathbb P^1_k\arrow[d,equal]\\
		C\arrow[r,"\rho"']&\mathbb P^1_k.
	\end{tikzcd}
\]
上边的公式在 \(Q\) 处同时为零，下边的局部公式 \([Y:Z]\) 却在 \(Q\) 附近有定义，并给出 \(\rho(Q)=[0:1]\)。同一个有理映射的某组齐次表达式有共同零点，并不决定该点一定无法延拓；应检查映射在所研究簇上的其他局部表达。

更一般地，\(\nu_d(\mathbb P^1_k)\subseteq\mathbb P^d_k\) 写作
\[
	[s:t]\longmapsto[s^d:s^{d-1}t:\cdots:st^{d-1}:t^d],
\]
称为\term[youlizhengguiquxian]{有理正规曲线}[rational normal curve]。\cref{g5:prop:veronese-embedding}给出的逆图表明它同构于 \(\mathbb P^1_k\)，所以其函数域都是 \(k(t/s)\)；嵌入坐标和所处射影空间则随 \(d\) 改变。函数域记录双有理信息，而具体嵌入还记录方程如何放在周围空间中。下一节的消去运算将从参数式计算像的方程，并说明哪些方程只确定了像的闭包。

% !TeX root = ../../Book5.tex
% 纲要标题：### 1.4 消去与有效计算
% 纲要位置：几何学计划纲要.md，第 5054 行。
% 正文状态：初稿；尚未正式审核。

% 纲要内容开始
% 
% * 全局单项式序和多元除法
% * Buchberger 判据及终止性
% * 消去定理与投影闭包
% * 饱和和射影分量
% * 零维商代数与重数核验
% 
% 纲要内容结束

\section{消去与有效计算}
\label{g5:sec:0104}

方程组给定以后，怎样判定另一个方程是否由它推出？怎样从参数式得到隐式方程，又怎样保留相交时的重数？这些问题需要在理想中实际运算。单变量多项式的除法提供了起点：选出最高项，逐次消去。多变量中首先要决定什么叫“最高”，其次要处理不同消去路线之间的抵消。本节先建立 Gröbner 基方法，再把它用于投影、射影闭包和有限点集。\footnote{Wolfgang Gröbner（1899--1980），奥地利数学家，研究代数几何及 Lie 级数；Gröbner 基以其姓氏命名。}

多项式计算部分适用于任意域；讨论零点与闭包时仍取代数闭域 \(k\)。实际执行算法还要求能够精确完成系数的四则运算和判零，例如在有理数域或明确给定的代数数域中计算。单项式序、除法和 Gröbner 基算法的进一步材料见\cite[第2章第2--7节]{CoxLittleOShea2015}。

\subsection{全局单项式序和多元除法}
\label{g5:subsec:0104:01}
在 \(k[x_1,\ldots,x_n]\) 中，\(x_1^2\) 与 \(x_2^3\) 哪一个应先消去，并没有由总次数之外的结构自动决定的答案。为了使一次消去不会重新产生更高的项，我们需要与乘法相容的次序；为了使消去最终结束，还需要良序性。

\begin{definition}\label{g5:def:monomial-order}
	设 \(k\) 为域。多项式环 \(R=k[x_1,\ldots,x_n]\) 的单项式 \(x^\alpha=x_1^{\alpha_1}\cdots x_n^{\alpha_n}\) 上的全序 \(\prec\)，若满足每个非空单项式集合都有最小元，且
	\[
		x^\alpha\prec x^\beta\quad\Longrightarrow\quad x^{\alpha+\gamma}\prec x^{\beta+\gamma},
	\]
	就称为\term[quanjudangxiangshixu]{全局单项式序}[global monomial order]。其中 \(1\) 必为最小单项式：否则某个 \(m\prec1\) 会给出无限下降链 \(1\succ m\succ m^2\succ\cdots\)。
\end{definition}

最常用的两个例子是字典序与分次字典序。固定变量优先次序 \(x_1\succ\cdots\succ x_n\)。\term[zidianxu]{字典序}[lexicographic order]从左到右比较指数，在第一个不同的位置上指数较大者为大；\term[fencizidianxu]{分次字典序}[graded lexicographic order]先比较总次数，总次数相同时再比较字典序。两者都与乘法相容。字典序的良序性可用依次选择最小的第一、第二、……坐标证明；分次字典序则先选最小总次数，再在有限集合中取最小元。例如在 \(k[x,y]\) 中取 \(x\succ y\)，字典序有 \(x\succ y^{100}\)，分次字典序却有 \(y^{100}\succ x\)。这个差别将在消去变量时发挥作用。

对非零多项式 \(f\)，记最大的单项式为 \(\operatorname{LM}(f)\)，其系数为 \(\operatorname{LC}(f)\)，并记
\[
	\operatorname{LT}(f)=\operatorname{LC}(f)\operatorname{LM}(f).
\]
它们分别是\term[shouxiangdanshi]{首单项式}[leading monomial]、\term[shouxiangxishu]{首项系数}[leading coefficient]和\term[shouxiang]{首项}[leading term]。\symidx{\(\operatorname{LM}(f),\operatorname{LT}(f)\)：首单项式、首项}乘法相容性给出 \(\operatorname{LM}(fg)=\operatorname{LM}(f)\operatorname{LM}(g)\)：其他乘积项都严格小于两个首单项式之积，因而不能抵消它。

\begin{algorithm}[htbp]
	\caption{多元除法}\label{g5:alg:polynomial-division}
	\begin{algorithmsteps}{域 \(k\) 上的 \(f\in R\)、非零多项式的有序表 \(g_1,\ldots,g_s\) 及全局单项式序。}{多项式 \(q_1,\ldots,q_s,r\)，使 \(f=\sum_iq_ig_i+r\)，且 \(r\) 的任何单项式都不被任一 \(\operatorname{LM}(g_i)\) 整除。}
		\State 置 \(p=f\)，\(q_1=\cdots=q_s=r=0\)。
		\While{\(p\ne0\)}
			\If{存在 \(i\) 使 \(\operatorname{LM}(g_i)\mid\operatorname{LM}(p)\)}
				\State 取最小的这种 \(i\)，置 \(a=\operatorname{LT}(p)/\operatorname{LT}(g_i)\)。
				\State 以 \(q_i+a\) 替换 \(q_i\)，以 \(p-ag_i\) 替换 \(p\)。
			\Else
				\State 把 \(\operatorname{LT}(p)\) 加入 \(r\)，并从 \(p\) 中减去此项。
			\EndIf
		\EndWhile
		\State 返回 \(q_1,\ldots,q_s,r\)。
	\end{algorithmsteps}
\end{algorithm}

每次循环都消去了 \(p\) 的首项，且只留下更小的单项式，因此良序性保证算法终止。等式 \(f=\sum_iq_ig_i+p+r\) 在每一步保持成立，终止时就是所需分解。若 \(f\ne0\)，计算还保证每个非零 \(q_ig_i\) 的首单项式不超过 \(\operatorname{LM}(f)\)。这个上界将用来控制后面的抵消。

余式通常依赖除式的排列。例如在字典序 \(x\succ y\) 下，用 \(xy-1,y^2-1\) 除 \(xy^2\)，先用前者得到余式 \(y\)，先用后者则得到余式 \(x\)。两者之差 \(x-y\) 属于同一个理想，却不能再被这两个首单项式整除。因此，一般的生成元表还不能用非零余式判定理想成员关系。

\subsection{Buchberger 判据及终止性}
\label{g5:subsec:0104:02}
上面的障碍来自最高项的抵消：两个生成元的倍式相减后，可能产生新的首单项式，而旧的首单项式都除不尽它。Gröbner 基要求把这些新信息也纳入生成元。

\begin{definition}\label{g5:def:groebner-basis}
	设 \(I\subseteq R=k[x_1,\ldots,x_n]\) 为理想，并固定全局单项式序。由所有非零 \(f\in I\) 的 \(\operatorname{LM}(f)\) 生成的理想
	\[
		\operatorname{in}(I)=\bigl(\operatorname{LM}(f):0\ne f\in I\bigr)
	\]
	称为 \(I\) 的\term[chushilixiang]{初始理想}[initial ideal]。\symidx{\(\operatorname{in}(I)\)：初始理想}若有限集合 \(G=\{g_1,\ldots,g_s\}\subseteq I\setminus\{0\}\) 满足
	\[
		\operatorname{in}(I)=\bigl(\operatorname{LM}(g_1),\ldots,\operatorname{LM}(g_s)\bigr),
	\]
	则称 \(G\) 为 \(I\) 的\term[Groebnerji]{Gröbner 基}[Gröbner basis]。零理想允许取空基。
\end{definition}

由单项式生成的理想有一个直接判据：一个多项式属于它，当且仅当该多项式的每个单项式都被某个生成单项式整除。一个方向由展开乘积得到，另一个方向逐项写成生成元的倍式即可。因此，Gröbner 基条件等价于每个非零 \(f\in I\) 的首单项式都能被某个 \(\operatorname{LM}(g_i)\) 整除。

\begin{proposition}\label{g5:prop:groebner-remainder}
	设 \(k\) 为域，\(G\) 是 \(I\subseteq k[x_1,\ldots,x_n]\) 关于一个全局单项式序的 Gröbner 基。则 \(G\) 生成 \(I\)，每个多项式除以 \(G\) 的余式唯一，且 \(f\in I\) 当且仅当其余式为零。
\end{proposition}
\begin{proof}
	对 \(f\in I\) 作多元除法，所得余式 \(r=f-\sum_iq_ig_i\) 仍在 \(I\) 中。若 \(r\ne0\)，Gröbner 基条件要求它的首单项式被某个 \(\operatorname{LM}(g_i)\) 整除，这与余式的定义矛盾。故 \(r=0\)，同时得到生成性与成员判据。

	若同一个 \(f\) 有两个余式 \(r,r'\)，则 \(r-r'\in I\)。这个差中的每个单项式仍都不能被 \(G\) 的首单项式整除，所以刚才的论证给出 \(r-r'=0\)。
\end{proof}

要得到有限的判定条件，只需考察每对生成元的首项怎样相消。对非零 \(g,h\)，设 \(L\) 是 \(\operatorname{LM}(g)\) 与 \(\operatorname{LM}(h)\) 的最小公倍单项式，即每个变量的指数取两者中的较大值。定义其\term[Sduoxiangshi]{\(S\) 多项式}[S-polynomial]为
\[
	S(g,h)=\frac{L}{\operatorname{LT}(g)}g-\frac{L}{\operatorname{LT}(h)}h.
\]
两个最高项都是 \(L\)，因而被消去；若结果非零，它的首单项式严格小于 \(L\)。下面的 Buchberger 判据把所有可能的抵消归结为这些成对的抵消。\footnote{Bruno Buchberger（1942--），奥地利数学家，建立了 Gröbner 基的算法理论，并推动符号计算与自动推理的发展。}

\begin{theorem}[Buchberger 判据]\label{g5:thm:buchberger-criterion}
	\termidx[Buchbergerpanju]{Buchberger 判据}\foreignidx{Buchberger's criterion}设 \(k\) 为域，\(G=\{g_1,\ldots,g_s\}\) 为非零多项式的有限集合，并固定全局单项式序及除法时的排列。则 \(G\) 是 \((G)\) 的 Gröbner 基，当且仅当每个 \(S(g_i,g_j)\) 除以 \(G\) 的余式都是零。
\end{theorem}
\begin{proof}
	若 \(G\) 是 Gröbner 基，则 \(S(g_i,g_j)\in(G)\)，由\cref{g5:prop:groebner-remainder}，余式为零。

	反过来，把各 \(g_i\) 除以其首项系数，可设它们都首一；这不影响条件与结论。取 \(0\ne f\in(G)\)。把任意一个表达式 \(f=\sum_iq_ig_i\) 中的 \(q_i\) 展开成单项式之和，得到
	\[
		f=\sum_{\nu=1}^t c_\nu x^{\alpha_\nu}g_{i_\nu},\qquad c_\nu\in k\setminus\{0\}.
	\]
	在全部这种表示中，选取使最大单项式
	\[
		M=\max_\nu\bigl(x^{\alpha_\nu}\operatorname{LM}(g_{i_\nu})\bigr)
	\]
	最小的一种；良序性保证这种选择存在。若 \(M=\operatorname{LM}(f)\)，则某个 \(\operatorname{LM}(g_{i_\nu})\) 整除 \(\operatorname{LM}(f)\)，已得结论。我们证明 \(M\succ\operatorname{LM}(f)\) 不可能。

	把首单项式等于 \(M\) 的那些项写成
	\[
		H=\sum_{\nu=1}^r a_\nu\frac{M}{m_\nu}g_{j_\nu},
		\qquad m_\nu=\operatorname{LM}(g_{j_\nu}).
	\]
	由于 \(f\) 中没有单项式 \(M\)，而其他项都小于 \(M\)，有 \(\sum_\nu a_\nu=0\)。固定最后一个指标，将此等式用于消去 \(a_r\)，得到
	\[
		H=\sum_{\nu=1}^{r-1}a_\nu
		\left(\frac{M}{m_\nu}g_{j_\nu}-\frac{M}{m_r}g_{j_r}\right)
		=\sum_{\nu=1}^{r-1}a_\nu\frac{M}{L_\nu}S(g_{j_\nu},g_{j_r}),
	\]
	其中 \(L_\nu=\operatorname{lcm}(m_\nu,m_r)\)；相同生成元对应的差为零。由假设，每个非零 \(S(g_{j_\nu},g_{j_r})\) 经除法写成 \(\sum_i b_{\nu i}g_i\)，且除法的首单项式控制给出
	\[
		\operatorname{LM}(b_{\nu i}g_i)
		\preceq\operatorname{LM}\bigl(S(g_{j_\nu},g_{j_r})\bigr)\prec L_\nu
	\]
	（只对非零项写此式）。乘以 \(M/L_\nu\) 后，所有首单项式都严格小于 \(M\)。把这些表达式代回 \(H\)，就将原表示中全部最高项替换成了更低的项，从而得到最大单项式小于 \(M\) 的 \(f\) 的表示，违背最小选择。故 \(M=\operatorname{LM}(f)\)，判据得证。
\end{proof}

这个证明中特别需要“由除法得到的表示”：仅知道 \(S(g_i,g_j)\in(G)\) 毫无额外信息，因为它本来就是两个生成元倍式之差。除法保证表示中每个乘积的首单项式低于已经相消的最小公倍单项式，才使整个证明能够下降。

\begin{algorithm}[htbp]
	\caption{Buchberger 算法}\label{g5:alg:buchberger}
	\begin{algorithmsteps}{域 \(k\) 上的有限生成元表 \(f_1,\ldots,f_s\)，以及全局单项式序。}{理想 \((f_1,\ldots,f_s)\) 的一个 Gröbner 基。}
		\State 删去零生成元，将其余生成元首一化，组成 \(G\)；把 \(G\) 中全部不同元素对放入待处理表 \(P\)。
		\While{\(P\ne\varnothing\)}
			\State 从 \(P\) 取出并删除一对 \((g,h)\)，计算 \(S(g,h)\) 除以当前 \(G\) 的余式 \(r\)。
			\If{\(r\ne0\)}
				\State 将 \(r\) 首一化，把 \((r,g')\)（\(g'\in G\)）加入 \(P\)，再把 \(r\) 加入 \(G\)。
			\EndIf
		\EndWhile
		\State 返回 \(G\)。
	\end{algorithmsteps}
\end{algorithm}

\begin{proposition}\label{g5:prop:buchberger-termination}
	设 \(k\) 为域，并使用精确的域运算及判零。\cref{g5:alg:buchberger}在有限步后终止，返回输入理想关于指定单项式序的 Gröbner 基。
\end{proposition}
\begin{proof}
	加入的余式始终属于当前生成的理想，故生成理想不变。每次加入非零 \(r\)，它的首单项式都不属于旧的首单项式理想，因此后者严格增大。\cref{g5:thm:hilbert-basis}保证 \(k[x_1,\ldots,x_n]\) 的理想升链最终稳定，故只可能加入有限多个新元素。每次加入只增添有限多个待处理对，因而全部计算有限步结束。

	对任一已处理的对，当时的除法给出其 \(S\) 多项式关于当时生成元的表示；若余式非零，加入该余式后也得到表示。所有表示项的首单项式都低于相应最小公倍单项式。此后增添生成元不会破坏这个表示。因此最终 \(G\) 的每对元素都具有 Buchberger 判据证明所需的这种较低表示，该证明直接给出最终 \(G\) 为 Gröbner 基。
\end{proof}

终止性给出有限算法，却没有给出一个小的运行次数上界。变量次序和输入形式会显著影响中间多项式的大小；这里先保留可以逐步核验的基本算法。

\begin{example}\label{g5:exm:buchberger-three-points}
	在任意域 \(k\) 上取 \(I=(xy-1,x-y^2)\subseteq k[x,y]\)，使用字典序 \(x\succ y\)。令 \(g_1=xy-1\)、\(g_2=x-y^2\)，则
	\[
		S(g_1,g_2)=g_1-yg_2=y^3-1.
	\]
	这个余式不能被 \(xy\) 或 \(x\) 除尽，须加入 \(g_3=y^3-1\)。由于 \(g_1=yg_2+g_3\)，理想也由 \(g_2,g_3\) 生成。对这两个生成元，
	\[
		S(g_2,g_3)=y^3g_2-xg_3=x-y^5
		\xrightarrow{-g_2}y^2-y^5
		\xrightarrow{+y^2g_3}0.
	\]
	所以 \(\{x-y^2,y^3-1\}\) 是 \(I\) 的 Gröbner 基。这里多出的单变量方程已经把求解次序显露出来：先求 \(y^3=1\)，再令 \(x=y^2\)。
\end{example}

\subsection{消去定理与投影闭包}
\label{g5:subsec:0104:03}
若想消去 \(x_1,\ldots,x_r\)，应该让含这些变量的单项式大于所有其余单项式。字典序 \(x_1\succ\cdots\succ x_n\) 正好做到这一点。对理想 \(I\subseteq k[x_1,\ldots,x_n]\)，其\term[xiaoqulixiang]{消去理想}[elimination ideal]是
\[
	I_r=I\cap k[x_{r+1},\ldots,x_n].
\]
它收集原方程的全部多项式后果中已经没有前 \(r\) 个变量的那些方程。

\begin{theorem}[消去定理]\label{g5:thm:elimination}
	\termidx[xiaoqudingli]{消去定理}\foreignidx{elimination theorem}设 \(k\) 为域，\(I\subseteq k[x_1,\ldots,x_n]\)，\(0\le r<n\)。若 \(G\) 为 \(I\) 关于字典序 \(x_1\succ\cdots\succ x_n\) 的 Gröbner 基，则
	\[
		G_r=G\cap k[x_{r+1},\ldots,x_n]
	\]
	是 \(I_r\) 关于限制次序的 Gröbner 基。
\end{theorem}
\begin{proof}
	取 \(0\ne f\in I_r\)。存在 \(g\in G\) 使 \(\operatorname{LM}(g)\mid\operatorname{LM}(f)\)，因此 \(\operatorname{LM}(g)\) 不含 \(x_1,\ldots,x_r\)。若 \(g\) 的其他项含其中某个变量，按字典序它会大于所有只含后面变量的单项式，特别是大于 \(\operatorname{LM}(g)\)，矛盾。因此 \(g\in G_r\)，这正是 Gröbner 基判据。
\end{proof}

从几何上说，消去得到的方程必在投影像上成立。它们究竟刻画实际像还是闭包，需要另外论证。对 \(k^{r}\times k^{n-r}\) 的坐标投影记作 \(\pi\)。

\begin{theorem}\label{g5:thm:projection-closure}
	设 \(k\) 为代数闭域，\(I\subseteq k[x_1,\ldots,x_n]\) 为任意理想，\(I_r=I\cap k[x_{r+1},\ldots,x_n]\)。则
	\[
		\overline{\pi(V(I))}=V(I_r),
	\]
	其中闭包取在 \(\mathbb A_k^{n-r}\) 的 Zariski 拓扑中。
\end{theorem}
\begin{proof}
	关键是计算投影像上所有消失的多项式。一个 \(h\in k[x_{r+1},\ldots,x_n]\) 在 \(\pi(V(I))\) 上恒为零，当且仅当把它视为 \(n\) 元多项式时在 \(V(I)\) 上恒为零。强零点定理\cref{g5:thm:nullstellensatz}因而给出
	\[
		I\bigl(\pi(V(I))\bigr)
		=\sqrt I\cap k[x_{r+1},\ldots,x_n]
		=\sqrt{I_r}.
	\]
	最后一个等号是因为 \(h^N\in I\) 且 \(h\) 只含后面变量，等价于 \(h^N\in I_r\)。任意集合 \(S\) 的闭包都是 \(V(I(S))\)，取零点即得结论。
\end{proof}

\begin{example}\label{g5:exm:projection-not-closed}
	设 \(k\) 代数闭，取双曲线 \(X=V(xy-1)\subseteq\mathbb A_k^2\)，向 \(y\) 轴投影。实际像是 \(k\setminus\{0\}\)，因为 \(y\ne0\) 时才能取 \(x=y^{-1}\)。而字典序 \(x\succ y\) 下 \(\{xy-1\}\) 已是 Gröbner 基，消去理想为零，故投影闭包为整条 \(y\) 轴。新增的原点满足全部消去方程，却没有原方程组的解位于它上方。
\end{example}

参数式的隐式化可以统一按图像处理。例如参数式 \((x,y)=(t^2,t^3)\) 对应图像理想 \((x-t^2,y-t^3)\subseteq k[t,x,y]\)。消去 \(t\) 就得到参数像闭包的方程。先取图像再投影，使“求出关系式”与“关系式的每个解是否都来自参数”成为两个可以分别检查的问题。消去的几何说明见\cite[第3章第1--2节]{CoxLittleOShea2015}。

\subsection{饱和和射影分量}
\label{g5:subsec:0104:04}
把方程限制在 \(h\ne0\) 的区域时，\(hf=0\) 与 \(f=0\) 具有相同意义。返回原多项式环时，必须记录所有通过消去 \(h\) 的幂才能得到的后果。

\begin{definition}\label{g5:def:saturation}
	设 \(R=k[x_1,\ldots,x_n]\)、\(I\subseteq R\)，并取 \(h\in R\)。理想
	\[
		I:h^\infty=\{f\in R:\text{存在 }N\ge0\text{ 使 }h^Nf\in I\}
	\]
	称为 \(I\) 关于 \(h\) 的\term[baohe]{饱和}[saturation]。\symidx{\(I:h^\infty\)：关于 \(h\) 的饱和}它是升链 \(I:h^N\) 的并，因而是理想；多项式环的 Noether 性保证这个并等于某个 \(I:h^N\)。
\end{definition}

\begin{proposition}\label{g5:prop:saturation-elimination}
	设 \(k\) 为域，\(I\subseteq R=k[x_1,\ldots,x_n]\)，\(h\in R\)，并引入新变量 \(u\)。则
	\[
		I:h^\infty=(I,1-uh)\cap R.
	\]
	若 \(k\) 代数闭，则还有
	\[
		V(I:h^\infty)=\overline{V(I)\setminus V(h)}.
	\]
\end{proposition}
\begin{proof}
	若 \(h^Nf\in I\)，则 \((uh)^Nf\in I R[u]\)，而 \(f-(uh)^Nf\) 是 \(1-uh\) 的倍式，故 \(f\in(I,1-uh)\cap R\)。反之，若 \(h\ne0\)，将 \(f\in(I,1-uh)\cap R\) 的理想表达式代入 \(u=1/h\)，在有理函数域 \(k(x_1,\ldots,x_n)\) 中得到 \(f=\sum_i a_i(x,1/h)f_i\)，其中 \(f_i\in I\)。乘以足够高的 \(h\) 的幂即得 \(h^Nf\in I\)。若 \(h=0\)，两边都等于 \(R\)，结论同样成立。

	在代数闭情形，\(V(I,1-uh)\) 向原坐标空间的投影恰为 \(V(I)\setminus V(h)\)：每个这样的点都有唯一的 \(u=1/h\)。对这个投影应用\cref{g5:thm:projection-closure}，即得第二个等式。
\end{proof}

因此，只要采用先消去 \(u\) 的次序，就能用 Gröbner 基计算饱和。几何上它保留 \(h\ne0\) 部分的闭包；例如 \((xy):x^\infty=(y)\)，从两条坐标轴中去除了完全位于 \(x=0\) 上的那一支。关于饱和的进一步讨论见\cite[第4章第4节]{CoxLittleOShea2015}。

射影闭包需要同一种操作。以 \([x_0:x_1:\cdots:x_n]\) 为射影坐标，把仿射空间识别为 \(x_0\ne0\) 的标准图。对次数为 \(d\) 的非零多项式 \(f\)，记
\[
	 f^{\mathrm h}(x_0,\ldots,x_n)=x_0^d f(x_1/x_0,\ldots,x_n/x_0)
\]
为其齐次化，且规定 \(0^{\mathrm h}=0\)。对理想 \(I\)，以全部 \(f^{\mathrm h}\)（\(f\in I\)）生成的齐次理想记为 \(I^{\mathrm h}\)。\symidx{\(I^{\mathrm h}\)：理想的齐次化}关键在于“全部”：对一份任意生成元逐个齐次化，所得理想还可能缺少由最高次抵消产生的方程。

\begin{proposition}\label{g5:prop:projective-homogenization}
	设 \(k\) 为域，\(I=(f_1,\ldots,f_s)\subseteq k[x_1,\ldots,x_n]\)，且
	\[
		J=(f_1^{\mathrm h},\ldots,f_s^{\mathrm h})\subseteq k[x_0,\ldots,x_n].
	\]
	则 \(I^{\mathrm h}=J:x_0^\infty\)。若 \(k\) 代数闭、\(X\subseteq\mathbb A_k^n\) 为代数集并取 \(I=I(X)\)，则 \(X\) 在 \(\mathbb P_k^n\) 中的闭包是 \(V_+(I^{\mathrm h})\)。
\end{proposition}
\begin{proof}
	先说明 \(I^{\mathrm h}\subseteq J:x_0^\infty\)。对 \(f=\sum_i a_if_i\in I\)，取 \(D\) 不小于 \(\deg f\) 及各非零 \(a_if_i\) 的次数，将等式各项齐次化到同一次数 \(D\)，便得
	\[
		x_0^{D-\deg f}f^{\mathrm h}
		=\sum_i x_0^{D-\deg a_i-\deg f_i}a_i^{\mathrm h}f_i^{\mathrm h}\in J.
	\]
	故每个 \(f^{\mathrm h}\) 都属于该饱和。

	反向包含只需检查齐次多项式，因为 \(J\) 的齐次性保证 \(J:x_0^\infty\) 也齐次。若次数为 \(q\) 的非零齐次多项式 \(F\) 满足 \(x_0^NF\in J\)，令 \(x_0=1\) 得到 \(f=F(1,x_1,\ldots,x_n)\in I\)。又有 \(F=x_0^{q-\deg f}f^{\mathrm h}\)，故 \(F\in I^{\mathrm h}\)。

	最后设 \(I=I(X)\)。任意齐次多项式 \(F\) 在嵌入的 \(X\) 上消失，当且仅当 \(F(1,x)\in I(X)\)，这又等价于 \(F\in I^{\mathrm h}\)。因此 \(V_+(I^{\mathrm h})\) 是包含 \(X\) 的全部射影闭集的交，也就是射影闭包。
\end{proof}

\begin{example}\label{g5:exm:spurious-infinity-component}
	在代数闭域 \(k\) 上，曲线 \((x,y,z)=(t,t^2,t^3)\) 的仿射理想为 \(I=(y-x^2,z-x^3)\)：商环中的 \(y,z\) 分别由 \(x^2,x^3\) 代替，商环就是 \(k[x]\)。在射影坐标 \([w:x:y:z]\) 中，逐个齐次化得到
	\[
		J=(wy-x^2,w^2z-x^3).
	\]
	令 \(w=0\)，两式只要求 \(x=0\)，因而 \(V_+(J)\) 包含整条直线 \([0:0:y:z]\)。但仿射理想中还有
	\[
		xy-z=x(y-x^2)-(z-x^3),\qquad
		y^2-xz=(y+x^2)(y-x^2)-x(z-x^3).
	\]
	它们的齐次化要求 \(xy-wz=0\)、\(y^2-xz=0\)。三个二次式
	\[
		wy-x^2=0,\qquad wz-xy=0,\qquad xz-y^2=0
	\]
	在 \(w\ne0\) 上恢复原曲线；在 \(w=0\) 上则迫使 \(x=y=0\)，只留下 \([0:0:0:1]\)。这个点确实在闭包中：参数式
	\[
		[s:t]\longmapsto[s^3:s^2t:st^2:t^3]
	\]
	在 \(s\ne0\) 上给出原曲线，且在 \(t\ne0\) 的图中变为 \(u\mapsto[u^3:u^2:u:1]\)，\(u=s/t\)，原点是 \(u\ne0\) 部分的 Zariski 闭包中的点。于是上述三个二次式在集合意义上恰好刻画射影闭包。逐个齐次化产生的无穷远直线，是缺少饱和造成的多余分量。
\end{example}

这里的饱和是关于所选仿射图的坐标 \(x_0\)；它专门去掉完全落在 \(x_0=0\) 上的多余部分。有关射影闭包的更多计算见\cite[第8章第4节]{CoxLittleOShea2015}。

\subsection{零维商代数与重数核验}
\label{g5:subsec:0104:05}
求出零点以后，商代数还保留了不同方程在同一点处相遇的次数。最简单的例子是 \(k[t]/(t^2)\)：它只有原点这一个零点，但 \(1,t\) 是两个线性无关的剩余类。Gröbner 基把这种维数变成可以直接数出的单项式个数。

\begin{proposition}\label{g5:prop:standard-monomials}
	设 \(k\) 为域，\(I\subseteq R=k[x_1,\ldots,x_n]\)，并固定一个全局单项式序。所有不属于 \(\operatorname{in}(I)\) 的单项式称为\term[biaozhundanxiangshi]{标准单项式}[standard monomial]，它们在 \(R/I\) 中的剩余类构成一个 \(k\) 向量空间基。特别地，\(\dim_k(R/I)<\infty\) 当且仅当对每个 \(i\) 都有某个 \(N_i\ge0\) 使 \(x_i^{N_i}\in\operatorname{in}(I)\)。
\end{proposition}
\begin{proof}
	由 Buchberger 算法取得一个 Gröbner 基。多元除法把每个多项式模 \(I\) 化为标准单项式的线性组合，故这些类张成商空间。若一个非零的这种线性组合属于 \(I\)，其首单项式就必须属于 \(\operatorname{in}(I)\)，矛盾，因此这些类线性无关。

	若某个变量的所有纯幂都是标准单项式，商空间就有无限多个基向量。反之，若各 \(x_i^{N_i}\in\operatorname{in}(I)\)，则每个标准单项式的第 \(i\) 个指数小于 \(N_i\)，只有有限种可能。\(I=R\) 时标准单项式集合为空，商空间维数为零。
\end{proof}

当 \(k\) 代数闭时，这个有限维条件恰好等价于 \(V(I)\) 为有限集合。为核对两种有限性的关系，若 \(V(I)=\{a^{(1)},\ldots,a^{(s)}\}\) 非空，则对每个坐标 \(i\)，多项式 \(p_i(T)=\prod_{j=1}^s(T-a_i^{(j)})\) 满足 \(p_i(x_i)\in\sqrt I\)。由零点定理，某个幂 \(p_i(x_i)^{N_i}\) 属于 \(I\)；它是关于 \(x_i\) 的首一多项式，故每个坐标的高次幂都能约为较低次幂，商代数有限维。空集情形由弱零点定理给出 \(I=R\)。反过来，若商代数有限维，则 \(1,x_i,x_i^2,\ldots\) 的类线性相关，给出非零单变量关系 \(q_i(x_i)\in I\)。每个零点的第 \(i\) 个坐标只能是 \(q_i\) 的有限多个根之一，故零点总数有限。

本节将这种非零有限维商代数称为零维商代数。它的\term[changdu]{长度}[length]在代数闭基域上等于 \(\dim_k(R/I)\)；下面把这个数分配到各个点，说明它为什么是重数之和。

\begin{proposition}\label{g5:prop:finite-algebra-multiplicities}
	设 \(k\) 代数闭，\(A=k[x_1,\ldots,x_n]/I\) 是非零有限维 \(k\) 代数，\(V(I)=\{a^{(1)},\ldots,a^{(s)}\}\)。记 \(\mathfrak m_j\subseteq A\) 为在 \(a^{(j)}\) 处取值为零的剩余类组成的极大理想。则对充分大的整数 \(N\)，自然映射给出分解
	\[
		A\cong\prod_{j=1}^s A_j,\qquad A_j=A/\mathfrak m_j^N,
		\qquad \dim_kA=\sum_{j=1}^s\dim_kA_j.
	\]
	每个 \(A_j\) 只有一个极大理想，剩余域为 \(k\)。数 \(\dim_kA_j\) 是方程组 \(I\) 在点 \(a^{(j)}\) 处的\term[jubuzhongshu]{局部重数}[local multiplicity]。\(I\) 为根式理想，当且仅当每个局部重数均为 \(1\)。
\end{proposition}
\begin{proof}
	先把各点分开。不同点的求值理想两两互素：若第 \(i\) 个坐标不同，则 \(x_i-a_i^{(j)}\) 与 \(x_i-a_i^{(\ell)}\) 之差为非零常数，故 \(\mathfrak m_j+\mathfrak m_\ell=A\)。由零点定理，\(\mathfrak n=\bigcap_j\mathfrak m_j\) 恰是 \(A\) 中全部幂零元素的理想。它由有限多个元素生成，例如取它的一个 \(k\) 向量空间基 \(b_1,\ldots,b_q\)。若 \(b_i^{e_i}=0\)，则任意 \(N>\sum_i(e_i-1)\) 个生成元的乘积总包含某个 \(b_i^{e_i}\)，所以 \(\mathfrak n^N=0\)。

	互素理想 \(J,K\) 满足 \(J\cap K=JK\)：若 \(u+v=1\)，\(u\in J,v\in K\)，则对 \(a\in J\cap K\)，有 \(a=au+av\in JK\)。它们的 \(N\) 次幂也互素，因为展开 \((u+v)^{2N-1}=1\)，每项至少属于 \(J^N\) 或 \(K^N\)。对有限族反复应用，得到
	\[
		\bigcap_j\mathfrak m_j^N
		=\prod_j\mathfrak m_j^N
		=\left(\prod_j\mathfrak m_j\right)^N
		=\mathfrak n^N=0.
	\]
	\term[zhongguoshengyudingli]{中国剩余定理}[Chinese remainder theorem]于是给出所述分解。这里所用的定理也可直接核对：对两个互素理想，映射 \(A\to A/J\times A/K\) 的核为交，而 \(va+ub\) 同时代表指定的模 \(J\) 类 \(a\) 和模 \(K\) 类 \(b\)；有限多个理想逐次处理即可。

	在 \(A_j\) 中，\(\mathfrak m_j/\mathfrak m_j^N\) 是幂零理想，商去它得到 \(k\)。任何极大理想都包含幂零元素，因此这是 \(A_j\) 的唯一极大理想。若 \(I\) 根式，则 \(\mathfrak n=0\)，直接对 \(\mathfrak m_j\) 应用同一分解得到 \(A\cong k^s\)，各因子都是 \(k\)。反之，若各因子维数为 \(1\)，它们都等于 \(k\)，乘积没有非零幂零元素，故 \(I\) 根式。
\end{proof}

上述分解与充分大的 \(N\) 无关：在第 \(j\) 个因子中，\(\mathfrak m_j\) 的高次幂为零，在其他因子中它等于整个环；后一点由互素性及这些因子的幂零极大理想可见。因此继续增大 \(N\) 得到同一求值点对应的因子。其作为自身模的任何简单商都是 \(k\)，每一步下降都使向量空间维数至少减一，故有限维性保证存在组成列，而组成因子个数等于 \(k\) 维数。这也解释了这里“长度”与通常模长度术语的一致性。

\begin{example}\label{g5:exm:length-three-points}
	对\cref{g5:exm:buchberger-three-points}中的理想，标准单项式为 \(1,y,y^2\)，故商代数长度为 \(3\)，并有
	\[
		k[x,y]/(xy-1,x-y^2)\cong k[y]/(y^3-1).
	\]
	若 \(k\) 代数闭且特征不为 \(3\)，\(y^3-1\) 的导数 \(3y^2\) 在其根处不为零，故三个根不同，每点重数为 \(1\)。若特征为 \(3\)，则 \(y^3-1=(y-1)^3\)，唯一零点为 \((1,1)\)，局部重数为 \(3\)。同一份标准单项式计数在两种特征下都给出 \(3\)，但这 \(3\) 的几何分配不同。
\end{example}

\begin{example}\label{g5:exm:tangent-intersection-length}
	在任意代数闭域 \(k\) 上，抛物线 \(y=x^2\) 与直线 \(y=0\) 的共同零点只有原点，而
	\[
		k[x,y]/(y-x^2,y)=k[x,y]/(x^2,y)\cong k[x]/(x^2)
	\]
	的基为 \(1,x\)，长度为 \(2\)。若换成直线 \(x=0\)，商环变为 \(k\)，长度为 \(1\)。前一个商中的 \(x\) 非零而平方为零，它记录了取根式后丢掉的相切信息。这给出有限相交点的代数计数；一般相交理论将在后文讨论其在族中的行为。
\end{example}

商代数中的加法、乘法也可以用余式完成。选定标准单项式基后，逐个把 \(x_i\) 乘上基向量，再除以 Gröbner 基，就得到“乘以 \(x_i\)”的矩阵。这样，方程组的有限商代数可以作为一套具体的有限维线性运算保存，而不必只保存零点列表。商环算法的进一步例子见\cite[第5章第3节]{CoxLittleOShea2015}。
% 纲要条目排为编号小节标题，并进入目录与书签；说明段仍保留于纲要注释中。

% !TeX root = ../../Book5.tex
% 纲要标题：### 1.5 模型提出的结构问题
% 纲要位置：几何学计划纲要.md，第 5062 行。
% 正文状态：初稿；尚未正式审核。

% 纲要内容开始
% 
% * 结点、尖点与非约化点
% * 射影闭包中的无穷远分支
% * 为什么必须保留局部环和概形纤维
% 
% ---
% 
% 纲要内容结束

\section{模型提出的结构问题}
\label{g5:sec:0105}

消去给出映射像的闭包，参数化给出函数域，射影闭包补入无穷远点。这几种构造保留的信息不同。下面比较两个有理三次曲线，并考察两点重合时的方程；计算将说明，研究一个点附近的函数及保留方程中的幂次，怎样补充点集和函数域的描述。

\subsection{结点、尖点与非约化点}
\label{g5:subsec:0105:01}

先设 \(k\) 为代数闭域且 \(\operatorname{char}k\ne2\)，考虑
\[
	C_{\mathrm n}=V(y^2-x^2(x+1)),\qquad
	C_{\mathrm c}=V(y^2-x^3).
\]
第二条是\cref{g5:exm:cusp-birational}的尖点曲线。第一条曲线可用过原点的直线 \(y=tx\) 求参数：代入并取 \(x\ne0\)，得到
\[
	x=t^2-1,\qquad y=t(t^2-1).
\]
由此定义正则映射
\[
	\nu_{\mathrm n}:\mathbb A_k^1\longrightarrow C_{\mathrm n},
	\qquad t\longmapsto(t^2-1,t(t^2-1)).
\]
在 \(x\ne0\) 处，逆参数是 \(t=y/x\)；原点则有两个参数原像 \(t=1,-1\)。其坐标环确实为 \(k[t^2-1,t(t^2-1)]\)：将多项式除以关于 \(y\) 首一的方程后，余式为 \(a(x)+yb(x)\)。代入参数，第一项是 \(t\) 的偶多项式，第二项是奇多项式；两者之和为零迫使两者都为零，并进而迫使 \(a=b=0\)。所以代入同态的核恰好为所给主理想，它为素理想，曲线因而不可约。

\ProseParagraphBreak

在原点平移展开方程，最低次非零齐次项分别为
\[
	y^2-x^2=(y-x)(y+x),\qquad y^2.
\]
把最低次非零齐次方程给出的锥连同因子重数，称为这个平面曲线模型的\term[qiezhui]{切锥}[tangent cone]，第一式给两个不同方向，第二式给同一方向的二重方程。这里的两个方向还可以从参数的局部变化看出：在 \(t=1\) 附近令 \(t=1+s\)，则 \(x=2s+s^2\)、\(y=2s+3s^2+s^3\)；在 \(t=-1\) 附近令 \(t=-1+s\)，则 \(x=-2s+s^2\)、\(y=2s-3s^2+s^3\)。它们分别有方向 \(y=x\)、\(y=-x\)。这种模型的原点称为\term[jiedian]{结点}[node]；\(C_{\mathrm c}\) 的原点称为\term[jiandian]{尖点}[cusp]。参数化的比较也可参见\cite[第I章第4节，例4.9.1及习题4.10]{Hartshorne1977}。

\ProseParagraphBreak

在 \(k=\mathbb C\) 时，\cref{g5:fig:node-cusp}画出实参数所描绘的实点截面。结点的两个参数值对应同一个平面点，尖点则由一个参数值产生；完整的局部区分由随后的方程计算给出。
\begin{figure}[htbp]
	\centering
	\input{Figures/Chapter01SingularModels}
	\caption[结点与尖点的参数模型]{结点与尖点的实点截面。左图参数为 \((t^2-1,t(t^2-1))\)，右图参数为 \((t^2,t^3)\)。}
	\label{g5:fig:node-cusp}
\end{figure}
全局不可约曲线也能在一个点有两个局部分支。为把上面的说法落实到方程，可以允许形式幂级数，即按总次数逐项运算、不讨论数值收敛的级数。在 \(k[[x]]\) 中，方程 \(s(x)^2=1+x\)、\(s(0)=1\) 有唯一解：若 \(s=1+\sum_{r\ge1}c_rx^r\)，逐次比较系数给出
\[
	2c_r=\delta_{r1}-\sum_{i=1}^{r-1}c_i c_{r-i}.
\]
特征不为二保证每步可解。因此在 \(k[[x,y]]\) 中，
\[
	y^2-x^2(x+1)=(y-xs(x))(y+xs(x)).
\]
两个因子分别给出形式图像 \(y=\pm xs(x)\)，每个图像的商环同构于 \(k[[x]]\)，两个因子不是彼此的单位倍。这具体显示了两条形式分支，尽管多项式坐标环本身是整环。

\ProseParagraphBreak

尖点的形式方程则仍不可约。代入 \(x=t^2,y=t^3\) 将
\(k[[x,y]]/(y^2-x^3)\) 嵌入 \(k[[t]]\)：利用 \(y^2=x^3\)，每个类都可唯一写成 \(a(x)+yb(x)\)，其中 \(a,b\in k[[x]]\)。这一归约对幂级数有意义，因为把 \(y^{2r}\) 换成 \(x^{3r}\) 后，对任一固定次数只有有限项贡献。代入后偶次幂与大于等于三的奇次幂不相混，故核为零。商环是整环，因而形式方程有一条分支。这里把形式方程的不可约因子所确定的部分称为形式分支；一般局部环与其完备化的关系将在后续章节系统建立。

\begin{example}\label{g5:exm:nonreduced-point}
	在任意代数闭域 \(k\) 上，方程 \(x=0\) 与 \(x^2=0\) 有同一零点集，商环却分别为
	\[
		k[x]/(x)\simeq k,\qquad
		k[x]/(x^2)\simeq k[\varepsilon]/(\varepsilon^2).
	\]
	后一个环中 \(\varepsilon\ne0\) 而 \(\varepsilon^2=0\)。环中某个正整数次幂为零的元素称为幂零元；没有非零幂零元的环称为约化环。前者是一个约化点的函数环，后者是一个\term[feiyuehuadian]{非约化点}[nonreduced point]的代数模型。作为 \(k\) 向量空间，它们的维数分别为一和二，后一空间的基为 \(1,\varepsilon\)。

	每个到 \(k\) 的 \(k\) 代数同态都把 \(\varepsilon\) 送到零，因而普通 \(k\) 值点看不到这个区别。若允许取值于 \(D=k[\varepsilon]/(\varepsilon^2)\)，则把坐标 \(x\) 送到 \(a\varepsilon\)（\(a\in k\)）都满足 \(x^2=0\)，却只有 \(a=0\) 满足 \(x=0\)。保留这些测试环中的解，就能检测点集中消失的信息。
\end{example}

\subsection{射影闭包中的无穷远分支}
\label{g5:subsec:0105:02}

在 \(\mathbb P_k^2\) 中使用坐标 \([X:Y:Z]\)，把 \(Z\ne0\) 作为原仿射平面。前面两条曲线的射影闭包为
\[
	\overline C_{\mathrm n}=V_+(Y^2Z-X^2(X+Z)),\qquad
	\overline C_{\mathrm c}=V_+(Y^2Z-X^3).
\]
两者在 \(Z=0\) 上都只有 \(P_\infty=[0:1:0]\)。在 \(Y\ne0\) 的图上令 \(u=X/Y,v=Z/Y\)，方程变成
\[
	v-u^3-u^2v=0,\qquad v-u^3=0.
\]
结点模型中 \(1-u^2\) 在 \(P_\infty\) 附近不为零，所以 \(v=u^3/(1-u^2)\)；尖点模型中 \(v=u^3\)。每个模型在无穷远点附近都由一个正则参数 \(u\) 描述，只有一条这样的局部分支。

\ProseParagraphBreak

这也可由射影参数直接核验。将仿射参数写成 \(t=T/S\)，得到
\[
\begin{aligned}
	\nu_{\mathrm n}([S:T])
		&=[S(T^2-S^2):T(T^2-S^2):S^3],\\
	\nu_{\mathrm c}([S:T])
		&=[ST^2:T^3:S^3].
\end{aligned}
\]
两组坐标均为同次齐次多项式，而且在 \(\mathbb P^1\) 上没有公共零点：\(S\ne0\) 时最后一个坐标非零，\(S=0\) 时第二个坐标非零。代入曲线方程检验像确实落在闭包内，且 \(S=0\) 唯一地映到 \(P_\infty\)。这些射影像包含原仿射曲线及其无穷远点，故给出了闭包的全部点。

\begin{example}\label{g5:exm:two-infinity-branches}
	双曲线 \(xy=1\) 的射影闭包为 \(XY=Z^2\)，有两个无穷远点 \([1:0:0]\)、\([0:1:0]\)。在第一点附近用 \(u=Y/X,v=Z/X\)，方程为 \(u=v^2\)；另一点同理。两个点附近各有一条参数分支。参数 \(x=t,y=t^{-1}\) 中的 \(t=0\) 与 \(t=\infty\) 在射影闭包中成为两个不同点。
\end{example}

无穷远点的个数与局部分支的个数须分别判断。把一个结点移到无穷远处即可看到差别：在 \(\overline C_{\mathrm n}\) 中改以直线 \(X=0\) 为无穷远直线，原结点 \([0:0:1]\) 就成为无穷远点，仍带有两条形式分支。由此，齐次化后列出无穷远点只是第一步；还需要在各点的仿射图内查看方程或参数。

\subsection{为什么必须保留局部环和概形纤维}
\label{g5:subsec:0105:03}

一个有理函数能否在指定点延拓，取决于该点附近能否找到不消失的分母。为了把所有这样的函数放在同一对象中，设 \(X\) 为仿射代数集，\(p\in X\)。把包含 \(p\) 的开邻域 \(U\) 上的正则函数 \(f\) 与另一个 \((V,g)\) 视为相同，只要两者在某个更小的 \(p\) 的邻域上相等。所得等价类称为函数在 \(p\) 的\term[hanshuya]{芽}[germ]，其环记为
\[
	\mathcal O_{X,p}=\{\text{在 }p\text{ 附近的正则函数芽}\}.
\]
\symidx{\(\mathcal O_{X,p}\)：\(X\) 在 \(p\) 处的局部环}

\begin{proposition}\label{g5:prop:model-local-ring}
	设 \(k\) 为代数闭域，\(X\) 为仿射代数集，\(p\in X\)，\(A=k[X]\)。每个函数芽都可写成 \(a/b\)，其中 \(a,b\in A\)、\(b(p)\ne0\)。两个这样的分式代表同一芽，当且仅当存在 \(h\in A\)、\(h(p)\ne0\)，使
	\[
		h(ad-bc)=0
	\]
	成立，其中所比较的分式为 \(a/b\) 与 \(c/d\)。此环有唯一极大理想
	\[
		\mathfrak m_p=\{f\in\mathcal O_{X,p}:f(p)=0\},
		\qquad \mathcal O_{X,p}/\mathfrak m_p\simeq k.
	\]
\end{proposition}
\begin{proof}
	分式表达来自正则函数的定义。若两分式在一个邻域内相等，取包含 \(p\) 并位于该邻域内的主开集 \(D(h)\)。函数 \(ad-bc\) 在 \(D(h)\) 上为零，所以 \(h(ad-bc)\) 在整个 \(X\) 上为零。反之，上述等式在 \(D(hbd)\) 上可以除以不消失的分母，给出分式相等。

	求值到 \(k\) 是满同态，其核即 \(\mathfrak m_p\)，因此它是极大理想。若 \(f(p)\ne0\)，在较小邻域上 \(1/f\) 正则，所以 \(f\) 是单位；若 \(f(p)=0\)，求值表明它不可能为单位。因此所有非单位恰组成 \(\mathfrak m_p\)，任何真理想都包含在其中，极大理想唯一。
\end{proof}

只有一个极大理想的交换环称为\term[jubuhuan]{局部环}[local ring]。命题给出的分式环就是记号 \(A_{\mathfrak m}\) 所表示的局部化，其中 \(\mathfrak m=\{a\in A:a(p)=0\}\)。\cref{g5:exm:cusp-local-obstruction}证明了尖点处的局部环不含函数 \(t=y/x\)，而参数直线原点的局部环含有它；局部环由此记录了双有理对应中被遗漏的差别。

\ProseParagraphBreak

一阶测试可以同时检测局部函数及非约化信息。设 \(A=k[x_1,\ldots,x_n]/I\)，并固定一个零点 \(p=(p_1,\ldots,p_n)\)。到双数环 \(D=k[\varepsilon]/(\varepsilon^2)\) 的同态，若模 \(\varepsilon\) 后等于在 \(p\) 的求值，必把坐标送到
\[
	x_i\longmapsto p_i+v_i\varepsilon .
\]
对任意 \(F\in I\)，直接展开每个单项式得到
\[
	F(p+v\varepsilon)=F(p)+
	\varepsilon\sum_{i=1}^n\frac{\partial F}{\partial x_i}(p)v_i.
\]
因而这个赋值能通过商环，当且仅当每个 \(F\in I\) 的一次项都为零。这一解空间是该方程模型在 \(p\) 的\term[Zariskiqiekongjian]{Zariski 切空间}[Zariski tangent space]。计算时只需检查理想生成元：对 \(F=\sum_j a_jF_j\)，因 \(F_j(p)=0\)，一次项就是各 \(F_j\) 一次项的 \(a_j(p)\) 倍之和。

\ProseParagraphBreak

在结点和尖点原点，两方程的一次项都为零，故两者的切空间都是 \(k^2\)；切锥又分别保留了两个方向和一个二重方向。对约化点 \(A=k[x]/(x)\)，一次项条件为 \(v=0\)；对 \(A=k[x]/(x^2)\)，一次项没有约束，得到一维切空间。可见一阶线性化、最低次方程和完整局部环各自保留不同层次的信息。

\begin{example}\label{g5:exm:double-cover-fibres}
	设 \(k\) 为代数闭域且 \(\operatorname{char}k\ne2\)。考虑映射
	\[
		\pi:\mathbb A_k^1\longrightarrow\mathbb A_k^1,\qquad t\longmapsto s=t^2.
	\]
	在 \(s=a\) 上的集合论纤维是方程 \(t^2=a\) 的根集。若直接在坐标环中加入基底方程 \(s-a=0\)，得到纤维代数
	\[
		B_a=k[t]/(t^2-a).
	\]
	当 \(a=b^2\ne0\) 时，\(b\ne-b\)，求值给出同构
	\[
		B_a\longrightarrow k\times k,\qquad f\longmapsto(f(b),f(-b)).
	\]
	其核由 \((t-b)(t+b)\) 生成，满射性可用插值多项式 \((t+b)/(2b)\)、\((b-t)/(2b)\) 核验。当 \(a=0\) 时，
	\(B_0=k[t]/(t^2)\)，有一个二重点。每个 \(B_a\) 都有向量空间基 \(1,t\)，维数恒为二；只取根集，则点数从二变为一。

	特征条件在这里有直接作用：若 \(\operatorname{char}k=2\)，每个 \(a=b^2\) 都满足 \(t^2-a=(t-b)^2\)，所有纤维代数都有非零幂零元。
\end{example}

一般地，若映射由 \(k[s_1,\ldots,s_m]\rightarrow B\)、\(s_i\mapsto f_i\) 给出，在指定点 \(a=(a_1,\ldots,a_m)\) 上保留方程的纤维代数是
\[
	B/(f_1-a_1,\ldots,f_m-a_m).
\]
它满足明确的映射性质：对于任意 \(k\) 代数 \(R\)，从此商环到 \(R\) 的同态，恰对应从 \(B\) 到 \(R\) 且把每个 \(f_i\) 送到 \(a_i\) 的同态。这正是“在基底点上取纤维”对各种测试环应有的含义。把纤维理想替换为其根式再取商，只保留约化的点集结构，会丢掉\cref{g5:exm:double-cover-fibres}中的二重结构。

\ProseParagraphBreak

概形将允许非约化的交换环作为仿射模型，并把这些环的局部信息粘合起来；以上商环随后会成为仿射概形纤维的坐标环。接下来的局部化、有限模与整性理论，分别用于组织点附近的函数、控制代数的有限性，以及比较原曲线和参数曲线。此处各个模型已经显示了这些构造要保留的具体信息。


\begin{chaptersummary}
	多项式环的理想有限生成，使无限方程组可以化成有限方程组，也保证不可约分量只有有限个。零点定理进一步识别点集能够恢复的代数信息：闭集对应根式理想，不可约闭集对应素理想，点对应极大理想。非约化商环中的幂零元需要在点集之外保留。

	\ProseParagraphBreak

	仿射代数集的整体正则函数组成其坐标环，正则映射对应反向的坐标环同态。不可约对象的函数域允许在稠密开集上计算；函数域同构等价于双有理等价，也等价于适当非空开集同构。尖点参数说明，这些条件不能保证整个空间同构，延拓必须检查指定点附近的分母。

	\ProseParagraphBreak

	射影空间由仿射图及正则换图组成，齐次方程保持比例坐标的意义。乘积坐标、全部同次单项式与最大子式分别给出 Segre、Veronese 和 Plücker 嵌入；像的方程与恢复原对象的逆图共同落实了这些构造。射影曲线的参数还允许逐图检查无穷远点。

	\ProseParagraphBreak

	Gröbner 基用首项理想控制除法，Buchberger 判据把检验化为有限对的首项相消。消去理想决定投影像的闭包，饱和去除落在被排除超曲面内的部分；标准单项式则计算有限商代数的维数。局部环记录一个点附近的正则函数，保留方程的纤维代数记录重数及幂零方向，这些计算为后续局部代数和概形构造提供具体对象。
\end{chaptersummary}

\BookExercises

\begin{exercise}\label{g5:ex:radical-embedded-model}
	设 \(k\) 为代数闭域，\(I=(x^2,xy)\subseteq k[x,y]\)。求 \(V(I)\)、\(\sqrt I\)，并说明 \(k[x,y]/I\) 中 \(\overline x\) 非零而幂零。比较在 \(y\ne0\) 的部分与在原点附近，这个元素是否仍可被检测。前一种比较允许 \(y\) 的幂作分母；后一种允许所有 \(h(0,0)\ne0\) 的多项式作分母，元素 \(\overline f\) 变为零指存在这样的 \(h\) 使 \(hf\in I\)。
\end{exercise}
\begin{solution}
	\(x^2=0\) 迫使 \(x=0\)，故 \(V(I)=V(x)\)，零点定理给出 \(\sqrt I=(x)\)。\(\overline x^2=0\)，而 \(x\notin I\)：若 \(x=ax^2+bxy\)，在整环中约去 \(x\) 会得到 \(1=ax+by\)，代入原点矛盾。

	在允许 \(y\) 可逆的商环中，由 \(xy=0\) 得到 \(x=0\)，所以这个幂零元素消失。若在原点局部化后 \(x=0\)，就存在 \(h(0,0)\ne0\) 使 \(hx\in I\)。同样约去 \(x\)，得到 \(h\in(x,y)\)，矛盾。因此它在原点仍非零。这一比较直接使用分式等价关系，显示非约化信息可以集中在一个点附近。
\end{solution}

\begin{exercise}\label{g5:ex:punctured-plane-functions}
	设 \(k\) 为代数闭域，\(U=\mathbb A_k^2\setminus\{(0,0)\}\)。证明每个 \(U\) 上的正则函数都唯一延拓为 \(\mathbb A_k^2\) 上的多项式。进而证明 \(U\) 不同构于任何仿射代数集。
\end{exercise}
\begin{solution}
	\(U=D(x)\cup D(y)\)。由\cref{g5:prop:principal-open-affine}，一个正则函数在两片上分别写成 \(a/x^r\)、\(b/y^s\)，其中 \(a,b\in k[x,y]\)。它们在 \(D(xy)\) 上相同，故整环中有 \(y^sa=x^rb\)。比较关于 \(x\) 的系数，\(x^r\mid a\)，于是函数等于多项式 \(a/x^r\)。稠密性保证延拓唯一。

	若存在同构 \(Z\simeq U\)，其中 \(Z\) 仿射，则复合 \(Z\rightarrow U\hookrightarrow\mathbb A_k^2\) 在整体正则函数环上诱导同构，因为刚证得限制 \(k[x,y]\rightarrow\mathcal O(U)\) 是同构。\cref{g5:thm:affine-morphism-correspondence}遂迫使此复合是到整个 \(\mathbb A_k^2\) 的同构，但它的像遗漏原点，矛盾。
\end{solution}

\begin{exercise}\label{g5:ex:positive-characteristic-bijection}
	设 \(k\) 为特征 \(p>0\) 的代数闭域，考虑 \(F:\mathbb A_k^1\rightarrow\mathbb A_k^1\)、\(t\mapsto t^p\)。证明 \(F\) 在点集上双射，却不是同构，也不是双有理映射。
\end{exercise}
\begin{solution}
	每个 \(a\in k\) 都有 \(p\) 次根；若 \(u^p=v^p\)，则 \((u-v)^p=0\)，域中即有 \(u=v\)。因此 \(F\) 双射。拉回同态 \(k[s]\rightarrow k[t]\) 的像是 \(k[t^p]\)，不含 \(t\)，所以不是环同构，\(F\) 不是同构。

	函数域映射的像为 \(k(t^p)\)。若 \(t=P(t^p)/Q(t^p)\)，其中 \(Q\ne0\)，则 \(tQ(t^p)=P(t^p)\)；左端非零项次数均同余于 \(1\pmod p\)，右端均同余于 \(0\pmod p\)，不可能相等。因此该域嵌入也不满，\(F\) 不是双有理映射。这里结论针对指定映射 \(F\)，其源与靶作为空间当然同构。
\end{solution}

\begin{exercise}\label{g5:ex:dominant-nonsurjective-plane}
	设 \(k\) 为代数闭域，\(\varphi:\mathbb A_k^2\rightarrow\mathbb A_k^2\)、\((x,y)\mapsto(u,v)=(x,xy)\)。求实际像、各点纤维及像的闭包，并写出函数域映射的逆。
\end{exercise}
\begin{solution}
	若 \(u\ne0\)，唯一原像为 \((u,v/u)\)；若 \(u=0,v\ne0\)，没有原像；\((u,v)=(0,0)\) 的纤维则为整条 \(x=0\) 的直线。因此像为 \(D(u)\cup\{(0,0)\}\)，其闭包为整个平面。

	代数上，\(u^iv^j\) 拉回为 \(x^{i+j}y^j\)，不同 \((i,j)\) 得到不同单项式，所以 \(k[u,v]\rightarrow k[x,y]\) 单射。这也确认映射支配。函数域映射有逆 \(x\mapsto u,\ y\mapsto v/u\)，故映射双有理；在 \(x\ne0\) 与 \(u\ne0\) 的开集上实际互逆。支配性及双有理性均不要求所有纤维非空或只有一点。
\end{solution}

\begin{exercise}\label{g5:ex:projective-line-global-functions}
	设 \(k\) 为代数闭域。证明 \(\mathbb P_k^1\) 上每个整体正则函数都是常数，并推得每个正则映射 \(\mathbb P_k^1\rightarrow\mathbb A_k^n\) 都是常值映射。
\end{exercise}
\begin{solution}
	在两张仿射图上分别用 \(t\) 与 \(s=1/t\)。正则函数的限制由\cref{g5:thm:affine-global-functions}分别为 \(f(t)\in k[t]\)、\(g(s)\in k[s]\)。在交上 \(f(t)=g(t^{-1})\)，这在 Laurent 多项式环 \(k[t,t^{-1}]\) 中是恒等式。左端只有非负次数，右端只有非正次数，比较系数可知两者均为相同常数。到仿射空间的正则映射每个坐标函数都常值，结论随即得到。
\end{solution}

\begin{exercise}\label{g5:ex:segre-graph-square}
	设 \(k\) 为代数闭域，\(f:\mathbb P_k^1\rightarrow\mathbb P_k^1\)、\([s:t]\mapsto[s^2:t^2]\)。证明其图像是 \(\mathbb P_k^1\times\mathbb P_k^1\) 中的闭子集，同构于 \(\mathbb P_k^1\)。写出其 Segre 像，并将它与三次 Veronese 曲线比较。
\end{exercise}
\begin{solution}
	在乘积坐标 \(([s:t],[u:v])\) 中，图像由 \(t^2u-s^2v=0\) 给出。如果 \(s\ne0\)，方程强迫 \(u\ne0\) 且 \(v/u=t^2/s^2\)；若 \(s=0\)，则 \(t\ne0,u=0\)，同样唯一确定 \([u:v]\)。所以方程恰好给图像。它在每个仿射乘积图中由多项式方程定义，故闭；第一投影与 \([s:t]\mapsto([s:t],f([s:t]))\) 正则互逆。

	Segre 坐标为
	\[
		[su:sv:tu:tv]=[s^3:st^2:s^2t:t^3].
	\]
	交换中间两个坐标即得 \([s^3:s^2t:st^2:t^3]\)，所以它经射影线性坐标置换成为三次 Veronese 曲线。在特征二时 \(f\) 仍有这些图像与嵌入性质；图像同构由第一投影给出，不要求第二投影 \(f\) 为同构。
\end{solution}

\begin{exercise}\label{g5:ex:grassmann-incidence}
	设 \(k\) 为代数闭域，\(E=\operatorname{span}(e_1,e_2)\subseteq k^4\)。证明 \(\operatorname{Gr}(2,4)\) 中满足 \(W\cap E\ne0\) 的点恰好由 \(p_{34}=0\) 切出。在 \(p_{12}=1\) 的图上，把此条件写成四个矩阵参数的方程。
\end{exercise}
\begin{solution}
	限制到 \(W\) 的投影 \(k^4\rightarrow k^4/E\) 的核是 \(W\cap E\)。选 \(W\) 的一组行基，这个投影的矩阵就是第3、4列；存在非零核恰好等价于该二阶行列式为零，即 \(p_{34}=0\)。

	在指定图上写基矩阵
	\[
		\begin{pmatrix}1&0&a&b\\0&1&c&d\end{pmatrix}.
	\]
	条件为 \(ad-bc=0\)。这是一个四维仿射图中的行列式方程，且与全局 Plücker 超平面方程一致。
\end{solution}

\begin{exercise}\label{g5:ex:division-order}
	在域 \(k\) 上采用字典序 \(x>y\)，用有序列表 \((xy-1,y^2-1)\) 及反向列表分别对 \(xy^2\) 作多元除法。解释为什么余式不同，并求原理想的一个 Gröbner 基。
\end{exercise}
\begin{solution}
	先用 \(xy-1\)，有 \(xy^2-y(xy-1)=y\)，余式为 \(y\)；先用 \(y^2-1\)，有 \(xy^2-x(y^2-1)=x\)，余式为 \(x\)。两余式都不能再由原列表的首单项式整除。差 \(x-y\) 却属于原理想，因为
	\[
		x(y^2-1)-y(xy-1)=y-x.
	\]
	故原列表不是 Gröbner 基。

	理想等于 \((x-y,y^2-1)\)：上式给出一个包含方向，反向由 \(xy-1=y(x-y)+(y^2-1)\) 得到。对这两个首一多项式，
	\[
		S(x-y,y^2-1)=y^2(x-y)-x(y^2-1)=x-y^3
	\]
	先减去 \(x-y\)，再用 \(y^2-1\)，余式为零。因此它们构成 Gröbner 基。
\end{solution}

\begin{exercise}\label{g5:ex:groebner-triple-point}
	设 \(k\) 为代数闭域，\(I=(x^2-y,xy)\subseteq k[x,y]\)，仍取字典序 \(x>y\)。求 Gröbner 基、消去理想 \(I\cap k[y]\)、标准单项式和局部重数，并说明原方程组与消去方程在原点记录的长度为何不同。
\end{exercise}
\begin{solution}
	首项相消给出
	\[
		S(x^2-y,xy)=y(x^2-y)-x(xy)=-y^2.
	\]
	补入 \(y^2\) 后，\(S(xy,y^2)=0\)，而
	\(S(x^2-y,y^2)=-y^3\) 被 \(y^2\) 约为零，故
	\(G=(x^2-y,xy,y^2)\) 是 Gröbner 基。消去定理给出
	\(I\cap k[y]=(y^2)\)，标准单项式为 \(1,x,y\)，商代数长度为三，唯一零点是原点。

	在原商环中 \(y=x^2\)、\(x^3=0\)，所以商环同构于 \(k[x]/(x^3)\)；消去商环 \(k[y]/(y^2)\) 通过 \(y\mapsto x^2\) 嵌入其中，像由 \(1,x^2\) 张成，只有维数二。投影保留了 \(y\) 所能检测的信息，源端还有独立的 \(x\) 方向；两种长度因此不同。
\end{solution}

\begin{exercise}\label{g5:ex:saturation-components}
	设 \(k\) 为代数闭域，\(I=(xy,xz)\subseteq k[x,y,z]\)。证明
	\[
		I=(x)\cap(y,z),\qquad
		I:x^\infty=(y,z),\qquad I:y^\infty=(x).
	\]
	解释各个饱和对应保留了哪一个几何分量。
\end{exercise}
\begin{solution}
	右侧交中的多项式是 \(xf\)，且令 \(y=z=0\) 后为零，于是 \(xf(x,0,0)=0\)，故 \(f(x,0,0)=0\)，即 \(f\in(y,z)\)。反向包含直接成立，因此 \(I\) 为两个素理想的交，其零点集是平面 \(x=0\) 与直线 \(y=z=0\) 的并。

	若 \(x^Nf\in I\)，模 \((y,z)\) 后得到 \(x^N\overline f=0\) 于整环 \(k[x]\)，故 \(f\in(y,z)\)；反向对每个 \(f\in(y,z)\) 都有 \(xf\in I\)。若 \(y^Nf\in I\)，模 \((x)\) 后在 \(k[y,z]\) 中得 \(y^N\overline f=0\)，故 \(f\in(x)\)，反向由 \(xy=0\) 给出。关于 \(x\) 的饱和保留直线，关于 \(y\) 的饱和保留平面，各自去掉完全落在被排除超平面内的那个分量。
\end{solution}

\begin{exercise}\label{g5:ex:closure-intersection}
	设 \(k\) 为代数闭域。在仿射平面取 \(L_0=V(y)\)、\(L_1=V(y-1)\)。求它们在 \(\mathbb P_k^2\) 中的闭包及闭包的交，比较 \(\overline{L_0\cap L_1}\) 与 \(\overline L_0\cap\overline L_1\)。
\end{exercise}
\begin{solution}
	用仿射图 \(Z\ne0\) 的坐标 \(x=X/Z,y=Y/Z\)。两闭包为 \(Y=0\) 与 \(Y-Z=0\) 的射影直线，其交由 \(Y=Z=0\) 给出，是点 \([1:0:0]\)。仿射交为空，因为 \(y=0\) 与 \(y=1\) 不相容，所以 \(\overline{L_0\cap L_1}=\varnothing\)。闭包与有限并相容，却一般不与交相容；本题差别正是新增的无穷远交点。
\end{solution}

\begin{exercise}\label{g5:ex:second-order-directions}
	设 \(k\) 为代数闭域且特征不为二。取 \(R=k[\varepsilon]/(\varepsilon^3)\)，在结点 \(y^2=x^2(x+1)\) 和尖点 \(y^2=x^3\) 的原点分别代入
	\[
		x=a\varepsilon+b\varepsilon^2,\qquad
		y=c\varepsilon+d\varepsilon^2.
	\]
	求满足方程的条件，并解释哪些一阶方向 \((a,c)\) 可提升到这样的二阶解。
\end{exercise}
\begin{solution}
	因为 \(\varepsilon^3=0\)，有 \(x^2=a^2\varepsilon^2\)、\(y^2=c^2\varepsilon^2\)、\(x^3=0\)。结点方程给出 \(c^2=a^2\)，即 \(c=a\) 或 \(c=-a\)；\(b,d\) 任意。尖点方程给出 \(c^2=0\)，域中即为 \(c=0\)，而 \(a,b,d\) 任意。

	在模 \(\varepsilon^2\) 的一阶测试中，两个方程均没有约束，所以切空间同为 \(k^2\)。提升到本题的环后，结点只允许两个切向方向，尖点只允许一个切向方向。这个计算提取了最低次方程，并未声称每个有限阶解都能延拓到全部阶数。
\end{solution}

\begin{exercise}\label{g5:ex:hyperbola-family}
	设 \(k\) 为代数闭域，\(\pi:\mathbb A_k^2\rightarrow\mathbb A_k^1\)、\((x,y)\mapsto xy\)。计算点 \(a\in k\) 上保留方程的纤维代数，分别判断 \(a=0\) 与 \(a\ne0\) 时是否约化、是否不可约、是否连通。
\end{exercise}
\begin{solution}
	纤维代数为 \(B_a=k[x,y]/(xy-a)\)。\(a\ne0\) 时 \(x\) 可逆，\(y=a/x\)，所以
	\[
		B_a\simeq k[x,x^{-1}],
	\]
	它是整环，纤维约化且不可约，因此连通。\(a=0\) 时 \((xy)=(x)\cap(y)\) 为根式理想，纤维仍约化，却由两条坐标轴组成，因而可约。两轴各不可约且在原点相交，所以它们的并连通。本例各纤维都没有非零幂零元，特殊纤维发生的是分量分裂。
\end{solution}

\begin{exercise}\label{g5:ex:direction-map-graph}
	设 \(k\) 为代数闭域，考虑 \(\varphi:\mathbb A_k^2\mathrel{\cdots\!\rightarrow}\mathbb P_k^1\)、\((x,y)\mapsto[x:y]\)。证明它不能在原点正则延拓。证明其图像闭包为
	\[
		B=\{((x,y),[u:v])\in\mathbb A_k^2\times\mathbb P_k^1:xv=yu\},
	\]
	并求第一投影 \(B\rightarrow\mathbb A_k^2\) 在原点的纤维。
\end{exercise}
\begin{solution}
	若有原点邻域上的正则延拓，在 \(x\) 轴的非零部分恒取 \([1:0]\)，由连续性及该点在 \(\mathbb P^1\) 中闭，原点也必须取 \([1:0]\)。沿 \(y\) 轴同理必须取 \([0:1]\)，矛盾。这里原点的任意开邻域在每条轴上都是包含原点的非空开集，除去原点后仍稠密，所用极限结论是 Zariski 闭包意义的。

	图像在 \(B\) 内，且在 \((x,y)\ne(0,0)\) 上比例唯一。对任意 \([u:v]\)，曲线
	\[
		\lambda\longmapsto((\lambda u,\lambda v),[u:v])
	\]
	在 \(\lambda\ne0\) 时位于图像，在 \(\lambda=0\) 时给出 \(((0,0),[u:v])\)。非零参数在仿射直线上稠密，所以这一点也在图像闭包内，证明闭包就是 \(B\)。因此原点纤维为整个 \(\mathbb P^1\)。在 \(u\ne0\) 的图上令 \(w=v/u\)，方程为 \(y=xw\)；在 \(v\ne0\) 的图上同理为 \(x=y(u/v)\)，可以直接查看新增的方向参数。
\end{solution}

\begin{exercise}\label{g5:ex:quotient-multiplication}
	设 \(k\) 为代数闭域，\(I=(x^2-y,y^2)\subseteq k[x,y]\)。分别在字典序 \(x>y\) 和 \(y>x\) 下求 Gröbner 基及标准单项式；在有序基 \((1,x,y,xy)\) 上写出商代数中乘以 \(x\)、乘以 \(y\) 的矩阵。比较这个商代数与 \(k[x,y]/(x^2,y^2)\)：它们是否有相同长度，是否同构？
\end{exercise}
\begin{solution}
	在 \(x>y\) 时，\(G=(x^2-y,y^2)\) 的唯一一对 \(S\) 多项式为 \(-y^3\)，可约为零，标准单项式是 \(1,x,y,xy\)。在 \(y>x\) 时，可改写理想为 \((y-x^2,x^4)\)；其 \(S\) 多项式为 \(-x^6\)，同样可约为零，标准单项式为 \(1,x,x^2,x^3\)。

	以矩阵的第 \(j\) 列记录第 \(j\) 个基向量的像，乘法矩阵为
	\[
		M_x=\begin{pmatrix}
			0&0&0&0\\1&0&0&0\\0&1&0&0\\0&0&1&0
		\end{pmatrix},\qquad
		M_y=M_x^2=\begin{pmatrix}
			0&0&0&0\\0&0&0&0\\1&0&0&0\\0&1&0&0
		\end{pmatrix}.
	\]
	原商代数同构于 \(k[x]/(x^4)\)，唯一极大理想 \(\mathfrak m\) 满足 \(\mathfrak m^3\ne0,\mathfrak m^4=0\)。另一商代数也有基 \(1,x,y,xy\)，故两者长度均为四；但其唯一极大理想 \(\mathfrak n=(\overline x,\overline y)\) 满足 \(\mathfrak n^2=k\overline{xy}\ne0,\mathfrak n^3=0\)。环同构必须保持唯一极大理想及其各次幂，故二者不同构。长度相同并不确定乘法。
\end{solution}

% !TeX root = ../../Book5.tex
% 纲要标题：## 第2章　局部环、有限模与局部—整体原则
% 纲要位置：几何学计划纲要.md，第 5070 行。
% 正文状态：初稿；尚未正式审核。

% 纲要内容开始
% 
% > 本章目标与证明范围：局部化、模运算和有限性一起服务于函数芽及局部构造。完备化的证明移至独立章，不在这里把一个形式邻域定义扩充成未建立的忠实平坦性。
% 
% 纲要内容结束

\chapter{局部环、有限模与局部—整体原则}
\label{g5:ch:02}
\BookChapterNumber{g5:ch:02}

\begin{chapterabstract}
	一个点附近的函数由局部化给出，函数之间的线性关系则需要用模记录。本章建立环与模的局部化、有限呈示和正合性，证明 Nakayama 引理，使剩余域上的线性代数能够控制有限模的生成元。支集、关联素理想与准素分解进一步辨认模存在的位置、零因子及嵌入分支。最后证明模和同态的局部检测原则，并将有限投射模与有限呈示的局部自由模联系起来；纤维计算同时说明有限性与平坦性条件的作用。
\end{chapterabstract}

\begin{learninggoals}
	\begin{itemize}
		\item 从分式等价关系构造局部化，证明泛性质、素理想对应，并计算局部环及剩余域。
		\item 用生成元和关系矩阵描述有限呈示模，区分局部化、张量积与 Hom 对正合列的不同作用。
		\item 证明并使用 Nakayama 引理，由剩余向量空间寻找极小生成元，并判断生成元能否在邻域内延展。
		\item 计算支集、零化子和关联素理想，完成所需的准素分解，区分不可约分支与嵌入信息。
		\item 证明有限呈示局部自由与有限投射的等价条件，检验秩和纤维维数，并解释恒定纤维维数反例。
	\end{itemize}
\end{learninggoals}

在第一章的两条坐标轴 \(xy=0\) 上，把 \(x\) 变成单位就会迫使 \(y=0\)；在原点附近，却不能除以 \(x\) 或 \(y\)，两条轴都必须保留。同一个坐标环，因所允许的分母不同，可以描述一条分支的一般位置，也可以描述分支相遇的局部结构。局部化把这种变化变成统一的代数运算。

\ProseParagraphBreak

还有一种信息不能只用函数环本身说明。考虑 \(k[x,y]\) 中由 \(x,y\) 生成的理想：它是函数的集合，又可以视为由两个元素生成的模，关系 \(yx-xy=0\) 则说明这两个生成元不是一组基。在原点之外，至少有一个生成元成为单位，这个理想局部等于整个函数环；在原点，它仍需要两个生成元。要理解这类随点变化的线性数据，需要同时追踪生成元、关系和局部化。

\ProseParagraphBreak

本章从交换含幺环、理想、商环和有限维线性代数出发，引入所需的模、张量积与正合列。第一章的 Hilbert 基定理和函数芽模型将分别提供 Noether 环的主要例子与局部化的几何解释。一般论证不要求基域代数闭，也会逐项区分任意模、有限生成模及有限呈示模；这些区别决定能否选择共同分母，能否在开邻域内同时控制有限多条关系。主要代数构造可结合\cite[第2、4节]{Matsumura1989}阅读，支集与关联素理想的具体阅读位置列在相应节中。

% !TeX root = ../../Book5.tex
% 纲要标题：### 2.1 局部化与局部环
% 纲要位置：几何学计划纲要.md，第 5074 行。
% 正文状态：初稿；尚未正式审核。

% 纲要内容开始
% 
% * 乘法集、分式及泛性质
% * 素理想处的局部化
% * 单位、剩余域与局部同态
% * 理想扩缩和饱和
% 
% 纲要内容结束

\section{局部化与局部环}
\label{g5:sec:0201}

在\cref{g5:prop:model-local-ring}中，一个点附近的函数写成 \(a/b\)，其中分母在该点不为零。这个条件只规定哪些元素可以作分母，并不依赖多项式的具体形式。现在从任意交换环出发，构造同时把一批元素变成单位的环。环中可能有零因子，因而需要保留第一章已经出现的现象：允许更多分母，有时也会使非零元素变为零。

\ProseParagraphBreak

本章的环均交换且含幺，环同态保持单位元，模满足 \(1m=m\)。允许零环；局部环则按非零环理解。除非另作说明，基域不要求代数闭。局部化的基本构造和理想对应可对照\cite[第4节，第20--24页]{Matsumura1989}。

\subsection{乘法集、分式及泛性质}
\label{g5:subsec:0201:01}

设 \(A\) 为环。包含 \(1\)、且对乘法封闭的子集 \(S\subseteq A\) 称为\term[chengfaji]{乘法集}[multiplicative set]。典型例子是
\[
	\{1,f,f^2,\ldots\},\qquad A\setminus\mathfrak p
	\quad(\mathfrak p\text{ 为素理想}).
\]
后一集合封闭正是素性条件的另一种说法。这里允许 \(0\in S\)；这种情形最终得到零环。

\ProseParagraphBreak

希望把 \((a,s)\in A\times S\) 写成分式 \(a/s\)。若允许分母成为单位，则任何被 \(S\) 中某个元素杀掉的元素也必须成为零。因此适当的等价关系是
\[
	(a,s)\sim(b,t)
	\quad\Longleftrightarrow\quad
	u(ta-sb)=0\text{ 对某个 }u\in S.
\]
自反性和对称性直接成立。若 \(u(ta-sb)=0\) 且 \(v(rb-tc)=0\)，则
\[
	uvt(ra-sc)
	=uv\bigl[r(ta-sb)+s(rb-tc)\bigr]=0,
\]
而 \(uvt\in S\)，所以关系也传递。

\begin{definition}\label{g5:def:ring-localization}
	设 \(A\) 为交换含幺环，\(S\subseteq A\) 为乘法集。以上等价类组成的集合记作 \(S^{-1}A\)，定义
	\[
	\frac as+\frac bt=\frac{at+bs}{st},\qquad
	\frac as\frac bt=\frac{ab}{st}.
	\]
	这个环称为 \(A\) 关于 \(S\) 的\term[jubuhua]{局部化}[localization]，自然同态为 \(\iota:A\to S^{-1}A,\ a\mapsto a/1\)。 对 \(S=\{f^n:n\geq0\}\) 记作 \(A_f\)。
\end{definition}
\symidx{\(S^{-1}A,\ A_f\)：环的局部化}

说明运算不依赖代表元。例如 \(u(s'a-sa')=0\) 时，把 \(a/s\) 换成 \(a'/s'\)，加法所得两个分子的交叉差是 \(t^2(s'a-sa')\)，乘法所得的交叉差是 \(bt(s'a-sa')\)，都被 \(u\) 杀掉。再对第二个分式作同样比较即可。环公理由通分后的环公理给出，零元、单位元分别是 \(0/1,1/1\)，而 \(s/1\) 的逆为 \(1/s\)。

\begin{proposition}\label{g5:prop:localization-zero}
	设 \(A\) 为交换含幺环，\(S\) 为乘法集。则
	\[
	\frac as=0\ \Longleftrightarrow\
	ta=0\text{ 对某个 }t\in S,\qquad
	\ker\iota=\{a\in A:ta=0\text{ 对某个 }t\in S\}.
	\]
	\(S^{-1}A\) 为零环当且仅当 \(0\in S\)。自然同态 \(\iota\) 单射当且仅当对每个 \(s\in S\)，乘法 \(a\mapsto sa\) 在 \(A\) 上单射。
\end{proposition}
\begin{proof}
	第一式是 \(a/s=0/1\) 的定义，核的公式随之得到。取 \(a=s=1\)，可知 \(1/1=0/1\) 当且仅当某个 \(t\in S\) 为零。最后，核为零等价于没有非零元素被 \(S\) 中任何元素杀掉，这正是所述乘法映射的单射性。
\end{proof}

例如在 \(A=k[x,y]/(xy)\) 中，\(\overline y\ne0\)，但在 \(A_{\overline x}\) 中 \(\overline y/1=0\)，因为 \(\overline x\,\overline y=0\)。事实上
\[
	A_{\overline x}\cong k[x,x^{-1}],
\]
同构把 \(x\) 送到 \(x\)，把 \(y\) 送到零。局部化在这里保留 \(x\ne0\) 的轴上部分，也同时去掉了另一条轴的函数信息。

\begin{theorem}[局部化的泛性质]\label{g5:thm:localization-universal}
	设 \(A,B\) 为交换含幺环，\(S\subseteq A\) 为乘法集，\(\varphi:A\to B\) 为环同态。若每个 \(\varphi(s)\)、\(s\in S\) 都是 \(B\) 的单位，则存在唯一环同态
	\[
	\widetilde\varphi:S^{-1}A\longrightarrow B,\qquad
	\widetilde\varphi(a/s)=\varphi(a)\varphi(s)^{-1},
	\]
	满足 \(\widetilde\varphi\circ\iota=\varphi\)。
\end{theorem}
\begin{proof}
	若 \(u(ta-sb)=0\)，应用 \(\varphi\) 后，\(\varphi(u),\varphi(s),\varphi(t)\) 都可逆，所以 \(\varphi(a)\varphi(s)^{-1}=\varphi(b)\varphi(t)^{-1}\)。故公式良定义；通分公式说明它保持加法、乘法及单位元。另一方面，每个分式都是 \((a/1)(s/1)^{-1}\)，所以任何满足要求的同态都必须采用这个公式，唯一性成立。
\end{proof}

泛性质使不同的构造可直接比较。例如先把 \(f\) 变成单位，再把 \(g\) 变成单位，与一次把 \(fg\) 变成单位得到同构的环：
\[
	(A_f)_{g/1}\cong A_{fg}.
\]
原因是 \(fg\) 可逆当且仅当 \(f,g\) 都可逆；两边的同态由泛性质唯一确定，复合在 \(A\) 上为恒等，因而在局部化上也为恒等。

\subsection{素理想处的局部化}
\label{g5:subsec:0201:02}

把指定点之外不消失的函数都作分母，会得到局部环。对一般环，“指定点”的代数条件由一个素理想承担。先说明分式属于扩张理想的判据。若 \(I\subseteq A\) 是理想，则由 \(\iota(I)\) 生成的理想等于
\[
	S^{-1}I=\{a/s:a\in I,\ s\in S\}.
\]
有限个这种分式可通分，分子仍在 \(I\) 中，故右侧确实是理想。对任意代表元有
\begin{equation}\label{g5:eq:localized-ideal-membership}
	\frac as\in S^{-1}I
	\quad\Longleftrightarrow\quad
	ta\in I\text{ 对某个 }t\in S.
\end{equation}
若 \(a/s=b/u\)、\(b\in I\)，取 \(v\in S\) 使 \(v(ua-sb)=0\)，则 \(vua=vsb\in I\)。反过来，\(ta\in I\) 时写成 \(a/s=ta/(ts)\) 即可。

\begin{theorem}\label{g5:thm:prime-localization}
	设 \(A\) 为交换含幺环，\(S\subseteq A\) 为乘法集。扩张与收缩给出保包含关系的双射
	\[
	\{\mathfrak p\subseteq A:\mathfrak p\text{ 素且 }\mathfrak p\cap S=\varnothing\}
	\ \longleftrightarrow\ \operatorname{Spec}(S^{-1}A),
	\qquad \mathfrak p\longmapsto S^{-1}\mathfrak p.
	\]
	这里 \(\operatorname{Spec}R\) 暂表示环 \(R\) 的全部素理想所成的集合。 对 \(A\) 的任一素理想 \(\mathfrak p\)，环 \(A_{\mathfrak p}=(A\setminus\mathfrak p)^{-1}A\) 是局部环，其唯一极大理想为 \(\mathfrak pA_{\mathfrak p}\)。
\end{theorem}
\begin{proof}
	若 \(\mathfrak p\cap S=\varnothing\)，由\eqref{g5:eq:localized-ideal-membership}及素性， \(a/s\in S^{-1}\mathfrak p\) 等价于 \(a\in\mathfrak p\)。 因此扩张是真理想，且两个分式的乘积属于它时，至少一个分式属于它。它为素理想，收缩恰为 \(\mathfrak p\)。

	若 \(\mathfrak q\subset S^{-1}A\) 素，收缩 \(\mathfrak p=\iota^{-1}(\mathfrak q)\) 素，且与 \(S\) 不交，因为单位不属于真理想。 又因 \(s/1\) 可逆，\(a/s\in\mathfrak q\) 当且仅当 \(a/1\in\mathfrak q\)，即 \(a\in\mathfrak p\)。故 \(\mathfrak q=S^{-1}\mathfrak p\)，两种操作互逆。

	取 \(S=A\setminus\mathfrak p\)。若 \(a/s\notin\mathfrak pA_{\mathfrak p}\)，则 \(a\notin\mathfrak p\)，其逆就是 \(s/a\)。故该理想之外全是单位，一切真理想都包含在它里面，它是唯一极大理想。
\end{proof}
\symidx{\(A_{\mathfrak p}\)：环在素理想 \(\mathfrak p\) 处的局部化}

这里沿用第一章的定义：有唯一极大理想的非零环称为局部环，常写为 \((R,\mathfrak m)\)。任一非零环的真理想都包含于极大理想：包含该理想的真理想组成偏序集，每条链的并仍为真理想，因为若并含 \(1\)，链中某个理想就含 \(1\)；极大性原理遂给出极大元。因此在局部环中，非单位恰组成 \(\mathfrak m\)。特别地，\(1-a\) 在 \(a\in\mathfrak m\) 时必为单位，否则 \(1=(1-a)+a\in\mathfrak m\)。

\begin{example}\label{g5:exm:localizing-axis}
	设 \(k\) 为域，\(A=k[x,y]/(xy)\)，以下用 \(x,y\) 表示剩余类。 在素理想 \(\mathfrak p=(x)\) 处，\(y\) 可逆，故 \(x=0\)；所有非零 \(k[y]\) 元素也成为单位，得到
	\[
	A_{\mathfrak p}\cong k(y).
	\]
	在 \(\mathfrak m=(x,y)\) 处，\(x,y\) 都非零。例如若 \(x/1=0\)，则有 \(s\notin\mathfrak m\) 使 \(sx\in(xy)\) 于 \(k[x,y]\)；约去 \(x\) 得 \(s\in(y)\)，矛盾。对 \(y\) 同理。因此 \(A_{\mathfrak m}\) 仍有两个非零元素之积为零，能够看到两条轴在原点相遇。
\end{example}

整环 \(A\) 在零素理想处的局部化是分式域 \(\operatorname{Frac}A\)。对素数 \(p\)，\(\mathbb Z_{(p)}\) 的元素是分母不被 \(p\) 整除的有理数，其极大理想由 \(p\) 生成。这些例子表明，局部化既可取在几何点处，也可取在一个不可约分支所对应的素理想处。

\subsection{单位、剩余域与局部同态}
\label{g5:subsec:0201:03}

局部环把“不在指定素理想中”变成“可除”。在一般乘法集下，单位也有一个直接的判据。

\begin{proposition}\label{g5:prop:units-residue-fields}
	设 \(A\) 为交换含幺环，\(S\) 为乘法集。分式 \(a/s\in S^{-1}A\) 可逆当且仅当 \((a)\cap S\ne\varnothing\)。 对每个素理想 \(\mathfrak p\subset A\)，有自然同构
	\[
	\kappa(\mathfrak p):=A_{\mathfrak p}/\mathfrak pA_{\mathfrak p}
	\cong\operatorname{Frac}(A/\mathfrak p).
	\]
\end{proposition}
\begin{proof}
	若 \(ab\in S\)，则 \(sb/(ab)\) 是 \(a/s\) 的逆。反过来，若 \((a/s)(b/t)=1\)，存在 \(u\in S\) 使 \(u(ab-st)=0\)。于是 \(a(ub)=ust\in S\)，得到判据。

	映射
	\[
	A_{\mathfrak p}\longrightarrow\operatorname{Frac}(A/\mathfrak p),
	\qquad a/s\longmapsto\overline a/\overline s
	\]
	良定义，因为 \(s\notin\mathfrak p\)。分式域中的每个分母都可如此提升，故映射满。其核恰为分子属于 \(\mathfrak p\) 的分式，即 \(\mathfrak pA_{\mathfrak p}\)，商环同构成立。
\end{proof}

对局部环 \((R,\mathfrak m)\)，域 \(R/\mathfrak m\) 称为它的\term[shengyuyu]{剩余域}[residue field]。上式的 \(\kappa(\mathfrak p)\) 是点 \(\mathfrak p\) 的剩余域。\symidx{\(\kappa(\mathfrak p)\)：素理想 \(\mathfrak p\) 处的剩余域} 在有限型代数闭域代数的极大理想处，弱零点定理使剩余域等于基域；在 \(k[x,y]/(xy)\) 的 \((x)\) 处，剩余域却是 \(k(y)\)。两者记录的分别是闭点上的值和沿整条不可约分支的有理函数值。

\begin{definition}\label{g5:def:local-homomorphism}
	设 \((R,\mathfrak m)\)、\((T,\mathfrak n)\) 为局部环。满足 \(\varphi(\mathfrak m)\subseteq\mathfrak n\) 的环同态 \(\varphi:R\to T\) 称为\term[jubutongtai]{局部同态}[local homomorphism]。
\end{definition}

这个条件等价于 \(\varphi^{-1}(\mathfrak n)=\mathfrak m\)：\(\mathfrak m\) 外的元素是单位，其像仍为单位，因而不可能落在 \(\mathfrak n\) 中。局部同态于是诱导剩余域的嵌入 \(R/\mathfrak m\hookrightarrow T/\mathfrak n\)。 一般的局部环之间的环同态未必局部。例如
\[
	k[t]_{(t)}\hookrightarrow k(t)
\]
把非单位 \(t\) 送成单位，靶的极大理想是零，所以不是局部同态。

\ProseParagraphBreak

若 \(\varphi:A\to B\) 为环同态，\(\mathfrak q\subset B\) 素，令 \(\mathfrak p=\varphi^{-1}(\mathfrak q)\)，则有自然局部同态
\[
	A_{\mathfrak p}\longrightarrow B_{\mathfrak q},
	\qquad a/s\longmapsto\varphi(a)/\varphi(s).
\]
分母可用，是因为 \(s\notin\mathfrak p\) 意味着 \(\varphi(s)\notin\mathfrak q\)；分子在 \(\mathfrak p\) 时像在 \(\mathfrak q\)，故同态局部。这是以后几何映射在点附近诱导环同态的基本形式。

\subsection{理想扩缩和饱和}
\label{g5:subsec:0201:04}

对环同态 \(\varphi:A\to B\)，理想 \(I\subseteq A\) 的\term[kuozhanglixiang]{扩张}[extension]是 \(IB\)，由 \(\varphi(I)\) 生成；理想 \(J\subseteq B\) 的\term[shousuolixiang]{收缩}[contraction]是 \(\varphi^{-1}(J)\)。一般只有 \(I\subseteq(IB)^c\)、\(J^cB\subseteq J\)。局部化时，后一包含总是等号，前一包含的差别恰由可用分母决定。

\begin{proposition}\label{g5:prop:localization-ideals}
	设 \(A\) 为交换含幺环，\(S\subseteq A\) 为乘法集，\(\iota:A\to S^{-1}A\) 为自然同态。则
	\[
	\iota^{-1}(S^{-1}I)=
	I:S^\infty:=\{a\in A:sa\in I\text{ 对某个 }s\in S\},
	\qquad
	J=S^{-1}\iota^{-1}(J).
	\]
	因而扩缩给出 \(S^{-1}A\) 的理想与 \(A\) 中满足 \(I:S^\infty=I\) 的理想之间的双射。 若 \(\overline S\) 是 \(S\) 在 \(A/I\) 中的像，则
	\[
	S^{-1}A/S^{-1}I\cong\overline S^{-1}(A/I).
	\]
	若 \(A\) 是 Noether 环，则 \(S^{-1}A\) 也是 Noether 环。
\end{proposition}
\begin{proof}
	第一个公式是\eqref{g5:eq:localized-ideal-membership}在分母为 \(1\) 时的情形。 若 \(a/s\in J\)，乘以单位 \(s/1\) 得 \(a/1\in J\)，所以 \(a\in\iota^{-1}(J)\)；反向包含来自理想性。这证明第二个公式，进而证明双射。 映射 \(a/s\mapsto\overline a/\overline s\) 满，其核是满足某个 \(ta\in I\) 的分式，故核为 \(S^{-1}I\)，商环公式得证。 最后若 \(A\) Noether，每个收缩理想由有限个 \(a_1,\ldots,a_r\) 生成，其扩张由 \(a_1/1,\ldots,a_r/1\) 生成；第二个公式说明这已涵盖局部化中的每个理想。
\end{proof}

当 \(S=\{f^n:n\geq0\}\) 时，上式就是第一章消去计算中的饱和 \(I:f^\infty\)。例如 \((xy):x^\infty=(y)\) 于 \(k[x,y]\)，因为可逆的 \(x\) 不再对方程 \(xy=0\) 贡献一个分支。

\ProseParagraphBreak

现在给全部素理想一个拓扑，使“在某点附近”成为统一的说法。环 \(A\) 的\term[supuxiang]{素谱}[prime spectrum]是集合 \(\operatorname{Spec}A\)，其闭集定义为
\[
	V(I)=\{\mathfrak p\in\operatorname{Spec}A:I\subseteq\mathfrak p\},
	\qquad I\subseteq A\text{ 为理想}.
\]
素性给出 \(V(IJ)=V(I)\cup V(J)\)，且 \(\bigcap_\lambda V(I_\lambda)=V(\sum_\lambda I_\lambda)\)； 配合 \(V(0)=\operatorname{Spec}A\)、\(V(A)=\varnothing\)，得到 Zariski 拓扑。 记 \(D(f)=\{\mathfrak p:f\notin\mathfrak p\}\)。 这里的点包括非极大素理想；\(V(I)\) 表示含 \(I\) 的素理想集，与第一章的经典零点集须按所在空间区分。
\symidx{\(\operatorname{Spec}A\)：环 \(A\) 的素谱}

\begin{proposition}\label{g5:prop:radical-prime-intersection}
	设 \(A\) 为交换含幺环，\(I\subseteq A\) 为理想。则
	\[
	\sqrt I=\bigcap_{\mathfrak p\supseteq I,\ \mathfrak p\text{ 素}}\mathfrak p,
	\]
	其中空族理想的交约定为 \(A\)。因而 \(V(I)=V(J)\) 当且仅当 \(\sqrt I=\sqrt J\)。
\end{proposition}
\begin{proof}
	若 \(a^n\in I\subseteq\mathfrak p\)，素性给出 \(a\in\mathfrak p\)，所以左侧包含于右侧。 若 \(a\notin\sqrt I\)，在 \(A/I\) 中 \(\overline a\) 的任何幂都非零；故 \((A/I)_{\overline a}\ne0\)。取其一个极大理想，依次收缩到 \(A/I\) 和 \(A\)，得到一个含 \(I\) 而不含 \(a\) 的素理想。这证明反向包含。 最后，同一个闭集给出同一族素理想的交，所以给出同一根式；反方向由素理想都是根式理想立即得到。
\end{proof}

\begin{proposition}\label{g5:prop:spec-basic-opens}
	设 \(A\) 为交换含幺环。主开集 \(D(f)\) 构成 \(\operatorname{Spec}A\) 的开基， \(D(f)\cap D(g)=D(fg)\)，且每个 \(D(f)\) 都是拟紧的，即任意开覆盖均有有限子覆盖。 局部化的素理想对应给出同胚
	\[
	\operatorname{Spec}A_f\cong D(f).
	\]
	特别地，\(\operatorname{Spec}A=D(1)\) 拟紧。
\end{proposition}
\begin{proof}
	对任何理想 \(I\)，有 \(\operatorname{Spec}A\setminus V(I)=\bigcup_{a\in I}D(a)\)。 素性又给出交集公式，所以主开集构成开基。

	局部化的素理想对应把 \(\operatorname{Spec}A_f\) 双射到不含 \(f\) 的素理想集。 在此双射下，主开集 \(D(a/f^n)\) 对应 \(D(a)\cap D(f)\)，因此双射及其逆都连续，得到同胚。

	最后先证任意环 \(R\) 的素谱拟紧。把一个开覆盖细化为主开覆盖 \(\{D(a_\lambda)\}\)。 若 \((a_\lambda)_\lambda\) 是真理想，它包含于某个极大理想，该极大理想就不属于任何 \(D(a_\lambda)\)，矛盾。 所以 \(1=\sum_{i=1}^n b_i a_{\lambda_i}\) 对有限个指标成立。 这些 \(D(a_{\lambda_i})\) 已覆盖素谱，因为素理想不可能同时含所有 \(a_{\lambda_i}\)。 取回它们所属的原覆盖成员便得有限子覆盖。 再对 \(R=A_f\) 应用这一结论与刚证的同胚，得到 \(D(f)\) 拟紧。
\end{proof}

一点 \(\mathfrak p\) 的闭包是 \(V(\mathfrak p)\)：每个包含该点的闭集都含这个集合。它在 \(\mathfrak p\) 极大时才是闭点；对整环，零素理想的闭包是整个素谱。这样的点称为所在不可约闭集的\term[fanxingdian]{泛型点}[generic point]。确实，若 \(V(\mathfrak p)=V(I)\cup V(J)\)，点 \(\mathfrak p\) 必在其中一边，而该边既为闭集就含其闭包 \(V(\mathfrak p)\)，所以此闭集不可约。用一个点代表整条分支，与在该点取分式域是相容的：两者都让我们考察分支的一般位置。

% 纲要条目排为编号小节标题，并进入目录与书签；说明段仍保留于纲要注释中。

% !TeX root = ../../Book5.tex
% 纲要标题：### 2.2 有限模与正合性
% 纲要位置：几何学计划纲要.md，第 4787 行。
% 正文状态：未撰写；本文件目前是框架。

% 纲要内容开始
% 
% * 有限生成、有限呈示及关系矩阵
% * 核、余核和正合列
% * 张量积、Hom 及标量变换
% * 局部化的正合性
% 
% 纲要内容结束

\section{有限模与正合性}
\label{g5:sec:0202}

待撰写

\subsection{有限生成、有限呈示及关系矩阵}
\label{g5:subsec:0202:01}
待撰写

\subsection{核、余核和正合列}
\label{g5:subsec:0202:02}
待撰写

\subsection{张量积、Hom 及标量变换}
\label{g5:subsec:0202:03}
待撰写

\subsection{局部化的正合性}
\label{g5:subsec:0202:04}
待撰写
% 纲要条目排为编号小节标题，并进入目录与书签；说明段仍保留于纲要注释中。

% !TeX root = ../../Book5.tex
% 纲要标题：### 2.3 Nakayama 与局部生成
% 纲要位置：几何学计划纲要.md，第 4794 行。
% 正文状态：未撰写；本文件目前是框架。

% 纲要内容开始
% 
% * 行列式技巧
% * 局部 Nakayama 引理
% * 极小生成元和剩余向量空间
% * 生成元提升及有限性假设
% 
% 纲要内容结束

\section{Nakayama 与局部生成}
\label{g5:sec:0203}

待撰写

\subsection{行列式技巧}
\label{g5:subsec:0203:01}
待撰写

\subsection{局部 Nakayama 引理}
\label{g5:subsec:0203:02}
待撰写

\subsection{极小生成元和剩余向量空间}
\label{g5:subsec:0203:03}
待撰写

\subsection{生成元提升及有限性假设}
\label{g5:subsec:0203:04}
待撰写
% 纲要条目排为编号小节标题，并进入目录与书签；说明段仍保留于纲要注释中。

% !TeX root = ../../Book5.tex
% 纲要标题：### 2.4 支集与关联素理想
% 纲要位置：几何学计划纲要.md，第 5095 行。
% 正文状态：初稿；尚未正式审核。

% 纲要内容开始
% 
% * 支集和零化子
% * 关联素理想的存在和有限性
% * 零因子及嵌入分支
% * 准素分解的所需版本
% 
% 纲要内容结束

\section{支集与关联素理想}
\label{g5:sec:0204}

局部化能够使一个非零模变为零，也能够使模中的某个非零元素消失。于是有两个相近而不同的问题：这个模在哪些素理想处仍然非零？它的哪些部分能够被一个素理想精确地零化？前一个问题产生支集，后一个问题产生关联素理想。关联素理想还会记录集中在较小闭集上的代数信息，使我们能进一步把一个理想分解成保留幂次的准素分量。

\ProseParagraphBreak

以下在 \(\operatorname{Spec}A\) 上讨论支集，点是素理想，闭集 \(V(I)\) 由包含 \(I\) 的素理想组成。支集的基本定义适用于任意交换含幺环；关联素理想的存在、有限性和准素分解中的有限性条件将在各个结果中写明。可对照\cite[第4节，第25--26页；第6节，第38--42页]{Matsumura1989}阅读。

\subsection{支集和零化子}
\label{g5:subsec:0204:01}
例如，对 \(A/(f)\) 在 \(\mathfrak p\) 处局部化，如果 \(f\notin\mathfrak p\)，那么 \(f\) 成为单位，而它在这个模上作用为零，所以整个局部化模为零。若想同时处理多个关系以及一般的有限模，就要把“哪些标量把它杀掉”整理成一个理想。

\begin{definition}\label{g5:def:support-annihilator}
	设 \(A\) 为交换含幺环，\(M\) 为 \(A\) 模。元素 \(m\in M\) 的\term[linghuazi]{零化子}[annihilator]及模 \(M\) 的零化子分别是
	\[
		\operatorname{Ann}_A(m)=\{a\in A:am=0\},\qquad
		\operatorname{Ann}_A(M)=\{a\in A:aM=0\}
		=\bigcap_{m\in M}\operatorname{Ann}_A(m).
	\]
	它们都是 \(A\) 的理想。模 \(M\) 的\term[zhiji]{支集}[support]定义为
	\[
		\operatorname{Supp}_A(M)=\{\mathfrak p\in\operatorname{Spec}A:M_{\mathfrak p}\ne0\}.
	\]
	底环明确时省略下标 \(A\)。规定 \(\operatorname{Ann}_A(0)=A\)、\(\operatorname{Supp}_A(0)=\varnothing\)。
\end{definition}
\symidx{\(\operatorname{Ann}_A(M)\)：模的零化子}
\symidx{\(\operatorname{Supp}_A(M)\)：模的支集}

映射 \(A\to Am\)、\(a\mapsto am\) 的核是 \(\operatorname{Ann}_A(m)\)，所以
\[
	Am\cong A/\operatorname{Ann}_A(m).
\]
一个元素生成的循环子模，由它的零化子完全决定。对任意模则有 \(\operatorname{Ann}_A(M)\subseteq\operatorname{Ann}_A(m)\)；整个模的共同零化条件可能比任何一个元素的零化条件都少。

\begin{proposition}\label{g5:prop:support-annihilator}
	设 \(A\) 为交换含幺环，\(M\) 为 \(A\) 模。总有
	\[
		\operatorname{Supp}_A(M)\subseteq V(\operatorname{Ann}_A(M)).
	\]
	若 \(M\) 有限生成，则等号成立，支集因而是闭集。特别地，对任意理想 \(I\subseteq A\)，有 \(\operatorname{Supp}_A(A/I)=V(I)\)。
\end{proposition}
\begin{proof}
	若某个 \(a\in\operatorname{Ann}_A(M)\) 不属于 \(\mathfrak p\)，则 \(a/1\) 在 \(A_{\mathfrak p}\) 中可逆，却把 \(M_{\mathfrak p}\) 的每个元素送到零，所以 \(M_{\mathfrak p}=0\)。这证明一般的包含关系。

	现在设 \(M\) 由 \(m_1,\ldots,m_r\) 生成，且 \(M_{\mathfrak p}=0\)。由模局部化的零元判据\eqref{g5:eq:module-localization-zero}，对每个 \(i\) 存在 \(s_i\notin\mathfrak p\) 使 \(s_im_i=0\)。令 \(s=\prod_i s_i\)，则 \(s\notin\mathfrak p\)，且 \(s\) 杀掉全部生成元，从而 \(sM=0\)。故 \(\operatorname{Ann}_A(M)\) 不包含在 \(\mathfrak p\) 中。这给出有限生成情形的反向包含。最后，\(A/I\) 由 \(1+I\) 生成，且其零化子为 \(I\)，于是得到商模公式。
\end{proof}

证明中有限生成性的作用，是把各个元素所需的分母合成一个共同分母。对无限多个生成元，不能作这样的有限乘积。

\begin{example}\label{g5:exm:nonclosed-module-support}
	设 \(k\) 为代数闭域，\(A=k[t]\)，取无限直和
	\[
		M=\bigoplus_{a\in k}A/(t-a).
	\]
	一个多项式若杀掉整个 \(M\)，就在所有 \(a\in k\) 处取零，故 \(\operatorname{Ann}_A(M)=0\)。在 \((0)\) 处局部化时，每个直和分量都被可逆的 \(t-a\) 杀掉；直和中的每个元素只有有限个非零坐标，所以 \(M_{(0)}=0\)。在 \((t-a)\) 处，第 \(a\) 个分量仍非零。由于 \(k[t]\) 的非零素理想正是这些极大理想，
	\[
		\operatorname{Supp}_A(M)=\{(t-a):a\in k\},
		\qquad V(\operatorname{Ann}_A(M))=\operatorname{Spec}k[t].
	\]
	这个支集包含全部闭点而遗漏泛型点 \((0)\)。任何包含它的闭集 \(V(J)\) 都要求 \(J\) 中每个多项式在全部 \(a\) 处取零，因而 \(J=0\)。所以支集稠密，却不是闭集。
\end{example}

\begin{proposition}\label{g5:prop:support-exact}
	设 \(A\) 为交换含幺环，\(0\to L\to M\to N\to0\) 为 \(A\) 模的短正合列。则
	\[
		\operatorname{Supp}_A(M)=\operatorname{Supp}_A(L)\cup\operatorname{Supp}_A(N).
	\]
	这里不要求模有限生成。
\end{proposition}
\begin{proof}
	对任意 \(\mathfrak p\in\operatorname{Spec}A\)，\cref{g5:thm:localization-exact}给出短正合列
	\[
		0\longrightarrow L_{\mathfrak p}\longrightarrow M_{\mathfrak p}
		\longrightarrow N_{\mathfrak p}\longrightarrow0.
	\]
	中间项为零当且仅当两端都为零。对全部 \(\mathfrak p\) 应用这个等价关系，就得到支集公式。
\end{proof}

因此，通过一条有限滤过把模拆成较简单的商模时，原模的支集就是这些商模支集的并。接下来要构造一种商模形如 \(A/\mathfrak p\) 的滤过，使其支集成为素理想确定的不可约闭集。

\subsection{关联素理想的存在和有限性}
\label{g5:subsec:0204:02}
若一个非零元素 \(m\) 的零化子恰为素理想 \(\mathfrak p\)，则循环子模 \(Am\cong A/\mathfrak p\) 在整条 \(V(\mathfrak p)\) 上存活。它不仅指出了某个局部化非零，还在原模中找到了一个承载这条闭集的具体子模。

\begin{definition}\label{g5:def:associated-prime}
	设 \(A\) 为交换含幺环，\(M\) 为 \(A\) 模。若素理想 \(\mathfrak p\) 可以写成 \(\operatorname{Ann}_A(m)\)，其中 \(0\ne m\in M\)，则称 \(\mathfrak p\) 为 \(M\) 的\term[guanliansulixiang]{关联素理想}[associated prime ideal]。全部关联素理想组成的集合记为 \(\operatorname{Ass}_A(M)\)。等价地，\(\mathfrak p\in\operatorname{Ass}_A(M)\) 当且仅当存在单射 \(A/\mathfrak p\hookrightarrow M\)。零模的关联素理想集为空。
\end{definition}
\symidx{\(\operatorname{Ass}_A(M)\)：模的关联素理想集}

一般元素的零化子未必素。例如在 \(k[t]/(t^2)\) 中，\(1\) 的零化子为 \((t^2)\)，而 \(t\) 的零化子是 \((t)\)。把元素换成一个适当的非零倍式，可能增大它的零化子；Noether 条件保证这个过程不会无限严格上升。

\begin{lemma}\label{g5:lem:associated-prime-exists}
	设 \(A\) 为交换含幺 Noether 环，\(M\ne0\) 为任意 \(A\) 模。在集合
	\[
		\{\operatorname{Ann}_A(m):0\ne m\in M\}
	\]
	中按包含关系取的每个极大元都是素理想。因此 \(\operatorname{Ass}_A(M)\ne\varnothing\)。
\end{lemma}
\begin{proof}
	Noether 环的理想升链条件保证这个非空理想族有极大元。取 \(\mathfrak p=\operatorname{Ann}_A(m)\) 为其中一个；由于 \(m\ne0\)，有 \(\mathfrak p\ne A\)。若 \(ab\in\mathfrak p\)、\(b\notin\mathfrak p\)，则 \(bm\ne0\)，并有
	\[
		\mathfrak p\subseteq\operatorname{Ann}_A(bm).
	\]
	极大性迫使二者相等。由 \(abm=0\)，得到 \(a\in\operatorname{Ann}_A(bm)=\mathfrak p\)，所以 \(\mathfrak p\) 为素理想。
\end{proof}

这个引理只要求环 Noether，没有要求 \(M\) 有限生成。有限生成性将用于控制关联素理想的总数。为此先说明，关联素理想怎样经过短正合列。

\begin{proposition}\label{g5:prop:ass-exact-sequence}
	设 \(A\) 为交换含幺环，\(0\to L\to M\to N\to0\) 为短正合列，则
	\[
		\operatorname{Ass}_A(L)\subseteq\operatorname{Ass}_A(M)
		\subseteq\operatorname{Ass}_A(L)\cup\operatorname{Ass}_A(N).
	\]
	特别地，有限个模的直和满足
	\[
		\operatorname{Ass}_A(M_1\oplus\cdots\oplus M_r)
		=\bigcup_{i=1}^r\operatorname{Ass}_A(M_i).
	\]
\end{proposition}
\begin{proof}
	将 \(L\) 识别为 \(M\) 的子模。其非零元素的零化子不变，所以第一个包含成立。对 \(\mathfrak p\in\operatorname{Ass}_A(M)\)，取循环子模 \(C=Am\cong A/\mathfrak p\)。若 \(C\cap L=0\)，则 \(C\to N\) 为单射，故 \(\mathfrak p\in\operatorname{Ass}_A(N)\)。

	若 \(C\cap L\ne0\)，选 \(0\ne bm\in C\cap L\)。由于 \(\operatorname{Ann}_A(m)=\mathfrak p\)，有 \(b\notin\mathfrak p\)。对于 \(a\in A\)，
	\[
		abm=0\quad\Longleftrightarrow\quad ab\in\mathfrak p
		\quad\Longleftrightarrow\quad a\in\mathfrak p,
	\]
	最后一个等价使用了素性及 \(b\notin\mathfrak p\)。因此 \(\operatorname{Ann}_A(bm)=\mathfrak p\)，说明 \(\mathfrak p\in\operatorname{Ass}_A(L)\)。这证明第二个包含。对直和，短正合列给出向右的包含，每个直和分量嵌入整个直和又给出反向包含，有限次归纳即得公式。
\end{proof}

关联素理想在商模中的行为与子模不同。例如取域 \(k\) 上的 \(A=k[t]\)，短正合列
\[
	0\longrightarrow A\xrightarrow{\,t\,}A\longrightarrow A/(t)\longrightarrow0
\]
中，前两个非零模的关联素理想均为 \((0)\)，商模却出现 \((t)\)。因此不能把上个命题中的包含关系一概改成等号。

\begin{theorem}\label{g5:thm:prime-filtration}
	设 \(A\) 为交换含幺 Noether 环，\(M\) 为有限生成 \(A\) 模。则存在有限滤过
	\[
		0=M_0\subsetneq M_1\subsetneq\cdots\subsetneq M_r=M,
		\qquad M_i/M_{i-1}\cong A/\mathfrak p_i,
	\]
	其中各 \(\mathfrak p_i\) 为素理想；\(M=0\) 时允许 \(r=0\)。这样的滤过称为\term[sulvguo]{素滤过}[prime filtration]，并满足
	\[
		\operatorname{Ass}_A(M)\subseteq\{\mathfrak p_1,\ldots,\mathfrak p_r\},
		\qquad \operatorname{Supp}_A(M)=\bigcup_{i=1}^r V(\mathfrak p_i).
	\]
	特别地，\(\operatorname{Ass}_A(M)\) 为有限集合。
\end{theorem}
\begin{proof}
	若 \(M\ne0\)，由\cref{g5:lem:associated-prime-exists}选择 \(M_1\cong A/\mathfrak p_1\) 的非零循环子模。若 \(M_1\ne M\)，对非零商模 \(M/M_1\) 作同样选择，再把选出的循环子模拉回 \(M\)，得到 \(M_2\)。如此继续，每一步都严格增大子模。由\cref{g5:prop:noetherian-finite-modules}，有限生成 \(A\) 模满足子模升链条件，因此只能作有限步，并且终止时必须达到 \(M\)。

	对素理想 \(\mathfrak p\)，模 \(A/\mathfrak p\) 中每个非零元素的零化子都是 \(\mathfrak p\)，所以 \(\operatorname{Ass}_A(A/\mathfrak p)=\{\mathfrak p\}\)。对滤过中逐个短正合列使用\cref{g5:prop:ass-exact-sequence}，得到关联素理想的包含关系。同样使用\cref{g5:prop:support-exact,g5:prop:support-annihilator}，得到支集的有限并公式。
\end{proof}

素滤过中的 \(\mathfrak p_i\) 可能重复，也可能包含并非 \(M\) 的关联素理想的素理想；滤过是构造工具，\(\operatorname{Ass}_A(M)\) 才是由模本身确定的集合。现在考察局部化怎样从这个集合中选出仍然可见的素理想。

\begin{theorem}\label{g5:thm:associated-localization}
	设 \(A\) 为交换含幺 Noether 环，\(M\) 为 \(A\) 模，\(S\subseteq A\) 为乘法集。则
	\[
		\operatorname{Ass}_{S^{-1}A}(S^{-1}M)
		=\{S^{-1}\mathfrak p:\mathfrak p\in\operatorname{Ass}_A(M),\quad
		\mathfrak p\cap S=\varnothing\}.
	\]
	左右的素理想通过\cref{g5:thm:prime-localization}的扩张、收缩对应。
\end{theorem}
\begin{proof}
	若 \(\mathfrak p=\operatorname{Ann}_A(m)\) 且 \(\mathfrak p\cap S=\varnothing\)，则 \(m/1\ne0\)，否则某个 \(s\in S\) 满足 \(sm=0\)，与不交条件矛盾。又有
	\[
		\frac as\frac m1=0
		\quad\Longleftrightarrow\quad tam=0\text{ 对某个 }t\in S
		\quad\Longleftrightarrow\quad a\in\mathfrak p.
	\]
	最后一步使用 \(ta\in\mathfrak p\)、\(t\notin\mathfrak p\) 及素性。因此 \(m/1\) 的零化子恰为 \(S^{-1}\mathfrak p\)。

	反过来，设 \(\mathfrak q\) 是某个非零 \(m/s\in S^{-1}M\) 的零化子。因为 \(s/1\) 可逆，\(m/1\) 的零化子同样是 \(\mathfrak q\)。令 \(\mathfrak p\subseteq A\) 为 \(\mathfrak q\) 的收缩；它为素理想、与 \(S\) 不交，且 \(\mathfrak q=S^{-1}\mathfrak p\)。Noether 性使 \(\mathfrak p=(a_1,\ldots,a_\ell)\) 有限生成。每个 \(a_i\) 在局部化后杀掉 \(m/1\)，所以存在 \(t_i\in S\) 使 \(t_i a_im=0\)。取 \(t=\prod_i t_i\)，便有 \(\mathfrak p\cdot tm=0\)。

	元素 \(tm\) 非零，因为其局部化是单位 \(t/1\) 乘以非零的 \(m/1\)。若 \(b\cdot tm=0\)，局部化并约去这个单位，得到 \((b/1)(m/1)=0\)，故 \(b\in\mathfrak p\)。于是 \(\operatorname{Ann}_A(tm)=\mathfrak p\)，证明 \(\mathfrak p\in\operatorname{Ass}_A(M)\)。若 \(0\in S\)，局部化环和模都为零，右侧也没有与 \(S\) 不交的素理想，公式仍成立。
\end{proof}

这个反向证明解释了为何需要把局部化中的元素乘以一个共同分母：最初的 \(m\) 未必已被整个 \(\mathfrak p\) 杀掉，而 \(tm\) 可以做到。所用有限性来自素理想 \(\mathfrak p\) 的有限生成，因此这里仍不需要 \(M\) 有限生成。

\subsection{零因子及嵌入分支}
\label{g5:subsec:0204:03}
支集知道模在哪里存在，却没有直接告诉我们哪些标量的乘法会失去单射性。关联素理想恰好补足这一信息：它们的并就是所有能杀掉非零元素的标量。

\begin{definition}\label{g5:def:module-zero-divisor}
	设 \(A\) 为交换含幺环，\(M\) 为 \(A\) 模。若存在 \(0\ne m\in M\) 使 \(am=0\)，则称 \(a\in A\) 为 \(M\) 上的\term[lingyinzi]{零因子}[zero divisor]。其余元素的乘法映射 \(M\xrightarrow{a}M\) 为单射，称为 \(M\) 上的\term[zhengzeyuan]{正则元}[regular element]。特别地，\(M\ne0\) 时零元素也是零因子；零模上每个元素都是正则元。
\end{definition}

\begin{theorem}\label{g5:thm:zero-divisors-associated}
	设 \(A\) 为交换含幺 Noether 环，\(M\) 为任意 \(A\) 模。则 \(M\) 上的零因子集恰为
	\[
		\bigcup_{\mathfrak p\in\operatorname{Ass}_A(M)}\mathfrak p.
	\]
\end{theorem}
\begin{proof}
	若 \(a\in\mathfrak p=\operatorname{Ann}_A(m)\)，其中 \(m\ne0\)，则 \(am=0\)，所以 \(a\) 为零因子。反过来，若 \(a\) 为零因子，则子模 \(L=\{m\in M:am=0\}\) 非零。由\cref{g5:lem:associated-prime-exists}，\(L\) 中有一个非零元素 \(m\)，其零化子是素理想 \(\mathfrak p\)。这个零化子在 \(M\) 中仍相同，所以 \(\mathfrak p\in\operatorname{Ass}_A(M)\)，而 \(am=0\) 又给出 \(a\in\mathfrak p\)。
\end{proof}

在 Noether 环的有限模上，检验一个元素是否正则因而只需避开有限多个关联素理想。这种判定保留了支集中可能看不见的较小部分；先比较两个集合各自的极小元。

\begin{corollary}\label{g5:cor:minimal-support-associated}
	设 \(A\) 为交换含幺 Noether 环，\(M\) 为有限生成 \(A\) 模。则
	\[
		\operatorname{Ass}_A(M)\subseteq\operatorname{Supp}_A(M),\qquad
		\operatorname{Supp}_A(M)=
		\bigcup_{\mathfrak p\in\operatorname{Ass}_A(M)}V(\mathfrak p).
	\]
	按理想包含关系，\(\operatorname{Supp}_A(M)\) 与 \(\operatorname{Ass}_A(M)\) 有相同的极小元。它们为有限多个，其闭包恰是支集的不可约分支。
\end{corollary}
\begin{proof}
	若 \(\mathfrak p=\operatorname{Ann}_A(m)\)、\(m\ne0\)，则对每个包含 \(\mathfrak p\) 的素理想 \(\mathfrak q\)，元素 \(m/1\in M_{\mathfrak q}\) 非零：否则有 \(s\notin\mathfrak q\) 使 \(sm=0\)，就会有 \(s\in\mathfrak p\subseteq\mathfrak q\)，矛盾。因此 \(V(\mathfrak p)\subseteq\operatorname{Supp}_A(M)\)，也得到第一个包含。

	现在取 \(\mathfrak q\in\operatorname{Supp}_A(M)\)。由\cref{g5:prop:localization-ideals}，\(A_{\mathfrak q}\) 为 Noether 环，而 \(M_{\mathfrak q}\ne0\)。\cref{g5:lem:associated-prime-exists}保证它有一个关联素理想。再由\cref{g5:thm:associated-localization}，该素理想来自某个 \(\mathfrak p\in\operatorname{Ass}_A(M)\)，且 \(\mathfrak p\subseteq\mathfrak q\)。于是 \(\mathfrak q\in V(\mathfrak p)\)，证明并集公式。

	由于关联素理想集有限，其中每个元素都包含一个极小元。因此上式可以只对这些极小元取并，而任意支集素理想都包含其中一个。结合 \(\operatorname{Ass}_A(M)\subseteq\operatorname{Supp}_A(M)\)，两个集合的极小元恰好相同。每个 \(V(\mathfrak p)\) 都不可约，且理想之间互不包含使这些闭集也互不包含。有限个闭集的并中的任一不可约子集必落在其中一项内，否则把相交拆开就违反不可约性。故这些 \(V(\mathfrak p)\) 正是支集的不可约分支。
\end{proof}

对 Noether 环上的有限模，支集中的极小素理想称为\term[guliguanliansulixiang]{孤立关联素理想}[isolated associated prime]，也称极小关联素理想；其余关联素理想称为\term[qianrusulixiang]{嵌入素理想}[embedded prime]。若 \(\mathfrak p\) 为嵌入素理想，它包含一个更小的关联素理想 \(\mathfrak q\)，所以 \(V(\mathfrak p)\) 严格包含在 \(V(\mathfrak q)\) 中。它所记录的子模集中在已有不可约分支的一个较小闭集上。

\begin{example}\label{g5:exm:associated-embedded-point}
	设 \(k\) 为任意域，\(R=k[x,y]\)，\(I=(x^2,xy)\)，\(M=R/I\)。这是\cref{g5:ex:radical-embedded-model}中已出现的模型；现在把它作为 \(R\) 模计算。以 \(\overline x,\overline y\) 表示剩余类，每个元素唯一写成
	\[
		a(\overline y)+b\overline x,\qquad a(y)\in k[y],\quad b\in k.
	\]
	存在性来自 \(\overline x^2=\overline x\,\overline y=0\)；唯一性来自 \(I\) 中每个单项式都含 \(x^2\) 或 \(xy\)，不能出现非零的 \(a(y)+bx\)。

	因为 \(f(x,y)\overline x=f(0,0)\overline x\)，且 \(\overline x\ne0\)，有
	\[
		\operatorname{Ann}_R(\overline x)=(x,y).
	\]
	另一方面，\(f(x,y)\overline y=f(0,y)\overline y\)，而 \(k[y]\) 是整环，所以
	\[
		\operatorname{Ann}_R(\overline y)=(x).
	\]
	两者都是关联素理想。为了排除其他关联素理想，取子模 \(R\overline x\cong R/(x,y)\)，其商模为 \(R/(x)\)。短正合列
	\[
		0\longrightarrow R/(x,y)\xrightarrow{\,1\mapsto\overline x\,}M
		\longrightarrow R/(x)\longrightarrow0
	\]
	与\cref{g5:prop:ass-exact-sequence}给出
	\[
		\operatorname{Ass}_R(M)=\{(x),(x,y)\},\qquad
		\operatorname{Supp}_R(M)=V(x).
	\]
	支集只有一条不可约分支 \(V(x)\)，但 \(R\overline x\) 是支撑在原点 \((x,y)\) 上的非零子模。其元素 \(\overline x\) 在 \(y\) 可逆时消失，在原点局部化后仍非零，因为没有原点之外的分母能落入它的零化子 \((x,y)\)。因此 \((x)\) 为孤立关联素理想，\((x,y)\) 为嵌入素理想。零因子集为 \((x)\cup(x,y)=(x,y)\)；例如 \(y\) 是零因子，虽然它不属于唯一的极小关联素理想 \((x)\)。
\end{example}

只检查支集的不可约分支，会漏掉上例中 \(y\) 杀掉 \(\overline x\) 的现象。准素分解的目标就是在分解方程理想时保留这种信息，同时区分哪些分量由理想唯一决定。

\subsection{准素分解的所需版本}
\label{g5:subsec:0204:04}
第一章的根式理想与不可约分支对应，使我们能把代数集拆成不可约闭集。但取根式会丢掉 \(x^2\) 与 \(x\) 的区别；本小节直接分解理想，允许一个分量内部带有幂零元素。与素理想相比，适当放宽的条件是：乘积为零时，一个因子可能只在取幂后才变为零。

\begin{definition}\label{g5:def:primary-ideal}
	设 \(A\) 为交换含幺环。真理想 \(Q\subsetneq A\) 若满足
	\[
		ab\in Q,\quad b\notin Q
		\quad\Longrightarrow\quad a^n\in Q\text{ 对某个 }n\ge1,
	\]
	则称为\term[zhunsulixiang]{准素理想}[primary ideal]。若 \(\sqrt Q=\mathfrak p\)，则称 \(Q\) 为 \(\mathfrak p\) 准素理想。把真理想写成有限个准素理想的交，称为\term[zhunsufenjie]{准素分解}[primary decomposition]。
\end{definition}

等价地，\(Q\) 为准素理想，当且仅当商环 \(A/Q\) 的每个零因子都是幂零元。定义中的 \(b\notin Q\) 正保证 \(b+Q\) 是检验零因子的非零元素。素理想当然准素，此时可以取 \(n=1\)。

\begin{lemma}\label{g5:lem:primary-ideal-properties}
	设 \(A\) 为交换含幺环，\(Q\subsetneq A\) 为准素理想，\(\mathfrak p=\sqrt Q\)。则 \(\mathfrak p\) 是素理想。对 \(c\notin Q\)，理想
	\[
		Q:c=\{a\in A:ac\in Q\}
	\]
	仍为 \(\mathfrak p\) 准素理想。有限个 \(\mathfrak p\) 准素理想的交也为 \(\mathfrak p\) 准素理想。此外，若 \(\iota:A\to A_{\mathfrak p}\) 为自然同态，则
	\[
		\iota^{-1}(QA_{\mathfrak p})=Q.
	\]
\end{lemma}
\begin{proof}
	若 \(ab\in\sqrt Q\)、\(b\notin\sqrt Q\)，则对某个 \(r\) 有 \(a^rb^r\in Q\)，而 \(b^r\notin Q\)。准素性给出 \(a^{rn}\in Q\)，所以 \(a\in\sqrt Q\)。这证明 \(\mathfrak p\) 为素理想。

	由于 \(c\notin Q\)，\(Q:c\) 是真理想。若 \(ab\in Q:c\)、\(b\notin Q:c\)，则 \(abc\in Q\)、\(bc\notin Q\)，故某个 \(a^n\in Q\subseteq Q:c\)。所以 \(Q:c\) 准素。包含关系 \(Q\subseteq Q:c\) 给出 \(\mathfrak p\subseteq\sqrt{Q:c}\)；反之，若 \(a^r c\in Q\)，由 \(c\notin Q\) 得到 \(a\in\mathfrak p\)，所以两根式相同。

	设 \(Q_1,\ldots,Q_r\) 都是 \(\mathfrak p\) 准素理想。有限交的根式仍为 \(\mathfrak p\)：一个元素分别有幂属于各个 \(Q_i\)，取指数的最大值即可使同一个幂属于交。若 \(ab\in\bigcap_iQ_i\)、\(b\notin\bigcap_iQ_i\)，选择一个 \(b\notin Q_j\)，则 \(a\in\mathfrak p\)，进而某个 \(a^N\in\bigcap_iQ_i\)。因此该交准素。

	最后，若 \(a/1\in QA_{\mathfrak p}\)，由\cref{g5:prop:localization-ideals}，存在 \(s\notin\mathfrak p\) 使 \(sa\in Q\)。如果 \(a\notin Q\)，准素性将给出 \(s\in\sqrt Q=\mathfrak p\)，矛盾。故 \(a\in Q\)，证明收缩公式。
\end{proof}

准素性也可以完全用关联素理想表述。

\begin{proposition}\label{g5:prop:primary-associated-singleton}
	设 \(A\) 为交换含幺 Noether 环，\(Q\subsetneq A\) 为理想。则 \(Q\) 准素，当且仅当 \(\operatorname{Ass}_A(A/Q)\) 只有一个元素。在这种情形下，
	\[
		\operatorname{Ass}_A(A/Q)=\{\sqrt Q\}.
	\]
\end{proposition}
\begin{proof}
	若 \(Q\) 准素，令 \(\mathfrak p=\sqrt Q\)。取任意 \(\mathfrak q\in\operatorname{Ass}_A(A/Q)\)，它是某个非零剩余类的零化子。由于 \(Q\) 杀掉全部商模，有 \(Q\subseteq\mathfrak q\)，故 \(\mathfrak p\subseteq\mathfrak q\)。另一方面，\(\mathfrak q\) 中每个元素都在 \(A/Q\) 上是零因子，准素性使它模 \(Q\) 幂零，所以 \(\mathfrak q\subseteq\mathfrak p\)。关联素理想的存在性保证集合非空，于是它恰为 \(\{\mathfrak p\}\)。

	反过来，若 \(\operatorname{Ass}_A(A/Q)=\{\mathfrak p\}\)，则\cref{g5:cor:minimal-support-associated,g5:prop:support-annihilator}给出 \(V(Q)=V(\mathfrak p)\)。由\cref{g5:prop:radical-prime-intersection}，\(\sqrt Q=\mathfrak p\)。\cref{g5:thm:zero-divisors-associated}又说明 \(A/Q\) 上的零因子恰为 \(\mathfrak p\) 中的元素，它们模 \(Q\) 都幂零，故 \(Q\) 准素。
\end{proof}

为了证明每个理想都有这样的分解，先作一个较粗的拆分。一个真理想 \(J\) 若不能表示成两个严格包含它的理想之交，就称为\term[bukeyuelixiang]{不可约理想}[irreducible ideal]。这里不可约描述的是理想的交分解，与闭集的不可约性需要分别判断；后者只看根式。

\begin{lemma}\label{g5:lem:irreducible-ideal-primary}
	设 \(A\) 为交换含幺 Noether 环。每个真理想都是有限个不可约理想的交，而且每个不可约理想均准素。
\end{lemma}
\begin{proof}
	先证明有限分解。若存在不能作有限不可约分解的真理想，则由理想升链条件，可在这些反例中取一个按包含关系极大的 \(J\)。它本身不可能不可约，故有 \(J=J_1\cap J_2\)，其中两个 \(J_i\) 都严格包含 \(J\)。两者仍为真理想，否则这个交就等于另一项而无法严格增大。由 \(J\) 的极大性，每个 \(J_i\) 都是有限个不可约理想的交，将两份表达式合并，就给出 \(J\) 的有限不可约分解，矛盾。

	再设 \(Q\) 不可约。转到商环 \(B=A/Q\)，其中零理想不可约。若 \(ab=0\) 且 \(b\ne0\)，要证明 \(a\) 幂零。考虑升链
	\[
		(0:a)\subseteq(0:a^2)\subseteq\cdots,
		\qquad (0:a^n)=\{u\in B:a^nu=0\}.
	\]
	商环 \(B\) 的理想升链提升为 \(A\) 中含 \(Q\) 的理想升链，故 \(B\) 仍 Noether。所以可以取 \(N\ge1\) 使 \((0:a^N)=(0:a^{N+1})\)。现在若 \(z\in(a^N)\cap(b)\)，写成 \(z=a^Nu=bv\)。由 \(ab=0\)，有 \(a^{N+1}u=az=0\)；稳定性遂给出 \(a^Nu=0\)，也就是 \(z=0\)。因此
	\[
		(0)=(a^N)\cap(b).
	\]
	零理想不可约而 \((b)\ne0\)，迫使 \((a^N)=0\)。这就证明 \(a^N=0\)，从而 \(Q\) 准素。
\end{proof}

\begin{theorem}[准素分解定理]\label{g5:thm:primary-decomposition}
	\termidx[zhunsufenjiedingli]{准素分解定理}\foreignidx{primary decomposition theorem}设 \(A\) 为交换含幺 Noether 环，\(I\subsetneq A\) 为真理想。则存在有限分解
	\[
		I=Q_1\cap\cdots\cap Q_r,
	\]
	其中每个 \(Q_i\) 均准素。可以要求这份分解不可删去任何一项，且 \(\mathfrak p_i=\sqrt{Q_i}\) 两两不同。满足这两个条件的分解称为\term[jixiaozhunsufenjie]{极小准素分解}[minimal primary decomposition]。
\end{theorem}
\begin{proof}
	由\cref{g5:lem:irreducible-ideal-primary}，先将 \(I\) 写成有限个不可约理想的交，它们都是准素理想。将根式相同的项合并，由\cref{g5:lem:primary-ideal-properties}，合并后的交仍是同一个素理想对应的准素理想。再删去不影响交的冗余项，有限次后便得到所需分解。
\end{proof}

这个定理保证分解存在。分解中的根式与实际准素理想承担不同信息：根式指出分量所在的闭集，准素理想还保留闭集上的幂次。下一定理分别判断它们的唯一性；对应结论可参见\cite[第6节，定理6.8]{Matsumura1989}。

\begin{theorem}\label{g5:thm:primary-uniqueness}
	设 \(A\) 为交换含幺 Noether 环，\(I\subsetneq A\)，并取一个极小准素分解
	\[
		I=\bigcap_{i=1}^rQ_i,\qquad \mathfrak p_i=\sqrt{Q_i}.
	\]
	则
	\[
		\operatorname{Ass}_A(A/I)=\{\mathfrak p_1,\ldots,\mathfrak p_r\}.
	\]
	因此分解中的素理想集与分量数由 \(I\) 唯一决定。若 \(\mathfrak p_i\) 为其中的极小元，则对应的孤立准素分量也唯一，且
	\[
		Q_i=\{a\in A:a/1\in IA_{\mathfrak p_i}\}.
	\]
\end{theorem}
\begin{proof}
	自然的对角映射
	\[
		A/I\longrightarrow\bigoplus_{i=1}^r A/Q_i,\qquad
		a+I\longmapsto(a+Q_1,\ldots,a+Q_r)
	\]
	为单射，因为它的核由 \(\bigcap_iQ_i=I\) 决定。由\cref{g5:prop:ass-exact-sequence,g5:prop:primary-associated-singleton}，得到
	\[
		\operatorname{Ass}_A(A/I)\subseteq\bigcup_i\operatorname{Ass}_A(A/Q_i)
		=\{\mathfrak p_1,\ldots,\mathfrak p_r\}.
	\]
	为了证明每个 \(\mathfrak p_i\) 都实际出现，固定 \(i\)。不可删去 \(Q_i\) 意味着可选择
	\[
		c\in\bigcap_{j\ne i}Q_j\setminus Q_i,
	\]
	单项分解时把空交约定为 \(A\)。于是 \(c+I\) 非零，且
	\[
		\operatorname{Ann}_A(c+I)=I:c=Q_i:c.
	\]
	由\cref{g5:lem:primary-ideal-properties}，这个理想为 \(\mathfrak p_i\) 准素理想。因此循环子模 \(A(c+I)\cong A/(Q_i:c)\) 的关联素理想集恰为 \(\{\mathfrak p_i\}\)。子模的关联素理想属于原模，所以 \(\mathfrak p_i\in\operatorname{Ass}_A(A/I)\)，反向包含得证。

	现在设 \(\mathfrak p_i\) 极小。对每个 \(j\ne i\)，都有 \(\mathfrak p_j\nsubseteq\mathfrak p_i\)。取 \(s_j\in\mathfrak p_j\setminus\mathfrak p_i\)，则某个 \(s_j^{n_j}\in Q_j\)。它在 \(A_{\mathfrak p_i}\) 中为单位，故 \(Q_jA_{\mathfrak p_i}=A_{\mathfrak p_i}\)。有限交与局部化相容：若一个分式分别属于各扩张理想，对每一项选一个使分子落入该理想的分母，再乘起这有限多个分母，就使分子落入原来的交。因此
	\[
		IA_{\mathfrak p_i}
		=\bigcap_j Q_jA_{\mathfrak p_i}
		=Q_iA_{\mathfrak p_i}.
	\]
	最后由\cref{g5:lem:primary-ideal-properties}的收缩公式，\(Q_iA_{\mathfrak p_i}\) 收缩回 \(A\) 正是 \(Q_i\)。所给表达式仅涉及 \(I\) 与 \(\mathfrak p_i\)，所以孤立分量唯一。
\end{proof}

\begin{example}\label{g5:exm:nonunique-embedded-component}
	在任意域 \(k\) 上，令 \(R=k[x,y]\)、\(I=(x^2,xy)\)。同一理想有两份极小准素分解：
	\[
		I=(x)\cap(x^2,y)
		=(x)\cap(x^2,xy,y^2).
	\]
	这个例子也见\cite[第6节，习题6.6]{Matsumura1989}。逐一核验如下。\((x)\) 为素理想，而
	\[
		R/(x^2,y)\cong k[x]/(x^2),\qquad
		R/(x^2,xy,y^2)\cong k\oplus k\overline x\oplus k\overline y
	\]
	中，第二个商的理想 \((\overline x,\overline y)\) 平方为零。这两个商环中，常数项非零的元素都可用有限几何级数求逆；常数项为零的元素都幂零。因此它们的全部零因子均幂零，两个理想都为 \((x,y)\) 准素理想。

	再核对交。若 \(f=xg\in(x^2,y)\)，令 \(y=0\)，可知 \(g(x,0)\) 被 \(x\) 整除，所以 \(g\in(x,y)\)，于是 \(f\in(x^2,xy)\)。若 \(f=xg\in(x^2,xy,y^2)=(x,y)^2\)，则 \(f\) 没有一次项，迫使 \(g\) 的常数项为零，同样得到 \(g\in(x,y)\)。这证明两个交都包含于 \(I\)；反向包含从生成元直接可见。每份分解都不可删项：\(x\) 在第一项中却不在第二项中，而 \(y\) 或 \(y^2\) 分别给出第二项中不属于第一项的元素。

	两份分解的孤立分量均为 \((x)\)，关联素理想集均为 \(\{(x),(x,y)\}\)，但嵌入分量不同，因为 \(y\in(x^2,y)\) 而 \(y\notin(x^2,xy,y^2)\)。所以嵌入素理想由 \(I\) 确定，具体的嵌入准素分量可以随分解改变。
\end{example}

准素分解因而同时保留两种可用信息：由关联素理想定位零因子及嵌入部分，由孤立分量的局部化公式恢复各不可约分支上的方程。所有这些有限分解均建立在 Noether 条件上；证明中它分别保证零化子能够取极大元、素滤过能够终止，以及不可约理想证明中的零化子升链能够稳定。
% 纲要条目排为编号小节标题，并进入目录与书签；说明段仍保留于纲要注释中。

% !TeX root = ../../Book5.tex
% 纲要标题：### 2.5 局部检测与局部自由
% 纲要位置：几何学计划纲要.md，第 5102 行。
% 正文状态：初稿；尚未正式审核。

% 纲要内容开始
% 
% * 模消失和同态同构的检测
% * 有限生成投射模的局部刻画
% * 常秩、有限呈示及局部自由
% * 纤维维数不足以单独判定平坦的模型
% 
% ---
% 
% 纲要内容结束

\section{局部检测与局部自由}
\label{g5:sec:0205}

一个模可以在每个点附近用不同的生成元描述。要把这些描述用于整体问题，需要知道两件事：逐点为零是否足以推出整体为零，以及一点处的基能否在邻域内继续作为基。前一个问题由零化子和局部化解决；后一个问题还要控制生成元之间的关系，因此有限呈示条件十分重要。本节把这些局部信息与有限投射模联系起来，最后用一个带幂零元的基环说明只看纤维维数会丢掉什么。

\ProseParagraphBreak

以下环均为交换含幺环，模均为含幺模。用 \(\operatorname{Spec}A\)、\(D(f)\) 和 \(\kappa(\mathfrak p)\) 分别表示素谱、主开集和素理想处的剩余域，沿用\secref{g5:sec:0201}的约定。

\subsection{模消失和同态同构的检测}
\label{g5:subsec:0205:01}

若一个元素在 \(\mathfrak p\) 处变为零，就有某个不属于 \(\mathfrak p\) 的元素杀掉它。反过来，只要该元素原来不为零，其零化子就是一个真理想，可以放进某个极大理想中；在这个极大理想处，任何允许作分母的元素都不能杀掉它。

\begin{proposition}
\label{g5:prop:local-detection}
	设 \(A\) 为交换含幺环，\(M\) 为任意 \(A\) 模。下列条件等价：
	\begin{equivalences}
		\item \(M=0\)；
		\item 对每个素理想 \(\mathfrak p\subset A\)，有 \(M_{\mathfrak p}=0\)；
		\item 对每个极大理想 \(\mathfrak m\subset A\)，有 \(M_{\mathfrak m}=0\)。
	\end{equivalences}
	因此，对任意 \(A\) 线性同态 \(u:M\rightarrow N\)，它是单射、满射或同构，当且仅当对每个素理想处的局部化具有相应性质；也可以只在全部极大理想处检验。
\end{proposition}
\begin{proof}
	前两个蕴含由局部化定义及极大理想是素理想得到。反过来，若 \(m\in M\) 非零，则
	\[
		\operatorname{Ann}_A(m)=\{a\in A:am=0\}
	\]
	是真理想，包含在某个极大理想 \(\mathfrak m\) 中。如果 \(m/1=0\) 于 \(M_{\mathfrak m}\)，局部化的零元判据给出 \(s\notin\mathfrak m\) 且 \(sm=0\)，这与 \(s\in\operatorname{Ann}_A(m)\subseteq\mathfrak m\) 矛盾。所以一个非零模必有非零的极大理想局部化。

	由\cref{g5:thm:localization-exact}，局部化与核、余核相容。把模消失的结论分别用于 \(\ker u\) 和 \(\operatorname{coker}u\)，就得到单射和满射的判据；二者同时成立便是同构。
\end{proof}

这里没有有限生成假设。逐点局部化保留了每个元素被哪些分母杀掉的信息，足以检测整个模。同一方法还说明，两张同态相等可以逐点核验：对它们的差应用零映射判据即可。这些局部检测可对照\cite[第4节，定理4.6]{Matsumura1989}。

\ProseParagraphBreak

若把局部模再模去其极大理想，就得到纤维
\[
	M(\mathfrak p)\coloneqq M\otimes_A\kappa(\mathfrak p)
	\cong M_{\mathfrak p}/\mathfrak pM_{\mathfrak p}.
\]
\symidx{\(M(\mathfrak p)\)：模在 \(\mathfrak p\) 处的纤维}
对于有限生成 \(M\)，Nakayama 引理给出 \(M(\mathfrak p)=0\) 当且仅当 \(M_{\mathfrak p}=0\)，所以纤维也能检测消失，只检查极大理想处便够。这是\cite[第4节，定理4.8]{Matsumura1989}中的有限模版本。没有有限生成条件时则不同：作为 \(\mathbb Z\) 模的 \(\mathbb Q\) 非零，但对每个素数 \(p\)，有 \(\mathbb Q/p\mathbb Q=0\)。在同态的检测中，也应先检验局部模上的核与余核，不能把上述同构判据直接换成任意模的纤维同构判据。

\subsection{有限生成投射模的局部刻画}
\label{g5:subsec:0205:02}

自由模的一个有用性质是可以逐个提升基向量：若 \(E\rightarrow F\) 满射，则任意映射 \(A^n\rightarrow F\) 都能提升到 \(E\)。投射模把这一提升性质从选定的自由基中抽取出来。它允许模在整体上没有统一的基，而仍能在满射下分离出自身。

\begin{definition}
\label{g5:def:projective-module}
	设 \(A\) 为交换含幺环。若 \(A\) 模 \(P\) 满足：对每个满同态 \(q:E\rightarrow F\) 和每个同态 \(v:P\rightarrow F\)，存在同态 \(\widetilde v:P\rightarrow E\) 使 \(q\widetilde v=v\)，则称 \(P\) 为\term[toushemo]{投射模}[projective module]。
\end{definition}
\[
	\begin{tikzcd}[column sep=large,row sep=large]
		&P\arrow[d,"v"]\arrow[dl,dashed,"\widetilde v"']\\
		E\arrow[r,two heads,"q"']&F
	\end{tikzcd}
\]
图中所有箭头均为 \(A\) 线性映射。提升要求存在，通常不唯一。该定义与自由模直和项的描述可参看\cite[附录B，Projective and injective modules，第277页]{Matsumura1989}。

\begin{proposition}
\label{g5:prop:projective-dual-basis}
	设 \(A\) 为交换含幺环，\(P\) 为 \(A\) 模。下列条件等价：
	\begin{equivalences}
		\item \(P\) 有限生成且投射；
		\item 存在 \(A\) 模 \(Q\) 和非负整数 \(n\)，使 \(P\oplus Q\cong A^n\)；
		\item 存在有限多个 \(p_i\in P\) 和 \(\lambda_i\in\operatorname{Hom}_A(P,A)\)，使对所有 \(p\in P\) 都有
		\[
			p=\sum_{i=1}^n\lambda_i(p)p_i.
		\]
	\end{equivalences}
	此外，有限生成投射模必有限呈示。
\end{proposition}
\begin{proof}
	若 \(P\) 有限生成，选满同态 \(\pi:A^n\rightarrow P\)。投射性将 \(\operatorname{id}_P\) 提升为 \(s:P\rightarrow A^n\)，满足 \(\pi s=\operatorname{id}_P\)。于是
	\[
		A^n=s(P)\oplus\ker\pi,
	\]
	因为每个 \(a\in A^n\) 可写为 \(s\pi(a)+(a-s\pi(a))\)，两部分分别在这两个子模中，交为零。这证明第一条件蕴含第二条件。

	若 \(P\) 是 \(A^n\) 的直和项，给定 \(q:E\rightarrow F\) 和 \(v:P\rightarrow F\)，先用到 \(P\) 的投影把 \(v\) 扩成 \(A^n\rightarrow F\)，再逐个提升标准基的像，得到 \(A^n\rightarrow E\)。限制到 \(P\) 就是所需提升。\(P\) 也是 \(A^n\) 的商模，因而有限生成。

	在 \(\pi s=\operatorname{id}_P\) 中，令 \(p_i=\pi(e_i)\)，令 \(\lambda_i\) 为 \(s\) 的第 \(i\) 个坐标函数，就得到第三条件。反过来，由第三条件定义
	\[
		\pi:A^n\longrightarrow P,\quad e_i\longmapsto p_i,
		\qquad
		s:P\longrightarrow A^n,\quad
		p\longmapsto(\lambda_1(p),\ldots,\lambda_n(p)).
	\]
	则 \(\pi s=\operatorname{id}_P\)，重新得到直和分解。

	最后，\(Q\) 是 \(A^n\) 的商模，故有限生成。上述满同态 \(A^n\rightarrow P\) 的核同构于 \(Q\)，从而给出 \(P\) 的有限呈示。
\end{proof}

第三个条件称为\term[duioujibiaozhun]{对偶基标准}[dual basis criterion]。这里的 \(p_i\) 可以有关系，\(\lambda_i(p_j)\) 也未必在 \(i=j\) 时为 \(1\)、在 \(i\ne j\) 时为 \(0\)；所要求的是这些秩一算子的和等于恒等算子。若把 \(P\) 看成 \(A^n\) 的直和项，复合 \(s\pi:A^n\rightarrow A^n\) 的矩阵 \(E\) 满足 \(E^2=E\)，且其像就是 \(P\)。反过来，一个幂等矩阵也给出分解 \(A^n=\operatorname{im}E\oplus\operatorname{im}(1-E)\)。

\ProseParagraphBreak

局部环上的有限投射模可以选出真正的基。剩余向量空间先决定所需的基向量数；投射分裂随后使关系模也有限生成，从而能对关系再应用一次 Nakayama 引理。

\begin{proposition}
\label{g5:prop:projective-local-free}
	设 \((R,\mathfrak m)\) 为交换含幺局部环，\(P\) 为有限生成投射 \(R\) 模。则 \(P\) 是有限秩自由模，其秩为
	\[
		r=\dim_{R/\mathfrak m}(P/\mathfrak mP).
	\]
	因此，任意交换含幺环 \(A\) 上的有限生成投射模，在每个素理想处的局部化都是有限秩自由模。
\end{proposition}
\begin{proof}
	提升 \(P/\mathfrak mP\) 的一组基。由\cref{g5:cor:nakayama-generators}，这些提升给出满同态 \(\pi:R^r\rightarrow P\)。投射性使它分裂，所以 \(R^r\cong K\oplus P\)，其中 \(K=\ker\pi\) 是有限生成模。模去 \(\mathfrak m\) 后仍有直和分解
	\[
		(R/\mathfrak m)^r\cong K/\mathfrak mK\oplus P/\mathfrak mP.
	\]
	右侧第二项已为 \(r\) 维，故第一项为零。Nakayama 引理给出 \(K=0\)，于是 \(\pi\) 是同构。

	若 \(P\oplus Q\cong A^n\)，局部化得到 \(P_{\mathfrak p}\oplus Q_{\mathfrak p}\cong A_{\mathfrak p}^n\)，所以 \(P_{\mathfrak p}\) 仍有限生成且投射。再应用刚证的局部环结论即可。
\end{proof}

这里使用的是有限投射模的版本；其证明可对照\cite[第2节，定理2.5的有限模情形]{Matsumura1989}。下一小节将证明反向结论：有限呈示模若在每个点处自由，便是整体上的有限投射模。证明需要把点处的基扩展到开集，再将开集上的分裂拼成一个整体分裂。

\subsection{常秩、有限呈示及局部自由}
\label{g5:subsec:0205:03}

在一点选出的生成元可以扩展到邻域，这是\cref{g5:prop:generators-near-prime}的结论。但“生成”之外还要证明没有关系，才能得到自由基。下面利用有限呈示，把这项要求化为有限多个关系元同时消失。

\begin{lemma}
\label{g5:lem:local-basis-extension}
	设 \(A\) 为交换含幺环，\(M\) 为有限呈示 \(A\) 模，\(\mathfrak p\in\operatorname{Spec}A\)。若 \(M_{\mathfrak p}\) 是自由 \(A_{\mathfrak p}\) 模，则存在 \(f\notin\mathfrak p\) 使 \(M_f\) 是自由 \(A_f\) 模；给定的 \(M_{\mathfrak p}\) 的一组基可以在这样的邻域上扩展为一组基。
\end{lemma}
\begin{proof}
	因为 \(M\) 有限生成，\(M_{\mathfrak p}\) 的自由秩 \(r\) 有限。把给定基写成 \(\xi_i=m_i/s_i\)，其中 \(m_i\in M\)、\(s_i\notin\mathfrak p\)。乘以各自的单位 \(s_i\)，可知 \(m_i/1\) 也构成基。令 \(\varphi:A^r\rightarrow M\) 把标准基送到 \(m_i\)。由\cref{g5:prop:generators-near-prime}，存在 \(f_0\notin\mathfrak p\)，使 \(\varphi_{f_0}:A_{f_0}^r\rightarrow M_{f_0}\) 满射。

	有限呈示在局部化后仍有限呈示，因为可以逐项局部化一个有限关系矩阵。故由\cref{g5:prop:finite-presentation-kernel}，\(K=\ker\varphi_{f_0}\) 是有限生成 \(A_{f_0}\) 模。它在 \(\mathfrak pA_{f_0}\) 处的局部化为零，因为 \(\varphi_{\mathfrak p}\) 是基给出的同构。逐个杀掉 \(K\) 的有限生成元，再把所用分母相乘，得到 \(a\notin\mathfrak p\)，使 \(K\) 在进一步反演 \(a\) 后为零。具体说，这些分母在 \(A_{f_0}\) 中可写成 \(a_j/f_0^{e_j}\)，其中 \(a_j\notin\mathfrak p\)，取 \(a=\prod_j a_j\) 即可。由局部化正合性，\(\varphi_{f_0a}\) 是同构。

	最后令 \(f=f_0a\prod_i s_i\)。所有 \(s_i\) 在 \(A_f\) 中为单位，所以再将基 \(m_i/1\) 分别除以 \(s_i\)，就得到扩展原基 \(\xi_i\) 的自由基。
\end{proof}

这同时证明有限呈示模的自由点组成开集，参看\cite[第4节，定理4.10(ii)]{Matsumura1989}。证明中有限生成用于控制余核，有限呈示用于控制选定满射的核；两个有限性条件对应不同的步骤。

\ProseParagraphBreak

若模可以在覆盖素谱的主开集上分别选出有限自由基，就称其\term[jubuziyou]{局部自由}[locally free]。在本章的模语言中，这具体指一个主开覆盖 \(D(f_i)\) 和同构 \(M_{f_i}\cong A_{f_i}^{r_i}\)；秩 \(r_i\) 可以随开集改变。

\begin{theorem}
\label{g5:thm:finite-projective-local}
	设 \(A\) 为交换含幺环，\(M\) 为有限呈示 \(A\) 模。下列条件等价：
	\begin{equivalences}
		\item \(M\) 是有限生成投射模；
		\item 对每个素理想 \(\mathfrak p\)，\(M_{\mathfrak p}\) 是自由 \(A_{\mathfrak p}\) 模；
		\item 对每个极大理想 \(\mathfrak m\)，\(M_{\mathfrak m}\) 是自由 \(A_{\mathfrak m}\) 模；
		\item 存在有限主开覆盖 \(\operatorname{Spec}A=\bigcup_{i=1}^tD(f_i)\)，使每个 \(M_{f_i}\) 是有限秩自由 \(A_{f_i}\) 模。
	\end{equivalences}
\end{theorem}
\begin{proof}
	第一条件蕴含第二条件，已由\cref{g5:prop:projective-local-free}证明；第二条件自然蕴含第三条件。若第三条件成立，由\cref{g5:lem:local-basis-extension}，每个极大理想都有一个主开邻域，使 \(M\) 在其中自由。这些主开集的定义元素生成单位理想，否则某个极大理想会包含它们全部，与其已被覆盖矛盾。于是 \(1\) 是有限多个定义元素的线性组合，相应的有限多个主开集就覆盖整个素谱，得到第四条件。

	最后设第四条件成立，选满同态 \(\pi:F=A^n\rightarrow M\)。因为 \(M_{f_i}\) 自由，可提升其自由基，得到一个局部分裂
	\[
		s_i:M_{f_i}\longrightarrow F_{f_i},\qquad
		\pi_{f_i}s_i=\operatorname{id}_{M_{f_i}}.
	\]
	我们要把它变成整体映射，再组合各个局部分裂。由 \(M\) 有限呈示及\cref{g5:prop:hom-localization}，
	\[
		\operatorname{Hom}_A(M,F)_{f_i}
		\cong\operatorname{Hom}_{A_{f_i}}(M_{f_i},F_{f_i}),
	\]
	所以可以写 \(s_i=\tau_i/f_i^{a_i}\)，其中 \(\tau_i:M\rightarrow F\) 为整体同态。等式 \(\pi_{f_i}s_i=\operatorname{id}\) 表明
	\(\pi\tau_i-f_i^{a_i}\operatorname{id}_M\) 在反演 \(f_i\) 后为零。再对 \(\operatorname{Hom}_A(M,M)\) 使用同一局部化结论，找到 \(b_i\geq0\) 使
	\[
		\pi(f_i^{b_i}\tau_i)=f_i^{a_i+b_i}\operatorname{id}_M
		\quad\text{在整个 }M\text{ 上成立}.
	\]
	可再乘一次 \(f_i\)，使指数 \(N_i\) 为正。记相应调整后的同态为 \(\rho_i\)，于是 \(\pi\rho_i=f_i^{N_i}\operatorname{id}_M\)。

	由于 \(D(f_i^{N_i})=D(f_i)\) 仍覆盖素谱，存在 \(c_i\in A\) 使 \(\sum_i c_i f_i^{N_i}=1\)。因此
	\[
		s=\sum_i c_i\rho_i:M\longrightarrow F,\qquad
		\pi s=\sum_i c_i f_i^{N_i}\operatorname{id}_M=\operatorname{id}_M.
	\]
	这给出整体分裂。由\cref{g5:prop:projective-dual-basis}，\(M\) 是有限生成投射模。
\end{proof}

局部基不必在交集上逐项相同；上面拼合的是满射的分裂。有限呈示让局部同态可以清除分母成为整体同态，而主开覆盖让相应分母的幂产生单位理想。这一局部刻画也见于\cite[第7节，定理7.12]{Matsumura1989}。

\ProseParagraphBreak

在自由邻域上，纤维的维数就是自由基的项数。因此，有限投射模具有一个逐点可读的\term[zhihanshu]{秩函数}[rank function]：
\[
	\operatorname{rk}_M(\mathfrak p)
	\coloneqq\dim_{\kappa(\mathfrak p)}M(\mathfrak p).
\]
\symidx{\(\operatorname{rk}_M(\mathfrak p)\)：有限投射模在一点的秩}
若 \(M_f\cong A_f^r\)，则这个函数在 \(D(f)\) 上恒等于 \(r\)，所以秩函数局部常值。其每个取值的原像都是开集，补集也是其他取值的原像之并，因而也是闭集。特别地，若 \(\operatorname{Spec}A\) 非空且连通，秩在全空间恒定。

\begin{example}
\label{g5:exm:projective-varying-rank}
	设 \(k\) 为域，\(A=k\times k\)，并取幂等元 \(e=(1,0)\)。模 \(P=Ae\) 满足
	\[
		A=Ae\oplus A(1-e),
	\]
	所以它有限生成且投射。\(A\) 的两个素理想是 \(\mathfrak p_1=0\times k\) 和 \(\mathfrak p_2=k\times0\)：素理想必须包含 \(e,1-e\) 之一，因为它们的乘积为零，而不能同时包含二者。模 \(P\) 在第一点处为一维自由模，在第二点处为零模，秩分别是 \(1,0\)。这两个点都是开闭集，故秩函数仍局部常值，但不存在整数 \(r\) 使 \(P\cong A^r\)。
\end{example}

对于一般有限呈示模，纤维维数可能在特殊点增大。一个关系矩阵可以把这种变化的发生位置写成子式方程。

\begin{proposition}
\label{g5:prop:fiber-upper-semicontinuity}
	设 \(A\) 为交换含幺环，\(M\) 有有限呈示
	\[
		A^s\xrightarrow{H}A^n\longrightarrow M\longrightarrow0.
	\]
	记 \(H\) 的 \(q\) 阶子式生成的理想为 \(J_q(H)\)，不存在这种子式时规定 \(J_q(H)=0\)。则
	\[
		\nu_M(\mathfrak p)\coloneqq
		\dim_{\kappa(\mathfrak p)}M(\mathfrak p)
		=n-\operatorname{rank}_{\kappa(\mathfrak p)}H(\mathfrak p),
	\]
	且对 \(1\leq r\leq n\)，有
	\[
		\{\mathfrak p:\nu_M(\mathfrak p)\geq r\}
		=V(J_{n-r+1}(H)).
	\]
	因此纤维维数函数\term[shangbanlianxu]{上半连续}[upper semicontinuous]，即它大于等于给定整数的点集为闭集。
\end{proposition}
\begin{proof}
	由\cref{g5:prop:tensor-right-exact}，将呈示张量 \(\kappa(\mathfrak p)\) 后得到
	\[
		\kappa(\mathfrak p)^s\xrightarrow{H(\mathfrak p)}
		\kappa(\mathfrak p)^n\longrightarrow M(\mathfrak p)\longrightarrow0.
	\]
	向量空间的秩公式给出所述维数。它至少为 \(r\)，当且仅当矩阵的秩至多为 \(n-r\)，也就是所有 \(n-r+1\) 阶子式在 \(\kappa(\mathfrak p)\) 中为零。一个 \(A\) 中元素在 \(\kappa(\mathfrak p)\) 中为零，当且仅当它属于 \(\mathfrak p\)。所以这一条件恰是 \(J_{n-r+1}(H)\subseteq\mathfrak p\)。当 \(r\leq0\) 或 \(r>n\) 时，相应集合分别为全空间或空集，仍为闭集。
\end{proof}

例如 \(A=k[t]\)、\(M=A/(t)\) 有关系矩阵 \((t)\)。在 \(t\ne0\) 的点处，纤维为零；在 \((t)\) 处，纤维维数为 \(1\)。关系的系数在该点消失，使纤维中可见的关系减少，因此维数增大。有限生成模的纤维维数也上半连续：在一点选出至多 \(r\) 个生成元后，\cref{g5:prop:generators-near-prime}使它们在邻域内继续生成；结合 Nakayama 引理，便知 \(\{\nu_M\leq r\}\) 为开集。有限呈示则进一步给出上面的有限矩阵与子式公式。

\subsection{纤维维数不足以单独判定平坦的模型}
\label{g5:subsec:0205:04}

局部自由模在每个邻域都像固定数目的独立坐标，因此与它作张量积会保留单射。平坦性抽取的正是这个性质：变换系数时，原来嵌在另一个模中的关系不能突然被压成零。\secref{g5:sec:0202}已经说明张量积总是右正合，因而还需要控制的是左端的单射。

\begin{definition}
\label{g5:def:flat-module}
	设 \(A\) 为交换含幺环。若 \(A\) 模 \(M\) 满足：对每个单同态 \(L\hookrightarrow N\)，诱导映射
	\[
		L\otimes_A M\longrightarrow N\otimes_A M
	\]
	仍为单射，则称 \(M\) 为\term[pingtanmo]{平坦模}[flat module]。等价地，函子 \(-\otimes_A M\) 保持短正合列。
\end{definition}

自由模 \(A^r\) 平坦，因为张量后只是取 \(r\) 份原映射的直和。有限投射模也平坦：若 \(P\oplus Q\cong A^r\)，那么与 \(P\) 张量后的映射是与 \(A^r\) 张量后所得单射的一个直和分量。局部化 \(A_f\) 作为 \(A\) 模同样平坦，这是张量积与局部化相容以及局部化正合性的直接结果。

\begin{example}
\label{g5:exm:constant-fiber-not-flat}
	设 \(k\) 为域，\(A=k[\varepsilon]/(\varepsilon^2)\)，\(M=A/(\varepsilon)\)。素谱只有 \(\mathfrak m=(\varepsilon)\) 一个点，剩余域为 \(k\)，且
	\[
		M\otimes_A k\cong k.
	\]
	所以纤维维数在整个基上恒为 \(1\)。然而 \(M\) 不平坦。取理想 \(I=(\varepsilon)\) 的包含 \(I\hookrightarrow A\)，张量 \(M\) 后得到
	\[
		I\otimes_A M\cong I/\varepsilon I=I
		\longrightarrow A\otimes_A M\cong M.
	\]
	左端非零，映射却把 \(\varepsilon\otimes\overline1\) 送到 \(\varepsilon\overline1=0\)，因而不是单射。纤维只看到剩余域，丢掉了基环中的非零幂零方向；平坦性同时检验这个方向与模之间的关系。
\end{example}

这个模甚至有限呈示：\(A\xrightarrow{\varepsilon}A\rightarrow M\rightarrow0\) 就是一个有限呈示。因此在允许非约化基环时，有限呈示与纤维维数局部常值合在一起，仍不能替代平坦性。要理解平坦与局部自由的联系，应直接分析张量积对关系模的作用。

\begin{proposition}
\label{g5:prop:flat-local-detection}
	设 \(A\) 为交换含幺环，\(M\) 为任意 \(A\) 模。则 \(M\) 平坦，当且仅当每个 \(M_{\mathfrak p}\) 都是平坦 \(A_{\mathfrak p}\) 模。也可以只检查全部极大理想。
\end{proposition}
\begin{proof}
	若 \(M\) 平坦，取 \(A_{\mathfrak p}\) 模的单射 \(L\hookrightarrow N\)，将其视为 \(A\) 模的单射。对任意 \(A_{\mathfrak p}\) 模 \(L\)，有自然同构
	\[
		L\otimes_A M\cong L\otimes_{A_{\mathfrak p}}M_{\mathfrak p},
	\]
	因为 \(A\setminus\mathfrak p\) 中的元素已经在 \(L\) 上可逆，纯张量的分母可以移到 \(L\) 一侧。于是 \(M\) 的平坦性给出 \(M_{\mathfrak p}\) 的平坦性。

	反过来，设所有极大理想处的局部模都平坦。对任意单射 \(L\hookrightarrow N\)，记诱导映射 \(L\otimes_A M\rightarrow N\otimes_A M\) 的核为 \(K\)。将它在极大理想 \(\mathfrak m\) 处局部化，利用局部化的正合性及其与张量积相容，得到
	\[
		K_{\mathfrak m}
		=\ker\bigl(L_{\mathfrak m}\otimes_{A_{\mathfrak m}}M_{\mathfrak m}
		\longrightarrow N_{\mathfrak m}\otimes_{A_{\mathfrak m}}M_{\mathfrak m}\bigr)=0.
	\]
	最后一个等号来自局部化仍保持原单射，以及 \(M_{\mathfrak m}\) 平坦。由\cref{g5:prop:local-detection}，\(K=0\)，故 \(M\) 平坦。
\end{proof}

下面的证明用一个具体等式控制关系。若 \(0\rightarrow K\rightarrow F\rightarrow M\rightarrow0\) 正合，平坦性保证 \(K\) 与 \(IF\) 的交没有比 \(IK\) 更多的元素。取局部环的极大理想作为 \(I\)，就可以把“极小生成元”提升为“自由基”。

\begin{theorem}
\label{g5:thm:finite-presented-flat}
	设 \(A\) 为交换含幺环，\(M\) 为有限呈示 \(A\) 模。则以下条件等价：\(M\) 平坦；\(M\) 局部自由；\(M\) 是有限生成投射模。
\end{theorem}
\begin{proof}
	由\cref{g5:thm:finite-projective-local}，后两个条件等价；有限投射模平坦已经在定义后证明。因此只需证明，有限呈示平坦模在每个素理想处自由。

	先设 \(M\) 为任意环上的平坦模，且 \(0\rightarrow K\rightarrow F\xrightarrow{\pi}M\rightarrow0\) 正合。对任意理想 \(I\subseteq A\)，我们证明
	\begin{equation}
	\label{g5:eq:flat-kernel-intersection}
		K\cap IF=IK.
	\end{equation}
	包含 \(IK\subseteq K\cap IF\) 直接成立。若 \(x\in K\cap IF\)，写 \(x=\sum_j a_jv_j\)，其中 \(a_j\in I\)、\(v_j\in F\)。张量
	\(z=\sum_j a_j\otimes v_j\in I\otimes_A F\)
	在 \(I\otimes_A M\) 中的像，经乘法映射到 \(M\) 后为 \(\pi(x)=0\)。平坦性使 \(I\otimes_A M\rightarrow A\otimes_A M=M\) 为单射，所以 \(z\) 在 \(I\otimes_A M\) 中的像已经为零。再用张量积的右正合性，
	\[
		I\otimes_A K\longrightarrow I\otimes_A F
		\longrightarrow I\otimes_A M\longrightarrow0,
	\]
	可知 \(z\) 来自 \(I\otimes_A K\)。对这些张量作标量乘法，就得到 \(x\in IK\)，证明了\eqref{g5:eq:flat-kernel-intersection}。

	现在固定 \(\mathfrak p\)，令 \(R=A_{\mathfrak p}\)、\(\mathfrak m=\mathfrak pA_{\mathfrak p}\)、\(P=M_{\mathfrak p}\)。模 \(P\) 仍有限呈示，并由\cref{g5:prop:flat-local-detection}知它平坦。提升 \(P/\mathfrak mP\) 的一组基，得到满同态 \(R^r\rightarrow P\)。记其核为 \(K\)。由\cref{g5:prop:finite-presentation-kernel}，\(K\) 有限生成；由\cref{g5:prop:minimal-presentation-local}，\(K\subseteq\mathfrak mR^r\)。把\eqref{g5:eq:flat-kernel-intersection}用于 \(I=\mathfrak m\)、\(F=R^r\)，便得到
	\[
		K=K\cap\mathfrak mR^r=\mathfrak mK.
	\]
	Nakayama 引理使 \(K=0\)，因此 \(P\cong R^r\)。这对所有素理想成立，再由\cref{g5:thm:finite-projective-local}得到结论。
\end{proof}

平坦性的局部检测与有限呈示模的上述等价关系可对照\cite[第7节，定理7.1、7.10、7.12及其推论]{Matsumura1989}。在\cref{g5:exm:constant-fiber-not-flat}中，极小满射 \(A\rightarrow M\) 的核为 \(K=(\varepsilon)\)，虽然 \(K\subseteq\mathfrak mA\)，却有 \(\mathfrak mK=0\ne K\)。这恰好指出平坦证明在哪里失败：纤维中看不见的关系 \(\varepsilon\overline1=0\) 仍保留在局部模中。

% 纲要条目排为编号小节标题，并进入目录与书签；说明段仍保留于纲要注释中。


\begin{chaptersummary}
	局部化由可用分母决定，其等价关系保留零因子造成的消失现象。素理想处的局部化是局部环，剩余域把该点的函数值组织成一个域；理想的扩缩则把局部化与饱和联系起来。素谱上的主开集为这些运算提供了邻域语言。

	\ProseParagraphBreak

	有限生成模只要求有限个生成方向，有限呈示还要求有限条生成关系。局部化保持正合，张量积一般只保持右正合，Hom 一般只保持左正合。有限呈示使 Hom 与局部化相容，因而允许把局部同态表示为有共同分母的整体同态。

	\ProseParagraphBreak

	Nakayama 引理把有限模的局部生成化为剩余域上的线性代数。支集说明模在哪些素理想处非零，关联素理想通过元素的零化子识别零因子和嵌入分支；Noether 环上的有限模只有有限个关联素理想，理想则能分解为有限个准素理想的交。

	\ProseParagraphBreak

	模是否为零、同态是否为同构，都可以在所有极大理想处检测。对于有限呈示模，逐点自由可以扩展成有限个主开集上的自由表示，并拼出整体投射结构。秩在局部自由情形下局部常值；仅有恒定的纤维维数仍可能遗漏幂零方向中的关系，平坦性正要求张量积保留原有的单射。
\end{chaptersummary}

\BookExercises

\begin{exercise}\label{g5:ex:localize-twelve}
	设 \(A=\mathbb Z/12\mathbb Z\)。分别计算 \(A_{\overline2}\) 和 \(A_{\overline6}\)，写出自然同态的核，并说明在第一个环中为何 \(\overline3\) 变为零而 \(\overline4\) 没有变为零。
\end{exercise}
\begin{solution}
	模三同态把 \(\overline2\) 送到单位，因而诱导 \(A_{\overline2}\to\mathbb F_3\)。在 \(A_{\overline2}\) 中， \(4\cdot3=0\)，且 \(4\) 为单位，所以 \(3=0\)。 因而自然同态经 \(\mathbb F_3\) 分解，两个方向互逆，得到 \(A_{\overline2}\cong\mathbb F_3\)，其核为 \((\overline3)\)。 因此 \(\overline3\) 变为零，而 \(\overline4\) 对应 \(\overline1\in\mathbb F_3\)，没有变为零。 直接用消失判据也可检查：\(4\cdot3=0\)，但 \(2^n4\) 从不被 \(12\) 整除。

	又因 \(\overline6^{\,2}=0\)，允许的分母中含零， \(A_{\overline6}=0\)，自然同态的核为整个 \(A\)。
\end{solution}

\begin{exercise}\label{g5:ex:localization-finite-intersection}
	设 \(A\) 为交换含幺环，\(S\subseteq A\) 为乘法集，\(N_1,N_2\) 为 \(A\)-模 \(M\) 的子模。证明
	\[
	S^{-1}(N_1\cap N_2)
	=(S^{-1}N_1)\cap(S^{-1}N_2)
	\]
	于 \(S^{-1}M\) 中成立。 再以 \(A=\mathbb Z\)、\(S=\{2^n:n\geq0\}\)、\(N_j=2^j\mathbb Z\) 说明这个等式不能直接推广为无限交。
\end{exercise}
\begin{solution}
	左侧包含于右侧。若 \(m/s\) 属于右侧，分式等价关系给出 \(u,v\in S\) 使 \(um\in N_1\)、\(vm\in N_2\)。 于是 \(uvm\in N_1\cap N_2\)，且 \(m/s=uvm/(uvs)\)，得到反向包含。 对给定无限族，\(\bigcap_{j\geq0}2^j\mathbb Z=0\)，所以先取交再局部化得零； 但每个 \(S^{-1}N_j\) 都等于 \(\mathbb Z[1/2]\)，局部化后再取交仍为该非零环。
\end{solution}

\begin{exercise}\label{g5:ex:tensor-hom-cyclic}
	设 \(m,n\) 为正整数。计算
	\[
	(\mathbb Z/m\mathbb Z)\otimes_{\mathbb Z}(\mathbb Z/n\mathbb Z),
	\qquad
	\operatorname{Hom}_{\mathbb Z}(\mathbb Z/m\mathbb Z,\mathbb Z/n\mathbb Z),
	\]
	并说明两者虽有同构的抽象群，所表示的数据却不同。
\end{exercise}
\begin{solution}
	令 \(d=\gcd(m,n)\)。由 \((\mathbb Z/m)\otimes N\cong N/mN\)，张量积为 \(\mathbb Z/(m,n)\cong\mathbb Z/d\mathbb Z\)，由 \(\overline1\otimes\overline1\) 生成，记录双线性配对的泛对象。

	一个同态由 \(\overline1\) 的像 \(b\in\mathbb Z/n\mathbb Z\) 唯一决定，条件是 \(mb=0\)。 写 \(m=dm'\)、\(n=dn'\)，其中 \(\gcd(m',n')=1\)，则 \(n\mid mb\) 等价于 \(n'\mid b\)。 因而 Hom 群由 \(n'\) 的剩余类生成，共有 \(d\) 个元素，也同构于 \(\mathbb Z/d\mathbb Z\)。 前者是商模，后者是被 \(m\) 杀掉的子模；两者在自然性与正合性中的作用仍然不同。
\end{solution}

\begin{exercise}\label{g5:ex:ideal-two-variable-presentation}
	设 \(k\) 为域，\(A=k[x,y]\)，\(I=(x^2,xy,y^2)\)。证明下列序列正合：
	\[
	0\longrightarrow A^2
	\xrightarrow{\ D\ } A^3
	\xrightarrow{\ (x^2\;\;xy\;\;y^2)\ } I
	\longrightarrow0,\qquad
	D=\begin{pmatrix}-y&0\\x&-y\\0&x\end{pmatrix}.
	\]
	并计算 \(I\otimes_A A/(x,y)\) 的维数，比较它与 \(I\otimes_A k(x,y)\) 的维数。
\end{exercise}
\begin{solution}
	右侧满性来自 \(I\) 的定义，两个箭头复合为零。 若 \(ax^2+bxy+cy^2=0\)，令 \(x=0\)，得 \(c(0,y)y^2=0\)，所以 \(c=x\beta\)。 在整环中约去非零 \(x\)，得到 \(ax+by+\beta y^2=0\)； 再令 \(y=0\)，可知 \(a=-y\alpha\) 对某个 \(\alpha\in A\) 成立。 代回并约去 \(y\)，得到 \(b=x\alpha-y\beta\)。 因而全部关系都由 \(D\) 的两列生成。 若 \(D(\alpha,\beta)=0\)，第一和第三个坐标分别给出 \(y\alpha=0\)、\(x\beta=0\)，故 \(\alpha=\beta=0\)，左侧单射。

	张量积 \(I\otimes_A A/(x,y)\cong I/(x,y)I\) 的基是 \(x^2,xy,y^2\) 的像，维数为三；其独立性由比较二次齐次项得到。 在 \(k(x,y)\) 中 \(x^2\) 为单位，故扩张理想等于整个域，维数为一。 生成方向的数量在原点增加，说明有限呈示模的纤维维数可以跳变。
\end{solution}

\begin{exercise}\label{g5:ex:hom-nonsplit-cyclic}
	设整数 \(n\geq2\)。证明短正合列
	\[
	0\longrightarrow\mathbb Z\xrightarrow{\ \times n\ }\mathbb Z
	\longrightarrow\mathbb Z/n\mathbb Z\longrightarrow0
	\]
	不分裂。分别用 Hom 和张量积指出其中一个满射或单射在相应运算后丢失。
\end{exercise}
\begin{solution}
	若有截面 \(\mathbb Z/n\to\mathbb Z\)，\(\overline1\) 的像被 \(n\) 杀掉，故只能为零，不能映回 \(\overline1\)，所以不分裂。 应用 \(\operatorname{Hom}_{\mathbb Z}(\mathbb Z/n,-)\) 后， \(\operatorname{Hom}(\mathbb Z/n,\mathbb Z)=0\)，而 \(\operatorname{Hom}(\mathbb Z/n,\mathbb Z/n)\) 含非零恒等同态； 因此最后的箭头不满。 应用 \(-\otimes_{\mathbb Z}\mathbb Z/n\) 后，原来的乘 \(n\) 单射变成 \(\mathbb Z/n\xrightarrow{0}\mathbb Z/n\)，不再单射。
\end{solution}

\begin{exercise}\label{g5:ex:surjective-finite-endomorphism}
	设 \(A\) 为交换含幺环，\(M\) 为有限生成 \(A\)-模。 证明每个满的 \(A\)-线性同态 \(\varphi:M\to M\) 都是同构。 提示：令多项式变量 \(T\) 通过 \(\varphi\) 作用于 \(M\)，对 \(A[T]\) 和理想 \((T)\) 使用\cref{g5:prop:nakayama-ideal}。
\end{exercise}
\begin{solution}
	赋予 \(M\) 以 \(A[T]\)-模结构 \(f(T)m=f(\varphi)m\)。原来的有限个 \(A\)-生成元也生成这个 \(A[T]\)-模。 满性给出 \(TM=M\)，故存在 \(Tq(T)\in(T)\) 使 \((1-Tq(T))M=0\)。 换成算子就是 \(\operatorname{id}_M=\varphi q(\varphi)\)。 多项式 \(q(\varphi)\) 与 \(\varphi\) 交换，因此也是左逆； \(\varphi\) 为同构，逆为 \(q(\varphi)\)。 这个论证不要求 \(A\) Noether。
\end{solution}

\begin{exercise}\label{g5:ex:generators-after-residue}
	设 \((A,\mathfrak m)\) 为交换含幺局部环，\(k=A/\mathfrak m\)， \(u:M\to N\) 为有限生成 \(A\)-模之间的同态。 证明 \(u\) 满当且仅当 \(u\otimes_A k:M\otimes_Ak\to N\otimes_Ak\) 满。 再举例说明：剩余域上的同态为同构，不能保证 \(u\) 为同构。
\end{exercise}
\begin{solution}
	令 \(Q=\operatorname{coker}u\)。张量积的右正合性给出 \(\operatorname{coker}(u\otimes_Ak)\cong Q/\mathfrak mQ\)。 \(Q\) 为有限生成模，Nakayama 引理说明该商为零当且仅当 \(Q=0\)，即得满性判据。 取 \(A=k[t]_{(t)}\)、\(M=A\)、\(N=A/(t)\)，令 \(u\) 为商映射。 剩余域上的同态是 \(k\) 的恒等同态，然而 \(u\) 的核为非零理想 \((t)\)。
\end{solution}

\begin{exercise}\label{g5:ex:support-without-finiteness}
	作为 \(\mathbb Z\)-模，令 \(M=\bigoplus_{p\text{ 为素数}}\mathbb Z/p\mathbb Z\)。 求 \(\operatorname{Ann}_{\mathbb Z}M\) 和 \(\operatorname{Supp}_{\mathbb Z}M\)，证明其支集不是闭集，并指出有限模支集公式在这里为何不能使用。
\end{exercise}
\begin{solution}
	杀掉整个 \(M\) 的整数必须被所有素数整除，所以 \(\operatorname{Ann}M=0\)。 在 \((0)\) 处局部化时，每个元素仅有有限个非零分量，故被某个非零整数杀掉；因此 \(M_{(0)}=0\)。 在 \((p)\) 处，第 \(p\) 个分量仍为 \(\mathbb Z/p\)，而其余分量被可逆素数杀掉，故 \(M_{(p)}\cong\mathbb Z/p\)。 支集于是是所有非零素理想 \((p)\)，不含 \((0)\)。

	一个闭集若含所有 \((p)\)，其定义理想便包含于 \(\bigcap_p(p)=0\)，所以它只能是整个 \(\operatorname{Spec}\mathbb Z\)。 因而这个支集不闭，也不等于 \(V(\operatorname{Ann}M)=\operatorname{Spec}\mathbb Z\)。 \(M\) 不有限生成，因为任何有限个元素只能涉及有限多个直和分量；缺少有限性时，不能把逐个元素的消失分母统一起来。
\end{solution}

\begin{exercise}\label{g5:ex:primary-product-powers}
	设 \(k\) 为域，\(A=k[x,y]\)，\(I=(x^2y)\)。 证明
	\[
	I=(x^2)\cap(y)
	\]
	是不可删减的准素分解，求 \(\operatorname{Ass}_A(A/I)\)，并判断是否有嵌入素理想。
\end{exercise}
\begin{solution}
	若 \(x^2a\in(y)\)，令 \(y=0\) 得 \(x^2a(x,0)=0\)，所以 \(y\mid a\)，从而 \((x^2)\cap(y)=(x^2y)\)。 理想 \((y)\) 素。 对 \((x^2)\)，把非零多项式写成 \(x^rf\)、\(x\nmid f\)； 关于 \(x\) 的最低次数在乘法下相加，因为系数环 \(k[y]\) 是整环。 因而若 \(ab\in(x^2)\) 且 \(a\notin(x)\)，则 \(b\in(x^2)\)，这说明 \((x^2)\) 是 \((x)\)-准素的。

	两个分量都不可删，根式不同，故此分解最简。 由准素分解的关联素理想结论， \(\operatorname{Ass}_A(A/I)=\{(x),(y)\}\)。 也可直接看到 \(\operatorname{Ann}(\overline{xy})=(x)\)、 \(\operatorname{Ann}(\overline{x^2})=(y)\)。 两个素理想都极小地包含 \(I\)，所以没有嵌入素理想；\(x^2\) 的幂次记录沿一个分支的非约化结构，本身不产生嵌入分支。
\end{solution}

\begin{exercise}\label{g5:ex:all-fibres-zero-infinite}
	证明非零 \(\mathbb Z\)-模 \(M=\mathbb Q/\mathbb Z\) 满足
	\[
	M\otimes_{\mathbb Z}\kappa(\mathfrak p)=0
	\quad\text{对每个 }\mathfrak p\in\operatorname{Spec}\mathbb Z.
	\]
	再证明 \(M_{(p)}\ne0\) 对每个素数 \(p\) 成立，以区分纤维与局部化。
\end{exercise}
\begin{solution}
	零素理想的剩余域是 \(\mathbb Q\)。每个 \(M\) 的元素都被某个非零整数杀掉，因此 \(M\otimes\mathbb Q=0\)。 对 \((p)\)，剩余域是 \(\mathbb F_p\)，且 \(M\otimes\mathbb F_p\cong M/pM=0\)，因为任何有理数的剩余类都可除以 \(p\)。 但 \(1/p+\mathbb Z\) 在 \(M_{(p)}\) 中非零：若被某个 \(s\notin(p)\) 杀掉，就会有 \(s/p\in\mathbb Z\)，即 \(p\mid s\)，矛盾。 有限模的剩余域检测来自 Nakayama 引理，这个非有限生成模不满足其条件；对任意模成立的局部化检测仍然有效。
\end{solution}

\begin{exercise}\label{g5:ex:module-two-open-gluing}
	设 \(A\) 为交换含幺环，\(f,g\in A\) 满足 \((f,g)=A\)，\(M\) 为任意 \(A\)-模。 证明序列
	\[
	0\longrightarrow M\longrightarrow M_f\oplus M_g
	\xrightarrow{\ (u,v)\mapsto u|_{fg}-v|_{fg}\ }M_{fg}
	\longrightarrow0
	\]
	正合，其中第一个非零箭头是同时局部化。
\end{exercise}
\begin{solution}
	先注意 \((f^r,g^s)=A\) 对正整数 \(r,s\) 成立：若这个理想真包含于某个极大理想，素性便迫使该极大理想同时含 \(f,g\)，矛盾。 若 \(m\) 在 \(M_f,M_g\) 中都为零，存在 \(r,s\) 使 \(f^rm=g^sm=0\)；用 \(cf^r+dg^s=1\) 得 \(m=0\)，所以第一个箭头单射。

	若 \(u=m/f^r\)、\(v=n/g^s\) 在 \(M_{fg}\) 中相同，则某个 \(t\geq0\) 满足 \((fg)^t(g^sm-f^rn)=0\)。 取充分大 \(N\)，令 \(a=f^{N-r}m\)、\(b=g^{N-s}n\)，便有 \(u=a/f^N\)、\(v=b/g^N\) 且 \(g^Na=f^Nb\) 于 \(M\) 中成立。 选 \(c,d\) 使 \(cf^N+dg^N=1\)，置 \(w=ca+db\)。 在 \(M_f\) 中 \(b=g^Nu\)、\(a=f^Nu\)，所以 \(w=u\)；在 \(M_g\) 中同理 \(w=v\)。 这证明中间正合。

	任取 \(m/(fg)^N\in M_{fg}\)，仍选 \(cf^N+dg^N=1\)，有
	\[
	\frac{m}{f^Ng^N}=\frac{dm}{f^N}+\frac{cm}{g^N}.
	\]
	它是 \((dm/f^N,-cm/g^N)\) 的像，故最后一个箭头满。
\end{solution}

\begin{exercise}\label{g5:ex:idempotent-variable-rank}
	设 \(k\) 为域，\(A=k\times k\times k\)、\(e=(1,0,1)\)，\(P=Ae\)。 证明 \(P\) 是有限生成投射模，求其各点的秩，并证明它不是自由模。
\end{exercise}
\begin{solution}
	\(e^2=e\)，且 \(A=Ae\oplus A(1-e)\)，所以 \(P\) 是秩一有限自由模 \(A\) 的直和项，为有限生成投射模。 环 \(A\) 的素理想恰为三个坐标投影的核 \(\mathfrak m_1,\mathfrak m_2,\mathfrak m_3\)。 事实上三个坐标幂等元两两乘积为零、和为 \(1\)，在整环 \(A/\mathfrak p\) 中恰有一个像为 \(1\)，其余为零，从而给出上述三个核。 在三个点的局部化分别等于 \(k\)，而 \(e\) 的像分别为 \(1,0,1\)，故 \(P\) 的秩分别为 \(1,0,1\)。 每个点都是开点，所以秩局部常值；若 \(P\cong A^r\)，三个点的秩都应等于同一个 \(r\)，矛盾。
\end{solution}

\begin{exercise}\label{g5:ex:projective-cyclic-idempotent}
	设 \(A\) 为交换含幺环，\(I\subseteq A\) 为理想。 证明 \(A/I\) 为投射 \(A\)-模，当且仅当 \(I=Ae\) 对某个幂等元 \(e^2=e\) 成立。
\end{exercise}
\begin{solution}
	若 \(A/I\) 投射，商映射有截面，因而 \(A=I\oplus J\) 为 \(A\)-模的直和。令 \(\pi:A\to I\) 为对应投影，\(e=\pi(1)\)。 \(A\)-线性给出 \(\pi(a)=ae\)，故 \(I=Ae\)； 又因 \(\pi\) 在 \(I\) 上恒等，\(\pi(e)=e\)，即 \(e^2=e\)。

	反过来，若 \(e^2=e\)，则 \(A=Ae\oplus A(1-e)\)：两项之和含 \(1\)，且若 \(z=ae=b(1-e)\)，乘以 \(e\) 得 \(z=0\)。 因此 \(A/I\cong A(1-e)\) 是有限自由模 \(A\) 的直和项，为投射模。
\end{solution}

\begin{exercise}\label{g5:ex:reduced-constant-fibre}
	设 \(A\) 为约化交换含幺环，\(M\) 为有限呈示 \(A\)-模。 若函数 \(\mathfrak p\mapsto\dim_{\kappa(\mathfrak p)} (M\otimes_A\kappa(\mathfrak p))\) 局部常值，证明 \(M\) 局部自由。 说明本题与双数环上的恒定纤维维数反例为何不矛盾。
\end{exercise}
\begin{solution}
	固定 \(\mathfrak p\)，在一个主开邻域上令纤维维数恒为 \(r\)，并把环替换为这个主开集对应的局部化。 有限呈示写成 \(A^q\xrightarrow C A^n\to M\to0\)。 对每个素理想，矩阵 \(C\) 在剩余域上的秩都是 \(s=n-r\)。 在 \(\mathfrak p\) 处选一个不消失的 \(s\) 阶子式 \(d\)（\(s=0\) 时取 \(d=1\)），再局部化于 \(d\)。 可逆子块的行列消元把矩阵化为
	\[
	\begin{pmatrix}I_s&0\\0&B\end{pmatrix}.
	\]
	其构造是先把左上可逆块化为单位块，再减去其余行、列在这些主元上的分量。 在每个剩余域上总秩仍为 \(s\)，所以 \(B\) 的每个元素都落在每个素理想中。 由\cref{g5:prop:radical-prime-intersection}，这些元素是幂零元。 局部化保持约化性：若 \((a/t)^N=0\)，某个 \(u\) 使 \(ua^N=0\)，则 \((ua)^N=0\)，约化性给出 \(ua=0\)，故 \(a/t=0\)。 因而此处 \(B=0\)，得到 \(M_d\cong A_d^r\)。 每个点都有这样的邻域，故 \(M\) 局部自由。 双数环 \(k[\varepsilon]/(\varepsilon^2)\) 不约化；纤维中的消失只能说明剩余关系为幂零，不能迫使关系本身为零。
\end{solution}

\begin{exercise}\label{g5:ex:tensor-support-intersection}
	设 \(A\) 为交换含幺环，\(M,N\) 为有限生成 \(A\)-模。 证明
	\[
	\operatorname{Supp}_A(M\otimes_A N)
	=\operatorname{Supp}_A M\cap\operatorname{Supp}_A N.
	\]
\end{exercise}
\begin{solution}
	在每个素理想 \(\mathfrak p\) 处，局部化与张量积相容； 若 \(M_{\mathfrak p}\) 或 \(N_{\mathfrak p}\) 为零，张量积局部化也为零。 反过来，若二者都非零，Nakayama 引理保证 \(M\otimes_A\kappa(\mathfrak p)\) 和 \(N\otimes_A\kappa(\mathfrak p)\) 都是非零有限维向量空间。 它们的张量积维数为两个正维数之积，故非零。 又
	\[
	(M\otimes_A N)\otimes_A\kappa(\mathfrak p)
	\cong
	(M\otimes_A\kappa(\mathfrak p))
	\otimes_{\kappa(\mathfrak p)}
	(N\otimes_A\kappa(\mathfrak p)),
	\]
	所以 \((M\otimes_A N)_{\mathfrak p}\ne0\)。 两个包含合在一起得到等式。
\end{solution}

% !TeX root = ../../Book5.tex
% 纲要标题：## 第3章　整性、Noether 正规化与维数
% 纲要位置：几何学计划纲要.md，第 5111 行。
% 正文状态：初稿；尚未正式审核。

% 纲要内容开始
% 
% > 本章目标与证明范围：严格按整性工具、有限投影、维数比较的次序展开。普通正规化和 Noether 正规化的对象不同，章名相近不意味着同一过程。
% 
% 纲要内容结束

\chapter{整性、Noether 正规化与维数}
\label{g5:ch:03}
\BookChapterNumber{g5:ch:03}

\begin{chapterabstract}
	首一方程使一个新函数的高次幂能够用低次幂表示，因而把代数关系转化为模的有限性。本章从这个观察建立整扩张的基本性质，证明素理想的 lying-over、going-up、incomparability 及整闭基底上的 going-down 定理。随后在任意基域上证明 Noether 正规化，将有限型代数表示为多项式子代数上的有限模。素链定义的维数与代数无关参数的数目由此相接；主理想定理和广义高度定理进一步控制方程能够产生的高度。参数系、零维环与纤维计算把这些结论带回具体的局部模型。
\end{chapterabstract}

\begin{learninggoals}
	\begin{itemize}
		\item 用首一方程与有限模判据判断整性，证明其传递性，并区分整闭包与 Noether 正规化。
		\item 证明 lying-over、incomparability、going-up 和整闭基底上的 going-down，辨认各项假设的用途。
		\item 完成 Noether 正规化的消元证明，解释无限域上的线性代换及任意域上的幂次代换。
		\item 从素链定义高度和维数，证明主理想定理、广义高度定理、有限型整代数的维数公式及整扩张的维数保持。
		\item 构造局部参数系，比较维数、嵌入维数与有限长度，并计算有限映射及一般映射的纤维环。
	\end{itemize}
\end{learninggoals}

在平面曲线 \(y^2=x^3\) 上，给定 \(x\) 后，\(y\) 满足首一方程 \(T^2-x^3=0\)。因此坐标环中的每个函数都能写成 \(a(x)+b(x)y\)：两个生成方向就足以控制全部函数。这个现象既表达了投影到 \(x\) 轴的有限性，也提示了维数的来源。函数 \(x\) 可以自由变化，而 \(y\) 只在 \(x\) 上增加代数关系，所以曲线只有一个代数无关的参数。

\ProseParagraphBreak

换成双曲线 \(xy=1\)，投影到 \(x\) 轴时，每个非空纤维仍只有一个点，坐标环却是 \(k[x,x^{-1}]\)，不能由有限多个元素作为 \(k[x]\) 模生成。缺少 \(x=0\) 的纤维，正与无法用首一方程控制 \(x^{-1}\) 相呼应。若改取 \(t=x+y\)，则 \(x\) 满足 \(X^2-tX+1=0\)，有限性便恢复了。选择合适的投影方向，是 Noether 正规化要解决的问题。

\ProseParagraphBreak

本章先证明整扩张如何控制函数与素理想，再构造有限投影，最后比较参数数目与素链长度。第二章的局部化、有限模、行列式技巧和 Nakayama 引理将贯穿这些证明。一般代数论证允许任意基域；只有把闭点直接当作坐标取值时才另行假设基域代数闭。首一关系、素链以及局部切片分别记录不同的信息，后面的算例会把它们放在同一个环中比较。

% !TeX root = ../../Book5.tex
% 纲要标题：### 3.1 整扩张及有限性
% 纲要位置：几何学计划纲要.md，第 4821 行。
% 正文状态：未撰写；本文件目前是框架。

% 纲要内容开始
% 
% * 首一方程与有限模判据
% * 整性传递和整闭包
% * lying-over 与 incomparability
% * going-up 及整闭基底下的 going-down
% 
% 纲要内容结束

\section{整扩张及有限性}
\label{g5:sec:0301}

待撰写

\subsection{首一方程与有限模判据}
\label{g5:subsec:0301:01}
待撰写

\subsection{整性传递和整闭包}
\label{g5:subsec:0301:02}
待撰写

\subsection{lying-over 与 incomparability}
\label{g5:subsec:0301:03}
待撰写

\subsection{going-up 及整闭基底下的 going-down}
\label{g5:subsec:0301:04}
待撰写
% 纲要条目排为编号小节标题，并进入目录与书签；说明段仍保留于纲要注释中。

% !TeX root = ../../Book5.tex
% 纲要标题：### 3.2 Noether 正规化的证明
% 纲要位置：几何学计划纲要.md，第 4828 行。
% 正文状态：未撰写；本文件目前是框架。

% 纲要内容开始
% 
% * 代数依赖的变换消元
% * 无限基域的线性投影
% * 有限基域的非线性变换
% * 有限多项式子代数
% * 几何中的有限投影
% 
% 纲要内容结束

\section{Noether 正规化的证明}
\label{g5:sec:0302}

待撰写

\subsection{代数依赖的变换消元}
\label{g5:subsec:0302:01}
待撰写

\subsection{无限基域的线性投影}
\label{g5:subsec:0302:02}
待撰写

\subsection{有限基域的非线性变换}
\label{g5:subsec:0302:03}
待撰写

\subsection{有限多项式子代数}
\label{g5:subsec:0302:04}
待撰写

\subsection{几何中的有限投影}
\label{g5:subsec:0302:05}
待撰写
% 纲要条目排为编号小节标题，并进入目录与书签；说明段仍保留于纲要注释中。

% !TeX root = ../../Book5.tex
% 纲要标题：### 3.3 Krull 维数与高度
% 纲要位置：几何学计划纲要.md，第 5130 行。
% 正文状态：初稿；尚未正式审核。

% 纲要内容开始
% 
% * 素链、局部维数和高度
% * 主理想定理的证明
% * 广义高度定理
% * 有限型代数维数与超越次数
% * 整扩张的维数保持
% 
% 纲要内容结束

\section{Krull 维数与高度}
\label{g5:sec:0303}

一条不可约曲线可以包含点，一个不可约曲面又可以包含曲线。沿这样的严格包含逐次下降，最多能走多少步，是一种不依赖坐标选择的维数定义。在素谱中，不可约闭集写成 \(V(\mathfrak p)\)，而闭集的包含方向与素理想的包含方向相反。因此，维数问题可以转化为素理想链的长度问题。

\ProseParagraphBreak

上一节的正规化给出有限个代数无关的参数。本节将证明，这些参数的数目恰好等于素链所测出的维数。另一项核心问题是：用 \(r\) 个方程切出的不可约分支，能够具有多大的高度？回答它需要 \term[Krullweishu]{Krull 维数}[Krull dimension]及高度定理。\footnote{Wolfgang Krull（1899--1971），德国数学家，建立了交换环的维数理论，并对赋值论与理想理论作出重要贡献。} 维数的基本定义与有限型整代数的结论可对照\cite[第5节，第30--35页]{Matsumura1989}；高度定理的陈述可对照\cite[第13节，定理13.5]{Matsumura1989}。

\subsection{素链、局部维数和高度}
\label{g5:subsec:0303:01}

\begin{definition}\label{g5:def:krull-dimension}
	设 \(A\) 为交换含幺环。一条严格包含的有限\term[sulian]{素链}[chain of prime ideals]
	\[
	\mathfrak p_0\subsetneq\mathfrak p_1\subsetneq\cdots\subsetneq\mathfrak p_r
	\]
	的长度为 \(r\)，即严格包含的次数。所有素链长度的上确界称为 \(A\) 的 Krull 维数，记为 \(\dim A\)。对非零环，它属于 \(\{0,1,2,\ldots\}\cup\{\infty\}\)；约定零环的维数为 \(-\infty\)。
	对 \(\mathfrak p\in\operatorname{Spec}A\)，所有以 \(\mathfrak p\) 为末项的素链长度的上确界，称为 \(\mathfrak p\) 的\term[gaodu]{高度}[height]，记为 \(\operatorname{ht}_A\mathfrak p\)。对真理想 \(I\subsetneq A\)，再定义
	\[
	\operatorname{ht}_A I
	=\inf_{\mathfrak p\supseteq I}\operatorname{ht}_A\mathfrak p.
	\]
\end{definition}
\symidx{\(\dim A\)：环的 Krull 维数}
\symidx{\(\operatorname{ht}_A\mathfrak p\)：素理想的高度}

长度为零的链只有一个素理想。因此，非零环维数为零，恰表示它的素理想之间没有严格包含，也即每个素理想都是极大理想。对整环，这又等价于环本身为域：零理想是素理想，若它也是极大理想，商环就是原环。

\ProseParagraphBreak

这个定义也给出了 \(\operatorname{Spec}A\) 的拓扑维数。事实上，每个非空不可约闭集都是某个 \(V(\mathfrak p)\)。若 \(V(I)\) 不可约且 \(fg\in\sqrt I\)，则
\[
	V(I)\subseteq V(fg)=V(f)\cup V(g),
\]
不可约性使 \(V(I)\) 包含于其中一边，由\cref{g5:prop:radical-prime-intersection}便有 \(f\in\sqrt I\) 或 \(g\in\sqrt I\)，所以 \(\sqrt I\) 为素理想；反向已经在\secref{g5:subsec:0201:04}说明。再由根式与素理想的对应，素链恰好给出严格下降的非空不可约闭集链。

\begin{proposition}\label{g5:prop:dimension-localization}
	设 \(A\) 为非零交换含幺环，\(\mathfrak p\in\operatorname{Spec}A\)。则
	\begin{equation}\label{g5:eq:height-local-dimension}
	\begin{aligned}
	\operatorname{ht}_A\mathfrak p&=\dim A_{\mathfrak p},\\
	\dim A&=\sup_{\mathfrak q\in\operatorname{Spec}A}\operatorname{ht}_A\mathfrak q
	=\sup_{\mathfrak m\text{ 为极大理想}}\dim A_{\mathfrak m}.
	\end{aligned}
	\end{equation}
	若 \(I\subsetneq A\)，则 \(\dim(A/I)\leq\dim A\)。特别地，
	\[
	\dim\bigl(A/\sqrt{(0)}\bigr)=\dim A.
	\]
\end{proposition}
\begin{proof}
	由\cref{g5:thm:prime-localization}，\(A_{\mathfrak p}\) 的素理想与 \(A\) 中包含于 \(\mathfrak p\) 的素理想保序对应。\(A_{\mathfrak p}\) 是局部环，它的每条素链都能在末端接上极大理想 \(\mathfrak pA_{\mathfrak p}\)，所以其维数就是以该极大理想为末项的链长上确界。这给出第一个等式。\(A\) 中任意素链的末项都包含于某个极大理想，把这个极大理想接到链的末端便得到其余等式。
	商环 \(A/I\) 的素理想与包含 \(I\) 的素理想保序对应，故商环中的链也是原环中的链。每个素理想都包含 \(\sqrt{(0)}\)，所以商去幂零根时素链一条也没有删去，维数不变。
\end{proof}

\(\dim A_{\mathfrak p}\) 称为环在 \(\mathfrak p\) 处的\term[jubuweishu]{局部维数}[local dimension]。它计算的是到达 \(\mathfrak p\) 以前的素链；\(\dim(A/\mathfrak p)\) 则计算从 \(\mathfrak p\) 出发向上的素链，也就是闭集 \(V(\mathfrak p)\) 的维数。例如在 \(k[t]\) 中，零素理想处的局部环是 \(k(t)\)，局部维数为零，而其闭包是整个一维素谱。对 \(a\in k\)，闭点 \((t-a)\) 的高度为一，其闭包只有一个点，维数为零。

\ProseParagraphBreak

这里 \(\dim k[t]=1\) 只用了单变量多项式环是主理想整环：它的每个非零素理想都由不可约多项式生成，因此极大。对 \(k[x_1,\ldots,x_n]\)，目前先能写出长度为 \(n\) 的素链
\begin{equation}\label{g5:eq:coordinate-prime-chain}
	(0)\subsetneq(x_1)\subsetneq(x_1,x_2)
	\subsetneq\cdots\subsetneq(x_1,\ldots,x_n).
\end{equation}
每一步的商环仍是多项式整环，故这些理想确实为素理想。它给出维数的下界；证明没有更长的素链，则需要本节后面的定理。

\begin{proposition}\label{g5:prop:dimension-minimal-components}
	设 \(A\) 为非零 Noether 环，\(\mathfrak p_1,\ldots,\mathfrak p_s\) 为其全部极小素理想。则
	\[
	\dim A=\max_{1\leq i\leq s}\dim(A/\mathfrak p_i).
	\]
\end{proposition}
\begin{proof}
	\cref{g5:cor:minimal-support-associated}应用于有限模 \(A\)，说明极小素理想只有有限多个，而且每个素理想都包含其中一个。给定 \(A\) 中的任意素链，取包含于链首项的 \(\mathfrak p_i\)。整条链都包含 \(\mathfrak p_i\)，所以在 \(A/\mathfrak p_i\) 中给出同样长度的链。反之，每个商环中的素链都提升为 \(A\) 中的素链。对长度取上确界即得结论。
\end{proof}

\subsection{主理想定理的证明}
\label{g5:subsec:0303:02}

一个方程 \(f=0\) 的不可约分支对应于 \((f)\) 上的极小素理想。\term[zhulixiangdingli]{主理想定理}[principal ideal theorem]断言，在 Noether 环中，这些素理想的高度至多为一。证明的关键是先把所考察的分支局部化，使它成为唯一的闭点，再比较素理想的幂。下面先建立证明所需的下降链条件。

\begin{lemma}\label{g5:lem:local-radical-artinian-quotient}
	设 \((R,\mathfrak m)\) 为 Noether 局部环，\(J\subsetneq R\) 满足 \(\sqrt J=\mathfrak m\)。则存在 \(N\geq1\) 使 \(\mathfrak m^N\subseteq J\)，并且 \(R/J\) 的理想下降链稳定。
\end{lemma}
\begin{proof}
	取 \(\mathfrak m\) 的有限组生成元 \(u_1,\ldots,u_s\)。若 \(s=0\)，则 \(R\) 为域且 \(J=0\)，结论成立。否则，由 \(u_i\in\sqrt J\)，可取 \(e_i\geq1\) 使 \(u_i^{e_i}\in J\)。令
	\[
	N=1+\sum_{i=1}^s(e_i-1).
	\]
	任意总次数为 \(N\) 的 \(u_i\) 单项式都有一个指数至少为相应的 \(e_i\)，所以属于 \(J\)。这些单项式生成 \(\mathfrak m^N\)，故 \(\mathfrak m^N\subseteq J\)。

	令 \(\kappa=R/\mathfrak m\)，考察有限滤过
	\[
	F_0=R/J\supseteq F_1\supseteq\cdots\supseteq F_N=0,
	\qquad F_i=(\mathfrak m^i+J)/J.
	\]
	每个 \(F_i/F_{i+1}\) 都被 \(\mathfrak m\) 杀掉，又是有限 \(R\) 模，因而是有限维 \(\kappa\) 向量空间。给定 \(R/J\) 的理想下降链 \(L_1\supseteq L_2\supseteq\cdots\)，各
	\[
	\frac{(L_j\cap F_i)+F_{i+1}}{F_{i+1}}
	\subseteq F_i/F_{i+1}
	\]
	随 \(j\) 构成向量子空间下降链。有限维性及 \(i\) 的有限性保证这些链在同一个指标以后全部稳定。

	如果相邻两项 \(L'\supseteq L''\) 在每层具有相同的像，则它们相等。确实，由 \(F_N=0\) 开始向上归纳：若已知 \(L'\cap F_{i+1}=L''\cap F_{i+1}\)，对任意 \(v\in L'\cap F_i\)，同像性给出 \(w\in L''\cap F_i\) 使 \(v-w\in F_{i+1}\)，归纳假设又给出 \(v-w\in L''\)，于是 \(v\in L''\cap F_i\)。取 \(i=0\) 即得 \(L'=L''\)。因此原下降链也稳定。
\end{proof}

这个引理把一个环上的理想链，化成有限多个剩余域向量空间中的子空间链。以下证明会把它用于 \(J=(a)\)。为使模 \(a\) 后得到的关系能重新提升到原环，我们先把素理想的幂局部化再收缩；这样得到的理想对乘以 \(a\) 自动具有消去性质。

\begin{lemma}\label{g5:lem:principal-local-domain}
	设 \((R,\mathfrak m)\) 为 Noether 局部整环，\(a\in\mathfrak m\) 且 \(\sqrt{(a)}=\mathfrak m\)。则 \(\dim R\leq1\)。
\end{lemma}
\begin{proof}
	若 \(a=0\)，则整环中 \(\sqrt{(0)}=0\)，所以 \(\mathfrak m=0\)，\(R\) 为域。以下设 \(a\ne0\)。只需证明每个严格包含于 \(\mathfrak m\) 的素理想 \(\mathfrak p\) 都为零。这样的 \(\mathfrak p\) 不含 \(a\)，否则它包含 \(\sqrt{(a)}=\mathfrak m\)。对 \(n\geq1\)，令
	\[
	Q_n=\mathfrak p^nR_{\mathfrak p}\cap R.
	\]
	这些理想满足 \(Q_{n+1}\subseteq Q_n\)，而 \(a/1\) 在 \(R_{\mathfrak p}\) 中为单位，故
	\begin{equation}\label{g5:eq:symbolic-power-cancellation}
	ac\in Q_n\quad\Longrightarrow\quad c\in Q_n.
	\end{equation}
	由\cref{g5:lem:local-radical-artinian-quotient}，\(R/(a)\) 中的下降链
	\[
	\frac{Q_1+(a)}{(a)}\supseteq
	\frac{Q_2+(a)}{(a)}\supseteq\cdots
	\]
	稳定。因此存在 \(n\geq1\) 使 \(Q_n+(a)=Q_{n+1}+(a)\)。对 \(v\in Q_n\)，可写 \(v=w+ac\)，其中 \(w\in Q_{n+1}\)。于是 \(ac\in Q_n\)，由\eqref{g5:eq:symbolic-power-cancellation}得 \(c\in Q_n\)，所以
	\[
	Q_n=Q_{n+1}+aQ_n.
	\]

	商模 \(Q_n/Q_{n+1}\) 有限生成，因为 \(R\) Noether；它等于自身的 \(a\) 倍，从而等于自身的 \(\mathfrak m\) 倍。由\cref{g5:thm:nakayama-local}，这个商模为零，即 \(Q_n=Q_{n+1}\)。

	现在在 \(\mathfrak p\) 处局部化。局部化理想的收缩再扩张恢复原理想，故
	\[
	\mathfrak p^nR_{\mathfrak p}=Q_nR_{\mathfrak p}
	=Q_{n+1}R_{\mathfrak p}=\mathfrak p^{n+1}R_{\mathfrak p}.
	\]
	\(R_{\mathfrak p}\) 也是 Noether 局部环。对有限模 \(\mathfrak p^nR_{\mathfrak p}\) 再用一次 Nakayama 引理，得到 \(\mathfrak p^nR_{\mathfrak p}=0\)。若 \(b\in\mathfrak p\)，则 \((b/1)^n=0\)。局部化后的环仍是整环，而且 \(R\rightarrow R_{\mathfrak p}\) 单射，所以 \(b=0\)。这证明 \(\mathfrak p=0\)，从而素链长度至多为一。
\end{proof}

上述 \(Q_n\) 称为 \(\mathfrak p\) 的第 \(n\) 个\term[fuhaomi]{符号幂}[symbolic power]，通常记为 \(\mathfrak p^{(n)}\)。一般有 \(\mathfrak p^n\subseteq\mathfrak p^{(n)}\)，两者未必相等。这里需要符号幂，是因为 \(a\notin\mathfrak p\) 使 \(a\) 在局部化后成为单位，从而保证\eqref{g5:eq:symbolic-power-cancellation}。证明中的两次 Nakayama 引理分别消去 \(Q_n/Q_{n+1}\) 和 \(\mathfrak p^nR_{\mathfrak p}\)，作用的局部环也分别是 \(R\) 和 \(R_{\mathfrak p}\)。

\begin{theorem}[Krull 主理想定理]\label{g5:thm:principal-ideal}
	设 \(A\) 为 Noether 环，\(f\in A\)。若 \(\mathfrak p\) 是包含 \((f)\) 的极小素理想，则
	\[
	\operatorname{ht}_A\mathfrak p\leq1.
	\]
	特别地，若 \(A\) 为整环且 \(f\ne0\) 不是单位，则每个这样的 \(\mathfrak p\) 都有高度一。
\end{theorem}
\begin{proof}
	令 \(R=A_{\mathfrak p}\)、\(\mathfrak m=\mathfrak pA_{\mathfrak p}\)。局部化的素理想对应说明，\(\mathfrak m\) 是 \(R\) 中唯一包含 \(fR\) 的素理想，所以 \(\sqrt{fR}=\mathfrak m\)。若 \(\dim R\geq2\)，便存在素链
	\[
	\mathfrak q_0\subsetneq\mathfrak q_1\subsetneq\mathfrak m.
	\]
	商环 \(R/\mathfrak q_0\) 是 Noether 局部整环；\(f\) 的像 \(\overline f\) 所生成理想的根式等于 \(\mathfrak m/\mathfrak q_0\)，因为 \(fR\subseteq fR+\mathfrak q_0\subseteq\mathfrak m\)。\cref{g5:lem:principal-local-domain}给出 \(\dim(R/\mathfrak q_0)\leq1\)，却又有长度为二的素链
	\[
	(0)\subsetneq\mathfrak q_1/\mathfrak q_0
	\subsetneq\mathfrak m/\mathfrak q_0,
	\]
	矛盾。因此 \(\operatorname{ht}_A\mathfrak p=\dim R\leq1\)。在整环且 \(f\ne0\) 的情形，\(\mathfrak p\ne0\)，素链 \((0)\subsetneq\mathfrak p\) 又给出高度至少为一。
\end{proof}

极小性不可省去。例如 \(k[x,y]\) 中的 \((x,y)\) 包含主理想 \((x)\)，但它下方存在素链 \((0)\subsetneq(x)\subsetneq(x,y)\)，高度至少为二。主理想定理控制的是方程 \(x=0\) 的整个不可约分支，而非分支上每一个更小闭子集的高度。

\subsection{广义高度定理}
\label{g5:subsec:0303:03}

一个方程所产生的高度界，怎样推广到多个方程？关键是把任意一条素链变成商环中的素链，并控制能够损失的长度。下面的引理完成这一点；它也将说明局部环中逐次加入方程时，维数每次至多下降一。

\begin{lemma}\label{g5:lem:principal-dimension-drop}
	设 \((R,\mathfrak m)\) 为 Noether 局部环，\(a\in\mathfrak m\)。则
	\[
	\dim R\leq\dim\bigl(R/(a)\bigr)+1.
	\]
\end{lemma}
\begin{proof}
	给定一条以极大理想为末项的素链
	\[
	\mathfrak q_0\subsetneq\mathfrak q_1\subsetneq\cdots
	\subsetneq\mathfrak q_d=\mathfrak m.
	\]
	我们将构造一条全部包含 \(a\)、长度至少为 \(d-1\) 的素链。它在 \(R/(a)\) 中仍严格，便得到所需估计。

	由于 \(a\in\mathfrak m\)，存在最小指标 \(t\) 使 \(a\in\mathfrak q_t\)。如果 \(t=0\)，原链的每一项都包含 \(a\)，结论已经成立。以下设 \(t\geq1\)。令 \(\mathfrak p_t=\mathfrak q_t\)，再从 \(i=t-1\) 向 \(i=0\) 递归选择
	\[
	\mathfrak q_i+(a)\subseteq\mathfrak p_i
	\subseteq\mathfrak p_{i+1},
	\]
	其中 \(\mathfrak p_i\) 为 \(\mathfrak q_i+(a)\) 上的一个极小素理想。这样的选择总能完成：\(\mathfrak p_{i+1}\) 包含 \(\mathfrak q_i+(a)\)，而 Noether 商环 \(R/(\mathfrak q_i+(a))\) 中的每个素理想都包含一个极小素理想，见\cref{g5:cor:minimal-support-associated}。

	对 \(0\leq i\leq t-2\)，有 \(\mathfrak p_i\subsetneq\mathfrak p_{i+1}\)。若相等，因为 \(a\notin\mathfrak q_{i+1}\)，就有
	\[
	\mathfrak q_i\subsetneq\mathfrak q_{i+1}
	\subsetneq\mathfrak p_i.
	\]
	但 \(\mathfrak p_i/\mathfrak q_i\) 是 Noether 环 \(R/\mathfrak q_i\) 中主理想 \((\overline a)\) 上的极小素理想，\cref{g5:thm:principal-ideal}要求它的高度至多为一，与上面的长度二素链矛盾。

	因此
	\[
	\mathfrak p_0\subsetneq\cdots\subsetneq\mathfrak p_{t-1}
	\subseteq\mathfrak q_t\subsetneq\cdots\subsetneq\mathfrak q_d
	\]
	全部包含 \(a\)，且只有中间的 \(\mathfrak p_{t-1}\subseteq\mathfrak q_t\) 可能成为等号。删去可能的重复项后，所得素链的长度至少为 \(d-1\)。于是每个这样的 \(d\) 都满足
	\(d\leq\dim(R/(a))+1\)。任意素链都能接到 \(\mathfrak m\)，对链长取上确界即得结论；这个论证也涵盖维数为无穷的可能情形。
\end{proof}

证明中向下选择 \(\mathfrak p_i\)，是为了始终保留素理想之间的包含关系。主理想定理随后保证：原链在遇到包含 \(a\) 的项以前，相邻两层不会被合并。把这个逐链估计在商环中反复使用，就得到\term[Krullgaodudingli]{Krull 高度定理}[Krull's height theorem]。

\begin{theorem}[Krull 高度定理]\label{g5:thm:height-theorem}
	设 \(A\) 为 Noether 环，真理想 \(I=(f_1,\ldots,f_r)\) 由 \(r\) 个元素生成。若 \(\mathfrak p\) 是 \(I\) 上的极小素理想，则
	\[
	\operatorname{ht}_A\mathfrak p\leq r.
	\]
	因而 \(\operatorname{ht}_A I\leq r\)。允许 \(r=0\)，此时 \(I=(0)\)。
\end{theorem}
\begin{proof}
	在 \(\mathfrak p\) 处局部化，记 \(R=A_{\mathfrak p}\)、\(\mathfrak m=\mathfrak pA_{\mathfrak p}\)。由于 \(\mathfrak p\) 在 \(I\) 上极小，\(\mathfrak m\) 是 \(R\) 中唯一包含 \(IR\) 的素理想。因此 \(R/IR\) 的素谱只有一个点，维数为零。

	依次记
	\[
	R_j=R/(f_1,\ldots,f_j)R\qquad(0\leq j\leq r),
	\]
	其中 \(R_0=R\)。各 \(R_j\) 都是非零 Noether 局部环，因为这些生成元都属于 \(\mathfrak m\)。对 \(R_{j-1}\) 中 \(f_j\) 的像应用\cref{g5:lem:principal-dimension-drop}，得到
	\[
	\dim R_{j-1}\leq\dim R_j+1.
	\]
	从 \(\dim R_r=0\) 开始向前归纳，便有 \(\dim R\leq r\)。由\eqref{g5:eq:height-local-dimension}，这正是所求的高度界。\(r=0\) 时，\(R=R/IR\) 已为零维，同样成立。
\end{proof}

这个结论也称\term[guangyigaodudingli]{广义高度定理}[generalized principal ideal theorem]。主理想定理对应于 \(r=1\)，而 \(k[x_1,\ldots,x_r]\) 中的 \((x_1,\ldots,x_r)\) 表明上界可以取到：\eqref{g5:eq:coordinate-prime-chain}中的坐标素链给出高度至少为 \(r\)，高度定理给出高度至多为 \(r\)。

\begin{corollary}\label{g5:cor:noetherian-local-finite-dimension}
	设 \((R,\mathfrak m)\) 为 Noether 局部环，且 \(\mathfrak m\) 能由 \(r\) 个元素生成。则 \(\dim R\leq r<\infty\)。特别地，Noether 环的每个素理想都有有限高度。
\end{corollary}
\begin{proof}
	\(\mathfrak m\) 是自身上唯一的极小素理想，故高度定理给出
	\(\dim R=\operatorname{ht}_R\mathfrak m\leq r\)。对一般 Noether 环的素理想 \(\mathfrak p\)，局部化 \(A_{\mathfrak p}\) 仍为 Noether 环，其极大理想有限生成，再用同一结论即可。
\end{proof}

维数因此给出了方程数目的下界：若 \((R,\mathfrak m)\) 为维数 \(d\) 的 Noether 局部环，且 \(\sqrt{(f_1,\ldots,f_r)}=\mathfrak m\)，则 \(r\geq d\)。下一节将构造恰好 \(d\) 个这样的元素，并把它们称为参数系。

\subsection{有限型代数维数与超越次数}
\label{g5:subsec:0303:04}

上一节把有限型代数表示为多项式子代数上的有限模。对于整环，其参数数目由分式域的超越次数确定。为了把这个数与素链长度联系起来，先说明严格增加一个素理想会使超越次数严格下降。下面的论证适用于任意基域，所用的只是代数独立与有限生成。

\begin{lemma}\label{g5:lem:strict-prime-transcendence-drop}
	设 \(k\) 为域，\(B\) 为有限型 \(k\) 代数且为整环，\(\mathfrak q\ne(0)\) 为 \(B\) 的素理想。则
	\[
	\operatorname{trdeg}_k\operatorname{Frac}(B/\mathfrak q)
	<\operatorname{trdeg}_k\operatorname{Frac}(B).
	\]
\end{lemma}
\begin{proof}
	写 \(B=k[b_1,\ldots,b_s]\)。从 \(b_i\) 在 \(B/\mathfrak q\) 中的像里选取一个极大代数无关子集，重排指标后记为 \(\overline b_1,\ldots,\overline b_e\)。其余生成元都在 \(k(\overline b_1,\ldots,\overline b_e)\) 上代数，所以这些像构成 \(\operatorname{Frac}(B/\mathfrak q)\) 的一个超越基。\(b_1,\ldots,b_e\) 本身也代数无关：原环中的多项式关系在商环中仍成立。因此
	\[
	e\leq d=\operatorname{trdeg}_k\operatorname{Frac}(B).
	\]
	若 \(e=d\)，则 \(\operatorname{Frac}(B)\) 在 \(K=k(b_1,\ldots,b_e)\) 上代数；否则再加入一个超越于 \(K\) 的元素，就得到 \(d+1\) 个代数无关元，与\cref{g5:lem:transcendence-bound}矛盾。令
	\[
	S=k[b_1,\ldots,b_e]\setminus\{0\}.
	\]
	因为 \(\overline b_1,\ldots,\overline b_e\) 代数无关，\(S\cap\mathfrak q=\varnothing\)。另一方面，
	\[
	S^{-1}B=K[b_{e+1},\ldots,b_s]
	\]
	是由有限多个代数元生成的 \(K\) 代数，故是有限维 \(K\) 向量空间：对各生成元取代数方程并用其约去高次幂，就得到有限个单项式张成。它又是整环，所以是域；事实上，任意非零元素的乘法为有限维向量空间的单射，因而满射。
	但 \(S\cap\mathfrak q=\varnothing\) 使 \(S^{-1}\mathfrak q\) 为这个域的真理想，而 \(\mathfrak q\ne0\) 及整环的局部化单射性又使 \(S^{-1}\mathfrak q\ne0\)，矛盾。因此 \(e<d\)。
\end{proof}

该证明的要点是：若取素商后没有损失代数无关参数，就能同时把这些参数中的非零多项式变为单位。原环此时变成有限代数域扩张，原来非零的素理想却必须仍为真理想，两者无法同时成立。

\begin{theorem}[有限型代数的维数]\label{g5:thm:affine-dimension}
	设 \(k\) 为任意域，\(A\) 为非零有限型 \(k\) 代数。
	若 \(A\) 为整环，则
	\begin{equation}\label{g5:eq:affine-dimension-transcendence}
	\dim A=\operatorname{trdeg}_k\operatorname{Frac}(A).
	\end{equation}
	一般地，若 \(\mathfrak p_1,\ldots,\mathfrak p_s\) 为 \(A\) 的全部极小素理想，则
	\begin{equation}\label{g5:eq:affine-dimension-components}
	\dim A=\max_{1\leq i\leq s}
	\operatorname{trdeg}_k\operatorname{Frac}(A/\mathfrak p_i).
	\end{equation}
\end{theorem}
\begin{proof}
	先设 \(A\) 为整环，并记 \(d=\operatorname{trdeg}_k\operatorname{Frac}(A)\)。对任意素链
	\[
	\mathfrak q_0\subsetneq\mathfrak q_1\subsetneq\cdots\subsetneq\mathfrak q_r,
	\]
	各 \(A/\mathfrak q_i\) 都是有限型 \(k\) 代数，且都是整环。对相邻两项应用\cref{g5:lem:strict-prime-transcendence-drop}，得到一列严格下降的非负整数
	\[
	\operatorname{trdeg}_k\operatorname{Frac}(A/\mathfrak q_0)
	>\cdots>
	\operatorname{trdeg}_k\operatorname{Frac}(A/\mathfrak q_r)\geq0.
	\]
	首项至多为 \(d\)：若 \(\mathfrak q_0=0\) 则相等，若 \(\mathfrak q_0\ne0\) 则再用一次该引理。故 \(r\leq d\)，从而 \(\dim A\leq d\)。

	反向用\cref{g5:thm:noether-normalization,g5:cor:normalization-transcendence}：存在代数无关的 \(t_1,\ldots,t_d\in A\)，使
	\[
	C=k[t_1,\ldots,t_d]\subseteq A
	\]
	为有限整扩张。\(C\) 中有长度为 \(d\) 的坐标素链
	\[
	(0)\subsetneq(t_1)\subsetneq\cdots\subsetneq(t_1,\ldots,t_d).
	\]
	当 \(d=0\) 时，这条链只含 \((0)\) 一项。由\cref{g5:thm:lying-over}，可先在 \(A\) 中选取位于 \((0)\) 上方的素理想；再反复应用\cref{g5:thm:going-up}，便把整条链提升到 \(A\)。各提升理想的收缩互不相同，所以提升后的包含也都严格，长度仍为 \(d\)。因此 \(\dim A\geq d\)，证明\eqref{g5:eq:affine-dimension-transcendence}。

	一般的 \(A\) 由\cref{g5:thm:hilbert-basis,g5:prop:affine-noetherian}知为 Noether 环。\cref{g5:prop:dimension-minimal-components}给出
	\(\dim A=\max_i\dim(A/\mathfrak p_i)\)，再对每个整环 \(A/\mathfrak p_i\) 使用已证的有限型代数结论，就得到\eqref{g5:eq:affine-dimension-components}。
\end{proof}

\begin{corollary}\label{g5:cor:polynomial-dimension}
	对任意域 \(k\) 和整数 \(n\geq0\)，有
	\[
	\dim k[x_1,\ldots,x_n]=n.
	\]
\end{corollary}
\begin{proof}
	分式域 \(k(x_1,\ldots,x_n)\) 对 \(k\) 的超越次数为 \(n\)，应用\cref{g5:thm:affine-dimension}即可。
\end{proof}

代数闭域 \(k\) 上的非空仿射代数集 \(X\)，其维数定义为 \(\dim X=\dim k[X]\)。坐标环同构保持素链，因此这个定义与仿射实现无关。若 \(X\) 不可约，就得到 \(\dim X=\operatorname{trdeg}_k k(X)\)；一般情形则取各不可约分量维数的最大值。这样，坐标中独立参数的数目与不可约闭集链的长度给出了同一个维数。

\begin{example}\label{g5:exm:mixed-dimensional-components}
	设 \(k\) 为任意域，令
	\[
	A=k[x,y,z]/(xy,xz).
	\]
	在 \(k[x,y,z]\) 中有
	\[
	(xy,xz)=(x)\cap(y,z).
	\]
	右端包含左端。反之，若 \(xg\in(y,z)\)，模去 \((y,z)\) 得到 \(x\overline g=0\)；商环 \(k[x]\) 是整环，故 \(\overline g=0\)，即 \(g\in(y,z)\)。这证明交理想等式。两个素理想 \((x)\)、\((y,z)\) 互不包含，而且任何包含其交的素理想都包含其中一个，所以它们恰是交理想上的极小素理想。
	相应两个分支的商环分别为
	\[
	A/\bigl((x)/(xy,xz)\bigr)\cong k[y,z],
	\qquad
	A/\bigl((y,z)/(xy,xz)\bigr)\cong k[x].
	\]
	由\cref{g5:thm:affine-dimension}，两者维数分别为二和一，而 \(\dim A=2\)。当 \(k\) 为代数闭域时，相应代数集为平面 \(x=0\) 与 \(x\) 轴，交于原点。
	令 \(\mathfrak p=(y,z)/(xy,xz)\)。它是极小素理想，所以 \(\operatorname{ht}_A\mathfrak p=0\)，但 \(\dim(A/\mathfrak p)=1\)。于是
	\[
	\operatorname{ht}_A\mathfrak p+\dim(A/\mathfrak p)=1<2=\dim A.
	\]
	这个计算说明，整体维数取决于最高维的分支，某个素理想的高度则只计算到达该素理想的链。两者相减一般不能给出其闭包的维数。
\end{example}

\subsection{整扩张的维数保持}
\label{g5:subsec:0303:05}

有限型代数的维数证明已经用到：多项式子环中的素链可以沿整扩张提升。对一般整扩张，lying-over 与 going-up 保证底环中的链不会丢失，incomparability 保证上层的严格链收缩后仍然严格。这两点合起来给出维数相等。

\begin{theorem}\label{g5:thm:integral-dimension}
	设 \(A\subseteq B\) 为非零交换含幺环的整扩张。则
	\[
	\dim B=\dim A,
	\]
	其中允许维数为无穷。对任意 \(\mathfrak q\in\operatorname{Spec}B\)，令 \(\mathfrak p=\mathfrak q\cap A\)，还有
	\[
	\dim(B/\mathfrak q)=\dim(A/\mathfrak p).
	\]
	更一般地，对任意交换含幺环同态 \(\varphi:A\rightarrow B\)，若 \(B\) 在 \(\varphi(A)\) 上整，则
	\[
	\dim B=\dim(A/\ker\varphi).
	\]
\end{theorem}
\begin{proof}
	设 \(\mathfrak q_0\subsetneq\cdots\subsetneq\mathfrak q_r\) 为 \(B\) 中的素链。收缩得到 \(A\) 中的包含链。若相邻两项收缩相同，就与\cref{g5:thm:incomparability}矛盾。因此收缩后的每一步都严格，长度仍为 \(r\)，从而 \(\dim B\leq\dim A\)。

	反过来，给定 \(A\) 中任意素链 \(\mathfrak p_0\subsetneq\cdots\subsetneq\mathfrak p_r\)，先由\cref{g5:thm:lying-over}选取 \(\mathfrak q_0\in\operatorname{Spec}B\)，使 \(\mathfrak q_0\cap A=\mathfrak p_0\)。再由\cref{g5:thm:going-up}逐步得到
	\[
	\mathfrak q_0\subseteq\mathfrak q_1\subseteq\cdots\subseteq\mathfrak q_r,
	\qquad \mathfrak q_i\cap A=\mathfrak p_i.
	\]
	收缩互不相同，故这些包含全部严格。因此 \(\dim B\geq\dim A\)。对所有有限链的长度取上确界，两个不等式也涵盖无穷维情形。

	给定 \(\mathfrak q\) 后，诱导包含 \(A/\mathfrak p\subseteq B/\mathfrak q\) 仍为整扩张，对它使用已经证明的结论即得第二个等式。对于一般整同态，先用 \(A/\ker\varphi\cong\varphi(A)\) 代替底环，再对 \(\varphi(A)\subseteq B\) 使用结论。若 \(B=0\)，两边均为零环的维数 \(-\infty\)。
\end{proof}

这个维数等式适用于任意整扩张。整性通过素理想的上行与不可比性直接控制全部有限素链。若希望比较指定点处的高度，则还需要保留提升链的上端点；这正是 going-down 所处理的问题。

\begin{proposition}\label{g5:prop:integral-height-comparison}
	设 \(A\subseteq B\) 为非零交换含幺环的整扩张，\(\mathfrak q\in\operatorname{Spec}B\)，\(\mathfrak p=\mathfrak q\cap A\)。则
	\[
	\operatorname{ht}_B\mathfrak q\leq\operatorname{ht}_A\mathfrak p.
	\]
	若这个扩张还满足 going-down，则等号成立。特别地，若 \(A\) 是整闭整环、\(B\) 是整环，则等号成立。
\end{proposition}
\begin{proof}
	每条以 \(\mathfrak q\) 为末项的素链收缩后，依不可比性仍严格，且末项为 \(\mathfrak p\)，因此得到第一个不等式。若有 going-down，给定任意以 \(\mathfrak p\) 为末项的链
	\[
	\mathfrak p_0\subsetneq\cdots\subsetneq\mathfrak p_r=\mathfrak p,
	\]
	从指定的 \(\mathfrak q_r=\mathfrak q\) 开始，逐次向下选取 \(\mathfrak q_i\subseteq\mathfrak q_{i+1}\)，使其收缩为 \(\mathfrak p_i\)。收缩的严格性保证所选链也严格，于是 \(\operatorname{ht}_B\mathfrak q\geq r\)。取上确界即得反向不等式。最后的充分条件由\cref{g5:thm:going-down}给出。
\end{proof}

比较这两个证明可以看出端点的作用。证明整体维数相等时，先随意选定链首素理想的一个上方素理想，再向上提升；证明指定 \(\mathfrak q\) 的高度相等时，必须从这个已经给定的末项向下构造。虽然得到的都是同样长度的链，所需要的提升方向却不同。

\begin{example}\label{g5:exm:integral-dimension-cusp}
	设 \(k\) 为任意域。包含
	\[
	k[t^2,t^3]\subseteq k[t]
	\]
	为有限整扩张：\(t\) 满足首一方程 \(T^2-t^2=0\)，而 \(k[t]\) 由 \(1,t\) 作为 \(k[t^2,t^3]\) 模生成。因此
	\[
	\dim k[t^2,t^3]=\dim k[t]=1.
	\]
	把 \(u,v\) 分别映到 \(t^2,t^3\)，可得
	\[
	k[u,v]/(v^2-u^3)\cong k[t^2,t^3].
	\]
	确实，模去 \(v^2-u^3\) 后每个多项式可写为 \(a(u)+v b(u)\)；代入后，\(a(t^2)\) 只含偶次幂，\(t^3b(t^2)\) 只含奇次幂，两者之和为零便迫使 \(a=b=0\)。这既核对了核理想，也说明尖点曲线与其参数直线具有同一维数。
\end{example}

一般整同态中的核确实影响维数。例如满同态 \(k[x,y]\rightarrow k[y]\)、\(x\mapsto0\)、\(y\mapsto y\) 为整同态，每个目标元素都满足一个首一的一次方程。但源环维数为二、目标环维数为一，维数相等发生在
\[
	k[x,y]/(x)\cong k[y]
\]
与目标环之间。因而应用整扩张的维数保持时，先确定底环在目标环中的像，就确定了实际参与比较的空间。
% 纲要条目排为编号小节标题，并进入目录与书签；说明段仍保留于纲要注释中。

% !TeX root = ../../Book5.tex
% 纲要标题：### 3.4 维数的局部模型
% 纲要位置：几何学计划纲要.md，第 5138 行。
% 正文状态：初稿；尚未正式审核。

% 纲要内容开始
% 
% * 参数系及嵌入维数
% * 零维 Noether 环和有限长度
% * 超曲面与完全交的预告
% * 纤维维数问题的代数准备
% 
% ---
% 
% 纲要内容结束

\section{维数的局部模型}
\label{g5:sec:0304}

一条曲线在一个点附近需要几个参数，可能与描述这个点附近的一阶变化所需的坐标数不同。例如尖点的函数域只有一个超越方向，但原点处的两个坐标都不能从极大理想的极小生成元组中删去。维数、生成元数和方程数因此需要分别计算。本节先用参数系测量局部维数，再用长度记录零维部分的厚度，最后把这些量放进超曲面、完全交与变纤维的具体模型中。

\subsection{参数系及嵌入维数}
\label{g5:subsec:0304:01}

设 \((R,\mathfrak m,\kappa)\) 为非零 Noether 局部环。\(\mathfrak m\) 有限生成，由广义高度定理，它的高度 \(d=\dim R\) 有限。寻找参数的目标是选出尽量少的 \(x_i\in\mathfrak m\)，使商环 \(R/(x_1,\ldots,x_r)\) 的素谱只剩下闭点。这个条件等价于 \(\sqrt{(x_1,\ldots,x_r)}=\mathfrak m\)，因而允许商环仍有幂零元。

\begin{definition}
\label{g5:def:system-of-parameters}
	设 \((R,\mathfrak m,\kappa)\) 为非零 Noether 局部环，\(d=\dim R\)。若 \(x_1,\ldots,x_d\in\mathfrak m\) 满足
	\[
		\sqrt{(x_1,\ldots,x_d)}=\mathfrak m,
	\]
	就称它们为 \(R\) 的一组\term[canshuxi]{参数系}[system of parameters]。当 \(d=0\) 时允许空参数系。数
	\[
		\operatorname{edim}R\coloneqq
		\dim_\kappa(\mathfrak m/\mathfrak m^2)
	\]
	称为 \(R\) 的\term[qianruweishu]{嵌入维数}[embedding dimension]。
	\symidx{\(\operatorname{edim}R\)：局部环的嵌入维数}
\end{definition}

Nakayama 引理说明，嵌入维数是生成极大理想本身所需的最少元素数。参数系则只要求生成理想的根式等于极大理想。证明参数系存在时，我们需要选一个元素，同时避开有限多个不可约分支所对应的素理想。下面的\term[youxiansulixiangbikaiyinli]{有限素理想避开引理}[finite prime avoidance lemma]提供这一选择。

\begin{lemma}[有限素理想避开引理]
\label{g5:lem:finite-prime-avoidance}
	设 \(A\) 为交换含幺环，\(I\subseteq A\) 为理想，\(\mathfrak p_1,\ldots,\mathfrak p_s\) 为有限多个素理想。若 \(I\) 不包含在任何一个 \(\mathfrak p_i\) 中，则存在
	\[
		a\in I\setminus\bigcup_{i=1}^s\mathfrak p_i.
	\]
\end{lemma}
\begin{proof}
	删去包含在其他 \(\mathfrak p_j\) 中的素理想，不改变它们的并，所以可设所余素理想两两互不包含。对每个 \(i\)，选 \(a_i\in I\setminus\mathfrak p_i\)；对每个 \(j\ne i\)，选 \(b_{ij}\in\mathfrak p_j\setminus\mathfrak p_i\)。令
	\[
		c_i=a_i\prod_{j\ne i}b_{ij}.
	\]
	则 \(c_i\in I\)，且 \(c_i\notin\mathfrak p_i\)，因为 \(\mathfrak p_i\) 为素理想；对 \(j\ne i\)，又有 \(c_i\in\mathfrak p_j\)。所以 \(a=\sum_i c_i\) 在模 \(\mathfrak p_i\) 后等于非零的 \(c_i\)，从而避开全部 \(\mathfrak p_i\)。空集合情形任取 \(a\in I\) 即可。
\end{proof}

\begin{theorem}
\label{g5:thm:system-of-parameters}
	设 \((R,\mathfrak m)\) 为非零 Noether 局部环，\(d=\dim R\)。则参数系存在，并且
	\[
		d=\min\{r:\text{存在 }x_1,\ldots,x_r\in\mathfrak m,\quad
		\sqrt{(x_1,\ldots,x_r)}=\mathfrak m\}.
	\]
\end{theorem}
\begin{proof}
	先证明下界。若 \(r\) 个元素生成的理想根式为 \(\mathfrak m\)，则 \(\mathfrak m\) 是该理想的唯一极小素理想。\cref{g5:thm:height-theorem}给出
	\(\dim R=\operatorname{ht}\mathfrak m\leq r\)。

	下面对 \(d\) 归纳构造至多 \(d\) 个这样的元素。若 \(d=0\)，局部环的每个素理想都等于 \(\mathfrak m\)，故由根式等于所有包含它的素理想之交，得到 \(\sqrt{(0)}=\mathfrak m\)，空参数系满足要求。

	设 \(d>0\)。Noether 环只有有限多个极小素理想；这也可直接从\cref{g5:thm:prime-filtration}用于 \(R\) 时的有限支集并公式看出：每个极小素理想必须等于其中某个滤过素理想。取其中满足
	\(\dim R/\mathfrak p_i=d\)
	的那些。这样的素理想确实存在：取一条长度为 \(d\) 的素链，其最小端必为极小素理想，否则可在前面再添一项，违背 \(d\) 的定义。每个所选 \(\mathfrak p_i\) 都严格包含于 \(\mathfrak m\)，因为 \(\dim R/\mathfrak m=0\)。

	由\cref{g5:lem:finite-prime-avoidance}，可选 \(x\in\mathfrak m\) 避开这些 \(\mathfrak p_i\)。我们断言 \(\dim R/(x)\leq d-1\)。否则有一条长度为 \(d\) 的素链，其每项都包含 \(x\)。这条链的最小端 \(\mathfrak q\) 同样必为极小素理想，且 \(\dim R/\mathfrak q=d\)，于是 \(\mathfrak q\) 属于被避开的那一组；但 \(x\in\mathfrak q\)，矛盾。

	商环 \(R/(x)\) 仍为非零 Noether 局部环，设其维数为 \(d'\leq d-1\)。由归纳假设，在商环中可以选 \(d'\) 个元素生成以 \(\mathfrak m/(x)\) 为根式的理想。将它们提升回 \(R\)，连同 \(x\) 一共得到 \(1+d'\leq d\) 个元素，其生成理想的根式为 \(\mathfrak m\)。前面的下界又要求元素数至少为 \(d\)，故恰好为 \(d\)，得到参数系。
\end{proof}

参数系把素链定义的维数转换成了“切到只剩一个点所需的最少方程数”。这里保留商环的幂零结构，因此“只剩一个点”与“商环是域”并不相同。上述结论可与\cite[第13节，定理13.4--13.6；第14节开头]{Matsumura1989}对照；这里由已经证明的高度定理得到下界，存在性则使用有限素理想避开和归纳。

\begin{proposition}
\label{g5:prop:dimension-embedding-bound}
	设 \((R,\mathfrak m,\kappa)\) 为非零 Noether 局部环。则
	\[
		\dim R\leq\operatorname{edim}R<\infty.
	\]
\end{proposition}
\begin{proof}
	\(\mathfrak m\) 是有限生成模，所以 \(\mathfrak m/\mathfrak m^2\) 有有限维数 \(e\)。提升其一组基，由\cref{g5:cor:nakayama-generators}得到 \(\mathfrak m\) 的 \(e\) 个生成元。它们生成的理想根式当然为 \(\mathfrak m\)，故\cref{g5:thm:system-of-parameters}给出 \(\dim R\leq e\)。
\end{proof}

例如，设 \(k\) 为域，
\[
	R=\bigl(k[x,y]/(y^2-x^3)\bigr)_{(x,y)}.
\]
由\cref{g5:exm:cusp-birational}，这个环的分式域是 \(k(t)\)，其中 \(x=t^2,y=t^3\)。仿射维数定理给出整条尖点曲线维数为 \(1\)；局部环中有严格素链 \((0)\subsetneq(x,y)R\)，所以 \(\dim R=1\)。另一方面，关系 \(y^2-x^3\) 没有一次项，\(x,y\) 的像在 \(\mathfrak m/\mathfrak m^2\) 中线性无关：若有线性关系，清去常数项非零的局部分母后比较一次项，就会使两个系数都为零。因此 \(\operatorname{edim}R=2\)。单个 \(x\) 已是一组参数，因为 \(y^2=x^3\) 使 \(y\in\sqrt{(x)}\)。参数数目与极大理想生成元数的差别，在这里由同一个方程明确呈现。

\subsection{零维 Noether 环和有限长度}
\label{g5:subsec:0304:02}

参数切片虽然只剩一个点，仍可能保留不同厚度。例如 \(k[t]/(t^2)\) 与 \(k[t]/(t^3)\) 都只有一个素理想，维数都是零，但它们作为 \(k\) 向量空间的维数分别为 \(2,3\)。在不预先指定基域的一般环上，长度以简单模为单位记录这种差别。

\begin{definition}
\label{g5:def:module-length}
	设 \(A\) 为交换含幺环，\(M\) 为 \(A\) 模。所有有限严格子模链
	\[
		0=M_0\subsetneq M_1\subsetneq\cdots\subsetneq M_r=M
	\]
	的长度 \(r\) 的上确界，记为 \(\ell_A(M)\)，称为 \(M\) 的长度。若这个上确界有限，就称 \(M\) 为有限长度模；约定 \(\ell_A(0)=0\)。非零且没有非零真子模的模称为\term[jiandanmo]{简单模}[simple module]。
	\symidx{\(\ell_A(M)\)：模的长度}
\end{definition}

这延续\subsecref{g5:subsec:0104:05}中有限商代数的长度。有限长度模中，一条达到最大长度的链不能再插入子模，所以它的每个相邻商都是简单模；这种链称为\term[zuchenglie]{组成列}[composition series]。反过来，存在有限组成列会给出有限长度，这可由下面的加性逐段证明。

\begin{proposition}
\label{g5:prop:length-additivity}
	设 \(A\) 为交换含幺环，且 \(0\rightarrow L\rightarrow M\rightarrow N\rightarrow0\) 为短正合列。则 \(M\) 有限长度，当且仅当 \(L,N\) 都有限长度；在此情形，
	\[
		\ell_A(M)=\ell_A(L)+\ell_A(N).
	\]
\end{proposition}
\begin{proof}
	将 \(L\) 看成 \(M\) 的子模。对 \(U\subseteq V\subseteq M\)，若 \(U\cap L=V\cap L\)，且 \(U,V\) 在 \(N\) 中的像相同，则 \(U=V\)：对 \(v\in V\)，取 \(u\in U\) 具有相同的像，便有 \(v-u\in V\cap L=U\cap L\)。

	因此，\(M\) 的严格子模链每前进一步，或者与 \(L\) 的交严格增大，或者在 \(N\) 中的像严格增大。如果 \(L,N\) 长度有限，链长至多为 \(\ell_A(L)+\ell_A(N)\)，所以 \(M\) 长度有限。反过来，\(L\) 的子模链可看成 \(M\) 中的链，\(N\) 的链可拉回 \(M\)，所以 \(M\) 长度有限时二者都有限。最后，将 \(L\) 中的一条最长链与 \(N\) 中一条最长链的逆像拼接，得到长度恰为两者之和的链，证明等式。
\end{proof}

简单 \(A\) 模都是 \(A/\mathfrak m\) 的形式，其中 \(\mathfrak m\) 极大：任取一个非零元素，它生成整个简单模，核便是极大理想。因此，在局部环上，每个组成因子都同构于剩余域。若局部环还是剩余域为 \(k\) 的有限维 \(k\) 代数，则每个组成因子恰有一维，故长度等于 \(k\) 维数。长度有限也立即使子模升链和降链都终止。加性与链条件的关系可对照\cite[第3节，定理3.1及其后的讨论]{Matsumura1989}。

\ProseParagraphBreak

满足理想降链条件的环称为\term[Artinhuan]{Artin 环}[Artinian ring]。\footnote{Emil Artin（1898--1962），奥地利裔美国数学家，研究域论、代数数论与具有降链条件的环。}Noether 条件控制升链；对交换环，零维性恰好补足降链条件。

\begin{theorem}
\label{g5:thm:artinian-zero-dimension}
	设 \(A\ne0\) 为交换含幺环。以下条件等价：
	\begin{equivalences}
		\item \(A\) 为 Noether 环且 \(\dim A=0\)；
		\item \(A\) 为 Artin 环；
		\item \(A\) 作为自身的模具有有限长度。
	\end{equivalences}
	在这些条件下，\(\operatorname{Spec}A\) 是有限集，所有素理想均极大。零环同样为 Noether 环、Artin 环且长度为零，其空素谱按前文的空集约定处理。
\end{theorem}
\begin{proof}
	先设 \(A\) Noether 且维数为零。对 \(A\) 使用\cref{g5:thm:prime-filtration}，得到有限个相邻商 \(A/\mathfrak p_i\)。每个 \(\mathfrak p_i\) 都极大，所以这些商都是简单模。长度加性使 \(\ell_A(A)\) 有限。有限长度又立即蕴含理想降链条件。

	现在设 \(A\) Artin。对任意素理想 \(\mathfrak p\)，整环 \(D=A/\mathfrak p\) 仍满足降链条件。若 \(0\ne a\in D\)，链 \((a)\supseteq(a^2)\supseteq\cdots\) 终止，故某个 \(n\) 满足 \(a^n=a^{n+1}b\)。在整环中约去 \(a^n\) 得 \(ab=1\)。所以 \(D\) 是域，全部素理想均极大。

	极大理想也只能有有限多个。否则可取两两不同的 \(\mathfrak m_1,\mathfrak m_2,\ldots\)。链
	\(\mathfrak m_1\supseteq\mathfrak m_1\cap\mathfrak m_2\supseteq\cdots\)
	每一步都严格：对第 \(r+1\) 个极大理想，从每个 \(\mathfrak m_i\setminus\mathfrak m_{r+1}\) 中取元素，乘积属于前 \(r\) 个的交，却不属于素理想 \(\mathfrak m_{r+1}\)。这违背降链条件。记全部极大理想为 \(\mathfrak m_1,\ldots,\mathfrak m_s\)，并记
	\(\mathfrak n=\bigcap_i\mathfrak m_i=\sqrt{(0)}\)。

	我们证明 \(\mathfrak n\) 幂零，而不预先假设它有限生成。降链条件使某个 \(q\geq1\) 满足 \(\mathfrak n^q=\mathfrak n^{q+1}\)。令 \(I=\mathfrak n^q\)，则 \(I^2=I\)。若 \(I\ne0\)，在满足 \(IJ\ne0\) 的理想 \(J\) 中选一个极小者。存在 \(x\in J\) 使 \(Ix\ne0\)，故极小性使 \(J=Ax\)。另一方面，\(I(IJ)=IJ\ne0\)，极小性还给出 \(IJ=J\)。于是有限生成模 \(J\) 满足 \(IJ=J\)。由\cref{g5:prop:nakayama-ideal}，存在 \(a\in I\subseteq\mathfrak n\) 使 \((1-a)J=0\)。\(a\) 幂零，故 \(1-a\) 可逆，迫使 \(J=0\)，矛盾。因此 \(I=0\)。

	不同极大理想两两互素，中国剩余定理给出
	\[
		A/\mathfrak n\cong\prod_{i=1}^s A/\mathfrak m_i.
	\]
	每层 \(\mathfrak n^j/\mathfrak n^{j+1}\) 都是这个有限个域的乘积环上的模，并且继承降链条件。用乘积环的坐标幂等元分解它，得到有限个域上的向量空间。每个向量空间都必须有限维：无限基可以不断删去一个基向量，构造严格的下降子空间链。因此各层都是有限长度 \(A\) 模。由于 \(\mathfrak n^q=0\)，有限层的长度加性给出 \(\ell_A(A)<\infty\)。

	最后，有限长度使理想升链终止，所以 \(A\) Noether；又由刚才对 Artin 环素商的论证得 \(\dim A=0\)。全部等价关系得证。
\end{proof}

Artin 环推出有限长度的证明可对照\cite[第3节，定理3.2]{Matsumura1989}。在 Noether 局部环 \((R,\mathfrak m)\) 中，对真理想 \(J\) 因而有
\[
	\sqrt J=\mathfrak m
	\quad\Longleftrightarrow\quad
	\dim R/J=0
	\quad\Longleftrightarrow\quad
	\ell_R(R/J)<\infty.
\]
第一个等价来自局部环中的素理想对应，第二个来自上个定理；商环的子模就是它的理想，因此作为 \(R\) 模与作为 \(R/J\) 模计算的长度相同。特别地，参数系切出的局部商环总有有限长度。

\begin{example}
\label{g5:exm:parameter-lengths}
	设 \(k\) 为域，\(R=(k[x,y]/(y^2-x^3))_{(x,y)}\)。参数 \(x\) 与参数 \(y\) 分别给出
	\[
		R/(x)\cong k[y]/(y^2),\qquad
		R/(y)\cong k[x]/(x^3).
	\]
	这里局部化没有再改变右侧的环，因为它们已经是仅有一个极大理想的局部环。两个商的长度分别为 \(2,3\)。所以参数系的元素数只由维数决定，而参数切片的长度还取决于具体选择。

	另取任意域扩张 \(K/k\)。域 \(K\) 作为自身的模长度为 \(1\)，但 \(\dim_k K\) 可以无限。这说明有限长度不等于对任意外部基域都有限维；第一章有限点模型中的等同还使用了有限型及剩余域条件。
\end{example}

\subsection{超曲面与完全交的预告}
\label{g5:subsec:0304:03}

一个非零方程通常被期待削去一个维数。高度定理给出了一方面的控制：由 \(r\) 个元素生成的理想，其极小素理想高度至多为 \(r\)。要得到精确维数，还须核对方程所处的环及方程之间的关系。多项式环中的单个非恒定方程可以通过首一化得到一个完整模型。

\begin{proposition}
\label{g5:prop:polynomial-hypersurface-dimension}
	设 \(k\) 为任意域，\(n\geq1\)，\(f\in k[x_1,\ldots,x_n]\) 为非恒定多项式。则
	\[
		\dim k[x_1,\ldots,x_n]/(f)=n-1.
	\]
	更具体地，经过多项式坐标的可逆变换及乘以非零常数，可以使 \(f\) 关于最后一个变量首一；此时商环是其余 \(n-1\) 个变量组成的多项式子环上的有限自由模。
\end{proposition}
\begin{proof}
	由\cref{g5:lem:weighted-monic-change}，可以作一个对任意基域都有效的三角代换，使最后一个变量 \(T\) 的最高次系数为非零常数。除以该常数，写成
	\[
		f=T^e+a_1(Y)T^{e-1}+\cdots+a_e(Y),
		\qquad e\geq1,
	\]
	其中 \(Y=(Y_1,\ldots,Y_{n-1})\)。令 \(C=k[Y_1,\ldots,Y_{n-1}]\)。关于首一多项式作除法，每个剩余类都有唯一表达
	\[
		b_0(Y)+b_1(Y)T+\cdots+b_{e-1}(Y)T^{e-1}.
	\]
	唯一性来自非零首一多项式的倍式不可能有小于 \(e\) 的 \(T\) 次数。因此 \(C\) 单射到商环，且 \(1,T,\ldots,T^{e-1}\) 为其自由基。特别地，这是有限整扩张。由\cref{g5:thm:integral-dimension}及多项式环维数，商环维数为 \(n-1\)。
\end{proof}

在代数闭域上，由单个非恒定多项式定义的 \(V(f)\subseteq\mathbb A_k^n\) 称为\term[chaoqumian]{超曲面}[hypersurface]。取根式不会改变素谱的包含关系，所以 \(V(f)\) 的约化坐标环也有维数 \(n-1\)。商环 \(k[x_1,\ldots,x_n]/(f)\) 还保存方程的重数；例如 \(f=x_1^a\) 与 \(f=x_1\) 给出相同的点集，却一般给出不同的商环。

\ProseParagraphBreak

对多个方程，必须要求新方程在前面方程切出的商环中仍有独立作用。设 \(S\) 为交换含幺环，\(f_1,\ldots,f_r\in S\) 生成真理想。若对每个 \(i\)，乘以 \(f_i\) 在
\(S/(f_1,\ldots,f_{i-1})\)
上是单射，则称它们为一个\term[zhengzexulie]{正则序列}[regular sequence]。在给定的多项式环中，由正则序列生成的理想所定义的商代数是\term[wanquanjiao]{完全交}[complete intersection]的基本形式。后续将从局部方程研究其内在判别；这里先计算一类可以逐项验证的模型。

\begin{example}
\label{g5:exm:coordinate-complete-intersection}
	设 \(k\) 为域，\(S=k[x_1,\ldots,x_n]\)，\(0\leq r\leq n\)，并取整数 \(a_1,\ldots,a_r\geq1\)。序列
	\[
		x_1^{a_1},\ldots,x_r^{a_r}
	\]
	是正则序列。事实上，商去前 \(i-1\) 项后，所得环作为 \(k[x_i,\ldots,x_n]\) 模，有自由基
	\[
		x_1^{b_1}\cdots x_{i-1}^{b_{i-1}},
		\qquad 0\leq b_j<a_j.
	\]
	在这个自由模上乘以 \(x_i^{a_i}\) 单射，因为系数环是整环。最终商环
	\[
		B=S/(x_1^{a_1},\ldots,x_r^{a_r})
	\]
	是 \(k[x_{r+1},\ldots,x_n]\) 上的有限自由模，基为所有 \(0\leq b_i<a_i\) 的单项式 \(\prod_{i=1}^r x_i^{b_i}\)。该多项式子环单射到 \(B\)，所以整扩张保持维数给出
	\[
		\dim B=n-r.
	\]
	当 \(r=n\) 时，\(B\) 只支撑在原点，且
	\(\ell_B(B)=\dim_k B=\prod_i a_i\)。
	于是余维由被控制的坐标方向数决定，长度则还记住各方向的幂次。
\end{example}

两个最小生成方程也可能未使每个分支的余维都等于二。在代数闭域 \(k\) 上取
\[
	I=(xy,xz)\subseteq S=k[x,y,z].
\]
有
\[
	I=(x)\cap(y,z):
\]
若 \(xg\in(y,z)\)，模去 \((y,z)\) 后得到 \(x\overline g=0\) 于整环 \(k[x]\)，故 \(g\in(y,z)\)，证明了反向包含。因此零点集是平面 \(x=0\) 与直线 \(y=z=0\) 的并，两个分支的维数分别为 \(2,1\)，整个商环维数为 \(2\)。在原点的局部环中，\(xy,xz\) 仍是这个理想的极小生成元：其像在 \(I/\mathfrak mI\) 中线性无关，因为 \(\mathfrak mI\) 中没有非零二次齐次项。方程个数为二，并非因为多写了一个可删去的生成元。

\ProseParagraphBreak

这个序列失败在第二步。\(xy\) 在整环 \(S\) 中是非零因子，但 \(xz\) 在 \(S/(xy)\) 中被非零类 \(y\) 杀掉：
\(y(xz)=z(xy)=0\)。
新方程在某个已有分支上恒等于零，因而不能把那个分支再切掉一个维数。高度定理控制的是最小可能的方程数与素理想高度；具体商的维数还需要检查这些消失关系。

\subsection{纤维维数问题的代数准备}
\label{g5:subsec:0304:04}

在映射 \(A\rightarrow B\) 中，把基点固定在 \(\mathfrak p\in\operatorname{Spec}A\)，应当使 \(\mathfrak p\) 中的元素变为零，并使 \(A\setminus\mathfrak p\) 中的元素可逆。两个要求合起来给出\term[xianweihuan]{纤维环}[fibre ring]
\[
	B(\mathfrak p)\coloneqq B\otimes_A\kappa(\mathfrak p)
	\cong (A\setminus\mathfrak p)^{-1}(B/\mathfrak pB).
\]

\symidx{\(B(\mathfrak p)=B\otimes_A\kappa(\mathfrak p)\)：素点处的纤维环}
其素理想恰对应 \(B\) 中收缩为 \(\mathfrak p\) 的素理想：在商环中要求包含 \(\mathfrak pB\)，在局部化中要求避开 \(A\setminus\mathfrak p\)，所以收缩既包含又不大于 \(\mathfrak p\)。这只是商环及局部化的素理想对应，适用于任意环同态，纤维环也可能为零环。

\ProseParagraphBreak

若 \(A\) 是整环，零素理想的剩余域为 \(K=\operatorname{Frac}A\)。相应的 \(B\otimes_AK\) 称为\term[fanxianweihuan]{泛纤维环}[generic fibre ring]。它把所有非零基底函数都反演，因而适合测量映射在基底的泛点处还剩多少独立方向。纤维环与泛纤维的记法可对照\cite[第15节，Fibres，第116页]{Matsumura1989}。

\begin{proposition}
\label{g5:prop:generic-fiber-dimension}
	设 \(A\subseteq B\) 为整环的单射，且 \(B\) 是有限型 \(A\) 代数。令 \(K=\operatorname{Frac}A\)、\(L=\operatorname{Frac}B\)。则
	\[
		\dim(B\otimes_AK)=\operatorname{trdeg}_K L.
	\]
	若 \(A,B\) 还是某个域 \(k\) 上的有限型代数，则
	\[
		\dim(B\otimes_AK)=\dim B-\dim A.
	\]
\end{proposition}
\begin{proof}
	因为 \(A\subseteq B\) 单射且二者为整环，
	\[
		B\otimes_AK=(A\setminus\{0\})^{-1}B
	\]
	是非零有限型 \(K\) 整代数，其分式域仍为 \(L\)。对这个 \(K\) 代数应用\cref{g5:thm:affine-dimension}，得到第一个等式。

	在第二个情形，超越次数具有塔式加性
	\(\operatorname{trdeg}_kL=\operatorname{trdeg}_kK+\operatorname{trdeg}_K L\)。
	确实，取 \(K/k\) 的超越基和 \(L/K\) 的超越基，它们的并在 \(k\) 上代数独立；而 \(L\) 对它们生成的域代数，因为两段剩余扩张都代数。再将 \(A,B\) 的维数分别写成其分式域的超越次数，便得结论。
\end{proof}

这个等式计算泛纤维，不能直接把其中的 \(K\) 换成任意点的剩余域。有限映射则对全部纤维有统一的零维结论，且保留了点集以外的有限代数。

\begin{proposition}
\label{g5:prop:finite-fiber-algebra}
	设 \(A\rightarrow B\) 为交换含幺环同态，且 \(B\) 是有限 \(A\) 模。对每个 \(\mathfrak p\in\operatorname{Spec}A\)，纤维环 \(B(\mathfrak p)\) 是有限维 \(\kappa(\mathfrak p)\) 代数。若它非零，则为零维 Artin 环，只有有限多个素理想，且各剩余域都是 \(\kappa(\mathfrak p)\) 的有限扩张。若 \(A\subseteq B\) 单射，则每个纤维环均非零。
\end{proposition}
\begin{proof}
	取 \(B\) 的有限组 \(A\) 模生成元。张量 \(\kappa(\mathfrak p)\) 后，它们的像张成 \(B(\mathfrak p)\)，所以这是有限维向量空间。它的每个理想都是向量子空间，故理想降链终止。\cref{g5:thm:artinian-zero-dimension}说明非零纤维环维数为零且素理想只有有限多个；每个剩余域还是原有限维向量空间的商，因此为有限域扩张。

	若 \(A\subseteq B\) 单射，模有限性使扩张整。由\cref{g5:thm:lying-over}，存在收缩为 \(\mathfrak p\) 的素理想 \(\mathfrak q\subset B\)。商环与局部化的素理想对应使它给出 \(B(\mathfrak p)\) 的一个素理想，故纤维环不能是零环。
\end{proof}

这把\cref{g5:prop:finite-affine-closed-fibres}的闭点结论推广到了任意素点。例如 \(k[t]\rightarrow k[z]\)、\(t\mapsto z^2\) 是秩二的有限自由扩张。在 \(t=0\) 处的纤维环是 \(k[z]/(z^2)\)，只有一个点、维数为零，但长度为二。有限性限制纤维的维数和点数，并不要求纤维约化。

\begin{example}
\label{g5:exm:fiber-dimension-jump}
	设 \(k\) 为代数闭域，考虑
	\[
		\pi:\mathbb A_k^2\longrightarrow\mathbb A_k^2,\qquad
		(x,y)\longmapsto(u,v)=(x,xy).
	\]
	坐标同态 \(A=k[u,v]\rightarrow B=k[x,y]\) 单射，因为不同单项式 \(u^iv^j\) 的像 \(x^{i+j}y^j\) 仍有不同指数。分式域却相同，因为 \(y=v/u\)，所以泛纤维环是 \(k(u,v)\)，维数为零。

	在闭点 \((a,b)\) 处，纤维环为
	\[
		k[x,y]/(x-a,xy-b).
	\]
	代入 \(x=a\) 后，第二个方程成为 \(ay=b\)，得到\cref{g5:tab:fiber-dimension-jump}中的三种情形。
\end{example}

\begin{table}[htbp]
	\centering
	\caption{映射 \(\pi(x,y)=(x,xy)\) 的闭点纤维}
	\label{g5:tab:fiber-dimension-jump}
	\begin{tabular}{lll}
		\toprule
		基点条件 & 纤维环 & 纤维的形状与维数 \\
		\midrule
		\(a\ne0\) & \(k\) & 一个点，维数 \(0\) \\
		\(a=0,\ b=0\) & \(k[y]\) & 一条直线，维数 \(1\) \\
		\(a=0,\ b\ne0\) & 零环 & 空纤维 \\
		\bottomrule
	\end{tabular}
\end{table}

在 \(u\ne0\) 的开集上，逆公式为 \(x=u,y=v/u\)；而整条直线 \(x=0\) 被送到原点。这里 \(A,B\) 都是维数二的整环，泛纤维维数为零，指定点的纤维却可能是一维或为空。这说明维数差公式的泛纤维结论与逐点结论必须区分。

\ProseParagraphBreak

还可以直接检测这个例子不平坦。取理想 \(J=(u,v)\subset A\)。它有呈示
\[
	A\xrightarrow{\,(-v,u)\,}A^2\longrightarrow J\longrightarrow0:
\]
若 \(au+bv=0\)，模去 \(u\) 后由 \(k[v]\) 为整环得 \(u\mid b\)，写 \(b=uc\) 后便有 \(a=-vc\)。张量 \(B\) 得到
\[
	J\otimes_A B\cong B^2/B(-xy,x).
\]
由包含 \(J\hookrightarrow A\) 诱导的映射把 \((a,b)\) 的类送到 \(ax+bxy\)。类 \((-y,1)\) 在核中，但它非零：若它等于 \(c(-xy,x)\)，第二坐标就要求 \(1=cx\)，这在 \(k[x,y]\) 中不可能。因此 \(J\otimes_A B\rightarrow B\) 不是单射，\(B\) 不是平坦 \(A\) 模。这个判断使用了张量关系，未把纤维维数变化本身当作平坦性判据。

\ProseParagraphBreak

参数系还给出一个适用于指定点的局部维数估计。它的等号条件可直接用素链能否向下提升来表达，因而与\secref{g5:sec:0301}的 going-down 性质相接。

\begin{proposition}
\label{g5:prop:local-fiber-dimension-inequality}
	设 \((A,\mathfrak m)\rightarrow(B,\mathfrak n)\) 为非零 Noether 局部环之间的局部同态，即 \(\mathfrak mB\subseteq\mathfrak n\)。则
	\[
		\dim B\leq\dim A+\dim(B/\mathfrak mB).
	\]
	若该同态具有 going-down 性质，即对 \(\mathfrak p_1\subseteq\mathfrak p_2\) 及收缩为 \(\mathfrak p_2\) 的 \(\mathfrak q_2\in\operatorname{Spec}B\)，总有 \(\mathfrak q_1\subseteq\mathfrak q_2\) 收缩为 \(\mathfrak p_1\)，则上式取等号。
\end{proposition}
\begin{proof}
	记 \(d=\dim A\)、\(r=\dim(B/\mathfrak mB)\)。选 \(A\) 的参数系 \(x_1,\ldots,x_d\)，再选 \(B/\mathfrak mB\) 的参数系，并将其提升为 \(y_1,\ldots,y_r\in\mathfrak n\)。因 \(\sqrt{(x_1,\ldots,x_d)A}=\mathfrak m\)，每个 \(\mathfrak m\) 中元素的某个幂属于 \((x_1,\ldots,x_d)\)。因此
	\[
		\mathfrak mB\subseteq\sqrt{(x_1,\ldots,x_d)B}.
	\]
	而商环中参数系的条件给出
	\(\sqrt{\mathfrak mB+(y_1,\ldots,y_r)B}=\mathfrak n\)。
	两者结合说明 \(x_i,y_j\) 这 \(d+r\) 个元素在 \(B\) 中生成的理想根式为 \(\mathfrak n\)。由\cref{g5:thm:system-of-parameters}的最小性，得到 \(\dim B\leq d+r\)。

	若 going-down 成立，取 \(B/\mathfrak mB\) 中长度为 \(r\) 的素链，提升为
	\[
		\mathfrak q_0\subsetneq\cdots\subsetneq\mathfrak q_r=\mathfrak n.
	\]
	这些素理想都收缩为 \(\mathfrak m\)。再取 \(A\) 中长度为 \(d\)、终点为 \(\mathfrak m\) 的素链，从 \(\mathfrak q_0\) 开始反复使用 going-down，便在 \(\mathfrak q_0\) 下方添上 \(d\) 次严格包含；严格性由收缩理想严格不同保证。这样得到 \(B\) 中长度为 \(d+r\) 的素链，给出相反不等式。
\end{proof}

这正是\cite[第15节，定理15.1]{Matsumura1989}中局部维数估计及 going-down 等号的证明内容。对于一般有限型映射，哪些开集上纤维维数稳定、平坦性怎样控制维数及重数，还需要相对结构定理；本节已经把这些问题分成可计算的泛纤维环、指定点的纤维环及其局部关系。

% 纲要条目排为编号小节标题，并进入目录与书签；说明段仍保留于纲要注释中。


\begin{chaptersummary}
	元素整于基环，等价于它生成的代数在基环上为有限模。这个判据使有限个整元能够共同处理，并导出整性的传递性。整扩张的素理想满足 lying-over、going-up 与 incomparability；整环扩张在基环整闭时还满足 going-down。上行和下行的出发数据不同，这一区别决定能否控制某个指定素理想下方的链。

	\ProseParagraphBreak

	Noether 正规化在有限型代数内部选取代数无关的参数，使原代数成为所选多项式子代数上的有限模。无限域上可以使用线性代换，有限域上则可用彼此分离的幂次使关系首一化。普通正规化是在给定分式域中补入整元；两种构造的方向与目的各不相同。

	\ProseParagraphBreak

	Krull 维数计算素链的长度，素理想的高度等于相应局部环的维数。主理想定理把一个方程上极小素理想的高度限制为一，广义高度定理把这个界推广到方程数。有限型整代数的维数等于分式域的超越次数；一般有限型代数取各极小分支的最大值。整扩张保持整体维数，而指定素理想的高度比较还要核对下降性质。

	\ProseParagraphBreak

	Noether 局部环的参数系用恰好 \(\dim R\) 个元素切出有限长度的商环。嵌入维数计算极大理想所需的生成元数，商环长度记录切片的代数厚度，两者均不能与参数数目混同。有限映射的非空纤维为零维有限代数；一般有限型映射的泛纤维维数由超越次数给出，特殊点上的纤维则可能改变维数、点数或幂零结构。
\end{chaptersummary}

\BookExercises

\begin{exercise}\label{g5:ex:integral-two-square-roots}
	设 \(k\) 为域，\(A=k[u,v]\)，
	\[
		B=A[s,t]/(s^2-u,t^2-v).
	\]
	证明 \(B\) 是秩为四的自由 \(A\) 模，并为 \(z=s+t\) 写出一个首一整性方程。说明该方程在特征二时仍成立，但不能据此断言它是最小次数的方程。
\end{exercise}
\begin{solution}
	先模去 \(s^2-u\)，再模去 \(t^2-v\)，两次首一多项式除法都给出唯一余式。因此每个元素唯一写成
	\[
		a+bs+ct+dst,\qquad a,b,c,d\in A,
	\]
	故 \(1,s,t,st\) 是一组自由基。由
	\[
		z^2-u-v=2st
	\]
	得
	\[
		z^4-2(u+v)z^2+(u-v)^2=0.
	\]
	这是对所有特征都成立的整系数恒等式，所以在特征二时依然有效。但在特征二时已经有 \(z^2=u+v\)，上面的四次式只是这个二次式的平方。有限模判据产生的首一方程保证存在性，不保证次数极小。
\end{solution}

\begin{exercise}\label{g5:ex:integral-localization-units}
	设 \(A\) 为整环，\(S\subseteq A\setminus\{0\}\) 为乘法集。证明下列条件等价：\(S^{-1}A\) 整于 \(A\)；\(S^{-1}A\) 是有限 \(A\) 模；\(S\) 中每个元素都已是 \(A\) 的单位；自然同态 \(A\to S^{-1}A\) 为同构。
\end{exercise}
\begin{solution}
	有限模蕴含整性。若局部化整于 \(A\)，对任意 \(s\in S\)，存在
	\[
		(1/s)^n+a_1(1/s)^{n-1}+\cdots+a_n=0,\qquad a_i\in A.
	\]
	乘以 \(s^{n-1}\)，得到
	\[
		1/s=-a_1-a_2s-\cdots-a_ns^{n-1}\in A.
	\]
	整环假设使 \(A\) 与局部化都能视为同一分式域中的子环，因此这个等式表明 \(s\) 已在 \(A\) 中可逆。若每个分母都是单位，局部化的逆同态就是 \(a/s\mapsto as^{-1}\)。最后，同构于 \(A\) 的模当然由一个元素生成，完成循环。
\end{solution}

\begin{exercise}\label{g5:ex:integral-base-change}
	设 \(A\to B\) 和 \(A\to C\) 为交换含幺环同态。证明：若 \(B\) 是有限 \(A\) 模，则 \(B\otimes_A C\) 是有限 \(C\) 模；若 \(B\) 整于 \(A\) 的像，则 \(B\otimes_A C\) 整于 \(C\) 的像。再计算
	\(\mathbb Z[i]\otimes_{\mathbb Z}\mathbb F_2\)，说明换基后的代数可能出现幂零元。
\end{exercise}
\begin{solution}
	若 \(b_1,\ldots,b_r\) 生成 \(B\)，则 \(b_j\otimes1\) 生成张量积，因为
	\[
		b\otimes c=\sum_j(b_j\otimes1)c_jc
	\]
	对 \(b=\sum_j a_jb_j\) 成立，这里的 \(c_j\) 是 \(a_j\) 在 \(C\) 中的像。

	在整性情形，每个 \(b\otimes1\) 满足由 \(b\) 的首一方程换系数得到的首一方程。张量积中的任意元素都是有限和 \(\sum_{j=1}^r b_j\otimes c_j\)，属于由这些有限个整元生成的 \(C\) 子代数。该子代数是有限 \(C\) 模，因而其中每个元素都整。这也涵盖结构同态有核的情形：方程系数只取决于 \(C\) 的像。

	最后，
	\[
		\mathbb Z[i]\otimes_{\mathbb Z}\mathbb F_2
		\cong\mathbb F_2[T]/(T^2+1)
		=\mathbb F_2[T]/((T+1)^2)
		\cong\mathbb F_2[\varepsilon]/(\varepsilon^2).
	\]
	\(\varepsilon\) 的类非零而平方为零。因此有限自由扩张在换基后仍有限自由，却未必保持约化性。
\end{solution}

\begin{exercise}\label{g5:ex:cusp-two-normalizations}
	设 \(k\) 为域，\(B=k[t^2,t^3]\subset k[t]\)。分别证明 \(B\) 在 \(k[t^2]\) 和 \(k[t^3]\) 上有限自由，求出秩与一组基，并求 \(B\) 在其分式域中的整闭包。说明这三个环扩张对应哪些不同的正规化构造。
\end{exercise}
\begin{solution}
	\(B\) 中出现的单项式指数为 \(0,2,3,4,\ldots\)。按奇偶分类，偶次幂属于 \(k[t^2]\)，奇次幂从 \(t^3\) 开始，均为 \(t^3\) 乘以 \(t^2\) 的幂。因此
	\[
		B=k[t^2]\oplus k[t^2]t^3.
	\]
	按模三余数分类，指数分别从 \(0,4,2\) 开始，故
	\[
		B=k[t^3]\oplus k[t^3]t^2\oplus k[t^3]t^4.
	\]
	两处的独立性均来自单项式指数集合互不相交，所以秩分别为二和三。

	又有 \(t=t^3/t^2\in\operatorname{Frac}B\)，故 \(\operatorname{Frac}B=k(t)\)。元素 \(t\) 满足 \(X^2-t^2=0\)，整于 \(B\)，却不属于 \(B\)。因此 \(B\) 不整闭。\(k[t]\) 是唯一分解整环，按\cref{g5:prop:ufd-integrally-closed}整闭。若 \(h\in k(t)\) 整于 \(B\)，同一个首一方程也表明它整于 \(k[t]\)，所以 \(h\in k[t]\)。反之，\(k[t]=B[t]\) 整于 \(B\)。于是整闭包恰为 \(k[t]\)。

	前两个子环 \(k[t^2]\)、\(k[t^3]\) 都给出 Noether 正规化，参数数目均为一，但模秩不同；扩张 \(B\subset k[t]\) 则是普通正规化，它在原分式域中补入缺失的整元。
\end{solution}

\begin{exercise}\label{g5:ex:finite-field-rank-six}
	在
	\[
		B=\mathbb F_2[x,y]/(xy(x-y))
	\]
	中令 \(u=x^2+xy+y^2\)。证明 \(u\) 在 \(\mathbb F_2\) 上超越，且 \(B\) 是 \(\mathbb F_2[u]\) 上秩为六的自由模，给出一组基。将其与\cref{g5:exm:finite-field-nonlinear-normalization}的秩七投影比较。
\end{exercise}
\begin{solution}
	引入独立变量 \(U,X,Y\)，令
	\[
		g_1=Y^2+XY+X^2+U,\qquad f=X^2Y+XY^2.
	\]
	则
	\[
		B\cong\mathbb F_2[U,X,Y]/(g_1,f).
	\]
	在特征二下，\(f+Xg_1=X^3+UX\)，故
	\[
		(g_1,f)=(g_1,X^3+UX).
	\]
	先对首一三次式 \(X^3+UX\) 取商，再对关于 \(Y\) 的首一二次式 \(g_1\) 取商，得到
	\[
		B\cong
		\bigl(\mathbb F_2[U][X]/(X^3+UX)\bigr)[Y]/
		(Y^2+XY+X^2+U).
	\]
	两次首一除法说明 \(B\) 是 \(\mathbb F_2[U]\) 上的自由模，基为
	\[
		1,\ x,\ x^2,\ y,\ xy,\ x^2y.
	\]
	特别地，\(\mathbb F_2[U]\to B\) 单射，其像即 \(\mathbb F_2[u]\)，故 \(u\) 超越。这里与正文的参数 \(x-y^3\) 都只需一个代数无关元，却得到不同的自由秩。Noether 正规化确定参数数目，并不把所选有限投影的次数固定下来。
\end{solution}

\begin{exercise}\label{g5:ex:mixed-codimension-equations}
	设 \(k\) 为域，\(r\geq2\)，\(S=k[x_0,x_1,\ldots,x_r]\)，
	\[
		I=(x_0x_1,\ldots,x_0x_r).
	\]
	求 \(I\) 的极小素理想及其高度，计算 \(\dim(S/I)\)。再证明，在原点处的局部环中，这 \(r\) 个方程仍是 \(I\) 的极小生成元。
\end{exercise}
\begin{solution}
	有
	\[
		I=(x_0)\cap(x_1,\ldots,x_r).
	\]
	从左到右显然；若 \(x_0g\in(x_1,\ldots,x_r)\)，在商环 \(k[x_0]\) 中得到 \(x_0\overline g=0\)，故 \(g\in(x_1,\ldots,x_r)\)，证明反向包含。这两个素理想互不包含，因而正是 \(I\) 上的极小素理想。

	素链 \((0)\subsetneq(x_0)\) 与主理想定理给出 \(\operatorname{ht}(x_0)=1\)。坐标素链
	\[
		(0)\subsetneq(x_1)\subsetneq\cdots\subsetneq(x_1,\ldots,x_r)
	\]
	给出另一个高度至少为 \(r\)，高度定理给出至多为 \(r\)。两个分支的商环分别为 \(k[x_1,\ldots,x_r]\) 与 \(k[x_0]\)，所以
	\[
		\dim(S/I)=\max\{r,1\}=r.
	\]

	令 \(\mathfrak m=(x_0,\ldots,x_r)\)。在 \(I/\mathfrak mI\) 中，\(x_0x_1,\ldots,x_0x_r\) 的像线性无关，因为 \(\mathfrak mI\) 的每个非零齐次项次数至少为三，而这些二次单项式互不相同。它们已经张成这个商。局部化后剩余向量空间不变，第二章的极小生成元判据便说明至少需要 \(r\) 个生成元。由 \(r\) 个必要方程定义的集合仍然可以同时含余维一与余维 \(r\) 的分支。
\end{solution}

\begin{exercise}\label{g5:ex:parameter-zero-divisor}
	设 \(k\) 为域，
	\[
		R=\bigl(k[x,y]/(x^2,xy)\bigr)_{(x,y)},\qquad
		\mathfrak m=(x,y)R.
	\]
	计算 \(\dim R\) 与嵌入维数，证明 \(y\) 是一个参数，却是零因子，并计算 \(\ell_R(R/yR)\)。
\end{exercise}
\begin{solution}
	\(x\) 幂零，而商去 \((x)\) 得到约化整环 \(k[y]_{(y)}\)，故 \(R\) 的幂零根恰为 \((x)\)。素链在商去幂零根时不变，所以 \(\dim R=1\)。\(\mathfrak m/\mathfrak m^2\) 由 \(x,y\) 的像组成一组 \(k\) 基：原关系 \(x^2,xy\) 没有一次项，局部化不会增加一次线性关系。因此嵌入维数为二。

	又有
	\[
		R/yR\cong k[x]/(x^2),
	\]
	其唯一素理想是 \((x)\)，所以 \(\sqrt{yR}=\mathfrak m\)，\(y\) 构成一个参数系。这里 \(x\ne0\)：它在 \(k[x,y]/(x^2,xy)\) 中的零化子为 \((x,y)\)，故原点之外的分母不能把它杀掉。但 \(yx=0\)，所以 \(y\) 是零因子。商环有组成列
	\[
		0\subset (x)\subset k[x]/(x^2),
	\]
	两个相邻商均同构于 \(k\)，因而长度为二。参数系切出零维局部商，不自动成为正则序列。
\end{solution}

\begin{exercise}\label{g5:ex:artinian-four-dimensional}
	设 \(k\) 为域，\(E=k[x,y]/(x^3,xy,y^2)\)。求其全部素理想、Krull 维数、嵌入维数与长度，并写出极大理想的幂滤过。比较这些数各自测量的对象。
\end{exercise}
\begin{solution}
	首一单项式关系给出 \(k\) 基 \(1,x,x^2,y\)。\(x,y\) 都幂零，所以每个素理想都包含它们；又 \(E/(x,y)\cong k\)，故唯一素理想也是唯一极大理想
	\(\mathfrak m=(x,y)\)。于是 \(E\) 为零维局部环，且
	\[
		E\supset\mathfrak m\supset\mathfrak m^2=(x^2)
		\supset\mathfrak m^3=0.
	\]
	各层 \(E/\mathfrak m,\mathfrak m/\mathfrak m^2,\mathfrak m^2\) 的 \(k\) 维数依次为 \(1,2,1\)。因此嵌入维数为二，长度为 \(1+2+1=4\)。Krull 维数只看素链，此处为零；嵌入维数记录极大理想的极小生成元数；长度还累计了幂零方向保留的全部组成因子。
\end{solution}

\begin{exercise}\label{g5:ex:parameter-ideal-length-six}
	设 \(k\) 为域，\(R=k[x,y]_{(x,y)}\)，\(\mathfrak m=(x,y)R\)，\(J=(x^2,y^3)R\)。证明 \(x^2,y^3\) 为参数系，而 \(J\ne\mathfrak m\)。计算 \(\dim R\)、嵌入维数以及 \(\ell_R(R/J)\)。
\end{exercise}
\begin{solution}
	素链 \((0)\subsetneq(x)R\subsetneq\mathfrak m\) 给出维数至少为二，高度定理给出 \(\dim R\leq2\)，故维数为二。\(x,y\) 的像组成 \(\mathfrak m/\mathfrak m^2\) 的基，所以嵌入维数也是二。由 \(x^2,y^3\in J\)，有 \(\sqrt J=\mathfrak m\)，故两个元素构成参数系。但 \(x\notin J\)，所以参数理想不等于极大理想。

	多项式商环 \(k[x,y]/(x^2,y^3)\) 有基
	\[
		1,\ y,\ y^2,\ x,\ xy,\ xy^2.
	\]
	其中极大理想 \((x,y)\) 幂零，任何常数项非零的元素 \(c+n\) 都可逆：将其写成 \(c(1+c^{-1}n)\)，逆由有限几何级数给出。因此原点处的局部化不再改变这个商环，\(R/J\) 的 \(k\) 维数为六。它的剩余域为 \(k\)，故 \(\ell_R(R/J)=6\)。参数个数为二，切片长度为六，二者记录不同的信息。
\end{solution}

\begin{exercise}\label{g5:ex:finite-local-base-maximals}
	设 \((A,\mathfrak m)\) 为整闭 Noether 局部整环，\(A\subseteq B\) 为有限扩张，且 \(B\) 为整环。证明 \(B\) 只有有限多个极大理想，每个极大理想 \(\mathfrak n\) 都满足
	\[
		\dim B_{\mathfrak n}=\dim A.
	\]
	这是否意味着 \(B\) 本身为局部环？以
	\[
		A=\mathbb Z_{(5)},\qquad B=A[T]/(T^2+1)
	\]
	为例，求 \(B\) 的全部极大理想和相应局部维数。
\end{exercise}
\begin{solution}
	有限性蕴含整性，由\cref{g5:lem:integral-field-extension}，\(B\) 的每个极大理想都收缩为 \(A\) 的极大理想 \(\mathfrak m\)，因此都包含 \(\mathfrak mB\)。这些极大理想恰与有限维剩余域代数
	\[
		B/\mathfrak mB=B\otimes_A(A/\mathfrak m)
	\]
	的极大理想对应。由\cref{g5:prop:finite-fiber-algebra}，它们只有有限多个。又因 \(A\) 整闭且 \(B\) 为整环，going-down 成立；\cref{g5:prop:integral-height-comparison}于是给出
	\[
		\dim B_{\mathfrak n}
		=\operatorname{ht}_B\mathfrak n
		=\operatorname{ht}_A\mathfrak m
		=\dim A.
	\]
	这里既用有限性控制极大理想个数，又用下降性质控制每个指定极大理想的高度。

	在所给例子中，\(A\) 为主理想整环的局部化，故 Noether 且整闭。\(B\) 是秩二的自由 \(A\) 模，并可识别为 \(\mathbb Z[i]\) 在所有不被 \(5\) 整除的整数处的局部化，因此也是整环。模去 \(5\) 得
	\[
		B/5B\cong\mathbb F_5[T]/((T-2)(T+2))
		\cong\mathbb F_5\times\mathbb F_5.
	\]
	故全部极大理想恰为
	\[
		\mathfrak n_+=(5,T-2),\qquad
		\mathfrak n_-=(5,T+2).
	\]
	两个局部环的维数都为一，但 \(B\) 有两个极大理想，因而不是局部环。
\end{solution}

\begin{exercise}\label{g5:ex:finite-extra-point}
	设 \(k\) 为域，令 \(A=k[t]\)，\(B=A\times k\)，结构同态为
	\[
		A\longrightarrow B,\qquad a(t)\longmapsto(a(t),a(0)).
	\]
	证明这是有限单射，计算 \(\dim A,\dim B\)。令 \(\mathfrak q=A\times(0)\)，计算其收缩与高度。再计算各 \(a\in k\) 处的纤维环 \(B\otimes_A A/(t-a)\)，并证明 \(B\) 不是平坦 \(A\) 模。
\end{exercise}
\begin{solution}
	第一分量保证单射；元素 \((1,0),(0,1)\) 生成 \(B\) 作为 \(A\) 模，因为
	\[
		(f,c)=f\cdot(1,0)+c\cdot(0,1).
	\]
	因此扩张有限且整，\(\dim B=\dim A=1\)。

	\(\mathfrak q\) 的商环为 \(k\)，收缩是 \((t)\)。乘积环中的素理想必含两个正交幂等元 \((1,0),(0,1)\) 之一，因为它们的乘积为零；于是素理想只有形如 \(\mathfrak p\times k\) 或 \(A\times(0)\) 的两类。因此 \(\mathfrak q\) 下方没有更小的素理想，\(\operatorname{ht}\mathfrak q=0\)，而 \(\operatorname{ht}(t)=1\)。这不违背整体维数保持；前题所需的上环为整环这一条件在这里失败。

	纤维环就是 \(B/(t-a)B\)。在第一分量取商得到 \(k\)；第二分量上 \(t-a\) 作用为标量 \(-a\)。故
	\[
		B\otimes_A A/(t-a)\cong
		\begin{cases}
			k,&a\ne0,\\
			k\times k,&a=0.
		\end{cases}
	\]
	最后，\(A\xrightarrow{\times t}A\) 单射，张量 \(B\) 后变为 \(B\) 上乘以 \((t,0)\)。它杀掉非零元素 \((0,1)\)，所以不再单射，\(B\) 不平坦。
\end{solution}

\begin{exercise}\label{g5:ex:gaussian-prime-fibers}
	把 \(\mathbb Z[i]\) 视为 \(\mathbb Z\) 上的有限代数。分别计算素数 \(p=2,3,5\) 处的纤维环，求纤维中素理想的个数及其剩余域，并比较纤维的维数、\(\mathbb F_p\) 向量空间维数与约化性。
\end{exercise}
\begin{solution}
	\(\mathbb Z[i]\) 在 \(\mathbb Z\) 上有基 \(1,i\)，故每个纤维都有 \(\mathbb F_p\) 维数二，且
	\[
		\mathbb Z[i]\otimes_{\mathbb Z}\mathbb F_p
		\cong\mathbb F_p[T]/(T^2+1).
	\]
	当 \(p=2\) 时，\(T^2+1=(T+1)^2\)，纤维为
	\(\mathbb F_2[\varepsilon]/(\varepsilon^2)\)。它只有一个素理想 \((\varepsilon)\)，剩余域为 \(\mathbb F_2\)，但不是约化环。

	当 \(p=3\) 时，\(\mathbb F_3\) 中的平方只有 \(0,1\)，因此 \(T^2+1\) 无根且不可约。纤维是九元域 \(\mathbb F_9\)，只有零素理想，剩余域为 \(\mathbb F_9\)。

	当 \(p=5\) 时，\(T^2+1=(T-2)(T+2)\)，两个因子互素，中国剩余定理给出 \(\mathbb F_5\times\mathbb F_5\)。纤维有两个素理想，剩余域均为 \(\mathbb F_5\)。

	三个纤维的 Krull 维数都为零，向量空间维数都为二，素理想个数却依次为一、一、二；只有第一个不约化。顺便由整扩张的维数保持得 \(\dim\mathbb Z[i]=\dim\mathbb Z=1\)，其中 \(\mathbb Z\) 的非零素理想均为极大理想。
\end{solution}

\begin{exercise}\label{g5:ex:flat-hyperbola-family}
	设 \(k\) 为域，\(A=k[t]\)，\(B=k[x,y]\)，结构同态由 \(t\mapsto xy\) 给出。证明 \(B\) 是自由 \(A\) 模，且一组基为
	\[
		1,\quad x^i\ (i\geq1),\quad y^j\ (j\geq1).
	\]
	由此证明 \(B\) 平坦但不是有限 \(A\) 模。计算泛纤维及每个 \(t=a\)、\(a\in k\) 处的纤维，并说明平坦性与特殊纤维可约并不冲突。
\end{exercise}
\begin{solution}
	不同的 \(t\) 次幂映到不同的单项式 \(x^ny^n\)，所以结构同态单射。每个单项式 \(x^ay^b\) 都可以取出 \(t^{\min(a,b)}\)，剩下 \(1\)、正次 \(x\) 幂或正次 \(y\) 幂，证明所列元素张成 \(B\)。对于独立性，展开系数多项式后，三个族的单项式指数分别位于 \(a=b\)、\(a>b\)、\(a<b\)；在每一族内，指数又唯一确定系数中的 \(t\) 次幂及基元素。因此全部单项式互不重复，任何有限线性关系的系数都为零。

	张量这个自由模等于取相应指标集上的直和。直和逐坐标保留单射，故 \(B\) 平坦。无限基的自由模不能有限生成：有限组元素只涉及有限多个基方向，不能生成剩余方向。因此这个扩张不有限。

	令 \(K=k(t)\)。因 \(t\) 在 \(K\) 中非零且可逆，
	\[
		B\otimes_AK\cong K[x,y]/(xy-t)\cong K[x,x^{-1}],
	\]
	同构由 \(y=t/x\) 给出。对 \(a\ne0\)，闭点纤维同样为
	\[
		k[x,y]/(xy-a)\cong k[x,x^{-1}],
	\]
	而 \(a=0\) 的纤维为 \(k[x,y]/(xy)\)，其两个极小素理想为 \((x)\)、\((y)\)。各环的维数都为一，泛纤维和非零参数纤维为整环，零参数纤维则有两个分支。这里平坦性由自由模结构直接证明，并不要求每个纤维都不可约。
\end{solution}

\begin{exercise}\label{g5:ex:zero-dimensional-finite-type}
	设 \(B\ne0\) 为域 \(k\) 上的有限型代数。证明
	\[
		\dim B=0\quad\Longleftrightarrow\quad\dim_k B<\infty.
	\]
	说明允许 \(B\) 含非零幂零元，并用 \(k(t)\) 解释有限型假设为何必要。
\end{exercise}
\begin{solution}
	由 Noether 正规化，存在多项式子代数 \(C=k[u_1,\ldots,u_d]\subseteq B\)，使 \(B\) 为有限 \(C\) 模。整扩张保持维数，故
	\[
		\dim B=\dim C=d.
	\]
	若 \(\dim B=0\)，则 \(d=0\)，\(C=k\)，所以 \(B\) 为有限维 \(k\) 向量空间。反之，若 \(B\) 有限维，全部理想都是向量子空间，因此降链终止，\(B\) 为 Artin 环，由\cref{g5:thm:artinian-zero-dimension}得 \(\dim B=0\)。两向证明均未使用约化性；例如 \(k[\varepsilon]/(\varepsilon^n)\)、\(n\geq2\)，满足等价条件。

	域 \(k(t)\) 的 Krull 维数为零，但 \(1,t,t^2,\ldots\) 在 \(k\) 上线性无关，故其 \(k\) 维数无穷。由刚证的等价关系可知，它不能是有限型 \(k\) 代数。环作为模有限与作为代数有限型的区别，在这个例子中均不可略去。
\end{solution}

\begin{exercise}\label{g5:ex:nonnoetherian-principal-height}
	设 \(k\) 为域，
	\[
		A_0=k[x]_{(x)},\qquad K=k(x),\qquad V=K[y]_{(y)}.
	\]
	记 \(\varepsilon:V\to K\) 为 \(y=0\) 处的取值同态，定义
	\[
		R=\varepsilon^{-1}(A_0),\qquad
		\mathfrak p=yV,\qquad
		\mathfrak m=\varepsilon^{-1}(xA_0)\subset R.
	\]
	证明 \(R\) 为局部整环，\(\mathfrak m=xR\) 为其极大理想，且
	\[
		(0)\subsetneq\mathfrak p\subsetneq\mathfrak m
	\]
	为素链。再构造 \(R\) 中严格上升的理想链，说明主理想定理的 Noether 假设不能删去。
\end{exercise}
\begin{solution}
	\(V\) 是整环，故其子环 \(R\) 也是整环。映射 \(\varepsilon:R\to A_0\) 满射，因为每个 \(A_0\) 元素都可以视为与 \(y\) 无关的常数。其核为 \(yV=\mathfrak p\)，而 \(\mathfrak p\subset R\)，所以
	\[
		R/\mathfrak p\cong A_0,\qquad R/\mathfrak m\cong k.
	\]
	因此 \(\mathfrak p\) 为素理想，\(\mathfrak m\) 为极大理想。

	若 \(f\in R\setminus\mathfrak m\)，则 \(\varepsilon(f)\) 为 \(A_0\) 的单位，特别地非零。这使 \(f\) 在 \(V=K[y]_{(y)}\) 中可逆，且
	\[
		\varepsilon(f^{-1})=\varepsilon(f)^{-1}\in A_0,
	\]
	故 \(f^{-1}\in R\)。反之，\(R\) 的单位在 \(A_0\) 中的像也必须可逆。因此全部非单位恰为 \(\mathfrak m\)，\(R\) 是局部环。

	显然 \(xR\subseteq\mathfrak m\)。若 \(f\in\mathfrak m\)，则 \(f/x\in V\)，因为 \(x\) 是 \(K\) 中的非零常数；又
	\(\varepsilon(f/x)=\varepsilon(f)/x\in A_0\)，故 \(f/x\in R\)。这给出 \(\mathfrak m=xR\)。元素 \(y\ne0\) 属于 \(\mathfrak p\)，元素 \(x\) 属于 \(\mathfrak m\setminus\mathfrak p\)，所列素链严格。因此一个主理想本身就是高度至少为二的极大素理想。

	最后，每个 \(y/x^n\) 都属于 \(R\)，因为其取值为零。于是有上升链
	\[
		(y)R\subsetneq(y/x)R\subsetneq(y/x^2)R\subsetneq\cdots.
	\]
	包含关系来自 \(y/x^n=x(y/x^{n+1})\)。若某一步相等，则 \(y/x^{n+1}=r(y/x^n)\) 对某个 \(r\in R\) 成立。在整环的分式域中约去非零因子，得到 \(r=1/x\)，但 \(\varepsilon(1/x)=1/x\notin A_0\)，矛盾。故 \(R\) 不为 Noether 环，确实处在高度定理的假设之外。
\end{solution}

% !TeX root = ../../Book5.tex
% 纲要标题：## 第4章　正规化、离散赋值与曲线奇点
% 纲要位置：几何学计划纲要.md，第 4853 行。
% 正文状态：未撰写；本文件目前是框架。

% 纲要内容开始
% 
% > 本章目标与证明范围：把正规化对象、有限性证明和局部曲线计算独立成章。任何一般有限性结论都注明有限型或 Nagata 类条件，不能推广到任意 Noether 环。
% 
% 纲要内容结束

\chapter{正规化、离散赋值与曲线奇点}
\label{g5:ch:04}
\BookChapterNumber{g5:ch:04}

\begin{chapterabstract}
待撰写
\end{chapterabstract}

\begin{learninggoals}
待撰写
\end{learninggoals}

待撰写

% !TeX root = ../../Book5.tex
% 纲要标题：### 4.1 一维正规局部结构
% 纲要位置：几何学计划纲要.md，第 5151 行。
% 正文状态：未撰写；本文件目前是框架。

% 纲要内容开始
% 
% * 离散赋值环的等价刻画
% * 统一化元、阶和分式理想
% * Dedekind 整环的局部观点
% * 正规曲线的局部参数
% 
% 纲要内容结束

\section{一维正规局部结构}
\label{g5:sec:0401}

待撰写

\subsection{离散赋值环的等价刻画}
\label{g5:subsec:0401:01}
待撰写

\subsection{统一化元、阶和分式理想}
\label{g5:subsec:0401:02}
待撰写

\subsection{Dedekind 整环的局部观点}
\label{g5:subsec:0401:03}
待撰写

\subsection{正规曲线的局部参数}
\label{g5:subsec:0401:04}
待撰写
% 纲要条目排为编号小节标题，并进入目录与书签；说明段仍保留于纲要注释中。

% !TeX root = ../../Book5.tex
% 纲要标题：### 4.2 正规化的构造与有限性
% 纲要位置：几何学计划纲要.md，第 4864 行。
% 正文状态：未撰写；本文件目前是框架。

% 纲要内容开始
% 
% * 分式域及总分式环中的整闭包
% * Noether 正规化后的有限性归约
% * 有限可分扩张的迹配对与分母控制
% * 不可分情形的独立有限性处理
% 
%   只涉及有限个系数的纯不可分扩张，不能默设基域完美或 Frobenius 有限。
% 
% * 约化对象的各分支与粘合
% 
% 纲要内容结束

\section{正规化的构造与有限性}
\label{g5:sec:0402}

待撰写

\subsection{分式域及总分式环中的整闭包}
\label{g5:subsec:0402:01}
待撰写

\subsection{Noether 正规化后的有限性归约}
\label{g5:subsec:0402:02}
待撰写

\subsection{有限可分扩张的迹配对与分母控制}
\label{g5:subsec:0402:03}
待撰写

\subsection{不可分情形的独立有限性处理}
\label{g5:subsec:0402:04}
待撰写

\subsection{约化对象的各分支与粘合}
\label{g5:subsec:0402:05}
待撰写
% 纲要条目排为编号小节标题，并进入目录与书签；说明段仍保留于纲要注释中。

% !TeX root = ../../Book5.tex
% 纲要标题：### 4.3 导子理想及显式计算
% 纲要位置：几何学计划纲要.md，第 4875 行。
% 正文状态：未撰写；本文件目前是框架。

% 纲要内容开始
% 
% * 有限双有理扩张的导子
% * 尖点 k[t²,t³] 的正规化
% * 结点的两个原像和导子
% * 商模长度、δ 不变量和局部差异
% 
% 纲要内容结束

\section{导子理想及显式计算}
\label{g5:sec:0403}

待撰写

\subsection{有限双有理扩张的导子}
\label{g5:subsec:0403:01}
待撰写

\subsection{\texorpdfstring{尖点 \(k[t^2,t^3]\) 的正规化}{尖点 k[t²,t³] 的正规化}}
\label{g5:subsec:0403:02}
待撰写

\subsection{结点的两个原像和导子}
\label{g5:subsec:0403:03}
待撰写

\subsection{\texorpdfstring{商模长度、\(\delta\) 不变量和局部差异}{商模长度、\GeometryPdfDelta{} 不变量和局部差异}}
\label{g5:subsec:0403:04}
待撰写
% 纲要条目排为编号小节标题，并进入目录与书签；说明段仍保留于纲要注释中。

% !TeX root = ../../Book5.tex
% 纲要标题：### 4.4 有限双有理映射的限制
% 纲要位置：几何学计划纲要.md，第 4882 行。
% 正文状态：未撰写；本文件目前是框架。

% 纲要内容开始
% 
% * 正规化与同胚、同构的区别
% * 正规化和基域扩张
% * 非约化对象先取约化的约定
% * 有限性失效及一般理论的准确边界
% 
% ---
% 
% 纲要内容结束

\section{有限双有理映射的限制}
\label{g5:sec:0404}

待撰写

\subsection{正规化与同胚、同构的区别}
\label{g5:subsec:0404:01}
待撰写

\subsection{正规化和基域扩张}
\label{g5:subsec:0404:02}
待撰写

\subsection{非约化对象先取约化的约定}
\label{g5:subsec:0404:03}
待撰写

\subsection{有限性失效及一般理论的准确边界}
\label{g5:subsec:0404:04}
待撰写
% 纲要条目排为编号小节标题，并进入目录与书签；说明段仍保留于纲要注释中。


\begin{chaptersummary}
待撰写
\end{chaptersummary}

\BookExercises
待撰写

% !TeX root = ../../Book5.tex
% 纲要标题：## 第5章　切空间、Kähler 微分与局部方程
% 纲要位置：几何学计划纲要.md，第 4891 行。
% 正文状态：未撰写；本文件目前是框架。

% 纲要内容开始
% 
% > 本章目标与证明范围：把切向量、微分模、切锥及形式方程放在不同层次。正则性的定义在此出现，完整局部结构在后面证明；光滑性留待相对态射理论。
% 
% 纲要内容结束

\chapter{切空间、Kähler 微分与局部方程}
\label{g5:ch:05}
\BookChapterNumber{g5:ch:05}

\begin{chapterabstract}
待撰写
\end{chapterabstract}

\begin{learninggoals}
待撰写
\end{learninggoals}

待撰写

% !TeX root = ../../Book5.tex
% 纲要标题：### 5.1 普遍导子与微分模
% 纲要位置：几何学计划纲要.md，第 4895 行。
% 正文状态：未撰写；本文件目前是框架。

% 纲要内容开始
% 
% * 导子的基环线性
% * Kähler 微分的构造和泛性质
% * 多项式、局部化和张量的计算
% * 域扩张的可分与不可分例子
% 
% 纲要内容结束

\section{普遍导子与微分模}
\label{g5:sec:0501}

待撰写

\subsection{导子的基环线性}
\label{g5:subsec:0501:01}
待撰写

\subsection{Kähler 微分的构造和泛性质}
\label{g5:subsec:0501:02}
待撰写

\subsection{多项式、局部化和张量的计算}
\label{g5:subsec:0501:03}
待撰写

\subsection{域扩张的可分与不可分例子}
\label{g5:subsec:0501:04}
待撰写
% 纲要条目排为编号小节标题，并进入目录与书签；说明段仍保留于纲要注释中。

% !TeX root = ../../Book5.tex
% 纲要标题：### 5.2 余切正合列
% 纲要位置：几何学计划纲要.md，第 4902 行。
% 正文状态：未撰写；本文件目前是框架。

% 纲要内容开始
% 
% * 商环的余法序列
% * 复合映射的传递序列
% * 基变换相容
% * 左端未必单射的含义
% 
% 纲要内容结束

\section{余切正合列}
\label{g5:sec:0502}

待撰写

\subsection{商环的余法序列}
\label{g5:subsec:0502:01}
待撰写

\subsection{复合映射的传递序列}
\label{g5:subsec:0502:02}
待撰写

\subsection{基变换相容}
\label{g5:subsec:0502:03}
待撰写

\subsection{左端未必单射的含义}
\label{g5:subsec:0502:04}
待撰写
% 纲要条目排为编号小节标题，并进入目录与书签；说明段仍保留于纲要注释中。

% !TeX root = ../../Book5.tex
% 纲要标题：### 5.3 切空间和 Jacobian
% 纲要位置：几何学计划纲要.md，第 5203 行。
% 正文状态：未撰写；本文件目前是框架。

% 纲要内容开始
% 
% * 双数点及切向量
% * 极大理想模平方
% * Jacobian 矩阵的秩计算
% * 正则局部环的维数定义
% * 非完美域上正则与光滑的区别
% 
% 纲要内容结束

\section{切空间和 Jacobian}
\label{g5:sec:0503}

待撰写

\subsection{双数点及切向量}
\label{g5:subsec:0503:01}
待撰写

\subsection{极大理想模平方}
\label{g5:subsec:0503:02}
待撰写

\subsection{Jacobian 矩阵的秩计算}
\label{g5:subsec:0503:03}
待撰写

\subsection{正则局部环的维数定义}
\label{g5:subsec:0503:04}
待撰写

\subsection{非完美域上正则与光滑的区别}
\label{g5:subsec:0503:05}
待撰写
% 纲要条目排为编号小节标题，并进入目录与书签；说明段仍保留于纲要注释中。

% !TeX root = ../../Book5.tex
% 纲要标题：### 5.4 切锥与高阶局部信息
% 纲要位置：几何学计划纲要.md，第 4917 行。
% 正文状态：未撰写；本文件目前是框架。

% 纲要内容开始
% 
% * 伴随分次和最低次初始理想
% * 切锥与切空间
% * 尖点、结点和加厚点
% * 形式同构与有限阶截断的差异
% 
% ---
% 
% 纲要内容结束

\section{切锥与高阶局部信息}
\label{g5:sec:0504}

待撰写

\subsection{伴随分次和最低次初始理想}
\label{g5:subsec:0504:01}
待撰写

\subsection{切锥与切空间}
\label{g5:subsec:0504:02}
待撰写

\subsection{尖点、结点和加厚点}
\label{g5:subsec:0504:03}
待撰写

\subsection{形式同构与有限阶截断的差异}
\label{g5:subsec:0504:04}
待撰写
% 纲要条目排为编号小节标题，并进入目录与书签；说明段仍保留于纲要注释中。


\begin{chaptersummary}
待撰写
\end{chaptersummary}

\BookExercises
待撰写

% !TeX root = ../../Book5.tex
% 纲要标题：## 第6章　平坦性、纤维与泛自由性
% 纲要位置：几何学计划纲要.md，第 5220 行。
% 正文状态：未撰写；本文件目前是框架。

% 纲要内容开始
% 
% > 本章目标与证明范围：把平坦性作为相对几何的核心工具独立建立，完整证明实际使用的判据和泛自由性。纤维维数、矩阵秩、局部自由及平坦性之间的假设逐项落实。
% 
% 纲要内容结束

\chapter{平坦性、纤维与泛自由性}
\label{g5:ch:06}
\BookChapterNumber{g5:ch:06}

\begin{chapterabstract}
待撰写
\end{chapterabstract}

\begin{learninggoals}
待撰写
\end{learninggoals}

待撰写

% !TeX root = ../../Book5.tex
% 纲要标题：### 6.1 平坦性的代数判据
% 纲要位置：几何学计划纲要.md，第 5224 行。
% 正文状态：未撰写；本文件目前是框架。

% 纲要内容开始
% 
% * 张量右正合与左端障碍
% * 有限关系判据
% * 有限生成理想张量判据
% * 局部化及局部检测
% * 无挠和平坦的不同范围
% 
% 纲要内容结束

\section{平坦性的代数判据}
\label{g5:sec:0601}

待撰写

\subsection{张量右正合与左端障碍}
\label{g5:subsec:0601:01}
待撰写

\subsection{有限关系判据}
\label{g5:subsec:0601:02}
待撰写

\subsection{有限生成理想张量判据}
\label{g5:subsec:0601:03}
待撰写

\subsection{局部化及局部检测}
\label{g5:subsec:0601:04}
待撰写

\subsection{无挠和平坦的不同范围}
\label{g5:subsec:0601:05}
待撰写
% 纲要条目排为编号小节标题，并进入目录与书签；说明段仍保留于纲要注释中。

% !TeX root = ../../Book5.tex
% 纲要标题：### 6.2 有限呈示与 Fitting 理想
% 纲要位置：几何学计划纲要.md，第 5232 行。
% 正文状态：未撰写；本文件目前是框架。

% 纲要内容开始
% 
% * 呈示子式和 Fitting 理想
% * 与呈示选择无关及基变换
% * 剩余域上的秩条件
% * 局部自由秩的理想判据
% * 幂零基底上恒定纤维维数的反例
% 
% 纲要内容结束

\section{有限呈示与 Fitting 理想}
\label{g5:sec:0602}

待撰写

\subsection{呈示子式和 Fitting 理想}
\label{g5:subsec:0602:01}
待撰写

\subsection{与呈示选择无关及基变换}
\label{g5:subsec:0602:02}
待撰写

\subsection{剩余域上的秩条件}
\label{g5:subsec:0602:03}
待撰写

\subsection{局部自由秩的理想判据}
\label{g5:subsec:0602:04}
待撰写

\subsection{幂零基底上恒定纤维维数的反例}
\label{g5:subsec:0602:05}
待撰写
% 纲要条目排为编号小节标题，并进入目录与书签；说明段仍保留于纲要注释中。

% !TeX root = ../../Book5.tex
% 纲要标题：### 6.3 泛自由性的证明
% 纲要位置：几何学计划纲要.md，第 4946 行。
% 正文状态：未撰写；本文件目前是框架。

% 纲要内容开始
% 
% * Noether 整基环上的有限型代数
% * 分式域上的标准单项式
% * 有限关系与清除分母
% * 局部化后的自由性
% * 自由秩可能无限的说明
% 
% 纲要内容结束

\section{泛自由性的证明}
\label{g5:sec:0603}

待撰写

\subsection{Noether 整基环上的有限型代数}
\label{g5:subsec:0603:01}
待撰写

\subsection{分式域上的标准单项式}
\label{g5:subsec:0603:02}
待撰写

\subsection{有限关系与清除分母}
\label{g5:subsec:0603:03}
待撰写

\subsection{局部化后的自由性}
\label{g5:subsec:0603:04}
待撰写

\subsection{自由秩可能无限的说明}
\label{g5:subsec:0603:05}
待撰写
% 纲要条目排为编号小节标题，并进入目录与书签；说明段仍保留于纲要注释中。

% !TeX root = ../../Book5.tex
% 纲要标题：### 6.4 忠实平坦与几何族
% 纲要位置：几何学计划纲要.md，第 5248 行。
% 正文状态：未撰写；本文件目前是框架。

% 纲要内容开始
% 
% * 消失和正合性的反映
% * 谱映射满射的判据
% * 泛平坦性及开放条件
% * 平坦退化的显式方程
% * 为下降准备的双重基变换
% 
% ---
% 
% 纲要内容结束

\section{忠实平坦与几何族}
\label{g5:sec:0604}

待撰写

\subsection{消失和正合性的反映}
\label{g5:subsec:0604:01}
待撰写

\subsection{谱映射满射的判据}
\label{g5:subsec:0604:02}
待撰写

\subsection{泛平坦性及开放条件}
\label{g5:subsec:0604:03}
待撰写

\subsection{平坦退化的显式方程}
\label{g5:subsec:0604:04}
待撰写

\subsection{为下降准备的双重基变换}
\label{g5:subsec:0604:05}
待撰写
% 纲要条目排为编号小节标题，并进入目录与书签；说明段仍保留于纲要注释中。


\begin{chaptersummary}
待撰写
\end{chaptersummary}

\BookExercises
待撰写

% !TeX root = ../../Book5.tex
% 纲要标题：## 第7章　滤过、Artin–Rees 与完备化
% 纲要位置：几何学计划纲要.md，第 4964 行。
% 正文状态：未撰写；本文件目前是框架。

% 纲要内容开始
% 
% > 本章目标与证明范围：以形式邻域为何保持有限模正合列为主问题，连续完成 Artin–Rees、张量表示和完备化忠实平坦性。逆极限的一般非正合性必须在证明中实际处理。
% 
% 纲要内容结束

\chapter{滤过、Artin–Rees 与完备化}
\label{g5:ch:07}
\BookChapterNumber{g5:ch:07}

\begin{chapterabstract}
待撰写
\end{chapterabstract}

\begin{learninggoals}
待撰写
\end{learninggoals}

待撰写

% !TeX root = ../../Book5.tex
% 纲要标题：### 7.1 滤过及伴随分次
% 纲要位置：几何学计划纲要.md，第 4968 行。
% 正文状态：未撰写；本文件目前是框架。

% 纲要内容开始
% 
% * 理想进拓扑和分离性
% * Rees 环及有限生成
% * 伴随分次与切锥
% * Hilbert–Samuel 函数的所需增长工具
% 
% 纲要内容结束

\section{滤过及伴随分次}
\label{g5:sec:0701}

待撰写

\subsection{理想进拓扑和分离性}
\label{g5:subsec:0701:01}
待撰写

\subsection{Rees 环及有限生成}
\label{g5:subsec:0701:02}
待撰写

\subsection{伴随分次与切锥}
\label{g5:subsec:0701:03}
待撰写

\subsection{Hilbert–Samuel 函数的所需增长工具}
\label{g5:subsec:0701:04}
待撰写
% 纲要条目排为编号小节标题，并进入目录与书签；说明段仍保留于纲要注释中。

% !TeX root = ../../Book5.tex
% 纲要标题：### 7.2 Artin–Rees 与交定理
% 纲要位置：几何学计划纲要.md，第 4975 行。
% 正文状态：未撰写；本文件目前是框架。

% 纲要内容开始
% 
% * 诱导滤过与本来滤过
% * Rees 模的 Noether 论证
% * Artin–Rees 引理
% * Krull 交定理的局部形式
% 
% 纲要内容结束

\section{Artin–Rees 与交定理}
\label{g5:sec:0702}

待撰写

\subsection{诱导滤过与本来滤过}
\label{g5:subsec:0702:01}
待撰写

\subsection{Rees 模的 Noether 论证}
\label{g5:subsec:0702:02}
待撰写

\subsection{Artin–Rees 引理}
\label{g5:subsec:0702:03}
待撰写

\subsection{Krull 交定理的局部形式}
\label{g5:subsec:0702:04}
待撰写
% 纲要条目排为编号小节标题，并进入目录与书签；说明段仍保留于纲要注释中。

% !TeX root = ../../Book5.tex
% 纲要标题：### 7.3 有限模的完备化
% 纲要位置：几何学计划纲要.md，第 4982 行。
% 正文状态：未撰写；本文件目前是框架。

% 纲要内容开始
% 
% * 逆系统和 Mittag–Leffler 的最低引理
% * 完备化的正合性
% * 有限模的张量表示
% * 完备化环的 Noether 性
% 
% 纲要内容结束

\section{有限模的完备化}
\label{g5:sec:0703}

待撰写

\subsection{逆系统和 Mittag–Leffler 的最低引理}
\label{g5:subsec:0703:01}
待撰写

\subsection{完备化的正合性}
\label{g5:subsec:0703:02}
待撰写

\subsection{有限模的张量表示}
\label{g5:subsec:0703:03}
待撰写

\subsection{完备化环的 Noether 性}
\label{g5:subsec:0703:04}
待撰写
% 纲要条目排为编号小节标题，并进入目录与书签；说明段仍保留于纲要注释中。

% !TeX root = ../../Book5.tex
% 纲要标题：### 7.4 忠实平坦及局部结构
% 纲要位置：几何学计划纲要.md，第 4989 行。
% 正文状态：未撰写；本文件目前是框架。

% 纲要内容开始
% 
% * 理想张量判据推出平坦
% * 剩余域不变与忠实性
% * 维数和伴随分次比较
% * 扩张收缩与局部检测
% * 完备化不是局部化或 Hensel 化
% 
% 纲要内容结束

\section{忠实平坦及局部结构}
\label{g5:sec:0704}

待撰写

\subsection{理想张量判据推出平坦}
\label{g5:subsec:0704:01}
待撰写

\subsection{剩余域不变与忠实性}
\label{g5:subsec:0704:02}
待撰写

\subsection{维数和伴随分次比较}
\label{g5:subsec:0704:03}
待撰写

\subsection{扩张收缩与局部检测}
\label{g5:subsec:0704:04}
待撰写

\subsection{完备化不是局部化或 Hensel 化}
\label{g5:subsec:0704:05}
待撰写
% 纲要条目排为编号小节标题，并进入目录与书签；说明段仍保留于纲要注释中。

% !TeX root = ../../Book5.tex
% 纲要标题：### 7.5 形式方程与有效近似
% 纲要位置：几何学计划纲要.md，第 5291 行。
% 正文状态：未撰写；本文件目前是框架。

% 纲要内容开始
% 
% * 有限阶邻域和相容族
% * 多项式、代数幂级数与任意形式级数
% * 收敛、代数化及逼近的不同问题
% 
% ---
% 
% 纲要内容结束

\section{形式方程与有效近似}
\label{g5:sec:0705}

待撰写

\subsection{有限阶邻域和相容族}
\label{g5:subsec:0705:01}
待撰写

\subsection{多项式、代数幂级数与任意形式级数}
\label{g5:subsec:0705:02}
待撰写

\subsection{收敛、代数化及逼近的不同问题}
\label{g5:subsec:0705:03}
待撰写
% 纲要条目排为编号小节标题，并进入目录与书签；说明段仍保留于纲要注释中。


\begin{chaptersummary}
待撰写
\end{chaptersummary}

\BookExercises
待撰写

% !TeX root = ../../Book5.tex
% 纲要标题：## 第8章　消解、Koszul 复形与深度
% 纲要位置：几何学计划纲要.md，第 5005 行。
% 正文状态：未撰写；本文件目前是框架。

% 纲要内容开始
% 
% > 本章目标与证明范围：局部同调工具自成一章，实际计算非横截交和深度。先建立 Ext、Tor 与维数平移，再使用 Koszul 正合性和深度引理，不把抽象函子当作已学黑箱。
% 
% 纲要内容结束

\chapter{消解、Koszul 复形与深度}
\label{g5:ch:08}
\BookChapterNumber{g5:ch:08}

\begin{chapterabstract}
待撰写
\end{chapterabstract}

\begin{learninggoals}
待撰写
\end{learninggoals}

待撰写

% !TeX root = ../../Book5.tex
% 纲要标题：### 8.1 自由消解与导出障碍
% 纲要位置：几何学计划纲要.md，第 5009 行。
% 正文状态：未撰写；本文件目前是框架。

% 纲要内容开始
% 
% * 自由消解的构造
% * 比较定理和链同伦
% * Ext、Tor 与长正合列
% * 维数平移和合冲模
% * 非平坦张量的显式 Tor
% 
% 纲要内容结束

\section{自由消解与导出障碍}
\label{g5:sec:0801}

待撰写

\subsection{自由消解的构造}
\label{g5:subsec:0801:01}
待撰写

\subsection{比较定理和链同伦}
\label{g5:subsec:0801:02}
待撰写

\subsection{Ext、Tor 与长正合列}
\label{g5:subsec:0801:03}
待撰写

\subsection{维数平移和合冲模}
\label{g5:subsec:0801:04}
待撰写

\subsection{非平坦张量的显式 Tor}
\label{g5:subsec:0801:05}
待撰写
% 纲要条目排为编号小节标题，并进入目录与书签；说明段仍保留于纲要注释中。

% !TeX root = ../../Book5.tex
% 纲要标题：### 8.2 正则序列和 Koszul
% 纲要位置：几何学计划纲要.md，第 5017 行。
% 正文状态：未撰写；本文件目前是框架。

% 纲要内容开始
% 
% * 逐次非零因子及末端非零条件
% * 外代数上的微分
% * 正则序列给出消解
% * 局部有限模上的反向正合判据
% * 交换次序和基变换的限制
% 
% 纲要内容结束

\section{正则序列和 Koszul}
\label{g5:sec:0802}

待撰写

\subsection{逐次非零因子及末端非零条件}
\label{g5:subsec:0802:01}
待撰写

\subsection{外代数上的微分}
\label{g5:subsec:0802:02}
待撰写

\subsection{正则序列给出消解}
\label{g5:subsec:0802:03}
待撰写

\subsection{局部有限模上的反向正合判据}
\label{g5:subsec:0802:04}
待撰写

\subsection{交换次序和基变换的限制}
\label{g5:subsec:0802:05}
待撰写
% 纲要条目排为编号小节标题，并进入目录与书签；说明段仍保留于纲要注释中。

% !TeX root = ../../Book5.tex
% 纲要标题：### 8.3 深度的不同刻画
% 纲要位置：几何学计划纲要.md，第 5025 行。
% 正文状态：未撰写；本文件目前是框架。

% 纲要内容开始
% 
% * 极大正则序列长度
% * Ext 检测与等价证明
% * 深度不超过维数
% * 深度引理
% * Serre Sᵣ 条件的局部表述
% 
% 纲要内容结束

\section{深度的不同刻画}
\label{g5:sec:0803}

待撰写

\subsection{极大正则序列长度}
\label{g5:subsec:0803:01}
待撰写

\subsection{Ext 检测与等价证明}
\label{g5:subsec:0803:02}
待撰写

\subsection{深度不超过维数}
\label{g5:subsec:0803:03}
待撰写

\subsection{深度引理}
\label{g5:subsec:0803:04}
待撰写

\subsection{\texorpdfstring{Serre \(S_r\) 条件的局部表述}{Serre Sᵣ 条件的局部表述}}
\label{g5:subsec:0803:05}
待撰写
% 纲要条目排为编号小节标题，并进入目录与书签；说明段仍保留于纲要注释中。

% !TeX root = ../../Book5.tex
% 纲要标题：### 8.4 极小自由消解和投射维数
% 纲要位置：几何学计划纲要.md，第 5033 行。
% 正文状态：未撰写；本文件目前是框架。

% 纲要内容开始
% 
% * 局部极小生成和极小微分
% * Betti 数及剩余域张量
% * 有限投射维数的检测
% * Auslander–Buchsbaum 公式的证明
% 
% 纲要内容结束

\section{极小自由消解和投射维数}
\label{g5:sec:0804}

待撰写

\subsection{局部极小生成和极小微分}
\label{g5:subsec:0804:01}
待撰写

\subsection{Betti 数及剩余域张量}
\label{g5:subsec:0804:02}
待撰写

\subsection{有限投射维数的检测}
\label{g5:subsec:0804:03}
待撰写

\subsection{Auslander–Buchsbaum 公式的证明}
\label{g5:subsec:0804:04}
待撰写
% 纲要条目排为编号小节标题，并进入目录与书签；说明段仍保留于纲要注释中。

% !TeX root = ../../Book5.tex
% 纲要标题：### 8.5 交与嵌入的模型
% 纲要位置：几何学计划纲要.md，第 5040 行。
% 正文状态：未撰写；本文件目前是框架。

% 纲要内容开始
% 
% * 横截和非横截的 Tor 对照
% * 超曲面与完全交
% * 嵌入分支导致的深度下降
% * 同一支集的不同局部同调信息
% 
% ---
% 
% 纲要内容结束

\section{交与嵌入的模型}
\label{g5:sec:0805}

待撰写

\subsection{横截和非横截的 Tor 对照}
\label{g5:subsec:0805:01}
待撰写

\subsection{超曲面与完全交}
\label{g5:subsec:0805:02}
待撰写

\subsection{嵌入分支导致的深度下降}
\label{g5:subsec:0805:03}
待撰写

\subsection{同一支集的不同局部同调信息}
\label{g5:subsec:0805:04}
待撰写
% 纲要条目排为编号小节标题，并进入目录与书签；说明段仍保留于纲要注释中。


\begin{chaptersummary}
待撰写
\end{chaptersummary}

\BookExercises
待撰写

% !TeX root = ../../Book5.tex
% 纲要标题：## 第9章　正则、Cohen–Macaulay 与正规局部结构
% 纲要位置：几何学计划纲要.md，第 5049 行。
% 正文状态：未撰写；本文件目前是框架。

% 纲要内容开始
% 
% > 本章目标与证明范围：把原来一个小节承载的结构定理展开为独立章。完整证明主范围内的正则性判据、CM 性和 Serre 正规性；除子章实际需要的余维一结构和阶乘性也有证明落点。
% 
% 纲要内容结束

\chapter{正则、Cohen–Macaulay 与正规局部结构}
\label{g5:ch:09}
\BookChapterNumber{g5:ch:09}

\begin{chapterabstract}
待撰写
\end{chapterabstract}

\begin{learninggoals}
待撰写
\end{learninggoals}

待撰写

% !TeX root = ../../Book5.tex
% 纲要标题：### 9.1 正则局部环的基本结构
% 纲要位置：几何学计划纲要.md，第 5053 行。
% 正文状态：未撰写；本文件目前是框架。

% 纲要内容开始
% 
% * 参数系和伴随分次
% * 维数与极小生成元
% * 正则参数与非零因子
% * 局部化、完备化及整环性
% 
% 纲要内容结束

\section{正则局部环的基本结构}
\label{g5:sec:0901}

待撰写

\subsection{参数系和伴随分次}
\label{g5:subsec:0901:01}
待撰写

\subsection{维数与极小生成元}
\label{g5:subsec:0901:02}
待撰写

\subsection{正则参数与非零因子}
\label{g5:subsec:0901:03}
待撰写

\subsection{局部化、完备化及整环性}
\label{g5:subsec:0901:04}
待撰写
% 纲要条目排为编号小节标题，并进入目录与书签；说明段仍保留于纲要注释中。

% !TeX root = ../../Book5.tex
% 纲要标题：### 9.2 正则性的同调刻画
% 纲要位置：几何学计划纲要.md，第 5354 行。
% 正文状态：未撰写；本文件目前是框架。

% 纲要内容开始
% 
% * 极小消解的剩余域模型
% * 剩余域有限投射维数推出正则的关键归纳
% * Auslander–Buchsbaum–Serre 判据
% * 整体维数和有限模消解
% 
% 纲要内容结束

\section{正则性的同调刻画}
\label{g5:sec:0902}

待撰写

\subsection{极小消解的剩余域模型}
\label{g5:subsec:0902:01}
待撰写

\subsection{剩余域有限投射维数推出正则的关键归纳}
\label{g5:subsec:0902:02}
待撰写

\subsection{Auslander–Buchsbaum–Serre 判据}
\label{g5:subsec:0902:03}
待撰写

\subsection{整体维数和有限模消解}
\label{g5:subsec:0902:04}
待撰写
% 纲要条目排为编号小节标题，并进入目录与书签；说明段仍保留于纲要注释中。

% !TeX root = ../../Book5.tex
% 纲要标题：### 9.3 Cohen–Macaulay 与完全交
% 纲要位置：几何学计划纲要.md，第 5361 行。
% 正文状态：未撰写；本文件目前是框架。

% 纲要内容开始
% 
% * 参数系和正则序列
% * 商去正则元素及局部化
% * 正则蕴含 CM
% * 完全交及 Koszul 消解
% * CM 而非完全交的模型
% 
% 纲要内容结束

\section{Cohen–Macaulay 与完全交}
\label{g5:sec:0903}

待撰写

\subsection{参数系和正则序列}
\label{g5:subsec:0903:01}
待撰写

\subsection{商去正则元素及局部化}
\label{g5:subsec:0903:02}
待撰写

\subsection{正则蕴含 CM}
\label{g5:subsec:0903:03}
待撰写

\subsection{完全交及 Koszul 消解}
\label{g5:subsec:0903:04}
待撰写

\subsection{CM 而非完全交的模型}
\label{g5:subsec:0903:05}
待撰写
% 纲要条目排为编号小节标题，并进入目录与书签；说明段仍保留于纲要注释中。

% !TeX root = ../../Book5.tex
% 纲要标题：### 9.4 Serre 正规性判据
% 纲要位置：几何学计划纲要.md，第 5075 行。
% 正文状态：未撰写；本文件目前是框架。

% 纲要内容开始
% 
% * R₁ 与 S₂ 分别控制什么
% * 非零因子和深度的关键引理
% * R₁+S₂ 与正规性的等价
% * 高度一局部化的交表示
% * 一般约化与整环情形的区别
% 
% 纲要内容结束

\section{Serre 正规性判据}
\label{g5:sec:0904}

待撰写

\subsection{\texorpdfstring{\(R_1\) 与 \(S_2\) 分别控制什么}{R₁ 与 S₂ 分别控制什么}}
\label{g5:subsec:0904:01}
待撰写

\subsection{非零因子和深度的关键引理}
\label{g5:subsec:0904:02}
待撰写

\subsection{\texorpdfstring{\(R_1+S_2\) 与正规性的等价}{R₁+S₂ 与正规性的等价}}
\label{g5:subsec:0904:03}
待撰写

\subsection{高度一局部化的交表示}
\label{g5:subsec:0904:04}
待撰写

\subsection{一般约化与整环情形的区别}
\label{g5:subsec:0904:05}
待撰写
% 纲要条目排为编号小节标题，并进入目录与书签；说明段仍保留于纲要注释中。

% !TeX root = ../../Book5.tex
% 纲要标题：### 9.5 正则局部环阶乘性的证明
% 纲要位置：几何学计划纲要.md，第 5083 行。
% 正文状态：未撰写；本文件目前是框架。

% 纲要内容开始
% 
% * 高度一素理想主性的归约
% * 有限自由消解与行列式线
% * 主开集上的可逆模控制
% * 沿正则参数的维数归纳
% * 局部为主与全局平凡的区别
% * 由高度一主性得到唯一分解
% 
% 纲要内容结束

\section{正则局部环阶乘性的证明}
\label{g5:sec:0905}

待撰写

\subsection{高度一素理想主性的归约}
\label{g5:subsec:0905:01}
待撰写

\subsection{有限自由消解与行列式线}
\label{g5:subsec:0905:02}
待撰写

\subsection{主开集上的可逆模控制}
\label{g5:subsec:0905:03}
待撰写

\subsection{沿正则参数的维数归纳}
\label{g5:subsec:0905:04}
待撰写

\subsection{局部为主与全局平凡的区别}
\label{g5:subsec:0905:05}
待撰写

\subsection{由高度一主性得到唯一分解}
\label{g5:subsec:0905:06}
待撰写
% 纲要条目排为编号小节标题，并进入目录与书签；说明段仍保留于纲要注释中。

% !TeX root = ../../Book5.tex
% 纲要标题：### 9.6 秩一对象和奇点层次
% 纲要位置：几何学计划纲要.md，第 5386 行。
% 正文状态：未撰写；本文件目前是框架。

% 纲要内容开始
% 
% * 高度一局部化及反射包络
% * 可逆模与秩一反射模
% * 正规非正则和 CM 非正规模型
% * 几何光滑性的基域条件
% 
% ---
% 
% 纲要内容结束

\section{秩一对象和奇点层次}
\label{g5:sec:0906}

待撰写

\subsection{高度一局部化及反射包络}
\label{g5:subsec:0906:01}
待撰写

\subsection{可逆模与秩一反射模}
\label{g5:subsec:0906:02}
待撰写

\subsection{正规非正则和 CM 非正规模型}
\label{g5:subsec:0906:03}
待撰写

\subsection{几何光滑性的基域条件}
\label{g5:subsec:0906:04}
待撰写
% 纲要条目排为编号小节标题，并进入目录与书签；说明段仍保留于纲要注释中。


\begin{chaptersummary}
待撰写
\end{chaptersummary}

\BookExercises
待撰写


% !TeX root = ../Book5.tex
% 纲要标题：# 第二编　概形、态射与相对构造
% 纲要位置：几何学计划纲要.md，第 5101 行。
% 正文状态：未撰写；本文件目前是框架。

% 纲要内容开始
% 
% 纲要内容结束

\BookPart{概形、态射与相对构造}{g5:part:02}
\BookPartGuide
待撰写

% !TeX root = ../../Book5.tex
% 纲要标题：## 第10章　层、局部环空间与概形
% 纲要位置：几何学计划纲要.md，第 5397 行。
% 正文状态：未撰写；本文件目前是框架。

% 纲要内容开始
% 
% > 本章目标与证明范围：从局部化函数构造结构层，证明仿射反变等价，再处理粘合和点函子。概形不先被简化为闭点集合，也不预设约化、正规或光滑。
% 
% 纲要内容结束

\chapter{层、局部环空间与概形}
\label{g5:ch:10}
\BookChapterNumber{g5:ch:10}

\begin{chapterabstract}
待撰写
\end{chapterabstract}

\begin{learninggoals}
待撰写
\end{learninggoals}

待撰写

% !TeX root = ../../Book5.tex
% 纲要标题：### 10.1 层与局部数据
% 纲要位置：几何学计划纲要.md，第 5107 行。
% 正文状态：未撰写；本文件目前是框架。

% 纲要内容开始
% 
% * 预层、层化及茎
% * 基上的层和粘合引理
% * 局部环空间及局部同态
% * Zariski 空间的拟紧性与一般点
% 
% 纲要内容结束

\section{层与局部数据}
\label{g5:sec:1001}

待撰写

\subsection{预层、层化及茎}
\label{g5:subsec:1001:01}
待撰写

\subsection{基上的层和粘合引理}
\label{g5:subsec:1001:02}
待撰写

\subsection{局部环空间及局部同态}
\label{g5:subsec:1001:03}
待撰写

\subsection{Zariski 空间的拟紧性与一般点}
\label{g5:subsec:1001:04}
待撰写
% 纲要条目排为编号小节标题，并进入目录与书签；说明段仍保留于纲要注释中。

% !TeX root = ../../Book5.tex
% 纲要标题：### 10.2 仿射概形的结构层
% 纲要位置：几何学计划纲要.md，第 5408 行。
% 正文状态：未撰写；本文件目前是框架。

% 纲要内容开始
% 
% * 主开集上的局部化
% * 层条件的有限覆盖证明
% * 茎与局部环的识别
% * 整体截面恢复原环
% 
% 纲要内容结束

\section{仿射概形的结构层}
\label{g5:sec:1002}

待撰写

\subsection{主开集上的局部化}
\label{g5:subsec:1002:01}
待撰写

\subsection{层条件的有限覆盖证明}
\label{g5:subsec:1002:02}
待撰写

\subsection{茎与局部环的识别}
\label{g5:subsec:1002:03}
待撰写

\subsection{整体截面恢复原环}
\label{g5:subsec:1002:04}
待撰写
% 纲要条目排为编号小节标题，并进入目录与书签；说明段仍保留于纲要注释中。

% !TeX root = ../../Book5.tex
% 纲要标题：### 10.3 仿射态射和粘合
% 纲要位置：几何学计划纲要.md，第 5415 行。
% 正文状态：未撰写；本文件目前是框架。

% 纲要内容开始
% 
% * 环同态诱导的局部环空间态射
% * 仿射反变对应的证明
% * 开子概形和闭子概形
% * 沿开集粘合
% * 双原点及非约化模型
% 
% 纲要内容结束

\section{仿射态射和粘合}
\label{g5:sec:1003}

待撰写

\subsection{环同态诱导的局部环空间态射}
\label{g5:subsec:1003:01}
待撰写

\subsection{仿射反变对应的证明}
\label{g5:subsec:1003:02}
待撰写

\subsection{开子概形和闭子概形}
\label{g5:subsec:1003:03}
待撰写

\subsection{沿开集粘合}
\label{g5:subsec:1003:04}
待撰写

\subsection{双原点及非约化模型}
\label{g5:subsec:1003:05}
待撰写
% 纲要条目排为编号小节标题，并进入目录与书签；说明段仍保留于纲要注释中。

% !TeX root = ../../Book5.tex
% 纲要标题：### 10.4 点函子和相对观点
% 纲要位置：几何学计划纲要.md，第 5129 行。
% 正文状态：未撰写；本文件目前是框架。

% 纲要内容开始
% 
% * T 值点及函数的拉回
% * Yoneda 的所需形式
% * 从双数点到切空间
% * 同一底层空间上的不同概形结构
% 
% ---
% 
% 纲要内容结束

\section{点函子和相对观点}
\label{g5:sec:1004}

待撰写

\subsection{T 值点及函数的拉回}
\label{g5:subsec:1004:01}
待撰写

\subsection{Yoneda 的所需形式}
\label{g5:subsec:1004:02}
待撰写

\subsection{从双数点到切空间}
\label{g5:subsec:1004:03}
待撰写

\subsection{同一底层空间上的不同概形结构}
\label{g5:subsec:1004:04}
待撰写
% 纲要条目排为编号小节标题，并进入目录与书签；说明段仍保留于纲要注释中。


\begin{chaptersummary}
待撰写
\end{chaptersummary}

\BookExercises
待撰写

% !TeX root = ../../Book5.tex
% 纲要标题：## 第11章　纤维积、构造集与纤维维数
% 纲要位置：几何学计划纲要.md，第 5138 行。
% 正文状态：未撰写；本文件目前是框架。

% 纲要内容开始
% 
% > 本章目标与证明范围：把相对对象的构造和其像、纤维的几何控制相连。Chevalley 及纤维维数定理都有实际证明；点处局部纤维维数与整个纤维维数不能混为同一个半连续函数。
% 
% 纲要内容结束

\chapter{纤维积、构造集与纤维维数}
\label{g5:ch:11}
\BookChapterNumber{g5:ch:11}

\begin{chapterabstract}
待撰写
\end{chapterabstract}

\begin{learninggoals}
待撰写
\end{learninggoals}

待撰写

% !TeX root = ../../Book5.tex
% 纲要标题：### 11.1 纤维积和基变换
% 纲要位置：几何学计划纲要.md，第 5436 行。
% 正文状态：未撰写；本文件目前是框架。

% 纲要内容开始
% 
% * 仿射张量积及泛性质
% * 一般纤维积的粘合
% * 几何点和几何纤维
% * 不可分扩张及幂零纤维
% 
% 纲要内容结束

\section{纤维积和基变换}
\label{g5:sec:1101}

待撰写

\subsection{仿射张量积及泛性质}
\label{g5:subsec:1101:01}
待撰写

\subsection{一般纤维积的粘合}
\label{g5:subsec:1101:02}
待撰写

\subsection{几何点和几何纤维}
\label{g5:subsec:1101:03}
待撰写

\subsection{不可分扩张及幂零纤维}
\label{g5:subsec:1101:04}
待撰写
% 纲要条目排为编号小节标题，并进入目录与书签；说明段仍保留于纲要注释中。

% !TeX root = ../../Book5.tex
% 纲要标题：### 11.2 态射的有限性和稳定性
% 纲要位置：几何学计划纲要.md，第 5443 行。
% 正文状态：未撰写；本文件目前是框架。

% 纲要内容开始
% 
% * 有限型、有限呈示和拟紧
% * 有限与拟有限的纤维区别
% * 性质的局部性、复合及基变换
% * 相对维数的初步模型
% 
% 纲要内容结束

\section{态射的有限性和稳定性}
\label{g5:sec:1102}

待撰写

\subsection{有限型、有限呈示和拟紧}
\label{g5:subsec:1102:01}
待撰写

\subsection{有限与拟有限的纤维区别}
\label{g5:subsec:1102:02}
待撰写

\subsection{性质的局部性、复合及基变换}
\label{g5:subsec:1102:03}
待撰写

\subsection{相对维数的初步模型}
\label{g5:subsec:1102:04}
待撰写
% 纲要条目排为编号小节标题，并进入目录与书签；说明段仍保留于纲要注释中。

% !TeX root = ../../Book5.tex
% 纲要标题：### 11.3 Chevalley 定理
% 纲要位置：几何学计划纲要.md，第 5156 行。
% 正文状态：未撰写；本文件目前是框架。

% 纲要内容开始
% 
% * 构造集和局部闭分解
% * 支配有限型映射的像包含开集
% * Noether 归纳
% * 构造集像定理的证明
% * 投影像及消去闭包的实际差别
% 
% 纲要内容结束

\section{Chevalley 定理}
\label{g5:sec:1103}

待撰写

\subsection{构造集和局部闭分解}
\label{g5:subsec:1103:01}
待撰写

\subsection{支配有限型映射的像包含开集}
\label{g5:subsec:1103:02}
待撰写

\subsection{Noether 归纳}
\label{g5:subsec:1103:03}
待撰写

\subsection{构造集像定理的证明}
\label{g5:subsec:1103:04}
待撰写

\subsection{投影像及消去闭包的实际差别}
\label{g5:subsec:1103:05}
待撰写
% 纲要条目排为编号小节标题，并进入目录与书签；说明段仍保留于纲要注释中。

% !TeX root = ../../Book5.tex
% 纲要标题：### 11.4 维数定理及半连续性
% 纲要位置：几何学计划纲要.md，第 5458 行。
% 正文状态：未撰写；本文件目前是框架。

% 纲要内容开始
% 
% * 有限型整簇的维数差
% * 一般纤维与特殊纤维
% * 局部纤维维数的上半连续
% * 固有情形整体纤维维数的控制
% * 平坦性带来的额外结构
% 
% ---
% 
% 纲要内容结束

\section{维数定理及半连续性}
\label{g5:sec:1104}

待撰写

\subsection{有限型整簇的维数差}
\label{g5:subsec:1104:01}
待撰写

\subsection{一般纤维与特殊纤维}
\label{g5:subsec:1104:02}
待撰写

\subsection{局部纤维维数的上半连续}
\label{g5:subsec:1104:03}
待撰写

\subsection{固有情形整体纤维维数的控制}
\label{g5:subsec:1104:04}
待撰写

\subsection{平坦性带来的额外结构}
\label{g5:subsec:1104:05}
待撰写
% 纲要条目排为编号小节标题，并进入目录与书签；说明段仍保留于纲要注释中。


\begin{chaptersummary}
待撰写
\end{chaptersummary}

\BookExercises
待撰写

% !TeX root = ../../Book5.tex
% 纲要标题：## 第12章　分离性、固有性与 Zariski 主定理
% 纲要位置：几何学计划纲要.md，第 5174 行。
% 正文状态：未撰写；本文件目前是框架。

% 纲要内容开始
% 
% > 本章目标与证明范围：按“对角的唯一性—赋值的存在性—拟有限的有限化”展开。Zariski 主定理的环论关键引理独立证明；不把它解释成一般上同调长正合列的来源。
% 
% 纲要内容结束

\chapter{分离性、固有性与 Zariski 主定理}
\label{g5:ch:12}
\BookChapterNumber{g5:ch:12}

\begin{chapterabstract}
待撰写
\end{chapterabstract}

\begin{learninggoals}
待撰写
\end{learninggoals}

待撰写

% !TeX root = ../../Book5.tex
% 纲要标题：### 12.1 对角和分离性
% 纲要位置：几何学计划纲要.md，第 5178 行。
% 正文状态：未撰写；本文件目前是框架。

% 纲要内容开始
% 
% * 对角态射及分离判据
% * 仿射和开覆盖计算
% * 双原点的失效
% * 基变换和复合的稳定性
% 
% 纲要内容结束

\section{对角和分离性}
\label{g5:sec:1201}

待撰写

\subsection{对角态射及分离判据}
\label{g5:subsec:1201:01}
待撰写

\subsection{仿射和开覆盖计算}
\label{g5:subsec:1201:02}
待撰写

\subsection{双原点的失效}
\label{g5:subsec:1201:03}
待撰写

\subsection{基变换和复合的稳定性}
\label{g5:subsec:1201:04}
待撰写
% 纲要条目排为编号小节标题，并进入目录与书签；说明段仍保留于纲要注释中。

% !TeX root = ../../Book5.tex
% 纲要标题：### 12.2 固有态射和赋值判据
% 纲要位置：几何学计划纲要.md，第 5479 行。
% 正文状态：未撰写；本文件目前是框架。

% 纲要内容开始
% 
% * 有限型、分离与普遍闭
% * 赋值环及延拓图
% * 存在性和唯一性的不同部分
% * Noether 主版本的赋值判据证明
% * 有限态射的固有性
% 
% 纲要内容结束

\section{固有态射和赋值判据}
\label{g5:sec:1202}

待撰写

\subsection{有限型、分离与普遍闭}
\label{g5:subsec:1202:01}
待撰写

\subsection{赋值环及延拓图}
\label{g5:subsec:1202:02}
待撰写

\subsection{存在性和唯一性的不同部分}
\label{g5:subsec:1202:03}
待撰写

\subsection{Noether 主版本的赋值判据证明}
\label{g5:subsec:1202:04}
待撰写

\subsection{有限态射的固有性}
\label{g5:subsec:1202:05}
待撰写
% 纲要条目排为编号小节标题，并进入目录与书签；说明段仍保留于纲要注释中。

% !TeX root = ../../Book5.tex
% 纲要标题：### 12.3 拟有限态射的局部有限化
% 纲要位置：几何学计划纲要.md，第 5487 行。
% 正文状态：未撰写；本文件目前是框架。

% 纲要内容开始
% 
% * 拟有限纤维的代数控制
% * 整闭包和局部化的关键引理
% * 分离性如何粘合局部数据
% * 开浸入后接有限态射的分解
% 
% 纲要内容结束

\section{拟有限态射的局部有限化}
\label{g5:sec:1203}

待撰写

\subsection{拟有限纤维的代数控制}
\label{g5:subsec:1203:01}
待撰写

\subsection{整闭包和局部化的关键引理}
\label{g5:subsec:1203:02}
待撰写

\subsection{分离性如何粘合局部数据}
\label{g5:subsec:1203:03}
待撰写

\subsection{开浸入后接有限态射的分解}
\label{g5:subsec:1203:04}
待撰写
% 纲要条目排为编号小节标题，并进入目录与书签；说明段仍保留于纲要注释中。

% !TeX root = ../../Book5.tex
% 纲要标题：### 12.4 Zariski 主定理及应用
% 纲要位置：几何学计划纲要.md，第 5494 行。
% 正文状态：未撰写；本文件目前是框架。

% 纲要内容开始
% 
% * 所选拟紧分离拟有限版本
% * 正规目标上的双有理推论
% * 固有拟有限蕴含有限
% * 开浸入的判别及反例
% 
% 纲要内容结束

\section{Zariski 主定理及应用}
\label{g5:sec:1204}

待撰写

\subsection{所选拟紧分离拟有限版本}
\label{g5:subsec:1204:01}
待撰写

\subsection{正规目标上的双有理推论}
\label{g5:subsec:1204:02}
待撰写

\subsection{固有拟有限蕴含有限}
\label{g5:subsec:1204:03}
待撰写

\subsection{开浸入的判别及反例}
\label{g5:subsec:1204:04}
待撰写
% 纲要条目排为编号小节标题，并进入目录与书签；说明段仍保留于纲要注释中。

% !TeX root = ../../Book5.tex
% 纲要标题：### 12.5 后续的两个不同问题
% 纲要位置：几何学计划纲要.md，第 5207 行。
% 正文状态：未撰写；本文件目前是框架。

% 纲要内容开始
% 
% 射影修改与 Chow 引理在射影工具就绪后用于凝聚直像；连通纤维与 Stein 分解在形式函数之后处理，不在本章提前调用。
% 
% ---
% 
% 纲要内容结束

\section{后续的两个不同问题}
\label{g5:sec:1205}

待撰写
% 本节沿用纲要的连续叙述安排，不另造小节标题。


\begin{chaptersummary}
待撰写
\end{chaptersummary}

\BookExercises
待撰写

% !TeX root = ../../Book5.tex
% 纲要标题：## 第13章　拟凝聚层、凝聚层与向量丛
% 纲要位置：几何学计划纲要.md，第 5213 行。
% 正文状态：未撰写；本文件目前是框架。

% 纲要内容开始
% 
% > 本章目标与证明范围：在射影和上同调之前证明仿射模层对应、拉回及局部自由的性质。层的满射只在茎上检测，整体截面不因此自动满射。
% 
% 纲要内容结束

\chapter{拟凝聚层、凝聚层与向量丛}
\label{g5:ch:13}
\BookChapterNumber{g5:ch:13}

\begin{chapterabstract}
待撰写
\end{chapterabstract}

\begin{learninggoals}
待撰写
\end{learninggoals}

待撰写

% !TeX root = ../../Book5.tex
% 纲要标题：### 13.1 模到层的对应
% 纲要位置：几何学计划纲要.md，第 5217 行。
% 正文状态：未撰写；本文件目前是框架。

% 纲要内容开始
% 
% * 局部化构造模层
% * 仿射整体截面恢复模
% * 拟凝聚性及限制
% * 有限呈示和 Noether 凝聚性
% 
% 纲要内容结束

\section{模到层的对应}
\label{g5:sec:1301}

待撰写

\subsection{局部化构造模层}
\label{g5:subsec:1301:01}
待撰写

\subsection{仿射整体截面恢复模}
\label{g5:subsec:1301:02}
待撰写

\subsection{拟凝聚性及限制}
\label{g5:subsec:1301:03}
待撰写

\subsection{有限呈示和 Noether 凝聚性}
\label{g5:subsec:1301:04}
待撰写
% 纲要条目排为编号小节标题，并进入目录与书签；说明段仍保留于纲要注释中。

% !TeX root = ../../Book5.tex
% 纲要标题：### 13.2 层运算和函子
% 纲要位置：几何学计划纲要.md，第 5518 行。
% 正文状态：未撰写；本文件目前是框架。

% 纲要内容开始
% 
% * 核、余核、张量及内 Hom
% * 拉回、逆像和模结构
% * 拉回—直像伴随
% * 拟紧拟分离条件下的直像性质
% 
% 纲要内容结束

\section{层运算和函子}
\label{g5:sec:1302}

待撰写

\subsection{核、余核、张量及内 Hom}
\label{g5:subsec:1302:01}
待撰写

\subsection{拉回、逆像和模结构}
\label{g5:subsec:1302:02}
待撰写

\subsection{拉回—直像伴随}
\label{g5:subsec:1302:03}
待撰写

\subsection{拟紧拟分离条件下的直像性质}
\label{g5:subsec:1302:04}
待撰写
% 纲要条目排为编号小节标题，并进入目录与书签；说明段仍保留于纲要注释中。

% !TeX root = ../../Book5.tex
% 纲要标题：### 13.3 局部自由和投射
% 纲要位置：几何学计划纲要.md，第 5231 行。
% 正文状态：未撰写；本文件目前是框架。

% 纲要内容开始
% 
% * 有限局部自由层与向量丛
% * 秩、行列式和对偶
% * 仿射投射模的几何对应
% * 短正合列与整体不分裂
% 
% 纲要内容结束

\section{局部自由和投射}
\label{g5:sec:1303}

待撰写

\subsection{有限局部自由层与向量丛}
\label{g5:subsec:1303:01}
待撰写

\subsection{秩、行列式和对偶}
\label{g5:subsec:1303:02}
待撰写

\subsection{仿射投射模的几何对应}
\label{g5:subsec:1303:03}
待撰写

\subsection{短正合列与整体不分裂}
\label{g5:subsec:1303:04}
待撰写
% 纲要条目排为编号小节标题，并进入目录与书签；说明段仍保留于纲要注释中。

% !TeX root = ../../Book5.tex
% 纲要标题：### 13.4 支集和局部模型
% 纲要位置：几何学计划纲要.md，第 5532 行。
% 正文状态：未撰写；本文件目前是框架。

% 纲要内容开始
% 
% * 关联点与嵌入分支
% * 挠、无挠和反射
% * 局部自由点集及 Fitting 条件
% * 理想层、结构层和法层
% 
% 纲要内容结束

\section{支集和局部模型}
\label{g5:sec:1304}

待撰写

\subsection{关联点与嵌入分支}
\label{g5:subsec:1304:01}
待撰写

\subsection{挠、无挠和反射}
\label{g5:subsec:1304:02}
待撰写

\subsection{局部自由点集及 Fitting 条件}
\label{g5:subsec:1304:03}
待撰写

\subsection{理想层、结构层和法层}
\label{g5:subsec:1304:04}
待撰写
% 纲要条目排为编号小节标题，并进入目录与书签；说明段仍保留于纲要注释中。

% !TeX root = ../../Book5.tex
% 纲要标题：### 13.5 相对 Spec 与向量丛的总空间
% 纲要位置：几何学计划纲要.md，第 5245 行。
% 正文状态：未撰写；本文件目前是框架。

% 纲要内容开始
% 
% * 拟凝聚代数与仿射谱的局部粘合
% * 相对 Spec 的泛性质和仿射态射的反变对应
% 
%   完整构造和核验放在本节，不等待形式函数或固有直像理论。[R9]
% 
% * 结构代数的恢复与任意基变换相容
% * 对称代数、零截面及 $\operatorname{Spec}_S\operatorname{Sym}(E^\vee)$
% 
%   当 $E$ 表示有限局部自由截面层时，以 $\operatorname{Spec}_S\operatorname{Sym}(E^\vee)$ 表示其总空间；与第14章使用的一维商型射影丛约定分别说明。
% 
% * 有限代数与有限态射的局部检测
% * 为后续 Stein 分解预留的凝聚代数应用
% 
% ---
% 
% 纲要内容结束

\section{相对 Spec 与向量丛的总空间}
\label{g5:sec:1305}

待撰写

\subsection{拟凝聚代数与仿射谱的局部粘合}
\label{g5:subsec:1305:01}
待撰写

\subsection{相对 Spec 的泛性质和仿射态射的反变对应}
\label{g5:subsec:1305:02}
待撰写

\subsection{结构代数的恢复与任意基变换相容}
\label{g5:subsec:1305:03}
待撰写

\subsection{\texorpdfstring{对称代数、零截面及 \(\operatorname{Spec}_S\operatorname{Sym}(E^\vee)\)}{对称代数、零截面及 Spec\_S Sym(E∨)}}
\label{g5:subsec:1305:04}
待撰写

\subsection{有限代数与有限态射的局部检测}
\label{g5:subsec:1305:05}
待撰写

\subsection{为后续 Stein 分解预留的凝聚代数应用}
\label{g5:subsec:1305:06}
待撰写
% 纲要条目排为编号小节标题，并进入目录与书签；说明段仍保留于纲要注释中。


\begin{chaptersummary}
待撰写
\end{chaptersummary}

\BookExercises
待撰写

% !TeX root = ../../Book5.tex
% 纲要标题：## 第14章　Proj、射影态射与吹起
% 纲要位置：几何学计划纲要.md，第 5556 行。
% 正文状态：未撰写；本文件目前是框架。

% 纲要内容开始
% 
% > 本章目标与证明范围：构造在先，泛性质和稳定性在后。相对 Proj、Rees 代数和线丛已具备所需层语言；不提前使用曲面交数证明吹起的基础性质。
% 
% 纲要内容结束

\chapter{Proj、射影态射与吹起}
\label{g5:ch:14}
\BookChapterNumber{g5:ch:14}

\begin{chapterabstract}
待撰写
\end{chapterabstract}

\begin{learninggoals}
待撰写
\end{learninggoals}

待撰写

% !TeX root = ../../Book5.tex
% 纲要标题：### 14.1 分次环与 Proj
% 纲要位置：几何学计划纲要.md，第 5266 行。
% 正文状态：未撰写；本文件目前是框架。

% 纲要内容开始
% 
% * 齐次素理想及标准开集
% * 零次局部化的结构层
% * 分次模的层化和扭转
% * Veronese 子环及同一 Proj
% 
% 纲要内容结束

\section{分次环与 Proj}
\label{g5:sec:1401}

待撰写

\subsection{齐次素理想及标准开集}
\label{g5:subsec:1401:01}
待撰写

\subsection{零次局部化的结构层}
\label{g5:subsec:1401:02}
待撰写

\subsection{分次模的层化和扭转}
\label{g5:subsec:1401:03}
待撰写

\subsection{Veronese 子环及同一 Proj}
\label{g5:subsec:1401:04}
待撰写
% 纲要条目排为编号小节标题，并进入目录与书签；说明段仍保留于纲要注释中。

% !TeX root = ../../Book5.tex
% 纲要标题：### 14.2 相对射影构造
% 纲要位置：几何学计划纲要.md，第 5273 行。
% 正文状态：未撰写；本文件目前是框架。

% 纲要内容开始
% 
% * 拟凝聚分次代数的相对 Proj
% * 射影丛和线丛的商约定
% * 射影态射及基变换
% * 射影空间的固有性
% 
% 纲要内容结束

\section{相对射影构造}
\label{g5:sec:1402}

待撰写

\subsection{拟凝聚分次代数的相对 Proj}
\label{g5:subsec:1402:01}
待撰写

\subsection{射影丛和线丛的商约定}
\label{g5:subsec:1402:02}
待撰写

\subsection{射影态射及基变换}
\label{g5:subsec:1402:03}
待撰写

\subsection{射影空间的固有性}
\label{g5:subsec:1402:04}
待撰写
% 纲要条目排为编号小节标题，并进入目录与书签；说明段仍保留于纲要注释中。

% !TeX root = ../../Book5.tex
% 纲要标题：### 14.3 Rees 代数与吹起
% 纲要位置：几何学计划纲要.md，第 5574 行。
% 正文状态：未撰写；本文件目前是框架。

% 纲要内容开始
% 
% * 理想幂、Rees 环及例外除子
% * 标准仿射图的显式计算
% * 严格变换与全变换
% * 光滑点及非主理想的模型
% 
% 纲要内容结束

\section{Rees 代数与吹起}
\label{g5:sec:1403}

待撰写

\subsection{理想幂、Rees 环及例外除子}
\label{g5:subsec:1403:01}
待撰写

\subsection{标准仿射图的显式计算}
\label{g5:subsec:1403:02}
待撰写

\subsection{严格变换与全变换}
\label{g5:subsec:1403:03}
待撰写

\subsection{光滑点及非主理想的模型}
\label{g5:subsec:1403:04}
待撰写
% 纲要条目排为编号小节标题，并进入目录与书签；说明段仍保留于纲要注释中。

% !TeX root = ../../Book5.tex
% 纲要标题：### 14.4 泛性质和相容性
% 纲要位置：几何学计划纲要.md，第 5581 行。
% 正文状态：未撰写；本文件目前是框架。

% 纲要内容开始
% 
% * 理想可逆化的泛性质
% * 基变换与 Rees 代数的条件
% * 平坦基变换的相容
% * 任意非平坦基变换的限制
% 
% ---
% 
% 纲要内容结束

\section{泛性质和相容性}
\label{g5:sec:1404}

待撰写

\subsection{理想可逆化的泛性质}
\label{g5:subsec:1404:01}
待撰写

\subsection{基变换与 Rees 代数的条件}
\label{g5:subsec:1404:02}
待撰写

\subsection{平坦基变换的相容}
\label{g5:subsec:1404:03}
待撰写

\subsection{任意非平坦基变换的限制}
\label{g5:subsec:1404:04}
待撰写
% 纲要条目排为编号小节标题，并进入目录与书签；说明段仍保留于纲要注释中。


\begin{chaptersummary}
待撰写
\end{chaptersummary}

\BookExercises
待撰写

% !TeX root = ../../Book5.tex
% 纲要标题：## 第15章　光滑、非分歧、平展与 Hensel 局部性
% 纲要位置：几何学计划纲要.md，第 5590 行。
% 正文状态：未撰写；本文件目前是框架。

% 纲要内容开始
% 
% > 本章目标与证明范围：把平方零提升、微分模和几何纤维三种判据建立并比较。尖点正规化处必须分别检验非分歧和平坦，两者均失败的局部计算保留。
% 
% 纲要内容结束

\chapter{光滑、非分歧、平展与 Hensel 局部性}
\label{g5:ch:15}
\BookChapterNumber{g5:ch:15}

\begin{chapterabstract}
待撰写
\end{chapterabstract}

\begin{learninggoals}
待撰写
\end{learninggoals}

待撰写

% !TeX root = ../../Book5.tex
% 纲要标题：### 15.1 形式提升性质
% 纲要位置：几何学计划纲要.md，第 5594 行。
% 正文状态：未撰写；本文件目前是框架。

% 纲要内容开始
% 
% * 平方零扩张及提升图
% * 形式光滑、非分歧和平展
% * 有限呈示把形式条件转为几何条件
% * 微分模检测提升差异
% 
% 纲要内容结束

\section{形式提升性质}
\label{g5:sec:1501}

待撰写

\subsection{平方零扩张及提升图}
\label{g5:subsec:1501:01}
待撰写

\subsection{形式光滑、非分歧和平展}
\label{g5:subsec:1501:02}
待撰写

\subsection{有限呈示把形式条件转为几何条件}
\label{g5:subsec:1501:03}
待撰写

\subsection{微分模检测提升差异}
\label{g5:subsec:1501:04}
待撰写
% 纲要条目排为编号小节标题，并进入目录与书签；说明段仍保留于纲要注释中。

% !TeX root = ../../Book5.tex
% 纲要标题：### 15.2 光滑态射的结构
% 纲要位置：几何学计划纲要.md，第 5601 行。
% 正文状态：未撰写；本文件目前是框架。

% 纲要内容开始
% 
% * 标准光滑代数
% * Jacobian 的相对判据
% * 平坦性和几何正则纤维
% * 光滑点集的开放性
% * 特征零的一般光滑性及正特征限制
% 
% 纲要内容结束

\section{光滑态射的结构}
\label{g5:sec:1502}

待撰写

\subsection{标准光滑代数}
\label{g5:subsec:1502:01}
待撰写

\subsection{Jacobian 的相对判据}
\label{g5:subsec:1502:02}
待撰写

\subsection{平坦性和几何正则纤维}
\label{g5:subsec:1502:03}
待撰写

\subsection{光滑点集的开放性}
\label{g5:subsec:1502:04}
待撰写

\subsection{特征零的一般光滑性及正特征限制}
\label{g5:subsec:1502:05}
待撰写
% 纲要条目排为编号小节标题，并进入目录与书签；说明段仍保留于纲要注释中。

% !TeX root = ../../Book5.tex
% 纲要标题：### 15.3 非分歧和平展
% 纲要位置：几何学计划纲要.md，第 5609 行。
% 正文状态：未撰写；本文件目前是框架。

% 纲要内容开始
% 
% * Ω 消失与非分歧的有限性条件
% * 平坦加非分歧的完整判据
% * 标准平展局部模型
% * 对角判据及有限平展覆盖
% 
% 纲要内容结束

\section{非分歧和平展}
\label{g5:sec:1503}

待撰写

\subsection{\texorpdfstring{\(\Omega\) 消失与非分歧的有限性条件}{\GeometryPdfOmega{} 消失与非分歧的有限性条件}}
\label{g5:subsec:1503:01}
待撰写

\subsection{平坦加非分歧的完整判据}
\label{g5:subsec:1503:02}
待撰写

\subsection{标准平展局部模型}
\label{g5:subsec:1503:03}
待撰写

\subsection{对角判据及有限平展覆盖}
\label{g5:subsec:1503:04}
待撰写
% 纲要条目排为编号小节标题，并进入目录与书签；说明段仍保留于纲要注释中。

% !TeX root = ../../Book5.tex
% 纲要标题：### 15.4 Hensel 引理及邻域
% 纲要位置：几何学计划纲要.md，第 5322 行。
% 正文状态：未撰写；本文件目前是框架。

% 纲要内容开始
% 
% * 单根与因式分解提升
% * Hensel 局部环的等价刻画
% * Hensel 化和严格 Hensel 化
% * 完备化与 Hensel 邻域的不同泛性质
% 
% 纲要内容结束

\section{Hensel 引理及邻域}
\label{g5:sec:1504}

待撰写

\subsection{单根与因式分解提升}
\label{g5:subsec:1504:01}
待撰写

\subsection{Hensel 局部环的等价刻画}
\label{g5:subsec:1504:02}
待撰写

\subsection{Hensel 化和严格 Hensel 化}
\label{g5:subsec:1504:03}
待撰写

\subsection{完备化与 Hensel 邻域的不同泛性质}
\label{g5:subsec:1504:04}
待撰写
% 纲要条目排为编号小节标题，并进入目录与书签；说明段仍保留于纲要注释中。

% !TeX root = ../../Book5.tex
% 纲要标题：### 15.5 必须算清的反例
% 纲要位置：几何学计划纲要.md，第 5329 行。
% 正文状态：未撰写；本文件目前是框架。

% 纲要内容开始
% 
% * 尖点的 Ω 和非约化纤维
% * 有限双有理泛秩与局部非平坦
% * 结点正规化和分支
% * 纯不可分映射与 Frobenius
% 
% ---
% 
% 纲要内容结束

\section{必须算清的反例}
\label{g5:sec:1505}

待撰写

\subsection{\texorpdfstring{尖点的 \(\Omega\) 和非约化纤维}{尖点的 \GeometryPdfOmega{} 和非约化纤维}}
\label{g5:subsec:1505:01}
待撰写

\subsection{有限双有理泛秩与局部非平坦}
\label{g5:subsec:1505:02}
待撰写

\subsection{结点正规化和分支}
\label{g5:subsec:1505:03}
待撰写

\subsection{纯不可分映射与 Frobenius}
\label{g5:subsec:1505:04}
待撰写
% 纲要条目排为编号小节标题，并进入目录与书签；说明段仍保留于纲要注释中。


\begin{chaptersummary}
待撰写
\end{chaptersummary}

\BookExercises
待撰写

% !TeX root = ../../Book5.tex
% 纲要标题：## 第16章*　忠实平坦下降、覆盖与代数空间
% 纲要位置：几何学计划纲要.md，第 5632 行。
% 正文状态：未撰写；本文件目前是框架。

% 纲要内容开始
% 
% > 本章目标与证明范围：模空间路线需要本章的下降、层化与代数空间最低语言，不必先学完整平展上同调。证明有限呈示模、局部自由层与仿射对象的有效下降，再以平展等价关系引出代数空间；一般概形、代数空间和叠的存在性范围分别说明。
% 
% 纲要内容结束

\OptionalChapter{忠实平坦下降、覆盖与代数空间}{g5:ch:16}
\BookChapterNumber{g5:ch:16}

\begin{chapterabstract}
待撰写
\end{chapterabstract}

\begin{learninggoals}
待撰写
\end{learninggoals}

待撰写

% !TeX root = ../../Book5.tex
% 纲要标题：### 16.1 下降数据
% 纲要位置：几何学计划纲要.md，第 5636 行。
% 正文状态：未撰写；本文件目前是框架。

% 纲要内容开始
% 
% * 双重基变换上的同构
% * 三重基变换的余循环
% * 态射下降和对象下降
% * 数据与选择无关性
% 
% 纲要内容结束

\section{下降数据}
\label{g5:sec:1601}

待撰写

\subsection{双重基变换上的同构}
\label{g5:subsec:1601:01}
待撰写

\subsection{三重基变换的余循环}
\label{g5:subsec:1601:02}
待撰写

\subsection{态射下降和对象下降}
\label{g5:subsec:1601:03}
待撰写

\subsection{数据与选择无关性}
\label{g5:subsec:1601:04}
待撰写
% 纲要条目排为编号小节标题，并进入目录与书签；说明段仍保留于纲要注释中。

% !TeX root = ../../Book5.tex
% 纲要标题：### 16.2 模的忠实平坦下降
% 纲要位置：几何学计划纲要.md，第 5349 行。
% 正文状态：未撰写；本文件目前是框架。

% 纲要内容开始
% 
% * Amitsur 复形的最低正合性
% * 等化子构造下降模
% * 基变换恢复及有限呈示
% * 局部自由、秩和线丛的下降
% 
% 纲要内容结束

\section{模的忠实平坦下降}
\label{g5:sec:1602}

待撰写

\subsection{Amitsur 复形的最低正合性}
\label{g5:subsec:1602:01}
待撰写

\subsection{等化子构造下降模}
\label{g5:subsec:1602:02}
待撰写

\subsection{基变换恢复及有限呈示}
\label{g5:subsec:1602:03}
待撰写

\subsection{局部自由、秩和线丛的下降}
\label{g5:subsec:1602:04}
待撰写
% 纲要条目排为编号小节标题，并进入目录与书签；说明段仍保留于纲要注释中。

% !TeX root = ../../Book5.tex
% 纲要标题：### 16.3 景、覆盖与层化
% 纲要位置：几何学计划纲要.md，第 5650 行。
% 正文状态：未撰写；本文件目前是框架。

% 纲要内容开始
% 
% * 覆盖族和预拓扑
% * Zariski、平展、fppf、fpqc 的条件
% * 表示函子的层性质
% * 层化改变朴素商函子的方式
% 
% 纲要内容结束

\section{景、覆盖与层化}
\label{g5:sec:1603}

待撰写

\subsection{覆盖族和预拓扑}
\label{g5:subsec:1603:01}
待撰写

\subsection{Zariski、平展、fppf、fpqc 的条件}
\label{g5:subsec:1603:02}
待撰写

\subsection{表示函子的层性质}
\label{g5:subsec:1603:03}
待撰写

\subsection{层化改变朴素商函子的方式}
\label{g5:subsec:1603:04}
待撰写
% 纲要条目排为编号小节标题，并进入目录与书签；说明段仍保留于纲要注释中。

% !TeX root = ../../Book5.tex
% 纲要标题：### 16.4 几何对象的有效性
% 纲要位置：几何学计划纲要.md，第 5657 行。
% 正文状态：未撰写；本文件目前是框架。

% 纲要内容开始
% 
% * 仿射概形的下降
% * 粘合开子对象及态射
% * 一般概形有效下降的限制
% * 向模空间传递的最低输入
% 
% 纲要内容结束

\section{几何对象的有效性}
\label{g5:sec:1604}

待撰写

\subsection{仿射概形的下降}
\label{g5:subsec:1604:01}
待撰写

\subsection{粘合开子对象及态射}
\label{g5:subsec:1604:02}
待撰写

\subsection{一般概形有效下降的限制}
\label{g5:subsec:1604:03}
待撰写

\subsection{向模空间传递的最低输入}
\label{g5:subsec:1604:04}
待撰写
% 纲要条目排为编号小节标题，并进入目录与书签；说明段仍保留于纲要注释中。

% !TeX root = ../../Book5.tex
% 纲要标题：### 16.5 平展等价关系与代数空间
% 纲要位置：几何学计划纲要.md，第 5370 行。
% 正文状态：未撰写；本文件目前是框架。

% 纲要内容开始
% 
% 采用 fppf 层、概形可表对角和概形给出的平展满射图册定义；平展等价关系的层商是本节的构造主线。一般模空间存在定理不由定义自动得到。[R10]
% 
% * 平展等价关系 $R\rightrightarrows U$ 与朴素商的层化
% * 代数空间的层、可表对角与平展满射图册
% * $U\times_XU$ 与图册等价关系的恢复
% * 态射、纤维积及基变换的最低工具
% * 由图册检测局部性质的条件
% * 概形商、代数空间商与非概形模型的区别
% 
% 纲要内容结束

\section{平展等价关系与代数空间}
\label{g5:sec:1605}

待撰写

\subsection{\texorpdfstring{平展等价关系 \(R\rightrightarrows U\) 与朴素商的层化}{平展等价关系 R⇉U 与朴素商的层化}}
\label{g5:subsec:1605:01}
待撰写

\subsection{代数空间的层、可表对角与平展满射图册}
\label{g5:subsec:1605:02}
待撰写

\subsection{\texorpdfstring{\(U\times_XU\) 与图册等价关系的恢复}{U×\_X U 与图册等价关系的恢复}}
\label{g5:subsec:1605:03}
待撰写

\subsection{态射、纤维积及基变换的最低工具}
\label{g5:subsec:1605:04}
待撰写

\subsection{由图册检测局部性质的条件}
\label{g5:subsec:1605:05}
待撰写

\subsection{概形商、代数空间商与非概形模型的区别}
\label{g5:subsec:1605:06}
待撰写
% 纲要条目排为编号小节标题，并进入目录与书签；说明段仍保留于纲要注释中。

% !TeX root = ../../Book5.tex
% 纲要标题：### 16.6 从代数空间到模问题的对象层次
% 纲要位置：几何学计划纲要.md，第 5675 行。
% 正文状态：未撰写；本文件目前是框架。

% 纲要内容开始
% 
% 本节为第33—34章建立共同语言；叠的实际定义和商叠模型在 V.34.5 展开，不提前要求导出叠或一般六函子理论。
% 
% * 集合值模函子与族的同构类
% * 概形可表和代数空间可表的不同结论
% * 保留自同构的群胚值对象及叠的动机
% * 通用族、刚化及粗模空间的区别
% * Hilbert、Picard 与商模问题的依赖清单
% 
% ---
% 
% 纲要内容结束

\section{从代数空间到模问题的对象层次}
\label{g5:sec:1606}

待撰写

\subsection{集合值模函子与族的同构类}
\label{g5:subsec:1606:01}
待撰写

\subsection{概形可表和代数空间可表的不同结论}
\label{g5:subsec:1606:02}
待撰写

\subsection{保留自同构的群胚值对象及叠的动机}
\label{g5:subsec:1606:03}
待撰写

\subsection{通用族、刚化及粗模空间的区别}
\label{g5:subsec:1606:04}
待撰写

\subsection{Hilbert、Picard 与商模问题的依赖清单}
\label{g5:subsec:1606:05}
待撰写
% 纲要条目排为编号小节标题，并进入目录与书签；说明段仍保留于纲要注释中。


\begin{chaptersummary}
待撰写
\end{chaptersummary}

\BookExercises
待撰写


% !TeX root = ../Book5.tex
% 纲要标题：# 第三编　除子、上同调与对偶
% 纲要位置：几何学计划纲要.md，第 5393 行。
% 正文状态：未撰写；本文件目前是框架。

% 纲要内容开始
% 
% 纲要内容结束

\BookPart{除子、上同调与对偶}{g5:part:03}
\BookPartGuide
待撰写

% !TeX root = ../../Book5.tex
% 纲要标题：## 第17章　除子、线性系与充足性
% 纲要位置：几何学计划纲要.md，第 5395 行。
% 正文状态：未撰写；本文件目前是框架。

% 纲要内容开始
% 
% > 本章目标与证明范围：从局部方程到可逆层，再到截面定义的映射。反射层、类群和数值正性的不同层次保留；Bertini 在实际构造一般截面时有独立证明。
% 
% 纲要内容结束

\chapter{除子、线性系与充足性}
\label{g5:ch:17}
\BookChapterNumber{g5:ch:17}

\begin{chapterabstract}
待撰写
\end{chapterabstract}

\begin{learninggoals}
待撰写
\end{learninggoals}

待撰写

% !TeX root = ../../Book5.tex
% 纲要标题：### 17.1 Cartier 与 Weil 除子
% 纲要位置：几何学计划纲要.md，第 5693 行。
% 正文状态：未撰写；本文件目前是框架。

% 纲要内容开始
% 
% * 总商层和局部有理方程
% * 正规簇的余维一赋值
% * Cartier 到 Weil 的比较
% * 主除子与线性等价
% * 局部阶乘假设的作用
% 
% 纲要内容结束

\section{Cartier 与 Weil 除子}
\label{g5:sec:1701}

待撰写

\subsection{总商层和局部有理方程}
\label{g5:subsec:1701:01}
待撰写

\subsection{正规簇的余维一赋值}
\label{g5:subsec:1701:02}
待撰写

\subsection{Cartier 到 Weil 的比较}
\label{g5:subsec:1701:03}
待撰写

\subsection{主除子与线性等价}
\label{g5:subsec:1701:04}
待撰写

\subsection{局部阶乘假设的作用}
\label{g5:subsec:1701:05}
待撰写
% 纲要条目排为编号小节标题，并进入目录与书签；说明段仍保留于纲要注释中。

% !TeX root = ../../Book5.tex
% 纲要标题：### 17.2 可逆层与 Picard 群
% 纲要位置：几何学计划纲要.md，第 5701 行。
% 正文状态：未撰写；本文件目前是框架。

% 纲要内容开始
% 
% * 除子层的构造
% * 截面、零点和有效除子
% * 可逆层与 Cartier 类
% * 秩一反射层及除子类群
% 
% 纲要内容结束

\section{可逆层与 Picard 群}
\label{g5:sec:1702}

待撰写

\subsection{除子层的构造}
\label{g5:subsec:1702:01}
待撰写

\subsection{截面、零点和有效除子}
\label{g5:subsec:1702:02}
待撰写

\subsection{可逆层与 Cartier 类}
\label{g5:subsec:1702:03}
待撰写

\subsection{秩一反射层及除子类群}
\label{g5:subsec:1702:04}
待撰写
% 纲要条目排为编号小节标题，并进入目录与书签；说明段仍保留于纲要注释中。

% !TeX root = ../../Book5.tex
% 纲要标题：### 17.3 线性系与射影态射
% 纲要位置：几何学计划纲要.md，第 5414 行。
% 正文状态：未撰写；本文件目前是框架。

% 纲要内容开始
% 
% * 基点和固定部分
% * 无基点截面给出的映射
% * 很充足与充足的区别
% * 拉回 O(1) 和分离喷射
% * 典型射影嵌入的计算
% 
% 纲要内容结束

\section{线性系与射影态射}
\label{g5:sec:1703}

待撰写

\subsection{基点和固定部分}
\label{g5:subsec:1703:01}
待撰写

\subsection{无基点截面给出的映射}
\label{g5:subsec:1703:02}
待撰写

\subsection{很充足与充足的区别}
\label{g5:subsec:1703:03}
待撰写

\subsection{拉回 O(1) 和分离喷射}
\label{g5:subsec:1703:04}
待撰写

\subsection{典型射影嵌入的计算}
\label{g5:subsec:1703:05}
待撰写
% 纲要条目排为编号小节标题，并进入目录与书签；说明段仍保留于纲要注释中。

% !TeX root = ../../Book5.tex
% 纲要标题：### 17.4 Bertini 与一般位置
% 纲要位置：几何学计划纲要.md，第 5716 行。
% 正文状态：未撰写；本文件目前是框架。

% 纲要内容开始
% 
% * 关联簇和坏截面参数
% * 特征零光滑性的证明
% * 维数至少2的所选不可约性版本
% * 有基点线性系的限制
% * 正特征的额外问题
% 
% ---
% 
% 纲要内容结束

\section{Bertini 与一般位置}
\label{g5:sec:1704}

待撰写

\subsection{关联簇和坏截面参数}
\label{g5:subsec:1704:01}
待撰写

\subsection{特征零光滑性的证明}
\label{g5:subsec:1704:02}
待撰写

\subsection{维数至少2的所选不可约性版本}
\label{g5:subsec:1704:03}
待撰写

\subsection{有基点线性系的限制}
\label{g5:subsec:1704:04}
待撰写

\subsection{正特征的额外问题}
\label{g5:subsec:1704:05}
待撰写
% 纲要条目排为编号小节标题，并进入目录与书签；说明段仍保留于纲要注释中。


\begin{chaptersummary}
待撰写
\end{chaptersummary}

\BookExercises
待撰写

% !TeX root = ../../Book5.tex
% 纲要标题：## 第18章　层上同调与计算模型
% 纲要位置：几何学计划纲要.md，第 5432 行。
% 正文状态：未撰写；本文件目前是框架。

% 纲要内容开始
% 
% > 本章目标与证明范围：本卷独立建立所用导出工具。Čech 计算的成立条件、仿射无上同调和高阶直像各有证明位置；不将光滑单位分解机制移植到 Zariski 凝聚层。
% 
% 纲要内容结束

\chapter{层上同调与计算模型}
\label{g5:ch:18}
\BookChapterNumber{g5:ch:18}

\begin{chapterabstract}
待撰写
\end{chapterabstract}

\begin{learninggoals}
待撰写
\end{learninggoals}

待撰写

% !TeX root = ../../Book5.tex
% 纲要标题：### 18.1 截面函子与消解
% 纲要位置：几何学计划纲要.md，第 5730 行。
% 正文状态：未撰写；本文件目前是框架。

% 纲要内容开始
% 
% * 层的短正合列
% * 足够多内射或明确无上同调消解
% * 导出函子的独立性
% * 连接映射及自然长正合列
% 
% 纲要内容结束

\section{截面函子与消解}
\label{g5:sec:1801}

待撰写

\subsection{层的短正合列}
\label{g5:subsec:1801:01}
待撰写

\subsection{足够多内射或明确无上同调消解}
\label{g5:subsec:1801:02}
待撰写

\subsection{导出函子的独立性}
\label{g5:subsec:1801:03}
待撰写

\subsection{连接映射及自然长正合列}
\label{g5:subsec:1801:04}
待撰写
% 纲要条目排为编号小节标题，并进入目录与书签；说明段仍保留于纲要注释中。

% !TeX root = ../../Book5.tex
% 纲要标题：### 18.2 Čech 与导出比较
% 纲要位置：几何学计划纲要.md，第 5443 行。
% 正文状态：未撰写；本文件目前是框架。

% 纲要内容开始
% 
% * 覆盖复形和细化
% * 双复形的最低谱序列工具
% * 无上同调覆盖的比较
% * 非分离情形交集的注意事项
% 
% 纲要内容结束

\section{Čech 与导出比较}
\label{g5:sec:1802}

待撰写

\subsection{覆盖复形和细化}
\label{g5:subsec:1802:01}
待撰写

\subsection{双复形的最低谱序列工具}
\label{g5:subsec:1802:02}
待撰写

\subsection{无上同调覆盖的比较}
\label{g5:subsec:1802:03}
待撰写

\subsection{非分离情形交集的注意事项}
\label{g5:subsec:1802:04}
待撰写
% 纲要条目排为编号小节标题，并进入目录与书签；说明段仍保留于纲要注释中。

% !TeX root = ../../Book5.tex
% 纲要标题：### 18.3 仿射无上同调
% 纲要位置：几何学计划纲要.md，第 5744 行。
% 正文状态：未撰写；本文件目前是框架。

% 纲要内容开始
% 
% * 主开覆盖和局部化复形
% * 清除分母的正合论证
% * 拟凝聚层的高阶消失
% * 仿射性及全局截面的范围
% 
% 纲要内容结束

\section{仿射无上同调}
\label{g5:sec:1803}

待撰写

\subsection{主开覆盖和局部化复形}
\label{g5:subsec:1803:01}
待撰写

\subsection{清除分母的正合论证}
\label{g5:subsec:1803:02}
待撰写

\subsection{拟凝聚层的高阶消失}
\label{g5:subsec:1803:03}
待撰写

\subsection{仿射性及全局截面的范围}
\label{g5:subsec:1803:04}
待撰写
% 纲要条目排为编号小节标题，并进入目录与书签；说明段仍保留于纲要注释中。

% !TeX root = ../../Book5.tex
% 纲要标题：### 18.4 局部到整体与高阶直像
% 纲要位置：几何学计划纲要.md，第 5751 行。
% 正文状态：未撰写；本文件目前是框架。

% 纲要内容开始
% 
% * Rᵠf_* 的构造
% * Leray 谱序列及低阶序列
% * 局部到整体 Ext
% * 支集上同调的最低入口
% 
% 纲要内容结束

\section{局部到整体与高阶直像}
\label{g5:sec:1804}

待撰写

\subsection{\texorpdfstring{\(R^q f_*\) 的构造}{Rᵠf\_* 的构造}}
\label{g5:subsec:1804:01}
待撰写

\subsection{Leray 谱序列及低阶序列}
\label{g5:subsec:1804:02}
待撰写

\subsection{局部到整体 Ext}
\label{g5:subsec:1804:03}
待撰写

\subsection{支集上同调的最低入口}
\label{g5:subsec:1804:04}
待撰写
% 纲要条目排为编号小节标题，并进入目录与书签；说明段仍保留于纲要注释中。

% !TeX root = ../../Book5.tex
% 纲要标题：### 18.5 可复核的初步计算
% 纲要位置：几何学计划纲要.md，第 5758 行。
% 正文状态：未撰写；本文件目前是框架。

% 纲要内容开始
% 
% * 射影线的 O(d)
% * 闭子概形的理想层序列
% * 非分裂扩张与 H¹
% * 同一几何族的 Čech 代表
% 
% ---
% 
% 纲要内容结束

\section{可复核的初步计算}
\label{g5:sec:1805}

待撰写

\subsection{射影线的 O(d)}
\label{g5:subsec:1805:01}
待撰写

\subsection{闭子概形的理想层序列}
\label{g5:subsec:1805:02}
待撰写

\subsection{\texorpdfstring{非分裂扩张与 \(H^1\)}{非分裂扩张与 H¹}}
\label{g5:subsec:1805:03}
待撰写

\subsection{同一几何族的 Čech 代表}
\label{g5:subsec:1805:04}
待撰写
% 纲要条目排为编号小节标题，并进入目录与书签；说明段仍保留于纲要注释中。


\begin{chaptersummary}
待撰写
\end{chaptersummary}

\BookExercises
待撰写

% !TeX root = ../../Book5.tex
% 纲要标题：## 第19章　射影上同调、Serre 定理与有限性
% 纲要位置：几何学计划纲要.md，第 5473 行。
% 正文状态：未撰写；本文件目前是框架。

% 纲要内容开始
% 
% > 本章目标与证明范围：完整证明射影空间扭转层的上同调、整体生成和 Serre 消失，再处理一般固有凝聚直像所需的额外步骤。Hilbert 多项式不是提前无条件使用的平坦判据。
% 
% 纲要内容结束

\chapter{射影上同调、Serre 定理与有限性}
\label{g5:ch:19}
\BookChapterNumber{g5:ch:19}

\begin{chapterabstract}
待撰写
\end{chapterabstract}

\begin{learninggoals}
待撰写
\end{learninggoals}

待撰写

% !TeX root = ../../Book5.tex
% 纲要标题：### 19.1 射影空间的 Čech 计算
% 纲要位置：几何学计划纲要.md，第 5477 行。
% 正文状态：未撰写；本文件目前是框架。

% 纲要内容开始
% 
% * 标准仿射覆盖
% * 分次 Laurent 单项式的复形
% * 各次数上同调的确定
% * 顶次配对与对偶的准备
% 
% 纲要内容结束

\section{射影空间的 Čech 计算}
\label{g5:sec:1901}

待撰写

\subsection{标准仿射覆盖}
\label{g5:subsec:1901:01}
待撰写

\subsection{分次 Laurent 单项式的复形}
\label{g5:subsec:1901:02}
待撰写

\subsection{各次数上同调的确定}
\label{g5:subsec:1901:03}
待撰写

\subsection{顶次配对与对偶的准备}
\label{g5:subsec:1901:04}
待撰写
% 纲要条目排为编号小节标题，并进入目录与书签；说明段仍保留于纲要注释中。

% !TeX root = ../../Book5.tex
% 纲要标题：### 19.2 整体生成与消失
% 纲要位置：几何学计划纲要.md，第 5484 行。
% 正文状态：未撰写；本文件目前是框架。

% 纲要内容开始
% 
% * 高次扭转的有限截面生成
% * 由扭转自由层覆盖凝聚层
% * Serre 消失的维数或消解证明
% * 有限生成及 Noether 条件
% 
% 纲要内容结束

\section{整体生成与消失}
\label{g5:sec:1902}

待撰写

\subsection{高次扭转的有限截面生成}
\label{g5:subsec:1902:01}
待撰写

\subsection{由扭转自由层覆盖凝聚层}
\label{g5:subsec:1902:02}
待撰写

\subsection{Serre 消失的维数或消解证明}
\label{g5:subsec:1902:03}
待撰写

\subsection{有限生成及 Noether 条件}
\label{g5:subsec:1902:04}
待撰写
% 纲要条目排为编号小节标题，并进入目录与书签；说明段仍保留于纲要注释中。

% !TeX root = ../../Book5.tex
% 纲要标题：### 19.3 Hilbert 函数与 Euler 多项式
% 纲要位置：几何学计划纲要.md，第 5785 行。
% 正文状态：未撰写；本文件目前是框架。

% 纲要内容开始
% 
% * 分次增长和有理级数
% * Hilbert 多项式及维数
% * 次数、重数和首项系数
% * Euler 示性数的多项式表达
% 
% 纲要内容结束

\section{Hilbert 函数与 Euler 多项式}
\label{g5:sec:1903}

待撰写

\subsection{分次增长和有理级数}
\label{g5:subsec:1903:01}
待撰写

\subsection{Hilbert 多项式及维数}
\label{g5:subsec:1903:02}
待撰写

\subsection{次数、重数和首项系数}
\label{g5:subsec:1903:03}
待撰写

\subsection{Euler 示性数的多项式表达}
\label{g5:subsec:1903:04}
待撰写
% 纲要条目排为编号小节标题，并进入目录与书签；说明段仍保留于纲要注释中。

% !TeX root = ../../Book5.tex
% 纲要标题：### 19.4 固有凝聚直像
% 纲要位置：几何学计划纲要.md，第 5792 行。
% 正文状态：未撰写；本文件目前是框架。

% 纲要内容开始
% 
% * Chow 引理的射影修改构造
% * 支集维数的 Noether 归纳
% 
%   使用高阶直像的 Leray 比较、射影修改和支集递降建立有限性归纳；把修改上层的有限性与误差层支集变小分别论证，不把双有理本身当作上同调不变。
% 
% * 由射影情形推广有限性
% * 高阶直像的凝聚性
% * 不能仅凭每根纤维有限维推出凝聚
% 
% 纲要内容结束

\section{固有凝聚直像}
\label{g5:sec:1904}

待撰写

\subsection{Chow 引理的射影修改构造}
\label{g5:subsec:1904:01}
待撰写

\subsection{支集维数的 Noether 归纳}
\label{g5:subsec:1904:02}
待撰写

\subsection{由射影情形推广有限性}
\label{g5:subsec:1904:03}
待撰写

\subsection{高阶直像的凝聚性}
\label{g5:subsec:1904:04}
待撰写

\subsection{不能仅凭每根纤维有限维推出凝聚}
\label{g5:subsec:1904:05}
待撰写
% 纲要条目排为编号小节标题，并进入目录与书签；说明段仍保留于纲要注释中。

% !TeX root = ../../Book5.tex
% 纲要标题：### 19.5 射影模型与正合列
% 纲要位置：几何学计划纲要.md，第 5803 行。
% 正文状态：未撰写；本文件目前是框架。

% 纲要内容开始
% 
% * 超曲面及完全交
% * 曲线的次数和算术亏格
% * 生成、消失与嵌入的实际应用
% 
% ---
% 
% 纲要内容结束

\section{射影模型与正合列}
\label{g5:sec:1905}

待撰写

\subsection{超曲面及完全交}
\label{g5:subsec:1905:01}
待撰写

\subsection{曲线的次数和算术亏格}
\label{g5:subsec:1905:02}
待撰写

\subsection{生成、消失与嵌入的实际应用}
\label{g5:subsec:1905:03}
待撰写
% 纲要条目排为编号小节标题，并进入目录与书签；说明段仍保留于纲要注释中。


\begin{chaptersummary}
待撰写
\end{chaptersummary}

\BookExercises
待撰写

% !TeX root = ../../Book5.tex
% 纲要标题：## 第20章　上同调与基变换、上半连续性
% 纲要位置：几何学计划纲要.md，第 5517 行。
% 正文状态：未撰写；本文件目前是框架。

% 纲要内容开始
% 
% > 本章目标与证明范围：以局部有限复形为主工具，证明纤维上同调维数和基变换映射的关系。固有、Noether、凝聚及基底平坦性各自的用途明确列出。[R3]
% 
% 纲要内容结束

\chapter{上同调与基变换、上半连续性}
\label{g5:ch:20}
\BookChapterNumber{g5:ch:20}

\begin{chapterabstract}
待撰写
\end{chapterabstract}

\begin{learninggoals}
待撰写
\end{learninggoals}

待撰写

% !TeX root = ../../Book5.tex
% 纲要标题：### 20.1 基变换映射
% 纲要位置：几何学计划纲要.md，第 5815 行。
% 正文状态：未撰写；本文件目前是框架。

% 纲要内容开始
% 
% * 笛卡尔方块和自然映射
% * 平坦基变换的先行证明
% * 普通张量与导出张量的区别
% * 不交换的显式复形模型
% 
% 纲要内容结束

\section{基变换映射}
\label{g5:sec:2001}

待撰写

\subsection{笛卡尔方块和自然映射}
\label{g5:subsec:2001:01}
待撰写

\subsection{平坦基变换的先行证明}
\label{g5:subsec:2001:02}
待撰写

\subsection{普通张量与导出张量的区别}
\label{g5:subsec:2001:03}
待撰写

\subsection{不交换的显式复形模型}
\label{g5:subsec:2001:04}
待撰写
% 纲要条目排为编号小节标题，并进入目录与书签；说明段仍保留于纲要注释中。

% !TeX root = ../../Book5.tex
% 纲要标题：### 20.2 有限复形模型的构造
% 纲要位置：几何学计划纲要.md，第 5528 行。
% 正文状态：未撰写；本文件目前是框架。

% 纲要内容开始
% 
% * 仿射覆盖与有界平坦 Čech 复形
% * 凝聚直像和有限呈示替换
% * 基底平坦的凝聚层的有限自由模型
% * 任意基底模张量的计算
% * 局部构造与相容性
% 
% 纲要内容结束

\section{有限复形模型的构造}
\label{g5:sec:2002}

待撰写

\subsection{仿射覆盖与有界平坦 Čech 复形}
\label{g5:subsec:2002:01}
待撰写

\subsection{凝聚直像和有限呈示替换}
\label{g5:subsec:2002:02}
待撰写

\subsection{基底平坦的凝聚层的有限自由模型}
\label{g5:subsec:2002:03}
待撰写

\subsection{任意基底模张量的计算}
\label{g5:subsec:2002:04}
待撰写

\subsection{局部构造与相容性}
\label{g5:subsec:2002:05}
待撰写
% 纲要条目排为编号小节标题，并进入目录与书签；说明段仍保留于纲要注释中。

% !TeX root = ../../Book5.tex
% 纲要标题：### 20.3 上半连续与局部自由
% 纲要位置：几何学计划纲要.md，第 5536 行。
% 正文状态：未撰写；本文件目前是框架。

% 纲要内容开始
% 
% * 剩余域上微分矩阵的秩
% * 上同调维数的上半连续
% * 维数常定的核像分析
% * 局部自由及相邻次数条件
% 
% 纲要内容结束

\section{上半连续与局部自由}
\label{g5:sec:2003}

待撰写

\subsection{剩余域上微分矩阵的秩}
\label{g5:subsec:2003:01}
待撰写

\subsection{上同调维数的上半连续}
\label{g5:subsec:2003:02}
待撰写

\subsection{维数常定的核像分析}
\label{g5:subsec:2003:03}
待撰写

\subsection{局部自由及相邻次数条件}
\label{g5:subsec:2003:04}
待撰写
% 纲要条目排为编号小节标题，并进入目录与书签；说明段仍保留于纲要注释中。

% !TeX root = ../../Book5.tex
% 纲要标题：### 20.4 上同调与基变换定理
% 纲要位置：几何学计划纲要.md，第 5543 行。
% 正文状态：未撰写；本文件目前是框架。

% 纲要内容开始
% 
% * 满射性在邻域中的传播
% * 基变换同构的准确判据
% * Grauert 型结论的代数版本
% * 维数跳跃与失效例子
% 
% 纲要内容结束

\section{上同调与基变换定理}
\label{g5:sec:2004}

待撰写

\subsection{满射性在邻域中的传播}
\label{g5:subsec:2004:01}
待撰写

\subsection{基变换同构的准确判据}
\label{g5:subsec:2004:02}
待撰写

\subsection{Grauert 型结论的代数版本}
\label{g5:subsec:2004:03}
待撰写

\subsection{维数跳跃与失效例子}
\label{g5:subsec:2004:04}
待撰写
% 纲要条目排为编号小节标题，并进入目录与书签；说明段仍保留于纲要注释中。

% !TeX root = ../../Book5.tex
% 纲要标题：### 20.5 平坦族的 Hilbert 多项式
% 纲要位置：几何学计划纲要.md，第 5550 行。
% 正文状态：未撰写；本文件目前是框架。

% 纲要内容开始
% 
% * Euler 数的局部常性
% * 射影平坦族的 Hilbert 多项式
% * 反向平坦判据所需的具体假设
% * 退化曲线中的计算
% 
% ---
% 
% 纲要内容结束

\section{平坦族的 Hilbert 多项式}
\label{g5:sec:2005}

待撰写

\subsection{Euler 数的局部常性}
\label{g5:subsec:2005:01}
待撰写

\subsection{射影平坦族的 Hilbert 多项式}
\label{g5:subsec:2005:02}
待撰写

\subsection{反向平坦判据所需的具体假设}
\label{g5:subsec:2005:03}
待撰写

\subsection{退化曲线中的计算}
\label{g5:subsec:2005:04}
待撰写
% 纲要条目排为编号小节标题，并进入目录与书签；说明段仍保留于纲要注释中。


\begin{chaptersummary}
待撰写
\end{chaptersummary}

\BookExercises
待撰写

% !TeX root = ../../Book5.tex
% 纲要标题：## 第21章*　形式函数、连通纤维与 Stein 分解
% 纲要位置：几何学计划纲要.md，第 5559 行。
% 正文状态：未撰写；本文件目前是框架。

% 纲要内容开始
% 
% > 本章目标与证明范围：形式函数和 Stein 分解从基变换—对偶混合章中分离。采用 Noether 固有及凝聚的主版本，证明逆极限比较和几何连通纤维，不把纤维的普通截面张量式误当作全部证明。[R4]
% 
% 纲要内容结束

\OptionalChapter{形式函数、连通纤维与 Stein 分解}{g5:ch:21}
\BookChapterNumber{g5:ch:21}

\begin{chapterabstract}
待撰写
\end{chapterabstract}

\begin{learninggoals}
待撰写
\end{learninggoals}

待撰写

% !TeX root = ../../Book5.tex
% 纲要标题：### 21.1 无穷小邻域中的上同调
% 纲要位置：几何学计划纲要.md，第 5563 行。
% 正文状态：未撰写；本文件目前是框架。

% 纲要内容开始
% 
% * 理想幂作用和截断层
% * 两种诱导滤过的比较
% * 有限性及 Artin–Rees 型控制
% * Mittag–Leffler 条件的使用
% 
% 纲要内容结束

\section{无穷小邻域中的上同调}
\label{g5:sec:2101}

待撰写

\subsection{理想幂作用和截断层}
\label{g5:subsec:2101:01}
待撰写

\subsection{两种诱导滤过的比较}
\label{g5:subsec:2101:02}
待撰写

\subsection{有限性及 Artin–Rees 型控制}
\label{g5:subsec:2101:03}
待撰写

\subsection{Mittag–Leffler 条件的使用}
\label{g5:subsec:2101:04}
待撰写
% 纲要条目排为编号小节标题，并进入目录与书签；说明段仍保留于纲要注释中。

% !TeX root = ../../Book5.tex
% 纲要标题：### 21.2 形式函数定理的连续证明
% 纲要位置：几何学计划纲要.md，第 5864 行。
% 正文状态：未撰写；本文件目前是框架。

% 纲要内容开始
% 
% * 比较映射及主版本
% 
%   对 Noether 环 A、理想 I、固有 f:X→Spec(A)、凝聚 F，构造 H^q(X,F)^∧→lim_n H^q(X,F/I^nF)。
% 
% * 诱导滤过的等价
% 
%   证明由上同调映射核得到的滤过与 I 进滤过互相控制。
% 
% * 核的稳定与有限性
% 
%   用有限分次模及 Artin–Rees 型估计控制截断产生的核。
% 
% * 像系统的 Mittag–Leffler 性
% 
%   不把逆极限当作一般正合函子，单独证明所用系统的稳定像。
% 
% * 极限上的同构
% 
%   分别验证单射与满射，并识别两侧的拓扑。
% 
% * 局部化与纤维邻域
% 
%   得到高阶直像茎完备化和无穷小纤维上同调的比较。
% 
% 纲要内容结束

\section{形式函数定理的连续证明}
\label{g5:sec:2102}

待撰写

\subsection{比较映射及主版本}
\label{g5:subsec:2102:01}
待撰写

\subsection{诱导滤过的等价}
\label{g5:subsec:2102:02}
待撰写

\subsection{核的稳定与有限性}
\label{g5:subsec:2102:03}
待撰写

\subsection{像系统的 Mittag–Leffler 性}
\label{g5:subsec:2102:04}
待撰写

\subsection{极限上的同构}
\label{g5:subsec:2102:05}
待撰写

\subsection{局部化与纤维邻域}
\label{g5:subsec:2102:06}
待撰写
% 纲要条目排为编号小节标题，并进入目录与书签；说明段仍保留于纲要注释中。

% !TeX root = ../../Book5.tex
% 纲要标题：### 21.3 凝聚直像代数与 Stein 分解的有限部分
% 纲要位置：几何学计划纲要.md，第 5596 行。
% 正文状态：未撰写；本文件目前是框架。

% 纲要内容开始
% 
% 相对 Spec 的基础构造已在凝聚层章完成；本节只处理固有直像代数上的应用。几何连通纤维仍由第21.4节的独立论证得到。
% 
% * $f_*\mathcal O_X$ 的凝聚代数结构
% * 调用 V.13.5 构造 $Y^\prime=\operatorname{Spec}_Y(f_*\mathcal O_X)$
% * 自然因子分解及有限态射的验证
% * 第一段态射的结构层直像与恒等
% 
% 纲要内容结束

\section{凝聚直像代数与 Stein 分解的有限部分}
\label{g5:sec:2103}

待撰写

\subsection{\texorpdfstring{\(f_*\mathcal O_X\) 的凝聚代数结构}{f\_*O\_X 的凝聚代数结构}}
\label{g5:subsec:2103:01}
待撰写

\subsection{\texorpdfstring{调用 V.13.5 构造 \(Y^\prime=\operatorname{Spec}_Y(f_*\mathcal O_X)\)}{调用 V.13.5 构造 Y′=Spec\_Y(f\_*O\_X)}}
\label{g5:subsec:2103:02}
待撰写

\subsection{自然因子分解及有限态射的验证}
\label{g5:subsec:2103:03}
待撰写

\subsection{第一段态射的结构层直像与恒等}
\label{g5:subsec:2103:04}
待撰写
% 纲要条目排为编号小节标题，并进入目录与书签；说明段仍保留于纲要注释中。

% !TeX root = ../../Book5.tex
% 纲要标题：### 21.4 几何连通纤维
% 纲要位置：几何学计划纲要.md，第 5899 行。
% 正文状态：未撰写；本文件目前是框架。

% 纲要内容开始
% 
% * 纤维幂等元与连通分解
% * 形式邻域中的幂等元提升
% * 严格局部化及几何连通性
% * Stein 分解的唯一性
% * 正规性和有限双有理特例
% 
% ---
% 
% 纲要内容结束

\section{几何连通纤维}
\label{g5:sec:2104}

待撰写

\subsection{纤维幂等元与连通分解}
\label{g5:subsec:2104:01}
待撰写

\subsection{形式邻域中的幂等元提升}
\label{g5:subsec:2104:02}
待撰写

\subsection{严格局部化及几何连通性}
\label{g5:subsec:2104:03}
待撰写

\subsection{Stein 分解的唯一性}
\label{g5:subsec:2104:04}
待撰写

\subsection{正规性和有限双有理特例}
\label{g5:subsec:2104:05}
待撰写
% 纲要条目排为编号小节标题，并进入目录与书签；说明段仍保留于纲要注释中。


\begin{chaptersummary}
待撰写
\end{chaptersummary}

\BookExercises
待撰写

% !TeX root = ../../Book5.tex
% 纲要标题：## 第22章　Serre 对偶、典范层与对偶复形
% 纲要位置：几何学计划纲要.md，第 5615 行。
% 正文状态：未撰写；本文件目前是框架。

% 纲要内容开始
% 
% > 本章目标与证明范围：光滑射影纯维对象上的 Serre 对偶为完整主模块。一般对偶复形在有明确构造假设的版本下介绍，不能对任意奇异空间只写一个典范线丛公式。
% 
% 纲要内容结束

\chapter{Serre 对偶、典范层与对偶复形}
\label{g5:ch:22}
\BookChapterNumber{g5:ch:22}

\begin{chapterabstract}
待撰写
\end{chapterabstract}

\begin{learninggoals}
待撰写
\end{learninggoals}

待撰写

% !TeX root = ../../Book5.tex
% 纲要标题：### 22.1 射影空间的对偶
% 纲要位置：几何学计划纲要.md，第 5913 行。
% 正文状态：未撰写；本文件目前是框架。

% 纲要内容开始
% 
% * 顶次 Čech 配对
% * 迹映射和符号
% * 扭转线丛的非退化性
% * 有限消解及一般凝聚层
% 
% 纲要内容结束

\section{射影空间的对偶}
\label{g5:sec:2201}

待撰写

\subsection{顶次 Čech 配对}
\label{g5:subsec:2201:01}
待撰写

\subsection{迹映射和符号}
\label{g5:subsec:2201:02}
待撰写

\subsection{扭转线丛的非退化性}
\label{g5:subsec:2201:03}
待撰写

\subsection{有限消解及一般凝聚层}
\label{g5:subsec:2201:04}
待撰写
% 纲要条目排为编号小节标题，并进入目录与书签；说明段仍保留于纲要注释中。

% !TeX root = ../../Book5.tex
% 纲要标题：### 22.2 闭嵌入与典范层
% 纲要位置：几何学计划纲要.md，第 5626 行。
% 正文状态：未撰写；本文件目前是框架。

% 纲要内容开始
% 
% * 局部 Koszul 和 Ext
% * 光滑闭嵌入的余法行列式
% * 伴随公式
% * 由环境对偶传递到子簇
% 
% 纲要内容结束

\section{闭嵌入与典范层}
\label{g5:sec:2202}

待撰写

\subsection{局部 Koszul 和 Ext}
\label{g5:subsec:2202:01}
待撰写

\subsection{光滑闭嵌入的余法行列式}
\label{g5:subsec:2202:02}
待撰写

\subsection{伴随公式}
\label{g5:subsec:2202:03}
待撰写

\subsection{由环境对偶传递到子簇}
\label{g5:subsec:2202:04}
待撰写
% 纲要条目排为编号小节标题，并进入目录与书签；说明段仍保留于纲要注释中。

% !TeX root = ../../Book5.tex
% 纲要标题：### 22.3 Serre 对偶的完整证明
% 纲要位置：几何学计划纲要.md，第 5633 行。
% 正文状态：未撰写；本文件目前是框架。

% 纲要内容开始
% 
% * 光滑射影纯维的标准假设
% * Hⁱ 与 Extⁿ⁻ⁱ 的配对
% * 迹及自然性
% * 非退化性与长正合列相容
% * 向量丛上的经典形式
% 
% 纲要内容结束

\section{Serre 对偶的完整证明}
\label{g5:sec:2203}

待撰写

\subsection{光滑射影纯维的标准假设}
\label{g5:subsec:2203:01}
待撰写

\subsection{\texorpdfstring{\(H^i\) 与 \(\operatorname{Ext}^{n-i}\) 的配对}{Hⁱ 与 Extⁿ⁻ⁱ 的配对}}
\label{g5:subsec:2203:02}
待撰写

\subsection{迹及自然性}
\label{g5:subsec:2203:03}
待撰写

\subsection{非退化性与长正合列相容}
\label{g5:subsec:2203:04}
待撰写

\subsection{向量丛上的经典形式}
\label{g5:subsec:2203:05}
待撰写
% 纲要条目排为编号小节标题，并进入目录与书签；说明段仍保留于纲要注释中。

% !TeX root = ../../Book5.tex
% 纲要标题：### 22.4 奇异对象的对偶
% 纲要位置：几何学计划纲要.md，第 5641 行。
% 正文状态：未撰写；本文件目前是框架。

% 纲要内容开始
% 
% * CM 对象的典范模或层
% * Gorenstein 与可逆典范层
% * 非 CM 情形的多次数对偶复形
% * 一般 Grothendieck 对偶的精确外部输入
% 
% ---
% 
% 纲要内容结束

\section{奇异对象的对偶}
\label{g5:sec:2204}

待撰写

\subsection{CM 对象的典范模或层}
\label{g5:subsec:2204:01}
待撰写

\subsection{Gorenstein 与可逆典范层}
\label{g5:subsec:2204:02}
待撰写

\subsection{非 CM 情形的多次数对偶复形}
\label{g5:subsec:2204:03}
待撰写

\subsection{一般 Grothendieck 对偶的精确外部输入}
\label{g5:subsec:2204:04}
待撰写
% 纲要条目排为编号小节标题，并进入目录与书签；说明段仍保留于纲要注释中。


\begin{chaptersummary}
待撰写
\end{chaptersummary}

\BookExercises
待撰写


% !TeX root = ../Book5.tex
% 纲要标题：# 第四编　曲线、曲面与双有理模型
% 纲要位置：几何学计划纲要.md，第 5944 行。
% 正文状态：未撰写；本文件目前是框架。

% 纲要内容开始
% 
% 纲要内容结束

\BookPart{曲线、曲面与双有理模型}{g5:part:04}
\BookPartGuide
待撰写

% !TeX root = ../../Book5.tex
% 纲要标题：## 第23章　代数曲线、留数与 Riemann–Roch
% 纲要位置：几何学计划纲要.md，第 5652 行。
% 正文状态：未撰写；本文件目前是框架。

% 纲要内容开始
% 
% > 本章目标与证明范围：在光滑射影几何整曲线的范围内给出独立于复分析的代数证明，保留一般基域下的闭点次数。主部—坐标留数的证明先取完美基域；任意基域时在曲线层面建立所需的代数对偶配对，不对不可分闭点无条件使用系数域和剩余域迹公式。
% 
% 纲要内容结束

\chapter{代数曲线、留数与 Riemann–Roch}
\label{g5:ch:23}
\BookChapterNumber{g5:ch:23}

\begin{chapterabstract}
待撰写
\end{chapterabstract}

\begin{learninggoals}
待撰写
\end{learninggoals}

待撰写

% !TeX root = ../../Book5.tex
% 纲要标题：### 23.1 曲线和函数域
% 纲要位置：几何学计划纲要.md，第 5656 行。
% 正文状态：未撰写；本文件目前是框架。

% 纲要内容开始
% 
% * 一变量函数域的正规射影模型
% * 光滑、正则和基域条件
% * 闭点的离散赋值
% * 有限映射和局部参数
% 
% 纲要内容结束

\section{曲线和函数域}
\label{g5:sec:2301}

待撰写

\subsection{一变量函数域的正规射影模型}
\label{g5:subsec:2301:01}
待撰写

\subsection{光滑、正则和基域条件}
\label{g5:subsec:2301:02}
待撰写

\subsection{闭点的离散赋值}
\label{g5:subsec:2301:03}
待撰写

\subsection{有限映射和局部参数}
\label{g5:subsec:2301:04}
待撰写
% 纲要条目排为编号小节标题，并进入目录与书签；说明段仍保留于纲要注释中。

% !TeX root = ../../Book5.tex
% 纲要标题：### 23.2 除子与主部
% 纲要位置：几何学计划纲要.md，第 5957 行。
% 正文状态：未撰写；本文件目前是框架。

% 纲要内容开始
% 
% * 闭点次数和除子次数
% * 线丛及有理截面
% * 主部空间和局部有限支集
% * 长正合列的具体代表
% 
% 纲要内容结束

\section{除子与主部}
\label{g5:sec:2302}

待撰写

\subsection{闭点次数和除子次数}
\label{g5:subsec:2302:01}
待撰写

\subsection{线丛及有理截面}
\label{g5:subsec:2302:02}
待撰写

\subsection{主部空间和局部有限支集}
\label{g5:subsec:2302:03}
待撰写

\subsection{长正合列的具体代表}
\label{g5:subsec:2302:04}
待撰写
% 纲要条目排为编号小节标题，并进入目录与书签；说明段仍保留于纲要注释中。

% !TeX root = ../../Book5.tex
% 纲要标题：### 23.3 完美基域上的代数留数与对偶
% 纲要位置：几何学计划纲要.md，第 5670 行。
% 正文状态：未撰写；本文件目前是框架。

% 纲要内容开始
% 
% * 局部参数中的留数
% * 可分剩余域的迹及参数无关性
% * 整体留数定理
% * Serre 配对及非退化性
% * 与一般对偶章的独立比较
% 
% 纲要内容结束

\section{完美基域上的代数留数与对偶}
\label{g5:sec:2303}

待撰写

\subsection{局部参数中的留数}
\label{g5:subsec:2303:01}
待撰写

\subsection{可分剩余域的迹及参数无关性}
\label{g5:subsec:2303:02}
待撰写

\subsection{整体留数定理}
\label{g5:subsec:2303:03}
待撰写

\subsection{Serre 配对及非退化性}
\label{g5:subsec:2303:04}
待撰写

\subsection{与一般对偶章的独立比较}
\label{g5:subsec:2303:05}
待撰写
% 纲要条目排为编号小节标题，并进入目录与书签；说明段仍保留于纲要注释中。

% !TeX root = ../../Book5.tex
% 纲要标题：### 23.4 任意基域下的曲线对偶
% 纲要位置：几何学计划纲要.md，第 5972 行。
% 正文状态：未撰写；本文件目前是框架。

% 纲要内容开始
% 
% 本节保持任意基域的 Riemann–Roch 目标；对偶工具取本卷已经建立的代数方法，不引用复解析曲线结论。
% 
% * 光滑曲线的典范线丛与迹
% 
%   把第22章的代数构造专门化到曲线，重新写出配对和迹映射，明确其独立于所选射影嵌入。
% 
% * 曲线对偶的非退化性
% 
%   在曲线的有限消解和长正合列中核验非退化性及自然性，保留非代数闭基域的线性代数。
% 
% * 不可分闭点与坐标留数的限制
% 
%   几何光滑不意味着所有闭点的剩余域扩张都可分；不把分离变量后的系数域公式无条件用于这些点。
% 
% * 点处修改与闭点次数
% 
%   以商层的 k 维数得到 [κ(x):k] 的贡献，为下一节的统一次数公式准备。
% 
% 纲要内容结束

\section{任意基域下的曲线对偶}
\label{g5:sec:2304}

待撰写

\subsection{光滑曲线的典范线丛与迹}
\label{g5:subsec:2304:01}
待撰写

\subsection{曲线对偶的非退化性}
\label{g5:subsec:2304:02}
待撰写

\subsection{不可分闭点与坐标留数的限制}
\label{g5:subsec:2304:03}
待撰写

\subsection{点处修改与闭点次数}
\label{g5:subsec:2304:04}
待撰写
% 纲要条目排为编号小节标题，并进入目录与书签；说明段仍保留于纲要注释中。

% !TeX root = ../../Book5.tex
% 纲要标题：### 23.5 Riemann–Roch 及线性系
% 纲要位置：几何学计划纲要.md，第 5992 行。
% 正文状态：未撰写；本文件目前是框架。

% 纲要内容开始
% 
% * 点处修改与 Euler 数
% * Riemann–Roch 公式的推导
% * 高次数有效性、无基点及嵌入
% * 典范类和亏格
% * 特殊除子与 Clifford 型结论的所选证明
% 
% ---
% 
% 纲要内容结束

\section{Riemann–Roch 及线性系}
\label{g5:sec:2305}

待撰写

\subsection{点处修改与 Euler 数}
\label{g5:subsec:2305:01}
待撰写

\subsection{Riemann–Roch 公式的推导}
\label{g5:subsec:2305:02}
待撰写

\subsection{高次数有效性、无基点及嵌入}
\label{g5:subsec:2305:03}
待撰写

\subsection{典范类和亏格}
\label{g5:subsec:2305:04}
待撰写

\subsection{特殊除子与 Clifford 型结论的所选证明}
\label{g5:subsec:2305:05}
待撰写
% 纲要条目排为编号小节标题，并进入目录与书签；说明段仍保留于纲要注释中。


\begin{chaptersummary}
待撰写
\end{chaptersummary}

\BookExercises
待撰写

% !TeX root = ../../Book5.tex
% 纲要标题：## 第24章　曲线映射、分歧与 Jacobi 几何
% 纲要位置：几何学计划纲要.md，第 6002 行。
% 正文状态：未撰写；本文件目前是框架。

% 纲要内容开始
% 
% > 本章目标与证明范围：按可分、野分歧及纯不可分分别处理 Hurwitz。典范映射、超椭圆及椭圆群结构有完整模型；一般 Jacobi 可表性在模空间章落实。
% 
% 纲要内容结束

\chapter{曲线映射、分歧与 Jacobi 几何}
\label{g5:ch:24}
\BookChapterNumber{g5:ch:24}

\begin{chapterabstract}
待撰写
\end{chapterabstract}

\begin{learninggoals}
待撰写
\end{learninggoals}

待撰写

% !TeX root = ../../Book5.tex
% 纲要标题：### 24.1 有限曲线映射与分歧
% 纲要位置：几何学计划纲要.md，第 6006 行。
% 正文状态：未撰写；本文件目前是框架。

% 纲要内容开始
% 
% * 局部次数、剩余域和阶
% * 微分模及不同理想
% * 可分 Hurwitz 公式
% * 驯分歧的 e−1 和野分歧修正
% * 纯不可分映射的单独处理
% 
% 纲要内容结束

\section{有限曲线映射与分歧}
\label{g5:sec:2401}

待撰写

\subsection{局部次数、剩余域和阶}
\label{g5:subsec:2401:01}
待撰写

\subsection{微分模及不同理想}
\label{g5:subsec:2401:02}
待撰写

\subsection{可分 Hurwitz 公式}
\label{g5:subsec:2401:03}
待撰写

\subsection{\texorpdfstring{驯分歧的 \(e-1\) 和野分歧修正}{驯分歧的 e−1 和野分歧修正}}
\label{g5:subsec:2401:04}
待撰写

\subsection{纯不可分映射的单独处理}
\label{g5:subsec:2401:05}
待撰写
% 纲要条目排为编号小节标题，并进入目录与书签；说明段仍保留于纲要注释中。

% !TeX root = ../../Book5.tex
% 纲要标题：### 24.2 平面和典范曲线
% 纲要位置：几何学计划纲要.md，第 5720 行。
% 正文状态：未撰写；本文件目前是框架。

% 纲要内容开始
% 
% * 平面曲线的算术亏格
% * 正规化和 δ 修正
% * 典范线性系的映射
% * 超椭圆曲线与二重覆盖
% * 小亏格模型
% 
% 纲要内容结束

\section{平面和典范曲线}
\label{g5:sec:2402}

待撰写

\subsection{平面曲线的算术亏格}
\label{g5:subsec:2402:01}
待撰写

\subsection{\texorpdfstring{正规化和 \(\delta\) 修正}{正规化和 \GeometryPdfDelta{} 修正}}
\label{g5:subsec:2402:02}
待撰写

\subsection{典范线性系的映射}
\label{g5:subsec:2402:03}
待撰写

\subsection{超椭圆曲线与二重覆盖}
\label{g5:subsec:2402:04}
待撰写

\subsection{小亏格模型}
\label{g5:subsec:2402:05}
待撰写
% 纲要条目排为编号小节标题，并进入目录与书签；说明段仍保留于纲要注释中。

% !TeX root = ../../Book5.tex
% 纲要标题：### 24.3 椭圆曲线的几何
% 纲要位置：几何学计划纲要.md，第 5728 行。
% 正文状态：未撰写；本文件目前是框架。

% 纲要内容开始
% 
% * 基点和 Weierstrass 模型
% * 割线切线群律及正则性
% * Pic⁰ 的解释
% * 同源与特征 p 的现象
% 
% 纲要内容结束

\section{椭圆曲线的几何}
\label{g5:sec:2403}

待撰写

\subsection{基点和 Weierstrass 模型}
\label{g5:subsec:2403:01}
待撰写

\subsection{割线切线群律及正则性}
\label{g5:subsec:2403:02}
待撰写

\subsection{\texorpdfstring{\(\operatorname{Pic}^0\) 的解释}{Pic⁰ 的解释}}
\label{g5:subsec:2403:03}
待撰写

\subsection{同源与特征 p 的现象}
\label{g5:subsec:2403:04}
待撰写
% 纲要条目排为编号小节标题，并进入目录与书签；说明段仍保留于纲要注释中。

% !TeX root = ../../Book5.tex
% 纲要标题：### 24.4 Jacobi 簇及 Abel 映射
% 纲要位置：几何学计划纲要.md，第 5735 行。
% 正文状态：未撰写；本文件目前是框架。

% 纲要内容开始
% 
% * 对称积和有效除子
% * 除子类映射及其纤维
% * 代数闭域上的 Jacobi 构造入口
% * Poincaré 线丛所需的基点条件
% * 与第四卷周期模型的比较
% 
% ---
% 
% 纲要内容结束

\section{Jacobi 簇及 Abel 映射}
\label{g5:sec:2404}

待撰写

\subsection{对称积和有效除子}
\label{g5:subsec:2404:01}
待撰写

\subsection{除子类映射及其纤维}
\label{g5:subsec:2404:02}
待撰写

\subsection{代数闭域上的 Jacobi 构造入口}
\label{g5:subsec:2404:03}
待撰写

\subsection{Poincaré 线丛所需的基点条件}
\label{g5:subsec:2404:04}
待撰写

\subsection{与第四卷周期模型的比较}
\label{g5:subsec:2404:05}
待撰写
% 纲要条目排为编号小节标题，并进入目录与书签；说明段仍保留于纲要注释中。


\begin{chaptersummary}
待撰写
\end{chaptersummary}

\BookExercises
待撰写

% !TeX root = ../../Book5.tex
% 纲要标题：## 第25章　曲面交理论、Bézout 与 Riemann–Roch
% 纲要位置：几何学计划纲要.md，第 5745 行。
% 正文状态：未撰写；本文件目前是框架。

% 纲要内容开始
% 
% > 本章目标与证明范围：在光滑射影曲面上连续建立交数、伴随、Riemann–Roch 和 Hodge 指数。先局部计重再整体配对；公共分支不属于有限交点长度的直接适用范围。
% 
% 纲要内容结束

\chapter{曲面交理论、Bézout 与 Riemann–Roch}
\label{g5:ch:25}
\BookChapterNumber{g5:ch:25}

\begin{chapterabstract}
待撰写
\end{chapterabstract}

\begin{learninggoals}
待撰写
\end{learninggoals}

待撰写

% !TeX root = ../../Book5.tex
% 纲要标题：### 25.1 局部相交与重数
% 纲要位置：几何学计划纲要.md，第 6043 行。
% 正文状态：未撰写；本文件目前是框架。

% 纲要内容开始
% 
% * 二维正则局部环中的长度
% * 正则序列和有限交
% * 非横截相切的模型
% * 公共分支与不适当相交
% 
% 纲要内容结束

\section{局部相交与重数}
\label{g5:sec:2501}

待撰写

\subsection{二维正则局部环中的长度}
\label{g5:subsec:2501:01}
待撰写

\subsection{正则序列和有限交}
\label{g5:subsec:2501:02}
待撰写

\subsection{非横截相切的模型}
\label{g5:subsec:2501:03}
待撰写

\subsection{公共分支与不适当相交}
\label{g5:subsec:2501:04}
待撰写
% 纲要条目排为编号小节标题，并进入目录与书签；说明段仍保留于纲要注释中。

% !TeX root = ../../Book5.tex
% 纲要标题：### 25.2 全局除子交数
% 纲要位置：几何学计划纲要.md，第 5756 行。
% 正文状态：未撰写；本文件目前是框架。

% 纲要内容开始
% 
% * 限制线丛的次数
% * 双线性、对称性及线性等价
% * 射影平面和乘积曲面
% * 平面 Bézout 定理的完整证明
% 
% 纲要内容结束

\section{全局除子交数}
\label{g5:sec:2502}

待撰写

\subsection{限制线丛的次数}
\label{g5:subsec:2502:01}
待撰写

\subsection{双线性、对称性及线性等价}
\label{g5:subsec:2502:02}
待撰写

\subsection{射影平面和乘积曲面}
\label{g5:subsec:2502:03}
待撰写

\subsection{平面 Bézout 定理的完整证明}
\label{g5:subsec:2502:04}
待撰写
% 纲要条目排为编号小节标题，并进入目录与书签；说明段仍保留于纲要注释中。

% !TeX root = ../../Book5.tex
% 纲要标题：### 25.3 伴随公式与曲面 Riemann–Roch
% 纲要位置：几何学计划纲要.md，第 6057 行。
% 正文状态：未撰写；本文件目前是框架。

% 纲要内容开始
% 
% * 平滑曲线的典范层
% * 余法正合列和行列式
% * 增加曲线时 Euler 数的变化
% * 曲面 Riemann–Roch 的推导
% * 奇异除子的算术亏格公式
% 
% 纲要内容结束

\section{伴随公式与曲面 Riemann–Roch}
\label{g5:sec:2503}

待撰写

\subsection{平滑曲线的典范层}
\label{g5:subsec:2503:01}
待撰写

\subsection{余法正合列和行列式}
\label{g5:subsec:2503:02}
待撰写

\subsection{增加曲线时 Euler 数的变化}
\label{g5:subsec:2503:03}
待撰写

\subsection{曲面 Riemann–Roch 的推导}
\label{g5:subsec:2503:04}
待撰写

\subsection{奇异除子的算术亏格公式}
\label{g5:subsec:2503:05}
待撰写
% 纲要条目排为编号小节标题，并进入目录与书签；说明段仍保留于纲要注释中。

% !TeX root = ../../Book5.tex
% 纲要标题：### 25.4 Hodge 指数定理
% 纲要位置：几何学计划纲要.md，第 5771 行。
% 正文状态：未撰写；本文件目前是框架。

% 纲要内容开始
% 
% * 数值等价与 Néron–Severi 空间的所需有限性
% 
%   数值配对先在给定有限维除子空间上讨论；需要 Néron–Severi 的一般有限生成时，将其作为独立明确输入，不反向调用后置 Picard 可表性来填补证明。
% 
% * 充足类正交补上的负定性
% * 代数证明中的 Riemann–Roch 估计
% * 等号及数值平凡
% * 不能直接引用复 Kähler 结论处理任意基域
% 
% 纲要内容结束

\section{Hodge 指数定理}
\label{g5:sec:2504}

待撰写

\subsection{数值等价与 Néron–Severi 空间的所需有限性}
\label{g5:subsec:2504:01}
待撰写

\subsection{充足类正交补上的负定性}
\label{g5:subsec:2504:02}
待撰写

\subsection{代数证明中的 Riemann–Roch 估计}
\label{g5:subsec:2504:03}
待撰写

\subsection{等号及数值平凡}
\label{g5:subsec:2504:04}
待撰写

\subsection{不能直接引用复 Kähler 结论处理任意基域}
\label{g5:subsec:2504:05}
待撰写
% 纲要条目排为编号小节标题，并进入目录与书签；说明段仍保留于纲要注释中。

% !TeX root = ../../Book5.tex
% 纲要标题：### 25.5 计算与几何约束
% 纲要位置：几何学计划纲要.md，第 5782 行。
% 正文状态：未撰写；本文件目前是框架。

% 纲要内容开始
% 
% * 吹起平面的交矩阵
% * 曲线次数、重数和线性系
% * nef 与负曲线的初步对照
% 
% ---
% 
% 纲要内容结束

\section{计算与几何约束}
\label{g5:sec:2505}

待撰写

\subsection{吹起平面的交矩阵}
\label{g5:subsec:2505:01}
待撰写

\subsection{曲线次数、重数和线性系}
\label{g5:subsec:2505:02}
待撰写

\subsection{nef 与负曲线的初步对照}
\label{g5:subsec:2505:03}
待撰写
% 纲要条目排为编号小节标题，并进入目录与书签；说明段仍保留于纲要注释中。


\begin{chaptersummary}
待撰写
\end{chaptersummary}

\BookExercises
待撰写

% !TeX root = ../../Book5.tex
% 纲要标题：## 第26章　吹起、收缩与曲面的极小模型
% 纲要位置：几何学计划纲要.md，第 5790 行。
% 正文状态：未撰写；本文件目前是框架。

% 纲要内容开始
% 
% > 本章目标与证明范围：将正规化、消解、双有理分解和极小化作为不同操作。完整建立光滑点吹起和所选 Castelnuovo 收缩；高维 MMP 不反向承担曲面的基本证明。
% 
% 纲要内容结束

\chapter{吹起、收缩与曲面的极小模型}
\label{g5:ch:26}
\BookChapterNumber{g5:ch:26}

\begin{chapterabstract}
待撰写
\end{chapterabstract}

\begin{learninggoals}
待撰写
\end{learninggoals}

待撰写

% !TeX root = ../../Book5.tex
% 纲要标题：### 26.1 光滑点吹起的计算
% 纲要位置：几何学计划纲要.md，第 5794 行。
% 正文状态：未撰写；本文件目前是框架。

% 纲要内容开始
% 
% * 两个仿射图及例外曲线
% * 严格变换和重数
% * 例外曲线自交−1
% * 典范类与交数变换
% * 有理映射基点的消去
% 
% 纲要内容结束

\section{光滑点吹起的计算}
\label{g5:sec:2601}

待撰写

\subsection{两个仿射图及例外曲线}
\label{g5:subsec:2601:01}
待撰写

\subsection{严格变换和重数}
\label{g5:subsec:2601:02}
待撰写

\subsection{\texorpdfstring{例外曲线自交\(-1\)}{例外曲线自交−1}}
\label{g5:subsec:2601:03}
待撰写

\subsection{典范类与交数变换}
\label{g5:subsec:2601:04}
待撰写

\subsection{有理映射基点的消去}
\label{g5:subsec:2601:05}
待撰写
% 纲要条目排为编号小节标题，并进入目录与书签；说明段仍保留于纲要注释中。

% !TeX root = ../../Book5.tex
% 纲要标题：### 26.2 曲面收缩
% 纲要位置：几何学计划纲要.md，第 5802 行。
% 正文状态：未撰写；本文件目前是框架。

% 纲要内容开始
% 
% * 负曲线与线性系
% * Castelnuovo 判据的标准版本
% * 收缩态射的构造
% * 收缩后光滑性与唯一性
% 
% 纲要内容结束

\section{曲面收缩}
\label{g5:sec:2602}

待撰写

\subsection{负曲线与线性系}
\label{g5:subsec:2602:01}
待撰写

\subsection{Castelnuovo 判据的标准版本}
\label{g5:subsec:2602:02}
待撰写

\subsection{收缩态射的构造}
\label{g5:subsec:2602:03}
待撰写

\subsection{收缩后光滑性与唯一性}
\label{g5:subsec:2602:04}
待撰写
% 纲要条目排为编号小节标题，并进入目录与书签；说明段仍保留于纲要注释中。

% !TeX root = ../../Book5.tex
% 纲要标题：### 26.3 双有理分解与极小性
% 纲要位置：几何学计划纲要.md，第 5809 行。
% 正文状态：未撰写；本文件目前是框架。

% 纲要内容开始
% 
% * 光滑曲面双有理态射的分解
% * 共同消解及基点
% * 极小模型的存在
% * 有理、直纹和非直纹情形的唯一性差异
% 
% 纲要内容结束

\section{双有理分解与极小性}
\label{g5:sec:2603}

待撰写

\subsection{光滑曲面双有理态射的分解}
\label{g5:subsec:2603:01}
待撰写

\subsection{共同消解及基点}
\label{g5:subsec:2603:02}
待撰写

\subsection{极小模型的存在}
\label{g5:subsec:2603:03}
待撰写

\subsection{有理、直纹和非直纹情形的唯一性差异}
\label{g5:subsec:2603:04}
待撰写
% 纲要条目排为编号小节标题，并进入目录与书签；说明段仍保留于纲要注释中。

% !TeX root = ../../Book5.tex
% 纲要标题：### 26.4 典型曲面及分类接口
% 纲要位置：几何学计划纲要.md，第 5816 行。
% 正文状态：未撰写；本文件目前是框架。

% 纲要内容开始
% 
% * 射影平面与 Hirzebruch 曲面
% * 最小化过程和数值不变量
% * 特征条件及完整分类的外部输入
% 
% ---
% 
% 纲要内容结束

\section{典型曲面及分类接口}
\label{g5:sec:2604}

待撰写

\subsection{射影平面与 Hirzebruch 曲面}
\label{g5:subsec:2604:01}
待撰写

\subsection{最小化过程和数值不变量}
\label{g5:subsec:2604:02}
待撰写

\subsection{特征条件及完整分类的外部输入}
\label{g5:subsec:2604:03}
待撰写
% 纲要条目排为编号小节标题，并进入目录与书签；说明段仍保留于纲要注释中。


\begin{chaptersummary}
待撰写
\end{chaptersummary}

\BookExercises
待撰写


% !TeX root = ../Book5.tex
% 纲要标题：# 第五编　代数圈、环面模型与奇点正性
% 纲要位置：几何学计划纲要.md，第 5824 行。
% 正文状态：未撰写；本文件目前是框架。

% 纲要内容开始
% 
% 纲要内容结束

\BookPart{代数圈、环面模型与奇点正性}{g5:part:05}
\BookPartGuide
待撰写

% !TeX root = ../../Book5.tex
% 纲要标题：## 第27章*　Chow 群、精细相交与 Riemann–Roch
% 纲要位置：几何学计划纲要.md，第 6120 行。
% 正文状态：未撰写；本文件目前是框架。

% 纲要内容开始
% 
% > 本章目标与证明范围：从圈和有理等价建立函子性、Chern 类与正规锥，完成射影丛、非横截交和 Grassmann 枚举的代表计算。再构造 Chern 特征与 Todd 类，准确外引所用 Riemann–Roch 版本并完整计算低维推论，不另增一般 GRR 的完整长证明任务。
% 
% 纲要内容结束

\OptionalChapter{Chow 群、精细相交与 Riemann–Roch}{g5:ch:27}
\BookChapterNumber{g5:ch:27}

\begin{chapterabstract}
待撰写
\end{chapterabstract}

\begin{learninggoals}
待撰写
\end{learninggoals}

待撰写

% !TeX root = ../../Book5.tex
% 纲要标题：### 27.1 代数圈及有理等价
% 纲要位置：几何学计划纲要.md，第 5830 行。
% 正文状态：未撰写；本文件目前是框架。

% 纲要内容开始
% 
% * 维数、重数及基本圈
% * 有理函数的主除子
% * Chow 群和局部化序列
% * 仿射、射影空间的基本计算
% 
% 纲要内容结束

\section{代数圈及有理等价}
\label{g5:sec:2701}

待撰写

\subsection{维数、重数及基本圈}
\label{g5:subsec:2701:01}
待撰写

\subsection{有理函数的主除子}
\label{g5:subsec:2701:02}
待撰写

\subsection{Chow 群和局部化序列}
\label{g5:subsec:2701:03}
待撰写

\subsection{仿射、射影空间的基本计算}
\label{g5:subsec:2701:04}
待撰写
% 纲要条目排为编号小节标题，并进入目录与书签；说明段仍保留于纲要注释中。

% !TeX root = ../../Book5.tex
% 纲要标题：### 27.2 推前、拉回及投影公式
% 纲要位置：几何学计划纲要.md，第 5837 行。
% 正文状态：未撰写；本文件目前是框架。

% 纲要内容开始
% 
% * 固有推前与域次数
% * 平坦拉回及相对维数
% * 两种操作与有理等价相容
% * 投影公式和基变换
% 
% 纲要内容结束

\section{推前、拉回及投影公式}
\label{g5:sec:2702}

待撰写

\subsection{固有推前与域次数}
\label{g5:subsec:2702:01}
待撰写

\subsection{平坦拉回及相对维数}
\label{g5:subsec:2702:02}
待撰写

\subsection{两种操作与有理等价相容}
\label{g5:subsec:2702:03}
待撰写

\subsection{投影公式和基变换}
\label{g5:subsec:2702:04}
待撰写
% 纲要条目排为编号小节标题，并进入目录与书签；说明段仍保留于纲要注释中。

% !TeX root = ../../Book5.tex
% 纲要标题：### 27.3 向量丛的 Chern 运算
% 纲要位置：几何学计划纲要.md，第 5844 行。
% 正文状态：未撰写；本文件目前是框架。

% 纲要内容开始
% 
% * 射影丛公式
% * 零截面和顶 Chern 类
% * 分裂原理及 Whitney 公式
% * 零点计数和枚举模型
% 
% 纲要内容结束

\section{向量丛的 Chern 运算}
\label{g5:sec:2703}

待撰写

\subsection{射影丛公式}
\label{g5:subsec:2703:01}
待撰写

\subsection{零截面和顶 Chern 类}
\label{g5:subsec:2703:02}
待撰写

\subsection{分裂原理及 Whitney 公式}
\label{g5:subsec:2703:03}
待撰写

\subsection{零点计数和枚举模型}
\label{g5:subsec:2703:04}
待撰写
% 纲要条目排为编号小节标题，并进入目录与书签；说明段仍保留于纲要注释中。

% !TeX root = ../../Book5.tex
% 纲要标题：### 27.4 正规锥与精细 Gysin
% 纲要位置：几何学计划纲要.md，第 6145 行。
% 正文状态：未撰写；本文件目前是框架。

% 纲要内容开始
% 
% * 正规锥、正规丛和正则嵌入
% * 变形到正规锥
% * 精细拉回的构造
% * 超额相交丛和公式
% * 自交及非横截实例
% 
% 纲要内容结束

\section{正规锥与精细 Gysin}
\label{g5:sec:2704}

待撰写

\subsection{正规锥、正规丛和正则嵌入}
\label{g5:subsec:2704:01}
待撰写

\subsection{变形到正规锥}
\label{g5:subsec:2704:02}
待撰写

\subsection{精细拉回的构造}
\label{g5:subsec:2704:03}
待撰写

\subsection{超额相交丛和公式}
\label{g5:subsec:2704:04}
待撰写

\subsection{自交及非横截实例}
\label{g5:subsec:2704:05}
待撰写
% 纲要条目排为编号小节标题，并进入目录与书签；说明段仍保留于纲要注释中。

% !TeX root = ../../Book5.tex
% 纲要标题：### 27.5 Schubert 演算与四条直线的完整计数
% 纲要位置：几何学计划纲要.md，第 6153 行。
% 正文状态：未撰写；本文件目前是框架。

% 纲要内容开始
% 
% 完整枚举模块固定复数域；一般 Schubert 理论只展开本模型需要的生成元与乘法，不扩展为全部旗簇的表示论。
% 
% * 完全旗、秩条件及 Schubert 子簇的定义
% * 胞腔分解、Chow 群基与特殊 Schubert 类
% * 通用子丛、商丛及低阶 Pieri 乘法
% * $\operatorname{Gr}(2,4)\subset\mathbb P^5$ 的 Plücker 二次方程
% * 与一条固定直线相交的超平面条件
% * 四个一般位置条件的方程消去与交数核验
% 
%   完成 $\int_{\operatorname{Gr}(2,4)}\sigma_1^4=2$，并与 Plücker 二次超曲面和四个超平面条件的计算互相核验：三维复射影空间中与四条一般位置直线均相交的直线有两条。
% 
% * 一般位置、横截性与退化时的重数
% 
%   给出一个一般位置实例验证两解不同，并说明退化时计重的意义；实数域下的实解数不能直接由复交数替代。
% 
% * 与 II.21.5 的胞腔、杯积及复定向比较
% 
% 纲要内容结束

\section{Schubert 演算与四条直线的完整计数}
\label{g5:sec:2705}

待撰写

\subsection{完全旗、秩条件及 Schubert 子簇的定义}
\label{g5:subsec:2705:01}
待撰写

\subsection{胞腔分解、Chow 群基与特殊 Schubert 类}
\label{g5:subsec:2705:02}
待撰写

\subsection{通用子丛、商丛及低阶 Pieri 乘法}
\label{g5:subsec:2705:03}
待撰写

\subsection{\texorpdfstring{\(\operatorname{Gr}(2,4)\subset\mathbb P^5\) 的 Plücker 二次方程}{Gr(2,4)⊂P⁵ 的 Plücker 二次方程}}
\label{g5:subsec:2705:04}
待撰写

\subsection{与一条固定直线相交的超平面条件}
\label{g5:subsec:2705:05}
待撰写

\subsection{四个一般位置条件的方程消去与交数核验}
\label{g5:subsec:2705:06}
待撰写

\subsection{一般位置、横截性与退化时的重数}
\label{g5:subsec:2705:07}
待撰写

\subsection{与 II.21.5 的胞腔、杯积及复定向比较}
\label{g5:subsec:2705:08}
待撰写
% 纲要条目排为编号小节标题，并进入目录与书签；说明段仍保留于纲要注释中。

% !TeX root = ../../Book5.tex
% 纲要标题：### 27.6 陈特征、Todd 类与 Euler 示性数
% 纲要位置：几何学计划纲要.md，第 6172 行。
% 正文状态：未撰写；本文件目前是框架。

% 纲要内容开始
% 
% 以域上的纯维光滑射影概形及向量丛为计算主范围，使用有理 Chow 类；整数 Euler 数与含分母的特征类之间的关系需要核验。
% 
% * 向量丛的 Grothendieck 群与凝聚层的加性关系
% * 分裂原理下的陈特征及可加、可乘性
% * Todd 幂级数、低次数项及正规化
% * 光滑射影概形上的 Hirzebruch–Riemann–Roch 公式
% 
%   准确引用 $\chi(X,E)=\deg\bigl([\operatorname{ch}(E)\operatorname{td}(T_{X/k})]_{\dim X}\bigr)$；前面两种低维 Riemann–Roch 的证明不反向依赖本节。[R11]
% 
% * 曲线、曲面及射影空间的完整展开
% * 与代数曲线、曲面 Riemann–Roch 的独立核对
% 
% 纲要内容结束

\section{陈特征、Todd 类与 Euler 示性数}
\label{g5:sec:2706}

待撰写

\subsection{向量丛的 Grothendieck 群与凝聚层的加性关系}
\label{g5:subsec:2706:01}
待撰写

\subsection{分裂原理下的陈特征及可加、可乘性}
\label{g5:subsec:2706:02}
待撰写

\subsection{Todd 幂级数、低次数项及正规化}
\label{g5:subsec:2706:03}
待撰写

\subsection{光滑射影概形上的 Hirzebruch–Riemann–Roch 公式}
\label{g5:subsec:2706:04}
待撰写

\subsection{曲线、曲面及射影空间的完整展开}
\label{g5:subsec:2706:05}
待撰写

\subsection{与代数曲线、曲面 Riemann–Roch 的独立核对}
\label{g5:subsec:2706:06}
待撰写
% 纲要条目排为编号小节标题，并进入目录与书签；说明段仍保留于纲要注释中。

% !TeX root = ../../Book5.tex
% 纲要标题：### 27.7 推前形式与解析指标的比较
% 纲要位置：几何学计划纲要.md，第 6186 行。
% 正文状态：未撰写；本文件目前是框架。

% 纲要内容开始
% 
% 先准确使用高阶直像有限局部自由的固有光滑版本，再说明更一般表述需要的对象；不对任意奇异态射直接套用向量丛公式。[R11]
% 
% * 高阶直像的交错和与普通直像的区别
% * 固有光滑态射上相对 Todd 类的公式
% * 高阶直像不局部自由时的 K 群或完美复形修正
% * 一般 Grothendieck–Riemann–Roch 的外部理论位置
% * 复数域下代数、解析及拓扑 Chern 特征的比较
% 
%   解析指标比较只在复数域及已具备相应比较定理时使用；任意基域的代数公式不能由复解析证明不加说明地转移。
% 
% * 与 II.24.4、III.17、IV.20 的指标和 GAGA 路线衔接
% 
% ---
% 
% 纲要内容结束

\section{推前形式与解析指标的比较}
\label{g5:sec:2707}

待撰写

\subsection{高阶直像的交错和与普通直像的区别}
\label{g5:subsec:2707:01}
待撰写

\subsection{固有光滑态射上相对 Todd 类的公式}
\label{g5:subsec:2707:02}
待撰写

\subsection{高阶直像不局部自由时的 K 群或完美复形修正}
\label{g5:subsec:2707:03}
待撰写

\subsection{一般 Grothendieck–Riemann–Roch 的外部理论位置}
\label{g5:subsec:2707:04}
待撰写

\subsection{复数域下代数、解析及拓扑 Chern 特征的比较}
\label{g5:subsec:2707:05}
待撰写

\subsection{与 II.24.4、III.17、IV.20 的指标和 GAGA 路线衔接}
\label{g5:subsec:2707:06}
待撰写
% 纲要条目排为编号小节标题，并进入目录与书签；说明段仍保留于纲要注释中。


\begin{chaptersummary}
待撰写
\end{chaptersummary}

\BookExercises
待撰写

% !TeX root = ../../Book5.tex
% 纲要标题：## 第28章*　环面簇：从半群到除子和退化
% 纲要位置：几何学计划纲要.md，第 6202 行。
% 正文状态：未撰写；本文件目前是框架。

% 纲要内容开始
% 
% > 本章目标与证明范围：作为概形与除子之后的独立深度路线，不以前置 MMP 为条件。把组合判据证明到能够计算具体扇、局部奇点、线丛截面和交数的程度。
% 
% 纲要内容结束

\OptionalChapter{环面簇：从半群到除子和退化}{g5:ch:28}
\BookChapterNumber{g5:ch:28}

\begin{chapterabstract}
待撰写
\end{chapterabstract}

\begin{learninggoals}
待撰写
\end{learninggoals}

待撰写

% !TeX root = ../../Book5.tex
% 纲要标题：### 28.1 格、锥与仿射模型
% 纲要位置：几何学计划纲要.md，第 5912 行。
% 正文状态：未撰写；本文件目前是框架。

% 纲要内容开始
% 
% * 对偶格和有理多面锥
% * Gordan 引理与半群有限性
% * 半群代数的正规性
% * 轨道及面对应
% 
% 纲要内容结束

\section{格、锥与仿射模型}
\label{g5:sec:2801}

待撰写

\subsection{对偶格和有理多面锥}
\label{g5:subsec:2801:01}
待撰写

\subsection{Gordan 引理与半群有限性}
\label{g5:subsec:2801:02}
待撰写

\subsection{半群代数的正规性}
\label{g5:subsec:2801:03}
待撰写

\subsection{轨道及面对应}
\label{g5:subsec:2801:04}
待撰写
% 纲要条目排为编号小节标题，并进入目录与书签；说明段仍保留于纲要注释中。

% !TeX root = ../../Book5.tex
% 纲要标题：### 28.2 扇与整体性质
% 纲要位置：几何学计划纲要.md，第 5919 行。
% 正文状态：未撰写；本文件目前是框架。

% 纲要内容开始
% 
% * 仿射图的交与粘合
% * 分离性、完备性和扇支集
% * 光滑锥的格基判据
% * 细分与双有理态射
% 
% 纲要内容结束

\section{扇与整体性质}
\label{g5:sec:2802}

待撰写

\subsection{仿射图的交与粘合}
\label{g5:subsec:2802:01}
待撰写

\subsection{分离性、完备性和扇支集}
\label{g5:subsec:2802:02}
待撰写

\subsection{光滑锥的格基判据}
\label{g5:subsec:2802:03}
待撰写

\subsection{细分与双有理态射}
\label{g5:subsec:2802:04}
待撰写
% 纲要条目排为编号小节标题，并进入目录与书签；说明段仍保留于纲要注释中。

% !TeX root = ../../Book5.tex
% 纲要标题：### 28.3 不变除子和线丛
% 纲要位置：几何学计划纲要.md，第 5926 行。
% 正文状态：未撰写；本文件目前是框架。

% 纲要内容开始
% 
% * Weil 类群的格正合列
% * Cartier 数据和分片线性函数
% * 截面多面体及格点
% * 充足性和严格凸性
% * Picard 群与类群的区别
% 
% 纲要内容结束

\section{不变除子和线丛}
\label{g5:sec:2803}

待撰写

\subsection{Weil 类群的格正合列}
\label{g5:subsec:2803:01}
待撰写

\subsection{Cartier 数据和分片线性函数}
\label{g5:subsec:2803:02}
待撰写

\subsection{截面多面体及格点}
\label{g5:subsec:2803:03}
待撰写

\subsection{充足性和严格凸性}
\label{g5:subsec:2803:04}
待撰写

\subsection{Picard 群与类群的区别}
\label{g5:subsec:2803:05}
待撰写
% 纲要条目排为编号小节标题，并进入目录与书签；说明段仍保留于纲要注释中。

% !TeX root = ../../Book5.tex
% 纲要标题：### 28.4 交数和显式模型
% 纲要位置：几何学计划纲要.md，第 5934 行。
% 正文状态：未撰写；本文件目前是框架。

% 纲要内容开始
% 
% * 射影空间及 Hirzebruch 曲面
% * 混合体积和所选次数公式
% * 锥奇点的组合消解
% * 平面吹起的扇计算
% 
% 纲要内容结束

\section{交数和显式模型}
\label{g5:sec:2804}

待撰写

\subsection{射影空间及 Hirzebruch 曲面}
\label{g5:subsec:2804:01}
待撰写

\subsection{混合体积和所选次数公式}
\label{g5:subsec:2804:02}
待撰写

\subsection{锥奇点的组合消解}
\label{g5:subsec:2804:03}
待撰写

\subsection{平面吹起的扇计算}
\label{g5:subsec:2804:04}
待撰写
% 纲要条目排为编号小节标题，并进入目录与书签；说明段仍保留于纲要注释中。

% !TeX root = ../../Book5.tex
% 纲要标题：### 28.5 退化与辛比较
% 纲要位置：几何学计划纲要.md，第 6235 行。
% 正文状态：未撰写；本文件目前是框架。

% 纲要内容开始
% 
% * 多面体提升和环面退化
% * 平坦性及特殊纤维
% * 矩多面体和 Delzant 条件
% * 复射影模型与辛约化的接口
% 
% ---
% 
% 纲要内容结束

\section{退化与辛比较}
\label{g5:sec:2805}

待撰写

\subsection{多面体提升和环面退化}
\label{g5:subsec:2805:01}
待撰写

\subsection{平坦性及特殊纤维}
\label{g5:subsec:2805:02}
待撰写

\subsection{矩多面体和 Delzant 条件}
\label{g5:subsec:2805:03}
待撰写

\subsection{复射影模型与辛约化的接口}
\label{g5:subsec:2805:04}
待撰写
% 纲要条目排为编号小节标题，并进入目录与书签；说明段仍保留于纲要注释中。


\begin{chaptersummary}
待撰写
\end{chaptersummary}

\BookExercises
待撰写

% !TeX root = ../../Book5.tex
% 纲要标题：## 第29章*　奇点、对数对与消解
% 纲要位置：几何学计划纲要.md，第 5950 行。
% 正文状态：未撰写；本文件目前是框架。

% 纲要内容开始
% 
% > 本章目标与证明范围：以显式局部环和吹起计算为起点，建立差异及后续 MMP 所需奇点类型。代数分层、Whitney 分层和拓扑局部平凡性各在自己的语境中出现。 对数消解和 MMP 接口以代数闭特征零域为主范围；正特征版本不得无标记替换。
% 
% 纲要内容结束

\OptionalChapter{奇点、对数对与消解}{g5:ch:29}
\BookChapterNumber{g5:ch:29}

\begin{chapterabstract}
待撰写
\end{chapterabstract}

\begin{learninggoals}
待撰写
\end{learninggoals}

待撰写

% !TeX root = ../../Book5.tex
% 纲要标题：### 29.1 局部奇点的结构
% 纲要位置：几何学计划纲要.md，第 6248 行。
% 正文状态：未撰写；本文件目前是框架。

% 纲要内容开始
% 
% * 正规、CM 和 Gorenstein 回顾
% * 有理奇点的高阶直像条件
% * 商奇点和锥的模型
% * 同一切锥的不同奇点
% 
% 纲要内容结束

\section{局部奇点的结构}
\label{g5:sec:2901}

待撰写

\subsection{正规、CM 和 Gorenstein 回顾}
\label{g5:subsec:2901:01}
待撰写

\subsection{有理奇点的高阶直像条件}
\label{g5:subsec:2901:02}
待撰写

\subsection{商奇点和锥的模型}
\label{g5:subsec:2901:03}
待撰写

\subsection{同一切锥的不同奇点}
\label{g5:subsec:2901:04}
待撰写
% 纲要条目排为编号小节标题，并进入目录与书签；说明段仍保留于纲要注释中。

% !TeX root = ../../Book5.tex
% 纲要标题：### 29.2 消解的具体方法
% 纲要位置：几何学计划纲要.md，第 6255 行。
% 正文状态：未撰写；本文件目前是框架。

% 纲要内容开始
% 
% * 平面曲线连续吹起
% * 曲面简单奇点及例外图
% * 嵌入消解和正规交叉
% * 特征零一般消解定理的准确引用
% 
% 纲要内容结束

\section{消解的具体方法}
\label{g5:sec:2902}

待撰写

\subsection{平面曲线连续吹起}
\label{g5:subsec:2902:01}
待撰写

\subsection{曲面简单奇点及例外图}
\label{g5:subsec:2902:02}
待撰写

\subsection{嵌入消解和正规交叉}
\label{g5:subsec:2902:03}
待撰写

\subsection{特征零一般消解定理的准确引用}
\label{g5:subsec:2902:04}
待撰写
% 纲要条目排为编号小节标题，并进入目录与书签；说明段仍保留于纲要注释中。

% !TeX root = ../../Book5.tex
% 纲要标题：### 29.3 对数对和差异
% 纲要位置：几何学计划纲要.md，第 5968 行。
% 正文状态：未撰写；本文件目前是框架。

% 纲要内容开始
% 
% * Q-Cartier 的典范除子
% * 对数消解上的差异公式
% * 独立于消解的检验
% * klt、lc、canonical、terminal 的区别
% * 低维差异的完整计算
% 
% 纲要内容结束

\section{对数对和差异}
\label{g5:sec:2903}

待撰写

\subsection{Q-Cartier 的典范除子}
\label{g5:subsec:2903:01}
待撰写

\subsection{对数消解上的差异公式}
\label{g5:subsec:2903:02}
待撰写

\subsection{独立于消解的检验}
\label{g5:subsec:2903:03}
待撰写

\subsection{klt、lc、canonical、terminal 的区别}
\label{g5:subsec:2903:04}
待撰写

\subsection{低维差异的完整计算}
\label{g5:subsec:2903:05}
待撰写
% 纲要条目排为编号小节标题，并进入目录与书签；说明段仍保留于纲要注释中。

% !TeX root = ../../Book5.tex
% 纲要标题：### 29.4 分层和族中的奇点
% 纲要位置：几何学计划纲要.md，第 5976 行。
% 正文状态：未撰写；本文件目前是框架。

% 纲要内容开始
% 
% * 局部闭分层与构造性
% * 复解析 Whitney 条件
% * 等奇异族和局部拓扑
% * 奇点不变量的跳跃
% * 分层不能自动消除平坦性问题
% 
% ---
% 
% 纲要内容结束

\section{分层和族中的奇点}
\label{g5:sec:2904}

待撰写

\subsection{局部闭分层与构造性}
\label{g5:subsec:2904:01}
待撰写

\subsection{复解析 Whitney 条件}
\label{g5:subsec:2904:02}
待撰写

\subsection{等奇异族和局部拓扑}
\label{g5:subsec:2904:03}
待撰写

\subsection{奇点不变量的跳跃}
\label{g5:subsec:2904:04}
待撰写

\subsection{分层不能自动消除平坦性问题}
\label{g5:subsec:2904:05}
待撰写
% 纲要条目排为编号小节标题，并进入目录与书签；说明段仍保留于纲要注释中。


\begin{chaptersummary}
待撰写
\end{chaptersummary}

\BookExercises
待撰写

% !TeX root = ../../Book5.tex
% 纲要标题：## 第30章*　正性、曲线锥与双有理程序
% 纲要位置：几何学计划纲要.md，第 6280 行。
% 正文状态：未撰写；本文件目前是框架。

% 纲要内容开始
% 
% > 本章目标与证明范围：先建立截面增长和数值锥的实际工具，再陈述带准确假设的锥、收缩和翻转定理。深层 MMP 输入逐条区分；不把某个终止结果写成任意 MMP 的无条件终止。
% 
% 纲要内容结束

\OptionalChapter{正性、曲线锥与双有理程序}{g5:ch:30}
\BookChapterNumber{g5:ch:30}

\begin{chapterabstract}
待撰写
\end{chapterabstract}

\begin{learninggoals}
待撰写
\end{learninggoals}

待撰写

% !TeX root = ../../Book5.tex
% 纲要标题：### 30.1 除子的不同正性
% 纲要位置：几何学计划纲要.md，第 6284 行。
% 正文状态：未撰写；本文件目前是框架。

% 纲要内容开始
% 
% * 充足、nef、big 与伪有效
% * 数值配对及截面增长
% * Kodaira 引理的所选证明
% * 体积函数和典型例子
% 
% 纲要内容结束

\section{除子的不同正性}
\label{g5:sec:3001}

待撰写

\subsection{充足、nef、big 与伪有效}
\label{g5:subsec:3001:01}
待撰写

\subsection{数值配对及截面增长}
\label{g5:subsec:3001:02}
待撰写

\subsection{Kodaira 引理的所选证明}
\label{g5:subsec:3001:03}
待撰写

\subsection{体积函数和典型例子}
\label{g5:subsec:3001:04}
待撰写
% 纲要条目排为编号小节标题，并进入目录与书签；说明段仍保留于纲要注释中。

% !TeX root = ../../Book5.tex
% 纲要标题：### 30.2 曲线锥与截面环
% 纲要位置：几何学计划纲要.md，第 5997 行。
% 正文状态：未撰写；本文件目前是框架。

% 纲要内容开始
% 
% * N¹、N₁ 及有限维性
% * 锥的对偶及极端射线
% * 典范环和 Kodaira 维数
% * 有限生成与射影模型
% 
% 纲要内容结束

\section{曲线锥与截面环}
\label{g5:sec:3002}

待撰写

\subsection{\texorpdfstring{\(N^1\)、\(N_1\) 及有限维性}{N¹、N₁ 及有限维性}}
\label{g5:subsec:3002:01}
待撰写

\subsection{锥的对偶及极端射线}
\label{g5:subsec:3002:02}
待撰写

\subsection{典范环和 Kodaira 维数}
\label{g5:subsec:3002:03}
待撰写

\subsection{有限生成与射影模型}
\label{g5:subsec:3002:04}
待撰写
% 纲要条目排为编号小节标题，并进入目录与书签；说明段仍保留于纲要注释中。

% !TeX root = ../../Book5.tex
% 纲要标题：### 30.3 锥、收缩及无基点
% 纲要位置：几何学计划纲要.md，第 6298 行。
% 正文状态：未撰写；本文件目前是框架。

% 纲要内容开始
% 
% * 对数对假设的逐项使用
% * 锥定理的准确版本
% * 收缩与基点自由的外部输入
% * 除子收缩和小收缩
% * 曲面及环面模型的证明计算
% 
% 纲要内容结束

\section{锥、收缩及无基点}
\label{g5:sec:3003}

待撰写

\subsection{对数对假设的逐项使用}
\label{g5:subsec:3003:01}
待撰写

\subsection{锥定理的准确版本}
\label{g5:subsec:3003:02}
待撰写

\subsection{收缩与基点自由的外部输入}
\label{g5:subsec:3003:03}
待撰写

\subsection{除子收缩和小收缩}
\label{g5:subsec:3003:04}
待撰写

\subsection{曲面及环面模型的证明计算}
\label{g5:subsec:3003:05}
待撰写
% 纲要条目排为编号小节标题，并进入目录与书签；说明段仍保留于纲要注释中。

% !TeX root = ../../Book5.tex
% 纲要标题：### 30.4 翻转与极小模型
% 纲要位置：几何学计划纲要.md，第 6306 行。
% 正文状态：未撰写；本文件目前是框架。

% 纲要内容开始
% 
% * 翻转的定义及方向
% * 极小模型和 Mori 纤维空间
% * 带缩放的程序及特定终止结论
% * 有限生成和典范模型的适用范围
% * 一般问题与已引用版本的边界
% 
% ---
% 
% 纲要内容结束

\section{翻转与极小模型}
\label{g5:sec:3004}

待撰写

\subsection{翻转的定义及方向}
\label{g5:subsec:3004:01}
待撰写

\subsection{极小模型和 Mori 纤维空间}
\label{g5:subsec:3004:02}
待撰写

\subsection{带缩放的程序及特定终止结论}
\label{g5:subsec:3004:03}
待撰写

\subsection{有限生成和典范模型的适用范围}
\label{g5:subsec:3004:04}
待撰写

\subsection{一般问题与已引用版本的边界}
\label{g5:subsec:3004:05}
待撰写
% 纲要条目排为编号小节标题，并进入目录与书签；说明段仍保留于纲要注释中。


\begin{chaptersummary}
待撰写
\end{chaptersummary}

\BookExercises
待撰写


% !TeX root = ../Book5.tex
% 纲要标题：# 第六编　导出结构、形变、模与比较
% 纲要位置：几何学计划纲要.md，第 6316 行。
% 正文状态：未撰写；本文件目前是框架。

% 纲要内容开始
% 
% 纲要内容结束

\BookPart{导出结构、形变、模与比较}{g5:part:06}
\BookPartGuide
待撰写

% !TeX root = ../../Book5.tex
% 纲要标题：## 第31章*　导出交与余切复形
% 纲要位置：几何学计划纲要.md，第 6024 行。
% 正文状态：未撰写；本文件目前是框架。

% 纲要内容开始
% 
% > 本章目标与证明范围：把形变实际要用的复形语言前置。完整处理有界消解、完美复形和局部完全交的两项余切模型；一般高阶导出几何仍有明确边界。
% 
% 纲要内容结束

\OptionalChapter{导出交与余切复形}{g5:ch:31}
\BookChapterNumber{g5:ch:31}

\begin{chapterabstract}
待撰写
\end{chapterabstract}

\begin{learninggoals}
待撰写
\end{learninggoals}

待撰写

% !TeX root = ../../Book5.tex
% 纲要标题：### 31.1 从非横截交到导出张量
% 纲要位置：几何学计划纲要.md，第 6028 行。
% 正文状态：未撰写；本文件目前是框架。

% 纲要内容开始
% 
% * 局部自由消解和拟同构
% * Tor 层与相交的缺失信息
% * 映射锥及三角的长正合列
% * 有界上方导出范畴及消解计算
% 
% 纲要内容结束

\section{从非横截交到导出张量}
\label{g5:sec:3101}

待撰写

\subsection{局部自由消解和拟同构}
\label{g5:subsec:3101:01}
待撰写

\subsection{Tor 层与相交的缺失信息}
\label{g5:subsec:3101:02}
待撰写

\subsection{映射锥及三角的长正合列}
\label{g5:subsec:3101:03}
待撰写

\subsection{有界上方导出范畴及消解计算}
\label{g5:subsec:3101:04}
待撰写
% 纲要条目排为编号小节标题，并进入目录与书签；说明段仍保留于纲要注释中。

% !TeX root = ../../Book5.tex
% 纲要标题：### 31.2 完美复形与导出运算
% 纲要位置：几何学计划纲要.md，第 6329 行。
% 正文状态：未撰写；本文件目前是框架。

% 纲要内容开始
% 
% * 局部有限自由模型
% * 导出张量和 RHom
% * 投影公式及基变换条件
% * 普通函子和导出函子的差别
% 
% 纲要内容结束

\section{完美复形与导出运算}
\label{g5:sec:3102}

待撰写

\subsection{局部有限自由模型}
\label{g5:subsec:3102:01}
待撰写

\subsection{导出张量和 RHom}
\label{g5:subsec:3102:02}
待撰写

\subsection{投影公式及基变换条件}
\label{g5:subsec:3102:03}
待撰写

\subsection{普通函子和导出函子的差别}
\label{g5:subsec:3102:04}
待撰写
% 纲要条目排为编号小节标题，并进入目录与书签；说明段仍保留于纲要注释中。

% !TeX root = ../../Book5.tex
% 纲要标题：### 31.3 余切复形的具体构造
% 纲要位置：几何学计划纲要.md，第 6042 行。
% 正文状态：未撰写；本文件目前是框架。

% 纲要内容开始
% 
% * 多项式呈示与两项截断
% 
%   从 A→P→B 构造 [I/I²→Ω_{P/A}⊗_P B]；它是余切复形的低阶截断模型，不在非完全交情形直接等同于完整 L_{B/A}。
% 
% * 局部完全交与幅度
% 
%   对由正则序列定义的商证明两项模型计算完整余切复形，并核验其幅度为 [−1,0]。
% 
% * 单纯多项式消解
% 
%   建立自由多项式代数的单纯消解、面与退化映射及所用比较同伦；先说明其与模的自由消解不同。
% 
% * 一般余切复形的构造
% 
%   对消解逐项取微分模并张量到 B，再以正规化复形形成 L_{B/A}；统一链与上链次数。
% 
% * 消解独立性与函子性
% 
%   通过共同消解和同伦比较证明导出对象独立于呈示，不仅核对 H⁰ 相同。
% 
% * 传递三角
% 
%   从 A→B→C 的相容消解构造三角，并验证微分模传递序列是其低阶后果。
% 
% * 导出基变换及普通基变换的条件
% 
%   把导出交换图与普通张量模型分开，注明 Tor 独立等可去掉导出的条件。
% 
% 纲要内容结束

\section{余切复形的具体构造}
\label{g5:sec:3103}

待撰写

\subsection{多项式呈示与两项截断}
\label{g5:subsec:3103:01}
待撰写

\subsection{局部完全交与幅度}
\label{g5:subsec:3103:02}
待撰写

\subsection{单纯多项式消解}
\label{g5:subsec:3103:03}
待撰写

\subsection{一般余切复形的构造}
\label{g5:subsec:3103:04}
待撰写

\subsection{消解独立性与函子性}
\label{g5:subsec:3103:05}
待撰写

\subsection{传递三角}
\label{g5:subsec:3103:06}
待撰写

\subsection{导出基变换及普通基变换的条件}
\label{g5:subsec:3103:07}
待撰写
% 纲要条目排为编号小节标题，并进入目录与书签；说明段仍保留于纲要注释中。

% !TeX root = ../../Book5.tex
% 纲要标题：### 31.4 由复形读出形变信息
% 纲要位置：几何学计划纲要.md，第 6072 行。
% 正文状态：未撰写；本文件目前是框架。

% 纲要内容开始
% 
% * 导子和自同构
% * 一阶扩张与 Ext¹
% * 障碍类及 Ext²
% * 双数、超曲面和非完全交模型
% 
% ---
% 
% 纲要内容结束

\section{由复形读出形变信息}
\label{g5:sec:3104}

待撰写

\subsection{导子和自同构}
\label{g5:subsec:3104:01}
待撰写

\subsection{\texorpdfstring{一阶扩张与 \(\operatorname{Ext}^1\)}{一阶扩张与 Ext¹}}
\label{g5:subsec:3104:02}
待撰写

\subsection{\texorpdfstring{障碍类及 \(\operatorname{Ext}^2\)}{障碍类及 Ext²}}
\label{g5:subsec:3104:03}
待撰写

\subsection{双数、超曲面和非完全交模型}
\label{g5:subsec:3104:04}
待撰写
% 纲要条目排为编号小节标题，并进入目录与书签；说明段仍保留于纲要注释中。


\begin{chaptersummary}
待撰写
\end{chaptersummary}

\BookExercises
待撰写

% !TeX root = ../../Book5.tex
% 纲要标题：## 第32章*　平坦形变与逐阶障碍
% 纲要位置：几何学计划纲要.md，第 6375 行。
% 正文状态：未撰写；本文件目前是框架。

% 纲要内容开始
% 
% > 本章目标与证明范围：以 Artin 基底的提升问题为组织轴。对嵌入、光滑对象和向量丛分别建立一阶分类及障碍；形式解、收敛解和实际代数族不能相互替代。
% 
% 纲要内容结束

\OptionalChapter{平坦形变与逐阶障碍}{g5:ch:32}
\BookChapterNumber{g5:ch:32}

\begin{chapterabstract}
待撰写
\end{chapterabstract}

\begin{learninggoals}
待撰写
\end{learninggoals}

待撰写

% !TeX root = ../../Book5.tex
% 纲要标题：### 32.1 形变函子及等价
% 纲要位置：几何学计划纲要.md，第 6379 行。
% 正文状态：未撰写；本文件目前是框架。

% 纲要内容开始
% 
% * 局部 Artin 环和小扩张
% * 指定特殊纤维的平坦族
% * 同构、标记及自同构
% * 切空间的加法与标量结构
% 
% 纲要内容结束

\section{形变函子及等价}
\label{g5:sec:3201}

待撰写

\subsection{局部 Artin 环和小扩张}
\label{g5:subsec:3201:01}
待撰写

\subsection{指定特殊纤维的平坦族}
\label{g5:subsec:3201:02}
待撰写

\subsection{同构、标记及自同构}
\label{g5:subsec:3201:03}
待撰写

\subsection{切空间的加法与标量结构}
\label{g5:subsec:3201:04}
待撰写
% 纲要条目排为编号小节标题，并进入目录与书签；说明段仍保留于纲要注释中。

% !TeX root = ../../Book5.tex
% 纲要标题：### 32.2 嵌入和丛的一阶形变
% 纲要位置：几何学计划纲要.md，第 6092 行。
% 正文状态：未撰写；本文件目前是框架。

% 纲要内容开始
% 
% * 理想层的 Hom 描述
% * 局部完全交与法层
% * 向量丛的 Ext¹
% * 平凡形变和坐标变换
% * 结点平滑化的显式方程
% 
% 纲要内容结束

\section{嵌入和丛的一阶形变}
\label{g5:sec:3202}

待撰写

\subsection{理想层的 Hom 描述}
\label{g5:subsec:3202:01}
待撰写

\subsection{局部完全交与法层}
\label{g5:subsec:3202:02}
待撰写

\subsection{\texorpdfstring{向量丛的 \(\operatorname{Ext}^1\)}{向量丛的 Ext¹}}
\label{g5:subsec:3202:03}
待撰写

\subsection{平凡形变和坐标变换}
\label{g5:subsec:3202:04}
待撰写

\subsection{结点平滑化的显式方程}
\label{g5:subsec:3202:05}
待撰写
% 纲要条目排为编号小节标题，并进入目录与书签；说明段仍保留于纲要注释中。

% !TeX root = ../../Book5.tex
% 纲要标题：### 32.3 障碍的构造
% 纲要位置：几何学计划纲要.md，第 6394 行。
% 正文状态：未撰写；本文件目前是框架。

% 纲要内容开始
% 
% * 小扩张上的局部提升
% * 双重、三重交上的差异
% * 具体障碍余圈
% * 障碍消失和解的扭子结构
% * 逐阶不唯一性
% 
% 纲要内容结束

\section{障碍的构造}
\label{g5:sec:3203}

待撰写

\subsection{小扩张上的局部提升}
\label{g5:subsec:3203:01}
待撰写

\subsection{双重、三重交上的差异}
\label{g5:subsec:3203:02}
待撰写

\subsection{具体障碍余圈}
\label{g5:subsec:3203:03}
待撰写

\subsection{障碍消失和解的扭子结构}
\label{g5:subsec:3203:04}
待撰写

\subsection{逐阶不唯一性}
\label{g5:subsec:3203:05}
待撰写
% 纲要条目排为编号小节标题，并进入目录与书签；说明段仍保留于纲要注释中。

% !TeX root = ../../Book5.tex
% 纲要标题：### 32.4 余切复形的统一解释
% 纲要位置：几何学计划纲要.md，第 6108 行。
% 正文状态：未撰写；本文件目前是框架。

% 纲要内容开始
% 
% * Ext 的自同构、切向及障碍次数
% * 局部到整体谱序列
% * 光滑对象的 H¹(T) 与 H²(T)
% * 半通用形式形变及 Schlessinger 条件的所选范围
% 
% 纲要内容结束

\section{余切复形的统一解释}
\label{g5:sec:3204}

待撰写

\subsection{Ext 的自同构、切向及障碍次数}
\label{g5:subsec:3204:01}
待撰写

\subsection{局部到整体谱序列}
\label{g5:subsec:3204:02}
待撰写

\subsection{\texorpdfstring{光滑对象的 \(H^1(\mathrm{T})\) 与 \(H^2(\mathrm{T})\)}{光滑对象的 H¹(T) 与 H²(T)}}
\label{g5:subsec:3204:03}
待撰写

\subsection{半通用形式形变及 Schlessinger 条件的所选范围}
\label{g5:subsec:3204:04}
待撰写
% 纲要条目排为编号小节标题，并进入目录与书签；说明段仍保留于纲要注释中。

% !TeX root = ../../Book5.tex
% 纲要标题：### 32.5 形式与实际模问题
% 纲要位置：几何学计划纲要.md，第 6409 行。
% 正文状态：未撰写；本文件目前是框架。

% 纲要内容开始
% 
% * 形式有效性及代数化
% * Artin 逼近的条件和用途
% * 自同构造成的可表性障碍
% * 有限阶计算不能直接推出模空间光滑
% 
% ---
% 
% 纲要内容结束

\section{形式与实际模问题}
\label{g5:sec:3205}

待撰写

\subsection{形式有效性及代数化}
\label{g5:subsec:3205:01}
待撰写

\subsection{Artin 逼近的条件和用途}
\label{g5:subsec:3205:02}
待撰写

\subsection{自同构造成的可表性障碍}
\label{g5:subsec:3205:03}
待撰写

\subsection{有限阶计算不能直接推出模空间光滑}
\label{g5:subsec:3205:04}
待撰写
% 纲要条目排为编号小节标题，并进入目录与书签；说明段仍保留于纲要注释中。


\begin{chaptersummary}
待撰写
\end{chaptersummary}

\BookExercises
待撰写

% !TeX root = ../../Book5.tex
% 纲要标题：## 第33章*　Hilbert、Picard、Quot 与通用族
% 纲要位置：几何学计划纲要.md，第 6124 行。
% 正文状态：未撰写；本文件目前是框架。

% 纲要内容开始
% 
% > 本章目标与证明范围：将两种可表性问题及其实际构造分开。Hilbert 的有限界和一般相对 Picard 存在性明确标记深层输入；有限维参数空间、方程及切空间计算仍完整展开。
% 
% 纲要内容结束

\OptionalChapter{Hilbert、Picard、Quot 与通用族}{g5:ch:33}
\BookChapterNumber{g5:ch:33}

\begin{chapterabstract}
待撰写
\end{chapterabstract}

\begin{learninggoals}
待撰写
\end{learninggoals}

待撰写

% !TeX root = ../../Book5.tex
% 纲要标题：### 33.1 模函子与代表对象
% 纲要位置：几何学计划纲要.md，第 6422 行。
% 正文状态：未撰写；本文件目前是框架。

% 纲要内容开始
% 
% * 通用族及 Yoneda
% * 细模与粗模
% * 下降和层化为何必要
% * 基底改变与对象的等价
% 
% 纲要内容结束

\section{模函子与代表对象}
\label{g5:sec:3301}

待撰写

\subsection{通用族及 Yoneda}
\label{g5:subsec:3301:01}
待撰写

\subsection{细模与粗模}
\label{g5:subsec:3301:02}
待撰写

\subsection{下降和层化为何必要}
\label{g5:subsec:3301:03}
待撰写

\subsection{基底改变与对象的等价}
\label{g5:subsec:3301:04}
待撰写
% 纲要条目排为编号小节标题，并进入目录与书签；说明段仍保留于纲要注释中。

% !TeX root = ../../Book5.tex
% 纲要标题：### 33.2 Hilbert 构造的有限维化
% 纲要位置：几何学计划纲要.md，第 6429 行。
% 正文状态：未撰写；本文件目前是框架。

% 纲要内容开始
% 
% * 固定 Hilbert 多项式
% * 正则性界的选定外部定理
% * Grassmann 参数空间
% * 乘法相容给出的闭条件
% * 通用商及平坦性的核验
% 
% 纲要内容结束

\section{Hilbert 构造的有限维化}
\label{g5:sec:3302}

待撰写

\subsection{固定 Hilbert 多项式}
\label{g5:subsec:3302:01}
待撰写

\subsection{正则性界的选定外部定理}
\label{g5:subsec:3302:02}
待撰写

\subsection{Grassmann 参数空间}
\label{g5:subsec:3302:03}
待撰写

\subsection{乘法相容给出的闭条件}
\label{g5:subsec:3302:04}
待撰写

\subsection{通用商及平坦性的核验}
\label{g5:subsec:3302:05}
待撰写
% 纲要条目排为编号小节标题，并进入目录与书签；说明段仍保留于纲要注释中。

% !TeX root = ../../Book5.tex
% 纲要标题：### 33.3 Hilbert 概形的局部几何
% 纲要位置：几何学计划纲要.md，第 6437 行。
% 正文状态：未撰写；本文件目前是框架。

% 纲要内容开始
% 
% * 切空间 Hom(I,O_Z)
% * 嵌入障碍及 Ext
% * 点、线和完全交的计算
% * Hilbᵈ(A¹) 的完整模型
% * 不可约分支及奇异点
% 
% 纲要内容结束

\section{Hilbert 概形的局部几何}
\label{g5:sec:3303}

待撰写

\subsection{\texorpdfstring{切空间 \(\operatorname{Hom}(I,\mathcal O_Z)\)}{切空间 Hom(I,O\_Z)}}
\label{g5:subsec:3303:01}
待撰写

\subsection{嵌入障碍及 Ext}
\label{g5:subsec:3303:02}
待撰写

\subsection{点、线和完全交的计算}
\label{g5:subsec:3303:03}
待撰写

\subsection{\texorpdfstring{\(\operatorname{Hilb}^d(\mathbb A^1)\) 的完整模型}{Hilbᵈ(A¹) 的完整模型}}
\label{g5:subsec:3303:04}
待撰写

\subsection{不可约分支及奇异点}
\label{g5:subsec:3303:05}
待撰写
% 纲要条目排为编号小节标题，并进入目录与书签；说明段仍保留于纲要注释中。

% !TeX root = ../../Book5.tex
% 纲要标题：### 33.4 Picard 函子、层化与代表对象
% 纲要位置：几何学计划纲要.md，第 6445 行。
% 正文状态：未撰写；本文件目前是框架。

% 纲要内容开始
% 
% * 相对线丛与来自基底的线丛之商
% * fppf 层化及局部线丛的下降障碍
% * 刚化、基点与 Poincaré 线丛
% * 一般可表性主版本及其代数空间结论
% 
%   明确外引如下版本：$f:X\to S$ 平坦、固有、有限呈示，且任意基变换 $T\to S$ 后 $\mathcal O_T\to(f_T)_*\mathcal O_{X_T}$ 为同构，则相对 Picard fppf 层由代数空间表示。是否为概形及是否存在未刚化通用线丛，需要分别给出附加条件。[R12]
% 
% * 代表空间、Picard 叠与通用线丛存在性的区别
% 
% 纲要内容结束

\section{Picard 函子、层化与代表对象}
\label{g5:sec:3304}

待撰写

\subsection{相对线丛与来自基底的线丛之商}
\label{g5:subsec:3304:01}
待撰写

\subsection{fppf 层化及局部线丛的下降障碍}
\label{g5:subsec:3304:02}
待撰写

\subsection{刚化、基点与 Poincaré 线丛}
\label{g5:subsec:3304:03}
待撰写

\subsection{一般可表性主版本及其代数空间结论}
\label{g5:subsec:3304:04}
待撰写

\subsection{代表空间、Picard 叠与通用线丛存在性的区别}
\label{g5:subsec:3304:05}
待撰写
% 纲要条目排为编号小节标题，并进入目录与书签；说明段仍保留于纲要注释中。

% !TeX root = ../../Book5.tex
% 纲要标题：### 33.5 曲线 Picard 与 Jacobi 簇
% 纲要位置：几何学计划纲要.md，第 6456 行。
% 正文状态：未撰写；本文件目前是框架。

% 纲要内容开始
% 
% * 高次对称积及除子类
% * 有基点曲线的通用线丛
% * Pic⁰ 的群结构与光滑性
% * 与周期构造的比较
% * 没有基点时的下降障碍
% 
% 纲要内容结束

\section{曲线 Picard 与 Jacobi 簇}
\label{g5:sec:3305}

待撰写

\subsection{高次对称积及除子类}
\label{g5:subsec:3305:01}
待撰写

\subsection{有基点曲线的通用线丛}
\label{g5:subsec:3305:02}
待撰写

\subsection{\texorpdfstring{\(\operatorname{Pic}^0\) 的群结构与光滑性}{Pic⁰ 的群结构与光滑性}}
\label{g5:subsec:3305:03}
待撰写

\subsection{与周期构造的比较}
\label{g5:subsec:3305:04}
待撰写

\subsection{没有基点时的下降障碍}
\label{g5:subsec:3305:05}
待撰写
% 纲要条目排为编号小节标题，并进入目录与书签；说明段仍保留于纲要注释中。

% !TeX root = ../../Book5.tex
% 纲要标题：### 33.6 Quot 函子与丛模问题的参数入口
% 纲要位置：几何学计划纲要.md，第 6170 行。
% 正文状态：未撰写；本文件目前是框架。

% 纲要内容开始
% 
% 本节给出参数语言与具体模型；一般 Quot 构造及丛模空间存在性不被压缩成一个未经证明的“标准构造”。深层可表性与有界性分别登记外部输入。
% 
% * 固定凝聚层的有限呈示商与商的等价关系
% * 商层对基底平坦及固定 Hilbert 多项式
% * Hilbert 函子作为结构层商的特例
% * 射影 Noether 范围内 Quot 可表性的准确外部输入
% * 点上的 Grassmann 商与射影线上的具体商层
% * 有界性、统一扭转及半稳定丛的参数化
% 
% ---
% 
% 纲要内容结束

\section{Quot 函子与丛模问题的参数入口}
\label{g5:sec:3306}

待撰写

\subsection{固定凝聚层的有限呈示商与商的等价关系}
\label{g5:subsec:3306:01}
待撰写

\subsection{商层对基底平坦及固定 Hilbert 多项式}
\label{g5:subsec:3306:02}
待撰写

\subsection{Hilbert 函子作为结构层商的特例}
\label{g5:subsec:3306:03}
待撰写

\subsection{射影 Noether 范围内 Quot 可表性的准确外部输入}
\label{g5:subsec:3306:04}
待撰写

\subsection{点上的 Grassmann 商与射影线上的具体商层}
\label{g5:subsec:3306:05}
待撰写

\subsection{有界性、统一扭转及半稳定丛的参数化}
\label{g5:subsec:3306:06}
待撰写
% 纲要条目排为编号小节标题，并进入目录与书签；说明段仍保留于纲要注释中。


\begin{chaptersummary}
待撰写
\end{chaptersummary}

\BookExercises
待撰写

% !TeX root = ../../Book5.tex
% 纲要标题：## 第34章*　几何不变量理论、稳定性与模空间
% 纲要位置：几何学计划纲要.md，第 6477 行。
% 正文状态：未撰写；本文件目前是框架。

% 纲要内容开始
% 
% > 本章目标与证明范围：先完成可计算商，再讨论稳定和模对象。线性约化、约化及基域特征的差异说明清楚；叠通过下降和稳定子进入，不作为“更高级的空间”一句带过。
% 
% 纲要内容结束

\OptionalChapter{几何不变量理论、稳定性与模空间}{g5:ch:34}
\BookChapterNumber{g5:ch:34}

\begin{chapterabstract}
待撰写
\end{chapterabstract}

\begin{learninggoals}
待撰写
\end{learninggoals}

待撰写

% !TeX root = ../../Book5.tex
% 纲要标题：### 34.1 仿射不变量与商
% 纲要位置：几何学计划纲要.md，第 6187 行。
% 正文状态：未撰写；本文件目前是框架。

% 纲要内容开始
% 
% * 有限群轨道多项式
% * Reynolds 算子及线性约化假设
% * Hilbert 有限生成的证明模块
% * 仿射商的泛性质
% * 闭轨道和一般轨道的区别
% 
% 纲要内容结束

\section{仿射不变量与商}
\label{g5:sec:3401}

待撰写

\subsection{有限群轨道多项式}
\label{g5:subsec:3401:01}
待撰写

\subsection{Reynolds 算子及线性约化假设}
\label{g5:subsec:3401:02}
待撰写

\subsection{Hilbert 有限生成的证明模块}
\label{g5:subsec:3401:03}
待撰写

\subsection{仿射商的泛性质}
\label{g5:subsec:3401:04}
待撰写

\subsection{闭轨道和一般轨道的区别}
\label{g5:subsec:3401:05}
待撰写
% 纲要条目排为编号小节标题，并进入目录与书签；说明段仍保留于纲要注释中。

% !TeX root = ../../Book5.tex
% 纲要标题：### 34.2 线性化与射影商
% 纲要位置：几何学计划纲要.md，第 6489 行。
% 正文状态：未撰写；本文件目前是框架。

% 纲要内容开始
% 
% * 线丛线性化
% * 不变量截面和半稳定点
% * 射影商的构造
% * 稳定点与几何商
% * 改变线性化的墙穿越模型
% 
% 纲要内容结束

\section{线性化与射影商}
\label{g5:sec:3402}

待撰写

\subsection{线丛线性化}
\label{g5:subsec:3402:01}
待撰写

\subsection{不变量截面和半稳定点}
\label{g5:subsec:3402:02}
待撰写

\subsection{射影商的构造}
\label{g5:subsec:3402:03}
待撰写

\subsection{稳定点与几何商}
\label{g5:subsec:3402:04}
待撰写

\subsection{改变线性化的墙穿越模型}
\label{g5:subsec:3402:05}
待撰写
% 纲要条目排为编号小节标题，并进入目录与书签；说明段仍保留于纲要注释中。

% !TeX root = ../../Book5.tex
% 纲要标题：### 34.3 数值稳定性
% 纲要位置：几何学计划纲要.md，第 6203 行。
% 正文状态：未撰写；本文件目前是框架。

% 纲要内容开始
% 
% * 一参数子群和权
% * Hilbert–Mumford 判据的所选版本
% * 对角环面作用的完整计算
% * 点构型和向量丛的不同稳定性
% 
% 纲要内容结束

\section{数值稳定性}
\label{g5:sec:3403}

待撰写

\subsection{一参数子群和权}
\label{g5:subsec:3403:01}
待撰写

\subsection{Hilbert–Mumford 判据的所选版本}
\label{g5:subsec:3403:02}
待撰写

\subsection{对角环面作用的完整计算}
\label{g5:subsec:3403:03}
待撰写

\subsection{点构型和向量丛的不同稳定性}
\label{g5:subsec:3403:04}
待撰写
% 纲要条目排为编号小节标题，并进入目录与书签；说明段仍保留于纲要注释中。

% !TeX root = ../../Book5.tex
% 纲要标题：### 34.4 实际模空间
% 纲要位置：几何学计划纲要.md，第 6504 行。
% 正文状态：未撰写；本文件目前是框架。

% 纲要内容开始
% 
% * 射影直线上的点构型
% * 稳定曲线的参数与商
% * 半稳定丛的参数、S-等价与粗模问题
% 
%   用 V.33.6 的 Quot 语言说明参数空间及 GIT 商；统一有界性、正则性界和模空间存在分别准确外引。一般半稳定对象按 S-等价而非逐个同构类得到粗参数，不宣称本章已完整构造所有丛模空间。
% 
% * 有界性、分离性和固有性分别需要的定理
% 
% 纲要内容结束

\section{实际模空间}
\label{g5:sec:3404}

待撰写

\subsection{射影直线上的点构型}
\label{g5:subsec:3404:01}
待撰写

\subsection{稳定曲线的参数与商}
\label{g5:subsec:3404:02}
待撰写

\subsection{半稳定丛的参数、S-等价与粗模问题}
\label{g5:subsec:3404:03}
待撰写

\subsection{有界性、分离性和固有性分别需要的定理}
\label{g5:subsec:3404:04}
待撰写
% 纲要条目排为编号小节标题，并进入目录与书签；说明段仍保留于纲要注释中。

% !TeX root = ../../Book5.tex
% 纲要标题：### 34.5 自同构和叠
% 纲要位置：几何学计划纲要.md，第 6514 行。
% 正文状态：未撰写；本文件目前是框架。

% 纲要内容开始
% 
% * 纤维范畴与下降数据
% * 商叠 [X/G] 的对象和态射
% * 粗模空间与稳定子信息
% * Deligne–Mumford 与 Artin 叠的图册条件
% 
%   对角的代数空间可表性与图册条件分别陈述：Artin 叠使用光滑满射概形图册，Deligne–Mumford 叠使用平展满射概形图册。稳定子的有限性、非分歧性和群概形的非约化现象分别检查，尤其不能在正特征只凭“有限自同构”推出全部 Deligne–Mumford 条件。
% 
% * 一般存在定理的边界
% 
% 纲要内容结束

\section{自同构和叠}
\label{g5:sec:3405}

待撰写

\subsection{纤维范畴与下降数据}
\label{g5:subsec:3405:01}
待撰写

\subsection{商叠 [X/G] 的对象和态射}
\label{g5:subsec:3405:02}
待撰写

\subsection{粗模空间与稳定子信息}
\label{g5:subsec:3405:03}
待撰写

\subsection{Deligne–Mumford 与 Artin 叠的图册条件}
\label{g5:subsec:3405:04}
待撰写

\subsection{一般存在定理的边界}
\label{g5:subsec:3405:05}
待撰写
% 纲要条目排为编号小节标题，并进入目录与书签；说明段仍保留于纲要注释中。

% !TeX root = ../../Book5.tex
% 纲要标题：### 34.6 四点标记有理曲线的稳定紧化
% 纲要位置：几何学计划纲要.md，第 6525 行。
% 正文状态：未撰写；本文件目前是框架。

% 纲要内容开始
% 
% 将一个紧化问题做完：主模型可在任意域上构造；复解析参数与长颈解释见 IV.6.5。先完成这个低维代表对象，再把一般稳定曲线模空间作为具有额外存在定理的研究出口。
% 
% * 有序标记、同构及交比归一化
% * $\mathcal M_{0,4}\simeq\mathbb P^1\setminus\{0,1,\infty\}$ 与光滑通用族
% * 四个截面的碰撞与三个吹起中心
% 
%   从 $\mathbb P^1_z\times\mathbb P^1_\lambda$ 及截面 $z=0,1,\infty,\lambda$ 出发，在 $(0,0)$、$(1,1)$、$(\infty,\infty)$ 吹起，取严格变换后的标记截面。
% 
% * 结点纤维、每个分支上的特殊点和稳定性
% * 三个边界点与标记集的三种二二分割
% * 平坦性、截面不交及稳定族的局部核验
% * 交比的族版本与代表对象的通用性
% 
%   除几何点分类外，还需在非约化测试基底上构造拉回对应，并验证局部平凡化的下降；不能仅凭交比给出几何点双射就宣布细模可表。
% 
% * $\overline{\mathcal M}_{0,4}\simeq\mathbb P^1$ 的固有性和粗细模辨别
% 
%   光滑四点有理曲线有四个标记点；边界纤维的每个分支恰有三个特殊点，逐支检查自同构平凡；置换标记后的商问题另行处理。
% 
% * 标记置换、商构造与更高亏格所需的新输入
% 
% ---
% 
% 纲要内容结束

\section{四点标记有理曲线的稳定紧化}
\label{g5:sec:3406}

待撰写

\subsection{有序标记、同构及交比归一化}
\label{g5:subsec:3406:01}
待撰写

\subsection{\texorpdfstring{\(\mathcal M_{0,4}\simeq\mathbb P^1\setminus\{0,1,\infty\}\) 与光滑通用族}{M\_\{0,4\}≃P¹∖\{0,1,∞\} 与光滑通用族}}
\label{g5:subsec:3406:02}
待撰写

\subsection{四个截面的碰撞与三个吹起中心}
\label{g5:subsec:3406:03}
待撰写

\subsection{结点纤维、每个分支上的特殊点和稳定性}
\label{g5:subsec:3406:04}
待撰写

\subsection{三个边界点与标记集的三种二二分割}
\label{g5:subsec:3406:05}
待撰写

\subsection{平坦性、截面不交及稳定族的局部核验}
\label{g5:subsec:3406:06}
待撰写

\subsection{交比的族版本与代表对象的通用性}
\label{g5:subsec:3406:07}
待撰写

\subsection{\texorpdfstring{\(\overline{\mathcal M}_{0,4}\simeq\mathbb P^1\) 的固有性和粗细模辨别}{\={M}\_\{0,4\}≃P¹ 的固有性和粗细模辨别}}
\label{g5:subsec:3406:08}
待撰写

\subsection{标记置换、商构造与更高亏格所需的新输入}
\label{g5:subsec:3406:09}
待撰写
% 纲要条目排为编号小节标题，并进入目录与书签；说明段仍保留于纲要注释中。


\begin{chaptersummary}
待撰写
\end{chaptersummary}

\BookExercises
待撰写

% !TeX root = ../../Book5.tex
% 纲要标题：## 第35章*　平展拓扑、基本群与上同调
% 纲要位置：几何学计划纲要.md，第 6550 行。
% 正文状态：未撰写；本文件目前是框架。

% 纲要内容开始
% 
% > 本章目标与证明范围：有限平展覆盖的 Galois 理论与平展上同调分阶段建立。Kummer 的系数可逆条件和复数域比较的对象范围写清楚；不把经典覆盖的全部性质无条件移植。
% 
% 纲要内容结束

\OptionalChapter{平展拓扑、基本群与上同调}{g5:ch:35}
\BookChapterNumber{g5:ch:35}

\begin{chapterabstract}
待撰写
\end{chapterabstract}

\begin{learninggoals}
待撰写
\end{learninggoals}

待撰写

% !TeX root = ../../Book5.tex
% 纲要标题：### 35.1 有限平展覆盖和纤维函子
% 纲要位置：几何学计划纲要.md，第 6260 行。
% 正文状态：未撰写；本文件目前是框架。

% 纲要内容开始
% 
% * 几何基点及纤维
% * 覆盖的分解与 Galois 闭包
% * 投射有限群的重建
% * 基本群正合列的条件
% * 有限域、环面及曲线模型
% 
% 纲要内容结束

\section{有限平展覆盖和纤维函子}
\label{g5:sec:3501}

待撰写

\subsection{几何基点及纤维}
\label{g5:subsec:3501:01}
待撰写

\subsection{覆盖的分解与 Galois 闭包}
\label{g5:subsec:3501:02}
待撰写

\subsection{投射有限群的重建}
\label{g5:subsec:3501:03}
待撰写

\subsection{基本群正合列的条件}
\label{g5:subsec:3501:04}
待撰写

\subsection{有限域、环面及曲线模型}
\label{g5:subsec:3501:05}
待撰写
% 纲要条目排为编号小节标题，并进入目录与书签；说明段仍保留于纲要注释中。

% !TeX root = ../../Book5.tex
% 纲要标题：### 35.2 平展层与上同调
% 纲要位置：几何学计划纲要.md，第 6562 行。
% 正文状态：未撰写；本文件目前是框架。

% 纲要内容开始
% 
% * 覆盖的局部截面
% * 常值层和局部常值层
% * 导出截面及低阶解释
% * 挠系数和有限性条件
% 
% 纲要内容结束

\section{平展层与上同调}
\label{g5:sec:3502}

待撰写

\subsection{覆盖的局部截面}
\label{g5:subsec:3502:01}
待撰写

\subsection{常值层和局部常值层}
\label{g5:subsec:3502:02}
待撰写

\subsection{导出截面及低阶解释}
\label{g5:subsec:3502:03}
待撰写

\subsection{挠系数和有限性条件}
\label{g5:subsec:3502:04}
待撰写
% 纲要条目排为编号小节标题，并进入目录与书签；说明段仍保留于纲要注释中。

% !TeX root = ../../Book5.tex
% 纲要标题：### 35.3 Kummer 与线丛
% 纲要位置：几何学计划纲要.md，第 6569 行。
% 正文状态：未撰写；本文件目前是框架。

% 纲要内容开始
% 
% * n 可逆时的根提升
% * Kummer 正合列
% * Pic/n、μₙ 上同调及 Brauer 的接口
% * 特征整除 n 时的 fppf 区别
% 
% 纲要内容结束

\section{Kummer 与线丛}
\label{g5:sec:3503}

待撰写

\subsection{n 可逆时的根提升}
\label{g5:subsec:3503:01}
待撰写

\subsection{Kummer 正合列}
\label{g5:subsec:3503:02}
待撰写

\subsection{\texorpdfstring{Pic/n、\(\mu_n\) 上同调及 Brauer 的接口}{Pic/n、\GeometryPdfMu{}ₙ 上同调及 Brauer 的接口}}
\label{g5:subsec:3503:03}
待撰写

\subsection{特征整除 n 时的 fppf 区别}
\label{g5:subsec:3503:04}
待撰写
% 纲要条目排为编号小节标题，并进入目录与书签；说明段仍保留于纲要注释中。

% !TeX root = ../../Book5.tex
% 纲要标题：### 35.4 局部和经典比较
% 纲要位置：几何学计划纲要.md，第 6282 行。
% 正文状态：未撰写；本文件目前是框架。

% 纲要内容开始
% 
% * 严格 Hensel 邻域
% * 形式邻域的不同信息
% * 复代数簇的有限覆盖比较
% * 所选挠系数上同调比较
% * 一般 ℓ-adic 理论所需的逆系统
% 
% ---
% 
% 纲要内容结束

\section{局部和经典比较}
\label{g5:sec:3504}

待撰写

\subsection{严格 Hensel 邻域}
\label{g5:subsec:3504:01}
待撰写

\subsection{形式邻域的不同信息}
\label{g5:subsec:3504:02}
待撰写

\subsection{复代数簇的有限覆盖比较}
\label{g5:subsec:3504:03}
待撰写

\subsection{所选挠系数上同调比较}
\label{g5:subsec:3504:04}
待撰写

\subsection{\texorpdfstring{一般 \(\ell\)-adic 理论所需的逆系统}{一般 \GeometryPdfEll{}-adic 理论所需的逆系统}}
\label{g5:subsec:3504:05}
待撰写
% 纲要条目排为编号小节标题，并进入目录与书签；说明段仍保留于纲要注释中。


\begin{chaptersummary}
待撰写
\end{chaptersummary}

\BookExercises
待撰写

% !TeX root = ../../Book5.tex
% 纲要标题：## 第36章*　解析比较、旗簇与几何表示
% 纲要位置：几何学计划纲要.md，第 6594 行。
% 正文状态：未撰写；本文件目前是框架。

% 纲要内容开始
% 
% > 本章目标与证明范围：保留原有比较和表示主题，但分别给明确模型。完整 D-模和一般导出几何不作为本卷的隐含承诺；本章通过已经建好的丛、上同调和导出函子进入它们。
% 
% 纲要内容结束

\OptionalChapter{解析比较、旗簇与几何表示}{g5:ch:36}
\BookChapterNumber{g5:ch:36}

\begin{chapterabstract}
待撰写
\end{chapterabstract}

\begin{learninggoals}
待撰写
\end{learninggoals}

待撰写

% !TeX root = ../../Book5.tex
% 纲要标题：### 36.1 解析化与代数化
% 纲要位置：几何学计划纲要.md，第 6598 行。
% 正文状态：未撰写；本文件目前是框架。

% 纲要内容开始
% 
% * 复代数局部环及收敛局部环
% * 射影 GAGA 的对象和使用
% * 代数、解析和拓扑上同调的区别
% * 非固有和非代数化模型
% 
% 纲要内容结束

\section{解析化与代数化}
\label{g5:sec:3601}

待撰写

\subsection{复代数局部环及收敛局部环}
\label{g5:subsec:3601:01}
待撰写

\subsection{射影 GAGA 的对象和使用}
\label{g5:subsec:3601:02}
待撰写

\subsection{代数、解析和拓扑上同调的区别}
\label{g5:subsec:3601:03}
待撰写

\subsection{非固有和非代数化模型}
\label{g5:subsec:3601:04}
待撰写
% 纲要条目排为编号小节标题，并进入目录与书签；说明段仍保留于纲要注释中。

% !TeX root = ../../Book5.tex
% 纲要标题：### 36.2 旗簇与线丛表示
% 纲要位置：几何学计划纲要.md，第 6303 行。
% 正文状态：未撰写；本文件目前是框架。

% 纲要内容开始
% 
% * Grassmann 及旗簇的坐标模型
% * 权和齐性线丛
% * P¹ 上的截面表示
% * Borel–Weil–Bott 的准确陈述
% * 消失、取向及权规范的核验
% 
% 纲要内容结束

\section{旗簇与线丛表示}
\label{g5:sec:3602}

待撰写

\subsection{Grassmann 及旗簇的坐标模型}
\label{g5:subsec:3602:01}
待撰写

\subsection{权和齐性线丛}
\label{g5:subsec:3602:02}
待撰写

\subsection{\texorpdfstring{\(\mathbb P^1\) 上的截面表示}{P¹ 上的截面表示}}
\label{g5:subsec:3602:03}
待撰写

\subsection{Borel–Weil–Bott 的准确陈述}
\label{g5:subsec:3602:04}
待撰写

\subsection{消失、取向及权规范的核验}
\label{g5:subsec:3602:05}
待撰写
% 纲要条目排为编号小节标题，并进入目录与书签；说明段仍保留于纲要注释中。

% !TeX root = ../../Book5.tex
% 纲要标题：### 36.3 D-模与解层
% 纲要位置：几何学计划纲要.md，第 6613 行。
% 正文状态：未撰写；本文件目前是框架。

% 纲要内容开始
% 
% * 微分算子环和滤过
% * 联络作为 D-模
% * 特征簇及完整的低维算例
% * 正则、holonomic 与全纯的不同术语
% * Riemann–Hilbert 的外部条件
% 
% 纲要内容结束

\section{D-模与解层}
\label{g5:sec:3603}

待撰写

\subsection{微分算子环和滤过}
\label{g5:subsec:3603:01}
待撰写

\subsection{联络作为 D-模}
\label{g5:subsec:3603:02}
待撰写

\subsection{特征簇及完整的低维算例}
\label{g5:subsec:3603:03}
待撰写

\subsection{正则、holonomic 与全纯的不同术语}
\label{g5:subsec:3603:04}
待撰写

\subsection{Riemann–Hilbert 的外部条件}
\label{g5:subsec:3603:05}
待撰写
% 纲要条目排为编号小节标题，并进入目录与书签；说明段仍保留于纲要注释中。

% !TeX root = ../../Book5.tex
% 纲要标题：### 36.4 不同几何化的接口
% 纲要位置：几何学计划纲要.md，第 6621 行。
% 正文状态：未撰写；本文件目前是框架。

% 纲要内容开始
% 
% * 导出推拉与对应
% * 卷积和表示的几何模型
% * 增强、叠及高阶结构的额外输入
% * 不以相似术语推断范畴等价
% 
% ---
% 
% 纲要内容结束

\section{不同几何化的接口}
\label{g5:sec:3604}

待撰写

\subsection{导出推拉与对应}
\label{g5:subsec:3604:01}
待撰写

\subsection{卷积和表示的几何模型}
\label{g5:subsec:3604:02}
待撰写

\subsection{增强、叠及高阶结构的额外输入}
\label{g5:subsec:3604:03}
待撰写

\subsection{不以相似术语推断范畴等价}
\label{g5:subsec:3604:04}
待撰写
% 纲要条目排为编号小节标题，并进入目录与书签；说明段仍保留于纲要注释中。


\begin{chaptersummary}
待撰写
\end{chaptersummary}

\BookExercises
待撰写


% !TeX root = ../Book5.tex
% 纲要标题：# 第五卷附录
% 纲要位置：几何学计划纲要.md，第 6328 行。
% 正文状态：未撰写；本文件目前是框架。

% 纲要内容开始
% 
% 纲要内容结束

\appendix
\renewcommand{\thechapter}{\Alph{chapter}}
\ctexset{chapter={number=\Alph{chapter}}}
\BookAppendixPart{g5:part:appendices}

% !TeX root = ../../Book5.tex
% 纲要标题：## 附录 A　算术几何与非 Archimedes 几何
% 纲要位置：几何学计划纲要.md，第 6330 行。
% 正文状态：未撰写；本文件目前是框架。

% 纲要内容开始
% 
% > 与数论系列形成接口，不重建完整类域论或 p 进 Hodge 理论。
% 
% * **整数谱上的几何**
%   * 素点、一般点及纤维
%   * 约化和整模型
%   * 好约化与坏约化的曲线算例
% 
% * **形式与非 Archimedes 空间**
%   * 形式概形和理想定义
%   * Tate 代数及 affinoid 模型
%   * 刚性与 Berkovich 的不同对象
%   * 一般纤维和特殊纤维
% 
% * **有理点与局部—整体**
%   * 点的下降和模解释
%   * 圆锥曲线的局部模型
%   * 高度及有限性的外部输入
% 
% ---
% 
% 纲要内容结束

\chapter{算术几何与非 Archimedes 几何}
\label{g5:app:A}
\BookChapterNumber{g5:app:A}

待撰写

\section{整数谱上的几何}
\label{g5:sec:A01}
待撰写

\subsection{素点、一般点及纤维}
\label{g5:subsec:A01:01}
待撰写

\subsection{约化和整模型}
\label{g5:subsec:A01:02}
待撰写

\subsection{好约化与坏约化的曲线算例}
\label{g5:subsec:A01:03}
待撰写

\section{形式与非 Archimedes 空间}
\label{g5:sec:A02}
待撰写

\subsection{形式概形和理想定义}
\label{g5:subsec:A02:01}
待撰写

\subsection{Tate 代数及 affinoid 模型}
\label{g5:subsec:A02:02}
待撰写

\subsection{刚性与 Berkovich 的不同对象}
\label{g5:subsec:A02:03}
待撰写

\subsection{一般纤维和特殊纤维}
\label{g5:subsec:A02:04}
待撰写

\section{有理点与局部—整体}
\label{g5:sec:A03}
待撰写

\subsection{点的下降和模解释}
\label{g5:subsec:A03:01}
待撰写

\subsection{圆锥曲线的局部模型}
\label{g5:subsec:A03:02}
待撰写

\subsection{高度及有限性的外部输入}
\label{g5:subsec:A03:03}
待撰写

% !TeX root = ../../Book5.tex
% 纲要标题：## 附录 B　代数叠、导出几何与几何化
% 纲要位置：几何学计划纲要.md，第 6352 行。
% 正文状态：未撰写；本文件目前是框架。

% 纲要内容开始
% 
% > 在正文的具体构造之后延伸，不再次仅重复名词定义。
% 
% * **商叠与层范畴**
%   * BG 和 [X/G] 的显式对象
%   * 拟凝聚层与等变模
%   * 稳定子和粗模的失真
% 
% * **导出局部模型**
%   * 自交的高阶 Tor
%   * 余切复形和切复形的对偶条件
%   * 完美障碍理论的动机
%   * 一般导出叠的新语言
% 
% * **几何表示的低维模型**
%   * 旗簇上的线丛
%   * 联络、单值和解层
%   * 推拉函子与卷积
%   * 一般等价定理的范围
% 
% ---
% 
% 纲要内容结束

\chapter{代数叠、导出几何与几何化}
\label{g5:app:B}
\BookChapterNumber{g5:app:B}

待撰写

\section{商叠与层范畴}
\label{g5:sec:B01}
待撰写

\subsection{BG 和 [X/G] 的显式对象}
\label{g5:subsec:B01:01}
待撰写

\subsection{拟凝聚层与等变模}
\label{g5:subsec:B01:02}
待撰写

\subsection{稳定子和粗模的失真}
\label{g5:subsec:B01:03}
待撰写

\section{导出局部模型}
\label{g5:sec:B02}
待撰写

\subsection{自交的高阶 Tor}
\label{g5:subsec:B02:01}
待撰写

\subsection{余切复形和切复形的对偶条件}
\label{g5:subsec:B02:02}
待撰写

\subsection{完美障碍理论的动机}
\label{g5:subsec:B02:03}
待撰写

\subsection{一般导出叠的新语言}
\label{g5:subsec:B02:04}
待撰写

\section{几何表示的低维模型}
\label{g5:sec:B03}
待撰写

\subsection{旗簇上的线丛}
\label{g5:subsec:B03:01}
待撰写

\subsection{联络、单值和解层}
\label{g5:subsec:B03:02}
待撰写

\subsection{推拉函子与卷积}
\label{g5:subsec:B03:03}
待撰写

\subsection{一般等价定理的范围}
\label{g5:subsec:B03:04}
待撰写

% !TeX root = ../../Book5.tex
% 纲要标题：## 附录 C　实、热带与组合代数几何
% 纲要位置：几何学计划纲要.md，第 6375 行。
% 正文状态：未撰写；本文件目前是框架。

% 纲要内容开始
% 
% > 以实不等式和赋值退化为两条独立入口，保持各自的对象、拓扑和重数。
% 
% * **实点与半代数结构**
%   * 实点集和 Zariski 闭包
%   * Tarski–Seidenberg 的准确陈述
%   * 实曲线的连通分支
%   * 正性证书和复几何的区别
% 
% * **初始退化与热带化**
%   * 赋值、权和初始理想
%   * Gröbner 退化与平坦族
%   * 热带超曲面的计算
%   * 平衡条件与重数
%   * 平凡及非平凡赋值的区别
% 
% * **组合和枚举**
%   * 格多面体与次数
%   * 热带曲线的有限模型
%   * 对应定理的具体假设
%   * 组合计数不能无条件替代代数重数
% 
% ---
% 
% 纲要内容结束

\chapter{实、热带与组合代数几何}
\label{g5:app:C}
\BookChapterNumber{g5:app:C}

待撰写

\section{实点与半代数结构}
\label{g5:sec:C01}
待撰写

\subsection{实点集和 Zariski 闭包}
\label{g5:subsec:C01:01}
待撰写

\subsection{Tarski–Seidenberg 的准确陈述}
\label{g5:subsec:C01:02}
待撰写

\subsection{实曲线的连通分支}
\label{g5:subsec:C01:03}
待撰写

\subsection{正性证书和复几何的区别}
\label{g5:subsec:C01:04}
待撰写

\section{初始退化与热带化}
\label{g5:sec:C02}
待撰写

\subsection{赋值、权和初始理想}
\label{g5:subsec:C02:01}
待撰写

\subsection{Gröbner 退化与平坦族}
\label{g5:subsec:C02:02}
待撰写

\subsection{热带超曲面的计算}
\label{g5:subsec:C02:03}
待撰写

\subsection{平衡条件与重数}
\label{g5:subsec:C02:04}
待撰写

\subsection{平凡及非平凡赋值的区别}
\label{g5:subsec:C02:05}
待撰写

\section{组合和枚举}
\label{g5:sec:C03}
待撰写

\subsection{格多面体与次数}
\label{g5:subsec:C03:01}
待撰写

\subsection{热带曲线的有限模型}
\label{g5:subsec:C03:02}
待撰写

\subsection{对应定理的具体假设}
\label{g5:subsec:C03:03}
待撰写

\subsection{组合计数不能无条件替代代数重数}
\label{g5:subsec:C03:04}
待撰写



\backmatter
\cleardoublepage
\pagenumbering{roman}
\bookmarksetup{startatroot}
% !TeX root = ../Book5.tex
\cleardoublepage
% !TeX root = ../Book5.tex
\chapter*{后记}
\addcontentsline{toc}{chapter}{后记}
待撰写

\cleardoublepage
% !TeX root = ../Book5.tex
\PrintBookBibliography
% 仅在核读并引用真实资料后补入书目。

\cleardoublepage
% !TeX root = ../Book5.tex
\PrintBookIndex{chinese}{中文索引}

\cleardoublepage
% !TeX root = ../Book5.tex
\PrintBookIndex{foreign}{外文索引}

\cleardoublepage
% !TeX root = ../Book5.tex
\PrintBookIndex{symbols}{符号索引}


\end{document}
